% Generated mechanically from frozen input p04/ko/varieties.tex.
% Do not edit this generated file; edit the bound locale source instead.

% OK, start here.
%

\setcounter{chapter}{32}
\chapter{다양체}



\phantomsection
\label{varieties-section-phantom}


\section{서론}
\label{varieties-section-introduction}

\noindent
이 장에서는 다양체와, 더 일반적으로는 체 위의 스킴을 연구하기 시작한다.
기본 참고문헌은 \cite{EGA}이다.








\section{표기}
\label{varieties-section-notation}

\noindent
이 장 전체에서 문자 $k$는 기초체를 나타낸다.








\section{다양체}
\label{varieties-section-varieties}

\noindent
Stacks 프로젝트에서는 다음을 다양체의 정의로 사용한다.

\begin{definition}
\label{varieties-definition-variety}
$k$를 체라 하자. {\it 다양체}란 스킴 $X$가 $k$ 위에 있고,
$X$가 정수적이고 구조 사상
$X \to \Spec(k)$가 분리이며 유한형인 것을 말한다.
\end{definition}

\noindent
이 정의에는 다음과 같은 단점이 있다. $k'/k$가 체의 확대라고 하자.
$X$가 $k$ 위의 다양체라고 하자. 그러면 기저 변환
$X_{k'} = X \times_{\Spec(k)} \Spec(k')$는 반드시 $k'$ 위의 다양체인 것은 아니다.
이 현상은 뒤의 절들에서 더 일반적인 형태로 자세히 논의한다.
두 다양체의 곱도 다양체가 아닐 수 있다(이는 사실 같은 현상이다).
다음은 그 예이다.

\begin{example}
\label{varieties-example-product-not-a-variety}
$k = \mathbf{Q}$라 하자. $X = \Spec(\mathbf{Q}(i))$이고
$Y = \Spec(\mathbf{Q}(i))$라 하자. 그러면 곱
$X \times_{\Spec(k)} Y$, 즉 다양체 $X$와 $Y$의 곱은 가약이므로 다양체가 아니다.
(이는 $X$ 두 개의 분리합과 동형이다.)
\end{example}

\noindent
그러나 기초체가 대수적으로 닫혀 있으면 다양체들의 곱은 다양체이다.
이는 대수학 장의 결과에서 따르지만, 그곳에서는 훨씬 더 일반적인 상황을 다룬다.
여기서는 간단한 직접 증명도 제시한다.

\begin{lemma}
\label{varieties-lemma-product-varieties}
\begin{slogan}
대수적으로 닫힌 체 위에서 다양체들의 곱은 다양체이다.
\end{slogan}
$k$를 대수적으로 닫힌 체라 하자.
$X$, $Y$를 $k$ 위의 다양체라 하자.
그러면 $X \times_{\Spec(k)} Y$는 $k$ 위의 다양체이다.
\end{lemma}

\begin{proof}
사상 $X \times_{\Spec(k)} Y \to \Spec(k)$는 분리이고 유한형인 사상들
$X \times_{\Spec(k)} Y \to Y \to \Spec(k)$의 합성이므로 유한형이고 분리이다.
Morphisms, Lemmas \ref{morphisms-lemma-base-change-finite-type}와
\ref{morphisms-lemma-composition-finite-type}, 그리고
Schemes, Lemma \ref{schemes-lemma-separated-permanence}를 보라.
증명을 끝내려면 $X \times_{\Spec(k)} Y$가 정수적임을 보이면 충분하다.
$X = \bigcup_{i = 1, \ldots, n} U_i$와
$Y = \bigcup_{j = 1, \ldots, m} V_j$를 유한 아핀 열린 덮개라 하자.
각 $U_i \times_{\Spec(k)} V_j$가 정수적임을 보일 수 있으면
Properties, Lemmas \ref{properties-lemma-characterize-reduced},
\ref{properties-lemma-characterize-irreducible}, 그리고
\ref{properties-lemma-characterize-integral}에 의해 증명이 끝난다.
따라서 아핀인 경우로 환원된다.

\medskip\noindent
아핀인 경우는 다음 대수적 명제로 옮겨진다. $A$, $B$가 정역이고 유한 생성
$k$-대수라고 하자. 그러면 $A \otimes_k B$는 정역이다. 모순을 얻기 위해
$$
(\sum\nolimits_{i = 1, \ldots, n} a_i \otimes b_i)
(\sum\nolimits_{j = 1, \ldots, m} c_j \otimes d_j) = 0
$$
이 $A \otimes_k B$에서 성립하고 두 인수가 모두 $A \otimes_k B$에서
영이 아니라고 하자. $b_1, \ldots, b_n$이 $k$ 위에서 선형 독립인 $B$의 원소이고,
$d_1, \ldots, d_m$이 $k$ 위에서 선형 독립인 $B$의 원소라고 가정해도 된다.
물론 $a_1$과 $c_1$이 $A$에서 영이 아니라고 가정해도 된다.
따라서 $D(a_1c_1) \subset \Spec(A)$는 공집합이 아니다. Hilbert 영점 정리
(Algebra, Theorem \ref{algebra-theorem-nullstellensatz})에 의해
극대 아이디얼 $\mathfrak m \subset A$를 찾을 수 있으며, 이는 $D(a_1c_1)$에 포함되고
$A/\mathfrak m = k$를 만족한다. 마지막 등식은 $k$가 대수적으로 닫혀 있기 때문이다.
$\overline{a}_i, \overline{c}_j$를 $a_i, c_j$의
$A/\mathfrak m = k$에서의 잉여류라 하자. 위 등식은
$$
(\sum\nolimits_{i = 1, \ldots, n} \overline{a}_i b_i)
(\sum\nolimits_{j = 1, \ldots, m} \overline{c}_j d_j) = 0
$$
이 된다. 이는 $\mathfrak m \in D(a_1c_1)$라는 사실, $b_1, \ldots, b_n$과
$d_1, \ldots, d_m$의 선형 독립성, 그리고 $B$가 정역이라는 사실과 모순이다.
\end{proof}






\section{다양체와 유리 사상}
\label{varieties-section-varieties-rational-maps}

\noindent
$k$를 체라 하자. $X$와 $Y$를 $k$ 위의 다양체라 하자.
{\it $X$에서 $Y$로 가는 다양체의 유리 사상}이라는 말은 Morphisms,
Definition \ref{morphisms-definition-rational-map}에서 정의된 대로
$\Spec(k)$-유리 사상, 즉 스킴 $X$에서 스킴 $Y$로 가는 유리 사상을 뜻한다.
관례상 “다양체의 유리 사상”이라는 말은 다양체들의 (공통) 기초체를 가리키지 않는다.
반면 일반 스킴에서는 유리 사상과 주어진 바탕 위의 유리 사상을 구별한다.

\medskip\noindent
이 절의 제목은 다음 기본 정리를 가리킨다.

\begin{theorem}
\label{varieties-theorem-varieties-rational-maps}
$k$를 체라 하자. 다양체와 우세 유리 사상의 범주는
유한 생성 체 확대 $K/k$의 범주와 동치이다.
\end{theorem}

\begin{proof}
$X$와 $Y$를 다양체라 하고, 그 일반점을 각각 $x \in X$와 $y \in Y$라 하자.
$X$에서 $Y$로 가는 우세 유리 사상은 정확히 $x$를 $y$로 보내는 유리 사상이다
(Morphisms, Definition \ref{morphisms-definition-dominant-rational}과
그 뒤의 논의를 보라). 따라서 우세 유리 사상 $X \supset U \to Y$가 주어지면
함수체의 사상
$$
k(Y) = \kappa(y) = \mathcal{O}_{Y, y}
\longrightarrow
\mathcal{O}_{X, x} = \kappa(x) = k(X)
$$
을 얻는다. 역으로, 그러한 $k$-대수 사상은 원천과 목표가 체이므로 자동으로
국소 사상이며, Morphisms, Lemma \ref{morphisms-lemma-rational-map-finite-presentation}에
의해 하나의 우세 유리 사상을 유일하게 결정한다. 이로써 완전 충실한 함자를 얻는다.
증명을 끝내려면 모든 유한 생성 체 확대 $K/k$가 본질적 상에 속함을 보이면 충분하다.
$K/k$가 유한 생성이므로 유한형 $k$-대수 $A \subset K$가 존재하여
$K$는 $A$의 분수체가 된다. 그러면 $X = \Spec(A)$는 함수체가 $K$인 다양체이다.
\end{proof}

\noindent
$k$를 체라 하자. $X$와 $Y$를 $k$ 위의 다양체라 하자.
{\it $X$와 $Y$가 쌍유리 다양체이다}라는 말은 Morphisms,
Definition \ref{morphisms-definition-birational}에서 정의된 대로
$X$와 $Y$가 $\Spec(k)$-쌍유리라는 뜻으로 사용한다.
관례상 “쌍유리 다양체”라는 말은 다양체들의 (공통) 기초체를 가리키지 않는다.
반면 일반 기약 스킴에서는 쌍유리인 것과 주어진 바탕 위에서 쌍유리인 것을 구별한다.

\begin{lemma}
\label{varieties-lemma-birational-varieties}
$X$와 $Y$를 체 $k$ 위의 다양체라 하자.
다음 조건들은 서로 동치이다.
\begin{enumerate}
\item $X$와 $Y$는 쌍유리 다양체이다.
\item 함수체 $k(X)$와 $k(Y)$는 동형이다.
\item $X$와 $Y$에 서로 동형인 공집합이 아닌 열린집합들이 다양체로서 존재한다.
\item 열린집합 $U \subset X$와 다양체의 쌍유리 사상 $U \to Y$가 존재한다.
\end{enumerate}
\end{lemma}

\begin{proof}
이는 Morphisms, Lemma
\ref{morphisms-lemma-criterion-birational-finite-presentation}의 특수한 경우이다.
\end{proof}







\section{체의 변경과 국소환}
\label{varieties-section-local-rings}

\noindent
기초체를 확대할 때 국소환에 일어나는 일에 관한 몇 가지 예비 결과를 제시한다.

\begin{lemma}
\label{varieties-lemma-change-fields-flat}
$K/k$를 체의 확대라 하자. $X$를 $k$ 위의 스킴이라 하고
$Y = X_K$로 두자. $y \in Y$의 상이 $x \in X$이면 다음이 성립한다.
\begin{enumerate}
\item $\mathcal{O}_{X, x} \to \mathcal{O}_{Y, y}$는 충실히 평탄한 국소환 준동형이다.
\item $\mathfrak p_0 = \Ker(\kappa(x) \otimes_k K \to \kappa(y))$로 두면
$\kappa(y) = \kappa(\mathfrak p_0)$이다.
\item $\mathcal{O}_{Y, y} = (\mathcal{O}_{X, x} \otimes_k K)_\mathfrak p$이며,
여기서 $\mathfrak p \subset \mathcal{O}_{X, x} \otimes_k K$는
$\mathfrak p_0$의 역상이다.
\item 다음이 성립한다.
$\mathcal{O}_{Y, y}/\mathfrak m_x\mathcal{O}_{Y, y} =
(\kappa(x) \otimes_k K)_{\mathfrak p_0}$
\end{enumerate}
\end{lemma}

\begin{proof}
$X = \Spec(A)$가 아핀이라고 가정해도 된다. 그러면 $Y = \Spec(A \otimes_k K)$이다.
$K$는 $k$ 위에서 평탄하므로 $A \to A \otimes_k K$가 평탄하다.
따라서 $Y \to X$는 평탄하고, Algebra, Lemma
\ref{algebra-lemma-local-flat-ff}도 사용하면 첫째 주장을 얻는다.
둘째 주장은 Schemes, Lemma \ref{schemes-lemma-points-fibre-product}에서 따른다.
이제 $y$는 소 아이디얼 $\mathfrak q \subset A \otimes_k K$에 대응하고,
$x$는 $\mathfrak r = A \cap \mathfrak q$에 대응한다. 그러면 $\mathfrak p_0$는
유도된 사상 $\kappa(\mathfrak r) \otimes_k K \to \kappa(\mathfrak q)$의 핵이다.
국소환 사이의 사상은
$$
A_\mathfrak r \longrightarrow (A \otimes_k K)_\mathfrak q
$$
이다. 이 사상을
$A_\mathfrak r \otimes_k K = (A \otimes_k K)_{\mathfrak r}$를 통하여 인수분해하면
$$
A_\mathfrak r \longrightarrow A_\mathfrak r \otimes_k K
\longrightarrow (A \otimes_k K)_\mathfrak q
$$
를 얻고, 둘째 화살표는 어떤 소 아이디얼에서의 국소화이다. 이 소 아이디얼은
$\mathfrak p_0$의 역상이며(세부 사항은 생략한다), 이로써 (3)이 증명된다.
(4)는 (3), 그리고 국소화와 $- \otimes_k K$가 완전 함자라는 사실에서 따른다.
\end{proof}

\begin{lemma}
\label{varieties-lemma-change-fields-algebraic-dim}
표기는 보조정리 \ref{varieties-lemma-change-fields-flat}과 같다고 하자.
$X$가 $k$ 위에서 국소 유한형이라고 가정하자. 그러면
$$
\dim(\mathcal{O}_{Y, y}/\mathfrak m_x\mathcal{O}_{Y, y}) =
\text{trdeg}_k(\kappa(x)) - \text{trdeg}_K(\kappa(y)) =
\dim(\mathcal{O}_{Y, y}) - \dim(\mathcal{O}_{X, x})
$$
이다.
\end{lemma}

\begin{proof}
이는 Algebra, Lemma
\ref{algebra-lemma-inequalities-under-field-extension}을 다시 서술한 것이다.
\end{proof}

\begin{lemma}
\label{varieties-lemma-change-fields-algebraic-unramified}
표기는 보조정리 \ref{varieties-lemma-change-fields-flat}과 같다고 하자.
$X$가 $k$ 위에서 국소 유한형이고,
$\dim(\mathcal{O}_{X, x}) = \dim(\mathcal{O}_{Y, y})$이며,
$\kappa(x) \otimes_k K$가 축약환이라고 가정하자
(예를 들어 $\kappa(x)/k$가 분리적이거나 $K/k$가 분리적인 경우이다).
그러면 $\mathfrak m_x \mathcal{O}_{Y, y} = \mathfrak m_y$이다.
\end{lemma}

\begin{proof}
(괄호 안의 주장은 Algebra, Lemma
\ref{algebra-lemma-separable-extension-preserves-reducedness}에서 따른다.)
보조정리 \ref{varieties-lemma-change-fields-flat}과
\ref{varieties-lemma-change-fields-algebraic-dim}을 결합하면
$\mathcal{O}_{Y, y}/\mathfrak m_x \mathcal{O}_{Y, y}$는 차원이 $0$이고
축약이다. 따라서 이는 체이다.
\end{proof}





\section{기하적으로 축약인 스킴}
\label{varieties-section-geometrically-reduced}

\noindent
$X$가 체 위의 축약 스킴이어도, 기초체를 확대하면 $X$가 축약이 아니게 될 수 있다.
기하적으로 축약인 스킴에는 이런 일이 일어나지 않는다.

\begin{definition}
\label{varieties-definition-geometrically-reduced}
$k$를 체라 하자.
$X$를 $k$ 위의 스킴이라 하자.
\begin{enumerate}
\item $x \in X$를 점이라 하자. $X$가 {\it $x$에서 기하적으로 축약}이라고 함은
모든 체 확대 $k'/k$와 모든 점 $x' \in X_{k'}$가 $x$ 위에 놓일 때
국소환 $\mathcal{O}_{X_{k'}, x'}$가 축약인 경우이다.
\item $X$가 $k$ 위에서 {\it 기하적으로 축약}이라고 함은
$X$가 $X$의 모든 점에서 기하적으로 축약인 경우이다.
\end{enumerate}
\end{definition}

\noindent
처음에는 이 개념이 다소 신비로워 보일 수 있지만, 실제로는 대수학 장에서
논의한 개념과 같은 것이다. 다음 기본 결과들이 그 관계를 설명한다.
보조정리 \ref{varieties-lemma-geometrically-reduced-dense-smooth-open}도 보라.
그 보조정리에 따르면 체 위의 다양체에 대하여 “기하적으로 축약”은
“일반점에서 매끄러움”과 동치이다.

\begin{lemma}
\label{varieties-lemma-geometrically-reduced-at-point}
\begin{slogan}
기하적 축약성은 국소환에서 확인할 수 있다.
\end{slogan}
$k$를 체라 하자.
$X$를 $k$ 위의 스킴이라 하자.
$x \in X$라 하자.
다음 조건들은 서로 동치이다.
\begin{enumerate}
\item $X$는 $x$에서 기하적으로 축약이다.
\item 환 $\mathcal{O}_{X, x}$는 $k$ 위에서 기하적으로 축약이다
(Algebra, Definition \ref{algebra-definition-geometrically-reduced}을 보라).
\end{enumerate}
\end{lemma}

\begin{proof}
(1)을 가정하자. 특히 $\mathcal{O}_{X, x}$는 축약이다.
$k'/k$를 유한 순수 비분리 체 확대라 하고 환
$\mathcal{O}_{X, x} \otimes_k k'$를 생각하자. Algebra, Lemma
\ref{algebra-lemma-p-ring-map}에 의해 그 스펙트럼은
$\mathcal{O}_{X, x}$의 스펙트럼과 같다. 따라서 이 환도 국소환이다
(Algebra, Lemma \ref{algebra-lemma-characterize-local-ring}).
그러므로 유일한 점 $x' \in X_{k'}$가 $x$ 위에 놓이고
$\mathcal{O}_{X_{k'}, x'} \cong \mathcal{O}_{X, x} \otimes_k k'$이다.
보조정리 \ref{varieties-lemma-change-fields-flat}을 보라. 가정에 의해 이 환은 축약이다.
따라서 Algebra, Lemma
\ref{algebra-lemma-geometrically-reduced-finite-purely-inseparable-extension}에 의해 (2)를 얻는다.

\medskip\noindent
(2)를 가정하자. $k'/k$를 체 확대라 하자.
$\Spec(k') \to \Spec(k)$가 전사이므로 $X_{k'} \to X$도 전사이다
(Morphisms, Lemma \ref{morphisms-lemma-base-change-surjective}).
임의의 점 $x' \in X_{k'}$를 $x$ 위에서 잡자.
보조정리 \ref{varieties-lemma-change-fields-flat}에 의해 국소환
$\mathcal{O}_{X_{k'}, x'}$는 환 $\mathcal{O}_{X, x} \otimes_k k'$의 국소화이다.
따라서 가정에 의해 축약이고, (1)이 증명된다.
\end{proof}

\noindent
표수가 영이면 이 개념은 특별한 내용을 갖지 않는다.

\begin{lemma}
\label{varieties-lemma-perfect-reduced}
$X$를 완전체 $k$ 위의 스킴이라 하자(예를 들어 $k$의 표수가 영인 경우이다).
$x \in X$라 하자. $\mathcal{O}_{X, x}$가 축약이면
$X$는 $x$에서 기하적으로 축약이다.
$X$가 축약이면 $X$는 $k$ 위에서 기하적으로 축약이다.
\end{lemma}

\begin{proof}
첫째 주장은 보조정리 \ref{varieties-lemma-geometrically-reduced-at-point},
Algebra, Lemma \ref{algebra-lemma-separable-extension-preserves-reducedness}, 그리고
완전체의 정의(Algebra, Definition \ref{algebra-definition-perfect})에서 따른다.
둘째 주장은 첫째 주장으로부터 따른다.
\end{proof}

\begin{lemma}
\label{varieties-lemma-geometrically-reduced}
$k$를 표수 $p > 0$인 체라 하고 $X$를 $k$ 위의 스킴이라 하자.
다음 조건들은 서로 동치이다.
\begin{enumerate}
\item $X$는 기하적으로 축약이다.
\item $X_{k'}$는 모든 체 확대 $k'/k$에 대해 축약이다.
\item $X_{k'}$는 모든 유한 순수 비분리 체 확대 $k'/k$에 대해 축약이다.
\item $X_{k^{1/p}}$는 축약이다.
\item $X_{k^{perf}}$는 축약이다.
\item $X_{\bar k}$는 축약이다.
\item 모든 아핀 열린집합 $U \subset X$에 대해 환 $\mathcal{O}_X(U)$는
기하적으로 축약이다(Algebra, Definition
\ref{algebra-definition-geometrically-reduced}을 보라).
\end{enumerate}
\end{lemma}

\begin{proof}
(1)을 가정하자. 모든 체 확대 $k'/k$와 모든 점 $x' \in X_{k'}$에 대해
$X_{k'}$의 $x'$에서의 국소환은 축약이다. 다시 말해 $X_{k'}$는 축약이다.
따라서 (2)가 성립한다.

\medskip\noindent
(2)를 가정하자. $U \subset X$를 아핀 열린집합이라 하자. 모든 체 확대
$k'/k$에 대해 스킴 $X_{k'}$가 축약이므로
$U_{k'} = \Spec(\mathcal{O}(U)\otimes_k k')$도 축약이고,
따라서 $\mathcal{O}(U)\otimes_k k'$는 축약이다
(Properties, Section \ref{properties-section-integral}을 보라).
다시 말해 $\mathcal{O}(U)$는 기하적으로 축약이므로 (7)이 성립한다.

\medskip\noindent
(7)을 가정하자. 임의의 체 확대 $k'/k$에 대해 기저 변환 $X_{k'}$는
환 $\mathcal{O}_X(U) \otimes_k k'$들의 스펙트럼을 이어붙여 얻는데,
여기서 $U$는 $X$의 아핀 열린집합들을 달린다
(Schemes, Section \ref{schemes-section-fibre-products}를 보라).
따라서 $X_{k'}$는 축약이고, (1)이 성립한다.

\medskip\noindent
이로써 (1), (2), (7)이 서로 동치임을 증명했다. 이 조건들은
환 $\mathcal{O}_X(U)$에 Algebra, Lemma
\ref{algebra-lemma-geometrically-reduced-finite-purely-inseparable-extension}을
적용할 수 있기 때문이다. 여기서 $U \subset X$는 아핀 열린집합을 달린다.
따라서 (3), (4), (5), (6)과도 동치이다.
\end{proof}

\begin{lemma}
\label{varieties-lemma-check-only-finite-inseparable-extensions}
$k$를 표수 $p > 0$인 체라 하고 $X$를 $k$ 위의 스킴이라 하자.
$x \in X$라 하자. 다음 조건들은 서로 동치이다.
\begin{enumerate}
\item $X$는 $x$에서 기하적으로 축약이다.
\item $\mathcal{O}_{X_{k'}, x'}$는 $k'$가 $k$의 모든 유한 순수 비분리
체 확대이고 유일한 점 $x' \in X_{k'}$가 $x$ 위에 놓일 때 축약이다.
\item $\mathcal{O}_{X_{k^{1/p}}, x'}$는 유일한 점
$x' \in X_{k^{1/p}}$가 $x$ 위에 놓일 때 축약이다.
\item $\mathcal{O}_{X_{k^{perf}}, x'}$는 유일한 점
$x' \in X_{k^{perf}}$가 $x$ 위에 놓일 때 축약이다.
\end{enumerate}
\end{lemma}

\begin{proof}
$k'/k$가 순수 비분리이면 $X_{k'} \to X$는 바탕 위상공간들 사이의
위상동형을 유도한다. Algebra, Lemma \ref{algebra-lemma-p-ring-map}을 보라.
따라서 진술에 나온 $x'$는 $x$ 위에 놓이는 점으로서 유일하다. 더욱이 이 경우
$\mathcal{O}_{X_{k'}, x'} = \mathcal{O}_{X, x} \otimes_k k'$이다.
그러므로 위의 보조정리 \ref{varieties-lemma-geometrically-reduced-at-point}과 Algebra, Lemma
\ref{algebra-lemma-geometrically-reduced-finite-purely-inseparable-extension}에서
결론이 따른다.
\end{proof}

\begin{lemma}
\label{varieties-lemma-geometrically-reduced-upstairs}
$k$를 체라 하자.
$X$를 $k$ 위의 스킴이라 하자.
$k'/k$를 체 확대라 하자.
$x \in X$를 점이라 하고, 점 $x' \in X_{k'}$를 $x$ 위에서 잡자.
다음 조건들은 서로 동치이다.
\begin{enumerate}
\item $X$는 $x$에서 기하적으로 축약이다.
\item $X_{k'}$는 $x'$에서 기하적으로 축약이다.
\end{enumerate}
특히 $X$가 $k$ 위에서 기하적으로 축약일 필요충분조건은
$X_{k'}$가 $k'$ 위에서 기하적으로 축약인 것이다.
\end{lemma}

\begin{proof}
(1)이 (2)를 함의함은 분명하다. (2)를 가정하자.
$k''/k$를 유한 순수 비분리 체 확대라 하고 점
$x'' \in X_{k''}$를 $x$ 위에서 잡자(실제로 유일하다). 공통 체 확대
$k'''/k$를 $k''/k$와 $k'/k$에 대해 택하자. Fields, Lemma
\ref{fields-lemma-common-extension-field}을 보라. 점 $x''' \in X_{k'''}$를
$x'$과 $x''$ 양쪽 위에서 잡자. 국소환의 사상
$$
\mathcal{O}_{X_{k''}, x''} \longrightarrow \mathcal{O}_{X_{k'''}, x'''}.
$$
을 생각하자. 이는 평탄한 국소환 준동형이므로 충실히 평탄하다.
(2)에 의해 오른쪽 국소환은 축약이다. 따라서 Algebra, Lemma
\ref{algebra-lemma-descent-reduced}에 의해
$\mathcal{O}_{X_{k''}, x''}$가 축약임을 알 수 있다.
그러므로 보조정리 \ref{varieties-lemma-check-only-finite-inseparable-extensions}에 의해
$X$는 $x$에서 기하적으로 축약이다.
\end{proof}

\begin{lemma}
\label{varieties-lemma-geometrically-reduced-any-base-change}
$k$를 체라 하고 $X$, $Y$를 $k$ 위의 스킴들이라 하자.
\begin{enumerate}
\item $X$가 $x$에서 기하적으로 축약이고 $Y$가 축약이면,
$X \times_k Y$는 $x$ 위에 놓이는 모든 점에서 축약이다.
\item $X$가 $k$ 위에서 기하적으로 축약이고 $Y$가 축약이면
$X \times_k Y$는 축약이다.
\item $X$와 $Y$가 $k$ 위에서 기하적으로 축약이면
$X \times_k Y$는 기하적으로 축약이다.
\item $k$가 완전체이고 $X$와 $Y$가 축약이면 $X \times_k Y$는 축약이다.
\item 여기에 더 추가할 것.
\end{enumerate}
\end{lemma}

\begin{proof}
(1)을 증명하려면 보조정리 \ref{varieties-lemma-geometrically-reduced-at-point}과 Algebra, Lemma
\ref{algebra-lemma-geometrically-reduced-any-reduced-base-change}을 결합한다.
(2)를 증명하려면 보조정리 \ref{varieties-lemma-geometrically-reduced}와 Algebra, Lemma
\ref{algebra-lemma-geometrically-reduced-any-reduced-base-change}을 결합한다.
(3)을 증명하려면
$(X \times_k Y)_{\overline{k}} =
X_{\overline{k}} \times_{\overline{k}} Y_{\overline{k}}$임을 관찰하고,
(2)와 보조정리 \ref{varieties-lemma-geometrically-reduced}을 사용한다.
(4)를 증명하려면 (3)과 보조정리 \ref{varieties-lemma-perfect-reduced}을 결합한다.
\end{proof}

\begin{lemma}
\label{varieties-lemma-generic-points-geometrically-reduced}
$k$를 체라 하자.
$X$를 $k$ 위의 스킴이라 하자.
\begin{enumerate}
\item $x' \leadsto x$가 특수화이고 $X$가 $x$에서 기하적으로 축약이면,
$X$는 $x'$에서도 기하적으로 축약이다.
\item $x \in X$가 다음을 만족한다고 하자. (a) $\mathcal{O}_{X, x}$는 축약이고,
(b) 모든 특수화 $x' \leadsto x$에 대해, $x'$가 $X$의 한 기약 성분의
일반점이면 $X$는 $x'$에서 기하적으로 축약이다.
그러면 $X$는 $x$에서 기하적으로 축약이다.
\item $X$가 축약이고 $X$의 기약 성분들의 모든 일반점에서 기하적으로 축약이면,
$X$는 기하적으로 축약이다.
\item $X$가 $k$ 위의 다양체이면, $X$가 기하적으로 축약일 필요충분조건은
그 함수체가 $k$의 분리 확대인 것이다.
\end{enumerate}
\end{lemma}

\begin{proof}
(1)은 보조정리 \ref{varieties-lemma-geometrically-reduced-at-point}과 다음 사실에서 따른다.
$A$가 기하적으로 축약인 $k$-대수이면, $S^{-1}A$도 기하적으로 축약인
$k$-대수이다. 여기서 $S$는 $A$의 임의의 곱셈적 부분집합이다. Algebra, Lemma
\ref{algebra-lemma-geometrically-reduced-permanence}를 보라.

\medskip\noindent
$A = \mathcal{O}_{X, x}$로 두자. (2)의 가정 (a), (b)는 $A$가 축약이고,
$A_{\mathfrak q}$가 $k$ 위에서 기하적으로 축약임을 뜻한다. 여기서
$\mathfrak q$는 $A$의 임의의 극소 소 아이디얼이다. 따라서 $A$는 $k$ 위에서 기하적으로 축약이다.
Algebra, Lemma \ref{algebra-lemma-generic-points-geometrically-reduced}을 보라.
그러므로 보조정리 \ref{varieties-lemma-geometrically-reduced-at-point}에 의해
$X$는 $x$에서 기하적으로 축약이다.

\medskip\noindent
(3)은 (2)에서 자명하게 따른다. 마지막으로 (4)는 (3)과 Algebra, Lemma
\ref{algebra-lemma-characterize-separable-field-extensions}에서 따른다.
\end{proof}

\begin{lemma}
\label{varieties-lemma-Noetherian-geometrically-reduced-at-point}
$k$를 체라 하자.
$X$를 $k$ 위의 스킴이라 하자.
$x \in X$라 하자.
$X$가 국소 뇌터이고 $x$에서 기하적으로 축약이라고 가정하자.
그러면 열린 근방 $U \subset X$가 $x$의 근방으로 존재하여 $k$ 위에서 기하적으로 축약이다.
\end{lemma}

\begin{proof}
$X$가 국소 뇌터이고 $x$에서 기하적으로 축약이라고 가정하자.
Properties, Lemma \ref{properties-lemma-ring-affine-open-injective-local-ring}에 의해
아핀 열린 근방 $U \subset X$를 $x$의 근방으로 찾아
$R = \mathcal{O}_X(U) \to \mathcal{O}_{X, x}$가 단사이게 할 수 있다.
보조정리 \ref{varieties-lemma-geometrically-reduced-at-point}에 의해 가정은
$\mathcal{O}_{X, x}$가 $k$ 위에서 기하적으로 축약임을 뜻한다.
Algebra, Lemma \ref{algebra-lemma-subalgebra-separable}에 의해
$R$도 $k$ 위에서 기하적으로 축약이고, 따라서 $U$도 기하적으로 축약이다.
\end{proof}

\begin{example}
\label{varieties-example-not-geometrically-reduced}
$k = \mathbf{F}_p(s, t)$, 즉 소체의 순수 초월 확대라 하자. 다양체
$X = \Spec(k[x, y]/(1 + sx^p + ty^p))$를 생각하자.
$k'/k$를 $s$와 $t$ 모두가 $p$제곱근을 $k'$에서 갖는 임의의 확대라 하자.
그러면 기저 변환 $X_{k'}$는 축약이 아니다. 실제로 환
$k'[x, y]/(1 + s x^p + ty^p)$는 원소
$1 + s^{1/p}x + t^{1/p}y$를 포함하며, 그 $p$제곱은 영이지만 그 자체는 영이 아니다
(아이디얼 $(1 + sx^p + ty^p)$는 차수가 $< p$인 영이 아닌 원소를 전혀 포함하지 않기 때문이다).
\end{example}

\begin{lemma}
\label{varieties-lemma-finite-extension-geometrically-reduced}
$k$를 체라 하자.
$X \to \Spec(k)$가 국소 유한형이라고 하자.
$X$가 유한 개의 기약 성분을 갖는다고 가정하자.
그러면 유한 순수 비분리 확대 $k'/k$가 존재하여
$(X_{k'})_{red}$는 $k'$ 위에서 기하적으로 축약이다.
\end{lemma}

\begin{proof}
이 보조정리를 증명할 때 $X$를 그 축약 $X_{red}$로 바꾸어도 된다.
따라서 $X$가 축약이고 $k$ 위에서 국소 유한형이라고 가정해도 된다.
$x_1, \ldots, x_n \in X$를 $X$의 기약 성분들의 일반점이라 하자.
모든 순수 비분리 대수적 확대 $k'/k$에 대해 사상
$(X_{k'})_{red} \to X$는 위상동형이다. Algebra, Lemma
\ref{algebra-lemma-p-ring-map}을 보라. 따라서 점들 $x'_1, \ldots, x'_n$은
$x_1, \ldots, x_n$ 위에 놓이고 $(X_{k'})_{red}$의 기약 성분들의 일반점이다.
$X$가 축약이므로 국소환 $K_i = \mathcal{O}_{X, x_i}$는 체이다.
Algebra, Lemma \ref{algebra-lemma-minimal-prime-reduced-ring}을 보라.
$X$가 $k$ 위에서 국소 유한형이므로 체 확대 $K_i/k$는 유한 생성 체 확대이다.
마지막으로 국소환 $\mathcal{O}_{(X_{k'})_{red}, x'_i}$는 체
$(K_i \otimes_k k')_{red}$이다. Algebra, Lemma
\ref{algebra-lemma-make-separable}에 의해 유한 순수 비분리 확대 $k'/k$를 찾아
$(K_i \otimes_k k')_{red}$가 $k'$의 분리 체 확대가 되게 할 수 있다.
특히 각각의 $(K_i \otimes_k k')_{red}$는 $k'$ 위에서 기하적으로 축약이다.
Algebra, Lemma \ref{algebra-lemma-characterize-separable-field-extensions}을 보라.
이제 보조정리 \ref{varieties-lemma-generic-points-geometrically-reduced}의 (3)에 의해
$(X_{k'})_{red}$는 기하적으로 축약이다.
\end{proof}







\section{기하적으로 연결인 스킴}
\label{varieties-section-geometrically-connected}

\noindent
$X$가 체 위의 연결 스킴이어도 기초체를 확대하면 $X$가 비연결이 될 수 있다.
기하적으로 연결인 스킴에서는 이런 일이 일어나지 않는다.

\begin{definition}
\label{varieties-definition-geometrically-connected}
$X$를 체 $k$ 위의 스킴이라 하자. $X$가 $k$ 위에서
{\it 기하적으로 연결}이라고 함은 스킴 $X_{k'}$가,
$k'$가 $k$의 임의의 체 확대일 때마다, 연결인 경우이다.
\end{definition}

\noindent
관례상 연결 위상공간은 공집합이 아니다. 따라서 특히 기하적으로 연결인
스킴은 공집합이 아니다. 다음은 기하적으로 연결이 아닌 다양체의 예이다.

\begin{example}
\label{varieties-example-not-geometrically-irreducible}
$k = \mathbf{Q}$라 하자. 스킴
$X = \Spec(\mathbf{Q}(i))$는 $\Spec(\mathbf{Q})$ 위의 다양체이다.
그러나 기저 변환 $X_{\mathbf{C}}$는
$\mathbf{C} \otimes_{\mathbf{Q}} \mathbf{Q}(i) \cong
\mathbf{C} \times \mathbf{C}$의 스펙트럼이며, 이는
$\Spec(\mathbf{C})$ 두 개의 서로소 합이다.
따라서 이것은 실제로 기하적으로 연결이 아닌 다양체의 예이다.
\end{example}

\begin{lemma}
\label{varieties-lemma-geometrically-connected-check-after-extension}
$X$를 체 $k$ 위의 스킴이라 하고 $k'/k$를 체 확대라 하자.
$X$가 $k$ 위에서 기하적으로 연결일 필요충분조건은
$X_{k'}$가 $k'$ 위에서 기하적으로 연결인 것이다.
\end{lemma}

\begin{proof}
$X$가 $k$ 위에서 기하적으로 연결이면 $X_{k'}$가 $k'$ 위에서
기하적으로 연결임은 분명하다. 역으로, 임의의 체 확대 $k''/k$에 대해
공통 체 확대 $k'''/k'$와 $k'''/k''$가 존재한다. Fields, Lemma
\ref{fields-lemma-common-extension-field}을 보라. 사상
$X_{k'''} \to X_{k''}$는 체들의 스펙트럼 사이 전사 사상의
기저 변환이므로 전사이다. 따라서 $X_{k'''}$의 연결성은
$X_{k''}$의 연결성을 함의한다. 그러므로 $X_{k'}$가 $k'$ 위에서
기하적으로 연결이면 $X$는 $k$ 위에서 기하적으로 연결이다.
\end{proof}

\begin{lemma}
\label{varieties-lemma-bijection-connected-components}
$k$를 체라 하자. $X$, $Y$를 $k$ 위의 스킴이라 하자.
$X$가 $k$ 위에서 기하적으로 연결이라고 가정하자. 그러면 사영 사상
$$
p : X \times_k Y \longrightarrow Y
$$
은 연결 성분들 사이의 전단사를 유도한다.
\end{lemma}

\begin{proof}
$p$의 스킴론적 올들은 기하적으로 연결인 스킴 $X$를 체 확대에 따라
기저 변환한 것이므로 연결이다. 또한 스킴론적 올은 집합론적 올과
위상동형이다. Schemes, Lemma
\ref{schemes-lemma-fibre-topological}을 보라.
Morphisms, Lemma
\ref{morphisms-lemma-scheme-over-field-universally-open}에 의해
사상 $p$는 열린 사상이다. 이제 Topology, Lemma
\ref{topology-lemma-connected-fibres-connected-components}를 적용하면 된다.
\end{proof}

\begin{lemma}
\label{varieties-lemma-affine-geometrically-connected}
$k$를 체라 하고 $A$를 $k$-대수라 하자. 그러면
$X = \Spec(A)$가 $k$ 위에서 기하적으로 연결일 필요충분조건은
$A$가 $k$ 위에서 기하적으로 연결인 것이다(Algebra, Definition
\ref{algebra-definition-geometrically-connected}을 보라).
\end{lemma}

\begin{proof}
정의에서 곧바로 따른다.
\end{proof}

\begin{lemma}
\label{varieties-lemma-separably-closed-field-connected-components}
$k'/k$를 체 확대라 하고 $X$를 $k$ 위의 스킴이라 하자.
$k$가 분리 대수적으로 닫혀 있다고 가정하자. 그러면 사상
$X_{k'} \to X$는 연결 성분들의 전단사를 유도한다. 특히
$X$가 $k$ 위에서 기하적으로 연결일 필요충분조건은 $X$가 연결인 것이다.
\end{lemma}

\begin{proof}
$k$가 분리 대수적으로 닫혀 있으므로 $k'$는 $k$ 위에서
기하적으로 연결이다. Algebra, Lemma
\ref{algebra-lemma-separably-closed-connected-implies-geometric}을 보라.
따라서 위의 보조정리 \ref{varieties-lemma-affine-geometrically-connected}에 의해
$Z = \Spec(k')$는 $k$ 위에서 기하적으로 연결이다.
$X_{k'} = Z \times_k X$이므로 결과는 보조정리
\ref{varieties-lemma-bijection-connected-components}의 특수한 경우이다.
\end{proof}

\begin{lemma}
\label{varieties-lemma-characterize-geometrically-connected}
$k$를 체라 하고 $X$를 $k$ 위의 스킴이라 하자.
$\overline{k}$를 $k$의 분리 대수적 폐포라 하자.
$X$가 기하적으로 연결일 필요충분조건은 기저 변환
$X_{\overline{k}}$가 연결인 것이다.
\end{lemma}

\begin{proof}
$X_{\overline{k}}$가 연결이라고 가정하자. $k'/k$를 체 확대라 하자.
체 확대 $\overline{k}'/\overline{k}$가 존재하여
$k'$가 $\overline{k}'$ 안에 $k$의 확대체로서 매장된다. 보조정리
\ref{varieties-lemma-separably-closed-field-connected-components}에 의해
$X_{\overline{k}'}$는 연결이다. 사상
$X_{\overline{k}'} \to X_{k'}$가 전사이므로 원하는 대로
$X_{k'}$가 연결임을 얻는다.
\end{proof}

\begin{lemma}
\label{varieties-lemma-descend-open}
$k$를 체라 하고 $X$를 $k$ 위의 스킴이라 하자.
$A$를 $k$-대수라 하고 $V \subset X_A$를 준콤팩트 열린집합이라 하자.
그러면 유한 생성 $k$-부분대수 $A' \subset A$와 준콤팩트 열린집합
$V' \subset X_{A'}$가 존재하여 $V = V'_A$이다.
\end{lemma}

\begin{proof}
$X$가 추가로 준분리이면 이는 Limits, Lemma
\ref{limits-lemma-descend-opens}에서 따른다는 점을 먼저 언급한다.
$U_1, \ldots, U_n$을 $X$의 유한 개 아핀 열린집합으로 택하여
$V \subset \bigcup U_{i, A}$가 되게 하자. $U_i = \Spec(R_i)$라 하자.
$V$가 준콤팩트이므로 유한 개의 원소
$f_{ij} \in R_i \otimes_k A$, $j = 1, \ldots, n_i$를 찾아
$V = \bigcup_i \bigcup_{j = 1, \ldots, n_i} D(f_{ij})$로 쓸 수 있다.
여기서 $D(f_{ij}) \subset U_{i, A}$는 대응하는
표준 열린집합이다. ($V \cap U_{i, A}$가
$D(f_{ij})$, $j = 1, \ldots, n_i$의 합집합이라고 주장하는 것은 아니다.)
유한 생성 $k$-부분대수 $A' \subset A$를 찾아 $f_{ij}$가 어떤
$f_{ij}' \in R_i \otimes_k A'$의 상이 되게 할 수 있음은 분명하다.
$V' = \bigcup D(f_{ij}')$로 두면 이는 $X_{A'}$의 준콤팩트 열린집합이다.
표준 사상을 $\pi : X_A \to X_{A'}$로 나타내자.
$\pi(V) \subset V'$인데, 이는 $\pi(D(f_{ij})) \subset D(f_{ij}')$이기 때문이다.
$x \in X_A$이고 $\pi(x) \in V'$이면
$\pi(x) \in D(f_{ij}')$인 어떤 $i, j$가 존재하며,
$x \in D(f_{ij})$임을 알 수 있다. 실제로 $f_{ij}'$가 $f_{ij}$로 간다.
따라서 원하는 대로
$V = \pi^{-1}(V')$이다.
\end{proof}

\noindent
$k$를 체라 하자. $\overline{k}/k$를 유한일 필요가 없는 갈루아 확대라 하자.
예를 들어 $\overline{k}$는 $k$의 분리 대수적 폐포일 수 있다.
임의의 $\sigma \in \text{Gal}(\overline{k}/k)$에 대응하여 자기동형
$
\Spec(\sigma) :
\Spec(\overline{k})
\longrightarrow
\Spec(\overline{k})
$
을 얻는다.
$\Spec(\sigma) \circ \Spec(\tau) = \Spec(\tau \circ \sigma)$임에 유의하자.
따라서 스킴 $\Spec(\overline{k})$ 위에 반대군의 작용
$$
\text{Gal}(\overline{k}/k)^{opp} \times \Spec(\overline{k})
\longrightarrow
\Spec(\overline{k})
$$
을 얻는다. $X$를 $k$ 위의 스킴이라 하자. 정의에 의해
$X_{\overline{k}} =
\Spec(\overline{k}) \times_{\Spec(k)} X$이므로
위 작용은 표준 작용
\begin{equation}
\label{varieties-equation-galois-action-base-change-kbar}
\text{Gal}(\overline{k}/k)^{opp} \times X_{\overline{k}}
\longrightarrow
X_{\overline{k}}.
\end{equation}
을 유도한다.

\begin{lemma}
\label{varieties-lemma-Galois-action-quasi-compact-open}
$k$를 체라 하고 $X$를 $k$ 위의 스킴이라 하자.
$\overline{k}$를 $k$의 유한일 필요가 없는 갈루아 확대라 하자.
$V \subset X_{\overline{k}}$를 준콤팩트 열린집합이라 하자. 그러면
\begin{enumerate}
\item 유한 부분확대 $\overline{k}/k'/k$와 준콤팩트 열린집합
$V' \subset X_{k'}$가 존재하여 $V = (V')_{\overline{k}}$이고,
\item 열린 부분군 $H \subset \text{Gal}(\overline{k}/k)$가 존재하여
$\sigma(V) = V$가 모든 $\sigma \in H$에 대해 성립한다.
\end{enumerate}
\end{lemma}

\begin{proof}
보조정리 \ref{varieties-lemma-descend-open}에 의해 유한 부분확대
$k'/k \subset \overline{k}$와 열린집합 $V' \subset X_{k'}$가 존재하여
후자를 당기면 $V$가 된다. 이것으로 (1)이 증명된다.
$\text{Gal}(\overline{k}/k')$가
$\text{Gal}(\overline{k}/k)$의 열린 부분군이므로 (2)도 분명하다.
\end{proof}

\begin{lemma}
\label{varieties-lemma-closed-fixed-by-Galois}
$k$를 체라 하고 $\overline{k}/k$를 유한일 필요가 없는 갈루아 확대라 하자.
$X$를 $k$ 위의 스킴이라 하자. 부분집합
$\overline{T} \subset X_{\overline{k}}$가 다음 성질들을 갖는다고 하자.
\begin{enumerate}
\item $\overline{T}$는 $X_{\overline{k}}$의 닫힌 부분집합이다.
\item 모든 $\sigma \in \text{Gal}(\overline{k}/k)$에 대해
$\sigma(\overline{T}) = \overline{T}$이다.
\end{enumerate}
그러면 닫힌 부분집합 $T \subset X$가 존재하여
$X_{\overline{k}}$에서의 역상이 $\overline{T}$이다.
\end{lemma}

\begin{proof}
이 보조정리는 $X = \Spec(A)$가 아핀인 경우로 즉시 환원된다.
이 경우 $\overline{I} \subset A \otimes_k \overline{k}$를
$\overline{T}$에 대응하는 근기 아이디얼이라 하자. 가정 (2)는
$\sigma(\overline{I}) = \overline{I}$가 모든
$\sigma \in \text{Gal}(\overline{k}/k)$에 대해 성립함을 뜻한다.
$x \in \overline{I}$를 택하자. 유한 갈루아 확대 $k'/k$가
$\overline{k}$ 안에 포함되고 $x \in A \otimes_k k'$가 되게 할 수 있다.
$G = \text{Gal}(k'/k)$로 두고
$$
P(T) = \prod\nolimits_{\sigma \in G} (T - \sigma(x)) \in (A \otimes_k k')[T]
$$
로 두자. $P(T)$가 일계수이고 실제로
$(A \otimes_k k')^G[T] = A[T]$의 원소임은 분명하다
(기본적인 갈루아 이론에 의한다). 또한
$P(T) = T^d + a_1T^{d - 1} + \ldots + a_d$로 쓰면
$a_i \in I := A \cap \overline{I}$임을 알 수 있다.
$P(x) = 0$과 $a_i \in I$를 합치면
$x^d = - a_1 x^{d - 1} - \ldots  - a_d \in I(A \otimes_k \overline{k})$를 얻는다. 따라서
$x$는 $I(A \otimes_k \overline{k})$의 근기에 속한다.
그러므로 $\overline{I}$는 $I(A \otimes_k \overline{k})$의 근기이고,
$T = V(I)$로 두면 원하는 해를 얻는다.
\end{proof}

\begin{lemma}
\label{varieties-lemma-characterize-geometrically-disconnected}
$k$를 체라 하고 $X$를 $k$ 위의 스킴이라 하자.
다음 조건들은 서로 동치이다.
\begin{enumerate}
\item $X$는 기하적으로 연결이다.
\item 모든 유한 분리 체 확대 $k'/k$에 대해 스킴 $X_{k'}$는 연결이다.
\end{enumerate}
\end{lemma}

\begin{proof}
정의에서 (1)이 (2)를 함의함은 곧바로 따른다.
$X$가 기하적으로 연결이 아니라고 가정하자.
$k \subset \overline{k}$를 $k$의 분리 대수적 폐포라 하자.
보조정리 \ref{varieties-lemma-characterize-geometrically-connected}에 의해
$X_{\overline{k}}$는 비연결이다. 따라서
$X_{\overline{k}} = \overline{U} \amalg \overline{V}$로 쓸 수 있으며,
$\overline{U}$와 $\overline{V}$는 열리고 닫혀 있고 공집합이 아니다.

\medskip\noindent
$W \subset X$를 임의의 준콤팩트 열린집합이라 하자.
$W_{\overline{k}} \cap \overline{U}$와
$W_{\overline{k}} \cap \overline{V}$는
$W_{\overline{k}}$에서 열리고 닫혀 있다. 특히
$W_{\overline{k}} \cap \overline{U}$와
$W_{\overline{k}} \cap \overline{V}$는 준콤팩트이고, 보조정리
\ref{varieties-lemma-Galois-action-quasi-compact-open}에 의해
$W_{\overline{k}} \cap \overline{U}$와
$W_{\overline{k}} \cap \overline{V}$는 모두 유한 부분확대 위에서
정의되며 $\text{Gal}(\overline{k}/k)$의 열린 부분군 아래 불변이다.
이하에서는 이를 따로 언급하지 않고 사용한다.

\medskip\noindent
준콤팩트 열린집합 $W_0 \subset X$를 택하여
$W_{0, \overline{k}} \cap \overline{U}$와
$W_{0, \overline{k}} \cap \overline{V}$가 모두 공집합이 아니게 하자.
유한 부분확대 $\overline{k}/k'/k$와 열린닫힌 부분집합들로의 분해
$W_{0, k'} = U_0' \amalg V_0'$를 택하여
$W_{0, \overline{k}} \cap \overline{U} = (U'_0)_{\overline{k}}$ 및
$W_{0, \overline{k}} \cap \overline{V} = (V'_0)_{\overline{k}}$가 되게 하자.
$H = \text{Gal}(\overline{k}/k') \subset \text{Gal}(\overline{k}/k)$로 두자. 특히
$\sigma(W_{0, \overline{k}} \cap \overline{U}) =
W_{0, \overline{k}} \cap \overline{U}$이고,
$\overline{V}$에 대해서도 마찬가지이다.

\medskip\noindent
위와 같이 $W_0$, $k'$를 택했을 때, 모든 준콤팩트 열린집합
$W \subset X$에 대해
$$
U_W =
\bigcap\nolimits_{\sigma \in H} \sigma(W_{\overline{k}} \cap \overline{U}),
\quad
V_W =
\bigcup\nolimits_{\sigma \in H} \sigma(W_{\overline{k}} \cap \overline{V}).
$$
로 두자. $W_{\overline{k}} \cap \overline{U}$와
$W_{\overline{k}} \cap \overline{V}$가
$\text{Gal}(\overline{k}/k)$의 열린 부분군에 의해 고정되므로
위의 합집합과 교집합은 유한하다. 따라서 $U_W$와 $V_W$는 모두
열리고 닫혀 있다. 또한 구성에 의해
$W_{\bar k} = U_W \amalg V_W$이다.

\medskip\noindent
$W \subset W' \subset X$가 준콤팩트 열린집합들이면
$W_{\overline{k}} \cap U_{W'} = U_W$이고
$W_{\overline{k}} \cap V_{W'} = V_W$라고 주장한다. 확인은 생략한다.
따라서 $U = \bigcup_{W \subset X} U_W$ 및
$V = \bigcup_{W \subset X} V_W$로 정의하면
$X_{\overline{k}} = U \amalg V$는 열린닫힌 부분집합들의 서로소 합이다.
$V$는 (국소적으로) 합집합을 취하여 구성했으므로 공집합이 아님이 분명하다.
한편 $U$는 구성에 의해 $W_0 \cap \overline{U}$를 포함하므로 공집합이 아니다.
마지막으로 $U, V \subset X_{\bar k}$는 구성에 의해 닫혀 있고
$H$-불변이다. 따라서 보조정리 \ref{varieties-lemma-closed-fixed-by-Galois}에 의해
$U = (U')_{\bar k}$이고 $V = (V')_{\bar k}$인 어떤 닫힌 부분집합
$U', V' \subset X_{k'}$가 존재한다.
분명히 $X_{k'} = U' \amalg V'$이므로, 원하는 대로
$X_{k'}$가 비연결임을 얻는다.
\end{proof}

\begin{lemma}
\label{varieties-lemma-tricky}
$k$를 체라 하고 $\overline{k}/k$를 유한일 필요가 없는 갈루아 확대라 하자.
$f : T \to X$를 $k$ 위의 스킴들의 사상이라 하자.
$T_{\overline{k}}$는 연결이고 $X_{\overline{k}}$는 비연결이라고 가정하자.
그러면 $X$는 비연결이다.
\end{lemma}

\begin{proof}
$X_{\overline{k}} = \overline{U} \amalg \overline{V}$로 쓰자.
여기서 $\overline{U}$와 $\overline{V}$는 열리고 닫혀 있다.
$\overline{f} : T_{\overline{k}} \to X_{\overline{k}}$를 $f$의 기저 변환이라 하자.
$T_{\overline{k}}$가 연결이므로 $T_{\overline{k}}$는
$\overline{f}^{-1}(\overline{U})$ 또는
$\overline{f}^{-1}(\overline{V})$ 중 하나에 포함된다.
$T_{\overline{k}} \subset \overline{f}^{-1}(\overline{U})$라고 하자.

\medskip\noindent
준콤팩트 열린집합 $W \subset X$를 고정하자. 유한 갈루아 부분확대
$\overline{k}/k'/k$가 존재하여
$\overline{U} \cap W_{\overline{k}}$와
$\overline{V} \cap W_{\overline{k}}$가 준콤팩트 열린집합
$U', V' \subset W_{k'}$로부터 온다. 그러면 또한
$W_{k'} = U' \amalg V'$이다. 다음을 생각하자.
$$
U'' = \bigcap\nolimits_{\sigma \in \text{Gal}(k'/k)} \sigma(U'),
\quad
V'' = \bigcup\nolimits_{\sigma \in \text{Gal}(k'/k)} \sigma(V').
$$
이들은 갈루아 불변인 열린닫힌 부분집합이며
$W_{k'} = U'' \amalg V''$이다.
보조정리 \ref{varieties-lemma-closed-fixed-by-Galois}에 의해 열린닫힌 부분집합
$U_W, V_W \subset W$를 얻어
$U'' = (U_W)_{k'}$, $V'' = (V_W)_{k'}$ 및
$W = U_W \amalg V_W$가 성립한다.

\medskip\noindent
$W \subset W' \subset X$가 준콤팩트 열린집합들이면
$W \cap U_{W'} = U_W$이고 $W \cap V_{W'} = V_W$라고 주장한다.
확인은 생략한다. 따라서 $U = \bigcup_{W \subset X} U_W$ 및
$V = \bigcup_{W \subset X} V_W$로 정의하면 $X = U \amalg V$를 얻는다.
$V$는 (국소적으로) 합집합을 취하여 구성했으므로 공집합이 아님이 분명하다.
한편 $U$는 구성에 의해 $f(T)$를 포함하므로 공집합이 아니다.
\end{proof}

\begin{lemma}
\label{varieties-lemma-geometrically-connected-criterion}
\begin{reference}
\cite[IV Corollary 4.5.13.1(i)]{EGA}
\end{reference}
$k$를 체라 하자. $T \to X$를 $k$ 위의 스킴들의 사상이라 하자.
$T$가 기하적으로 연결이고 $X$가 연결이라고 가정하자.
그러면 $X$는 기하적으로 연결이다.
\end{lemma}

\begin{proof}
이는 보조정리 \ref{varieties-lemma-tricky}를 다시 표현한 것이다.
\end{proof}

\begin{lemma}
\label{varieties-lemma-geometrically-connected-if-connected-and-point}
$k$를 체라 하고 $X$를 $k$ 위의 스킴이라 하자.
$X$가 연결이고 점 $x$를 가지며, $k$가 $\kappa(x)$ 안에서
대수적으로 닫혀 있다고 가정하자. 그러면 $X$는 기하적으로 연결이다.
특히 $X$가 $k$-유리점을 갖고 $X$가 연결이면
$X$는 기하적으로 연결이다.
\end{lemma}

\begin{proof}
$T = \Spec(\kappa(x))$로 두자. $\overline{k}$를 $k$의 분리 대수적 폐포라 하자.
$\kappa(x)/k$에 대한 가정은 $T_{\overline{k}}$가 기약임을 함의한다.
Algebra, Lemma \ref{algebra-lemma-field-extension-geometrically-irreducible}을 보라.
따라서 보조정리 \ref{varieties-lemma-geometrically-connected-criterion}에 의해
$X_{\overline{k}}$는 연결이다. 보조정리
\ref{varieties-lemma-characterize-geometrically-connected}에 의해
$X$가 기하적으로 연결임을 얻는다.
\end{proof}

\begin{lemma}
\label{varieties-lemma-inverse-image-connected-component}
$K/k$를 체 확대라 하고 $X$를 $k$ 위의 스킴이라 하자.
$T$를 $X$의 임의의 연결 성분이라 하면 역상
$T_K \subset X_K$는 $X_K$의 연결 성분들의 합집합이다.
\end{lemma}

\begin{proof}
이는 순수하게 위상적인 명제이다. 사영 사상을
$p : X_K \to X$로 나타내자. $T \subset X$를 $X$의 연결 성분이라 하자.
$t \in T_K = p^{-1}(T)$라 하고 $C \subset X_K$를 $t$를 포함하는
연결 성분이라 하자. 그러면 $p(C)$는 $X$의 연결 부분집합이고
$T$와 만나므로 $p(C) \subset T$이다. 따라서 $C \subset T_K$이다.
\end{proof}

\noindent
다음 보조정리는 아래의 더 강한 보조정리
\ref{varieties-lemma-image-connected-component}로 대체될 것이다.

\begin{lemma}
\label{varieties-lemma-image-connected-component-finite-extension}
$K/k$를 유한 체 확대라 하고 $X$를 $k$ 위의 스킴이라 하자.
사영 사상을 $p : X_K \to X$로 나타내자. 모든 연결 성분
$T$에 대해, 그것이 $X_K$의 연결 성분이면 상 $p(T)$는 $X$의 연결 성분이다.
\end{lemma}

\begin{proof}
상 $p(T)$는 어떤 연결 성분 $X'$에 포함되며, 이는 $X$의 연결 성분이다.
$X'$에 임의의 방식으로 닫힌 부분스킴 구조를 주어 $X$의 부분스킴으로 보자.
그러면 $T$는 $X'_K = p^{-1}(X')$의 연결 성분이기도 하므로
$X$가 연결인 경우를 가정해도 된다. 사상 $p$는 열린 사상이고
(Morphisms, Lemma
\ref{morphisms-lemma-scheme-over-field-universally-open}),
닫힌 사상이며
(Morphisms, Lemma
\ref{morphisms-lemma-integral-universally-closed})
$p$의 올들은 유한집합이다
(Morphisms, Lemma \ref{morphisms-lemma-finite-quasi-finite}).
따라서 Topology, Lemma
\ref{topology-lemma-finite-fibre-connected-components}를 적용하면 된다.
\end{proof}

\begin{lemma}[Gabber]
\label{varieties-lemma-image-connected-component}
\begin{reference}
2016년 6월 4일자 Ofer Gabber의 전자우편
\end{reference}
$K/k$를 체 확대라 하고 $X$를 $k$ 위의 스킴이라 하자.
사영 사상을 $p : X_K \to X$로 나타내자.
$\overline{T} \subset X_K$를 연결 성분이라 하자.
그러면 $p(\overline{T})$는 $X$의 연결 성분이다.
\end{lemma}

\begin{proof}
$K/k$가 유한이면 이는 보조정리
\ref{varieties-lemma-image-connected-component-finite-extension}이다.
일반적인 경우에는 증명이 더 어렵다.

\medskip\noindent
$T \subset X$를 $X$의 연결 성분 중 $\overline{T}$의 상을 포함하는 것으로 하자.
$X$를 $T$로 바꾸어도 된다(유도된 축약 부분스킴 구조를 준다).
따라서 $X$가 연결이라고 가정해도 된다.
$A = H^0(X, \mathcal{O}_X)$로 두자.
$L \subset A$를 극대 약 에탈 $k$-부분대수라 하자. More on Algebra, Lemma
\ref{more-algebra-lemma-max-weakly-etale-subalgebra}를 보라.
$A$에는 자명하지 않은 멱등원이 없으므로, More on Algebra, Lemma
\ref{more-algebra-lemma-class-weakly-etale-over-field}에 의해
$L$은 체이고 $k$의 분리 대수적 확대이다.
$L$은 또한 $L$-대수로 볼 때 $A$의 극대 약 에탈 부분대수이다
(More on Algebra, Lemma
\ref{more-algebra-lemma-composition-weakly-etale}에 의해 모든 약 에탈
$L$-대수는 $k$ 위에서도 약 에탈이기 때문이다).
Schemes, Lemma \ref{schemes-lemma-morphism-into-affine}에 의해
구조 사상의 분해 $X \to \Spec(L) \to \Spec(k)$를 얻는다.

\medskip\noindent
$L'/L$을 유한 분리 확대라 하자. Cohomology of Schemes, Lemma
\ref{coherent-lemma-finite-locally-free-base-change-cohomology}에 의해
$$
A \otimes_L L' =
H^0(X \times_{\Spec(L)} \Spec(L'), \mathcal{O}_{X \times_{\Spec(L)} \Spec(L')})
$$
이다. 극대 약 에탈 $L'$-부분대수로서 $A \otimes_L L'$ 안에 있는 것은
More on Algebra, Lemma
\ref{more-algebra-lemma-change-fields-max-weakly-etale-subalgebra}에 의해
$L \otimes_L L' = L'$이다. 특히 $A \otimes_L L'$에는 자명하지 않은
멱등원이 없다(그런 멱등원은 약 에탈 부분대수를 생성할 것이다).
따라서 $X \times_{\Spec(L)} \Spec(L')$가 연결임을 얻는다.
보조정리 \ref{varieties-lemma-characterize-geometrically-disconnected}에 의해
$X$는 $L$ 위에서 기하적으로 연결이다.

\medskip\noindent
$\overline{T}$에 유도된 축약 스킴 구조를 주고 합성
$$
\overline{T} \xrightarrow{i} X_K = X \times_{\Spec(k)} \Spec(K)
\xrightarrow{\pi}
\Spec(L \otimes_k K)
$$
을 생각하자. 그 상은 $\Spec(L \otimes_k K)$의 한 연결 성분에 포함된다.
$K \to L \otimes_k K$가 정수적이므로
$\Spec(L \otimes_k K)$의 연결 성분들은 점이고 모든 점이 닫혀 있다.
Algebra, Lemma \ref{algebra-lemma-integral-over-field}을 보라.
따라서 몫체 $L \otimes_k K \to E$가 존재하여
$\overline{T}$가 $\Spec(E) \subset \Spec(L \otimes_k K)$로 사상된다.
그러므로 $i(\overline{T}) \subset \pi^{-1}(\Spec(E))$이다. 한편
$$
\pi^{-1}(\Spec(E)) =
(X \times_{\Spec(k)} \Spec(K)) \times_{\Spec(L \otimes_k K)} \Spec(E) =
X \times_{\Spec(L)} \Spec(E)
$$
이고, $X$가 $L$ 위에서 기하적으로 연결이므로 이는 연결이다.
따라서 집합론적으로
$\overline{T} = X \times_{\Spec(L)} \Spec(E)$이고,
원하는 대로 $\overline{T} \to X$가 전사임을 얻는다.
\end{proof}

\noindent
$X$를 스킴이라 하자. $\pi_0(X)$로 $X$의 연결 성분들의 집합을 나타낸다.

\begin{lemma}
\label{varieties-lemma-galois-action-connected-components}
$k$를 체라 하고 그 분리 대수적 폐포를 $\overline{k}$라 하자.
$X$를 $k$ 위의 스킴이라 하자. 다음 작용이 존재한다.
$$
\text{Gal}(\overline{k}/k)^{opp} \times \pi_0(X_{\overline{k}})
\longrightarrow
\pi_0(X_{\overline{k}})
$$
이 작용은 다음 성질들을 갖는다.
\begin{enumerate}
\item 원소 $\overline{T} \in \pi_0(X_{\overline{k}})$가 이 작용에 의해
고정될 필요충분조건은 연결 성분 $T \subset X$가 $k$ 위에서
기하적으로 연결이고 $T_{\overline{k}} = \overline{T}$가 되도록 존재하는 것이다.
\item 임의의 체 확대 $k'/k$와 그 분리 대수적 폐포
$\overline{k}'$에 대해 도식
$$
\xymatrix{
\text{Gal}(\overline{k}'/k') \times \pi_0(X_{\overline{k}'})
\ar[r] \ar[d] &
\pi_0(X_{\overline{k}'}) \ar[d] \\
\text{Gal}(\overline{k}/k) \times \pi_0(X_{\overline{k}})
\ar[r] &
\pi_0(X_{\overline{k}})
}
$$
은 가환이다(오른쪽 수직 화살표는 보조정리
\ref{varieties-lemma-separably-closed-field-connected-components}에 의해 전단사이다).
\end{enumerate}
\end{lemma}

\begin{proof}
$\text{Gal}(\overline{k}/k)$의 $X_{\overline{k}}$ 위 작용
(\ref{varieties-equation-galois-action-base-change-kbar})은 그 연결 성분들 위의
작용을 유도한다. 연결 성분은 항상 닫혀 있다
(Topology, Lemma \ref{topology-lemma-connected-components}).
따라서 $\overline{T}$가 (1)과 같으면 보조정리
\ref{varieties-lemma-closed-fixed-by-Galois}에 의해 닫힌 부분집합
$T \subset X$가 존재하여 $\overline{T} = T_{\overline{k}}$이다.
$T$는 $k$ 위에서 기하적으로 연결임에 유의하자. 보조정리
\ref{varieties-lemma-characterize-geometrically-connected}을 보라.
$T$가 $X$의 연결 성분임을 보이기 위해, $T \subset T'$이고
$T \not = T'$이며 $T'$가 $X$의 연결 성분이라고 가정하자.
이 경우 $T'_{k'}$는 $\overline{T}$를 진부분집합으로 포함하므로 비연결이다.
보조정리 \ref{varieties-lemma-tricky}에 의해 이는 $T'$가 비연결임을 뜻한다!
모순이다.

\medskip\noindent
(2)의 함자성에 대한 증명은 생략한다.
\end{proof}

\begin{lemma}
\label{varieties-lemma-galois-action-connected-components-continuous}
$k$를 체라 하고 그 분리 대수적 폐포를 $\overline{k}$라 하자.
$X$를 $k$ 위의 스킴이라 하자. 다음을 가정한다.
\begin{enumerate}
\item $X$는 준콤팩트이다.
\item $X_{\overline{k}}$의 연결 성분들은 열려 있다.
\end{enumerate}
그러면
\begin{enumerate}
\item[(a)] $\pi_0(X_{\overline{k}})$는 유한하고,
\item[(b)] $\text{Gal}(\overline{k}/k)$의
$\pi_0(X_{\overline{k}})$ 위 작용은 연속이다.
\end{enumerate}
더욱이 $X$가 $k$ 위에서 유한형이면 가정 (1), (2)가 성립한다.
\end{lemma}

\begin{proof}
연결 성분들은 열려 있고 $X_{\overline{k}}$를 덮는다
(Topology, Lemma \ref{topology-lemma-connected-components}).
$X_{\overline{k}}$가 준콤팩트이므로 연결 성분은 유한 개뿐이다.
따라서 (a)가 성립한다. 보조정리 \ref{varieties-lemma-descend-open}에 의해
이 연결 성분들은 각각 $\overline{k}/k$의 유한 부분확대 위에서 정의되므로
(b)를 얻는다. $X$가 $k$ 위에서 유한형이면
$X_{\overline{k}}$는 $\overline{k}$ 위에서 유한형이다
(Morphisms, Lemma \ref{morphisms-lemma-base-change-finite-type}).
따라서 $X_{\overline{k}}$는 뇌터 스킴이다
(Morphisms, Lemma \ref{morphisms-lemma-finite-type-noetherian}).
그러므로 $X_{\overline{k}}$는 유한 개의 기약 성분을 갖고
(Properties, Lemma \ref{properties-lemma-Noetherian-irreducible-components}),
따라서 특히 연결 성분도 유한 개이다(그러므로 이들은 열려 있다).
\end{proof}








\section{기하적으로 기약인 스킴}
\label{varieties-section-geometrically-irreducible}

\noindent
$X$가 체 위의 기약 스킴이어도 기초체를 확대하면 $X$가 가약이 될 수 있다.
기하적으로 기약인 스킴에서는 이런 일이 일어나지 않는다.

\begin{definition}
\label{varieties-definition-geometrically-irreducible}
$X$를 체 $k$ 위의 스킴이라 하자. $X$가 $k$ 위에서
{\it 기하적으로 기약}이라고 함은 스킴
$X_{k'}$가 기약인 경우이다\footnote{기약 공간은 공집합이 아니다.}
여기서 $k'$는 $k$의 임의의 체 확대이다.
\end{definition}

\begin{lemma}
\label{varieties-lemma-geometrically-irreducible-check-after-extension}
$X$를 체 $k$ 위의 스킴이라 하고 $k'/k$를 체 확대라 하자.
$X$가 $k$ 위에서 기하적으로 기약일 필요충분조건은
$X_{k'}$가 $k'$ 위에서 기하적으로 기약인 것이다.
\end{lemma}

\begin{proof}
$X$가 $k$ 위에서 기하적으로 기약이면 $X_{k'}$가 $k'$ 위에서
기하적으로 기약임은 분명하다. 역으로 임의의 체 확대 $k''/k$에 대해
공통 체 확대 $k'''/k'$와 $k'''/k''$가 존재한다. Fields, Lemma
\ref{fields-lemma-common-extension-field}을 보라.
사상 $X_{k'''} \to X_{k''}$는 체들의 스펙트럼 사이 전사 사상의
기저 변환이므로 전사이다. 따라서 $X_{k'''}$의 기약성은
$X_{k''}$의 기약성을 함의한다. 그러므로 $X_{k'}$가 $k'$ 위에서
기하적으로 기약이면 $X$는 $k$ 위에서 기하적으로 기약이다.
\end{proof}

\begin{lemma}
\label{varieties-lemma-separably-closed-irreducible}
$X$를 분리 닫힌 체 $k$ 위의 스킴이라 하자. $X$가 기약이면
$X_K$는 임의의 체 확대 $K/k$에 대해 기약이다. 즉
$X$는 $k$ 위에서 기하적으로 기약이다.
\end{lemma}

\begin{proof}
Properties, Lemma \ref{properties-lemma-characterize-irreducible}과
Algebra, Lemma \ref{algebra-lemma-separably-closed-irreducible}을 사용하라.
\end{proof}

\begin{lemma}
\label{varieties-lemma-bijection-irreducible-components}
$k$를 체라 하고 $X$, $Y$를 $k$ 위의 스킴이라 하자.
$X$가 $k$ 위에서 기하적으로 기약이라고 가정하자. 그러면 사영 사상
$$
p : X \times_k Y \longrightarrow Y
$$
은 기약 성분들 사이의 전단사를 유도한다.
\end{lemma}

\begin{proof}
먼저 $p$의 스킴론적 올들은 기약이다. 실제로 이들은
기하적으로 기약인 스킴 $X$를 체 확대에 따라 기저 변환한 것이다.
또한 스킴론적 올들은 집합론적 올들과 위상동형이다. Schemes, Lemma
\ref{schemes-lemma-fibre-topological}을 보라.
Morphisms, Lemma \ref{morphisms-lemma-scheme-over-field-universally-open}에 의해
사상 $p$는 열린 사상이다. 따라서 Topology, Lemma
\ref{topology-lemma-irreducible-fibres-irreducible-components}를 적용하면 된다.
\end{proof}

\begin{lemma}
\label{varieties-lemma-geometrically-irreducible-local}
\begin{slogan}
기하적 기약성은 연결성을 법으로 하면 자리스키 국소적이다.
\end{slogan}
$k$를 체라 하고 $X$를 $k$ 위의 스킴이라 하자.
다음 조건들은 서로 동치이다.
\begin{enumerate}
\item $X$는 $k$ 위에서 기하적으로 기약이다.
\item 모든 공집합이 아닌 아핀 열린집합 $U$에 대해
$k$-대수 $\mathcal{O}_X(U)$는 $k$ 위에서 기하적으로 기약이다
(Algebra, Definition \ref{algebra-definition-geometrically-irreducible}을 보라).
\item $X$는 기약이고, 아핀 열린 덮개 $X = \bigcup U_i$가 존재하여
각 $k$-대수 $\mathcal{O}_X(U_i)$가 기하적으로 기약이다.
\item 열린 덮개 $X = \bigcup_{i \in I} X_i$가 존재하여
$I \not = \emptyset$이고, $X_i$가 각 $i$에 대해 기하적으로 기약이며,
$X_i \cap X_j \not = \emptyset$가 모든 $i, j \in I$에 대해 성립한다.
\end{enumerate}
더욱이 $X$가 기하적으로 기약이면 $X$의 모든 공집합이 아닌
열린 부분스킴도 기하적으로 기약이다.
\end{lemma}

\begin{proof}
$\Spec(A)$가 $k$ 위의 아핀 스킴일 때 기하적으로 기약일 필요충분조건은
$A$가 $k$ 위에서 기하적으로 기약인 것이다. 이는 정의에서 곧바로 따른다.
스킴이 기약이면 $X$의 모든 공집합이 아닌 열린 부분스킴도 기약이고,
공집합이 아닌 임의의 두 열린 부분집합이 서로 만남을 상기하자.
또한 모든 아핀 열린집합이 기약이면 스킴도 기약이다. Properties, Lemma
\ref{properties-lemma-characterize-irreducible}을 보라.
따라서 보조정리의 마지막 진술과 함의
(1) $\Rightarrow$ (2), (2) $\Rightarrow$ (3), (3) $\Rightarrow$ (4)는 분명하다.
(4)가 성립하면 임의의 체 확대 $k'/k$에 대해 스킴 $X_{k'}$는
쌍마다 서로 만나는 기약 열린집합들로 이루어진 덮개를 갖는다.
따라서 $X_{k'}$는 기약이고, (4)는 (1)을 함의한다.
\end{proof}

\begin{lemma}
\label{varieties-lemma-geometrically-irreducible-function-field}
$X$를 체 $k$ 위의 기약 스킴이라 하고 $\xi \in X$를 그 일반점이라 하자.
다음 조건들은 서로 동치이다.
\begin{enumerate}
\item $X$는 $k$ 위에서 기하적으로 기약이다.
\item $\kappa(\xi)$는 $k$ 위에서 기하적으로 기약이다.
\item $k$의 분리 대수적 폐포가 $\kappa(\xi)$ 안에서 $k$와 같다.
\end{enumerate}
\end{lemma}

\begin{proof}
Algebra, Lemma
\ref{algebra-lemma-geometrically-irreducible-separable-elements}에 의해
(2)와 (3)은 서로 동치이다.

\medskip\noindent
(1)을 가정하자. $\mathcal{O}_{X, \xi}$는 $\mathcal{O}_X(U)$들의
여과 쌍대극한임을 상기하자. 여기서 $U$는 $X$의 공집합이 아닌
아핀 열린 부분스킴들을 달린다.
보조정리 \ref{varieties-lemma-geometrically-irreducible-local}과
Algebra, Lemma \ref{algebra-lemma-subalgebra-geometrically-irreducible}을 합치면
$\mathcal{O}_{X, \xi}$가 $k$ 위에서 기하적으로 기약임을 알 수 있다.
$\mathcal{O}_{X, \xi} \to \kappa(\xi)$는 국소 멱영 핵을 갖는 전사이므로
(Algebra, Lemma \ref{algebra-lemma-minimal-prime-reduced-ring}을 보라)
$\kappa(\xi)$는 기하적으로 기약이다. Algebra, Lemma
\ref{algebra-lemma-p-ring-map}을 보라.

\medskip\noindent
(2)를 가정하자. $X$가 축약이라고 가정해도 된다.
$U \subset X$를 공집합이 아닌 아핀 열린집합이라 하자.
그러면 $U = \Spec(A)$이고 $A$는 분수체 $\kappa(\xi)$를 갖는 정역이다.
따라서 $A$는 기하적으로 기약인 $k$-대수의 $k$-부분대수이다.
그러므로 Algebra, Lemma
\ref{algebra-lemma-subalgebra-geometrically-irreducible}에 의해
$A$는 $k$ 위에서 기하적으로 기약이다.
보조정리 \ref{varieties-lemma-geometrically-irreducible-local}에 의해
$X$는 $k$ 위에서 기하적으로 기약이다.
\end{proof}

\begin{lemma}
\label{varieties-lemma-geometrically-irreducible-self-product}
스킴 $X$가 체 $k$ 위에 있을 때 $k$ 위에서 기하적으로 기약일 필요충분조건은
$X\times_k X$가 기약인 것이다.
\end{lemma}

\begin{proof}
$X$가 $k$ 위에서 기하적으로 기약이면 보조정리
\ref{varieties-lemma-bijection-irreducible-components}에 의해
$X\times_k X$는 기약이다.

\medskip\noindent
역으로 $X \times_k X$가 기약이라고 가정하자.
$X$를 그 바탕 축약 스킴으로 바꾸면 문제를 바꾸지 않고
$X$가 축약이라고 가정할 수 있다.
첫째 사영 $X \times_k X \to X$가 전사이므로 $X$도 기약이고, 따라서 정역적이다.
원소 $\alpha$가 함수체 $k(X)$ 안에 있고 $k$ 위에서 분리 대수적이지만
$k$에는 속하지 않는다고 가정하자. $E = k(\alpha)$로 두자.

\medskip\noindent
그러면 공집합이 아닌 아핀 열린 부분스킴 $U$가 $X$ 안에 존재하여
$k \subset E \subset O(U)$이고, 따라서
$E \otimes_k E \subset
O(U)\otimes_k O(U) = O(U\times_k U)$이다.
그러므로 우세 사상 $U \times_k U \to \Spec(E\otimes_k E)$를 얻는다.
$X\times_k X$가 기약이므로 $U \times_k U$도 기약이고,
따라서 $\Spec(E \otimes_k E)$는 연결이다.
$E$가 $k$보다 진정으로 크면 $E \otimes_k E$는 $E$를 한 인자로 포함하는
체들의 곱이다. 이는 차원이 $(\dim_k(E))^2>\dim_k(E)$인 $k$-벡터 공간이므로
$E\otimes_k E$는 적어도 두 체의 곱이다.
이는 $\Spec(E \otimes_k E)$가 연결이라는 사실과 모순이다.
따라서 실제로 $E$는 $k$와 같다.
보조정리 \ref{varieties-lemma-geometrically-irreducible-function-field}에 의해
$X$는 $k$ 위에서 기하적으로 기약이다.
\end{proof}

\begin{lemma}
\label{varieties-lemma-separably-closed-field-irreducible-components}
$k'/k$를 체 확대라 하고 $X$를 $k$ 위의 스킴이라 하자.
$X' = X_{k'}$로 두자. $k$가 분리 대수적으로 닫혀 있다고 가정하자.
그러면 사상 $X' \to X$는 기약 성분들의 전단사를 유도한다.
\end{lemma}

\begin{proof}
$k$가 분리 대수적으로 닫혀 있으므로 $k'$는 $k$ 위에서
기하적으로 기약이다. Algebra, Lemma
\ref{algebra-lemma-separably-closed-irreducible-implies-geometric}을 보라.
따라서 위의 보조정리 \ref{varieties-lemma-geometrically-irreducible-local}에 의해
$Z = \Spec(k')$는 $k$ 위에서 기하적으로 기약이다.
$X' = Z \times_k X$이므로 결과는 보조정리
\ref{varieties-lemma-bijection-irreducible-components}의 특수한 경우이다.
\end{proof}

\begin{lemma}
\label{varieties-lemma-characterize-geometrically-irreducible}
\begin{slogan}
기하적 기약성은 기초체의 분리 대수적 폐포 위에서 판정할 수 있다.
\end{slogan}
$k$를 체라 하고 $X$를 $k$ 위의 스킴이라 하자.
다음 조건들은 서로 동치이다.
\begin{enumerate}
\item $X$는 $k$ 위에서 기하적으로 기약이다.
\item 모든 유한 분리 체 확대 $k'/k$에 대해 스킴 $X_{k'}$는 기약이다.
\item $X_{\overline{k}}$는 기약이다. 여기서
$k \subset \overline{k}$는 $k$의 분리 대수적 폐포이다.
\end{enumerate}
\end{lemma}

\begin{proof}
$X_{\overline{k}}$가 기약이라고, 즉 (3)을 가정하자.
$k'/k$를 체 확대라 하자. 체 확대 $\overline{k}'/\overline{k}$가 존재하여
$k'$가 $\overline{k}'$ 안에 $k$의 확대체로서 매장된다.
보조정리 \ref{varieties-lemma-separably-closed-field-irreducible-components}에 의해
$X_{\overline{k}'}$는 기약이다.
$X_{\overline{k}'} \to X_{k'}$가 전사이므로 $X_{k'}$는 기약이다.
따라서 (1)이 성립한다.

\medskip\noindent
$k \subset \overline{k}$를 $k$의 분리 대수적 폐포라 하자.
(3)이 아니라고, 즉 $X_{\overline{k}}$가 가약이라고 가정하자.
목표는 $X_{k'}$가 어떤 유한 부분확대
$\overline{k}/k'/k$에 대해 가약임을 보이는 것이다.
$X = \bigcup_{i \in I} U_i$를 $U_i$들이 공집합이 아닌 아핀 열린 덮개라 하자.
어떤 $i$에 대해 스킴 $U_i$가 가약이거나 어떤 쌍
$i \not = j$에 대해 교집합 $U_i \cap U_j$가 공집합이면,
$X$는 가약이고(보조정리
\ref{properties-lemma-characterize-irreducible}) 증명이 끝난다.
특히 $U_{i, \overline{k}} \cap U_{j, \overline{k}}$가
모든 $i, j \in I$에 대해 공집합이 아니라고 가정할 수 있고,
그러면 $U_{i, \overline{k}}$가 어떤 $i$에 대해 가약이어야 한다.
Algebra, Lemma \ref{algebra-lemma-geometrically-irreducible}에 따르면
이는 $U_{i, k'}$가 어떤 유한 분리 체 확대 $k'/k$에 대해
가약이라는 뜻이다. 따라서 $X_{k'}$도 가약이고, (2)가 (3)을 함의함을 얻는다.

\medskip\noindent
함의 (1) $\Rightarrow$ (2)는 즉시 성립한다. 이것으로 보조정리가 증명된다.
\end{proof}

\begin{lemma}
\label{varieties-lemma-inverse-image-irreducible}
$K/k$를 체 확대라 하고 $X$를 $k$ 위의 스킴이라 하자.
$T$를 $X$의 임의의 기약 성분이라 하면 역상
$T_K \subset X_K$는 $X_K$의 기약 성분들의 합집합이다.
\end{lemma}

\begin{proof}
$T \subset X$를 $X$의 기약 성분이라 하자. 사상
$T_K \to T$는 평탄이므로 일반화는 $T_K \to T$를 따라 올라간다.
따라서 모든 $\xi \in T_K$가 $T_K$의 기약 성분의 일반점이면
일반점 $\eta$로 사상되며, 이는 $T$의 일반점이다.
$\xi' \leadsto \xi$가 $X_K$에서의 특수화이면
$\xi'$도 $\eta$로 사상된다. 실제로 $\eta$로 특수화되는 점은
$X$에 없기 때문이다.
따라서 $\xi' \in T_K$이고 $\xi = \xi'$임을 얻는다.
다시 말해 $\xi$는 $X_K$의 한 기약 성분의 일반점이다.
이는 $T_K$의 기약 성분들이 모두 $X_K$의 기약 성분임을 뜻한다.
\end{proof}

\noindent
스킴 $X$에 대해 $\text{IrredComp}(X)$로
$X$의 기약 성분들의 집합을 나타낸다.

\begin{lemma}
\label{varieties-lemma-image-irreducible}
$K/k$를 체 확대라 하고 $X$를 $k$ 위의 스킴이라 하자.
모든 기약 성분 $\overline{T} \subset X_K$에 대해
$\overline{T}$의 $X$에서의 상은 $X$의 기약 성분이다.
이로써 전사인 표준 사상
$$
\text{IrredComp}(X_K)
\longrightarrow
\text{IrredComp}(X)
$$
이 정의된다.
\end{lemma}

\begin{proof}
도식
$$
\xymatrix{
X_K \ar[d] & X_{\overline{K}} \ar[d] \ar[l] \\
X & X_{\overline{k}} \ar[l]
}
$$
을 생각하자. 여기서 $\overline{K}$는 $K$의 분리 대수적 폐포이고,
$\overline{k}$는 $k$의 분리 대수적 폐포이다. 보조정리
\ref{varieties-lemma-separably-closed-field-irreducible-components}에 의해
사상 $X_{\overline{K}} \to X_{\overline{k}}$는 기약 성분들의 전단사를 유도한다.
따라서 사상 $X_{\overline{k}} \to X$와
$X_{\overline{K}} \to X_K$에 대해 보조정리를 보이면 충분하다.
다시 말해 $K = \overline{k}$라고 가정해도 된다.

\medskip\noindent
사상 $p : X_{\overline{k}} \to X$는 정수적이고 평탄하며 전사이다.
평탄성에 의해 일반화는 $p$를 따라 올라간다. Morphisms, Lemma
\ref{morphisms-lemma-generalizations-lift-flat}을 보라.
따라서 $X_{\overline{k}}$의 기약 성분들의 일반점은
$X$의 기약 성분들의 일반점으로 사상된다. 정수성에 의해
$p$는 보편 닫힌 사상이다. Morphisms, Lemma
\ref{morphisms-lemma-integral-universally-closed}을 보라.
따라서 기약 성분의 상 $p(\overline{T})$는 $X$의 한 기약 성분의
일반점을 포함하는 닫힌 기약 부분집합이다. 그러므로
$p(\overline{T})$는 $X$의 기약 성분이다. 이것으로 첫째 주장이 증명된다.
$T \subset X$가 기약 성분이면
$p^{-1}(T) =T_K$는 공집합이 아닌 기약 성분들의 합집합이다.
보조정리 \ref{varieties-lemma-inverse-image-irreducible}을 보라.
첫째 부분에 의해 이 성분들은 각각 반드시 $T$ 위로 전사된다.
따라서 이 사상은 전사이다.
\end{proof}

\begin{lemma}
\label{varieties-lemma-irreducible-dense-rational-points-geometrically-irreducible}
$k$를 체라 하고 $X$를 $k$ 위의 스킴이라 하자.
$X$가 기약이고 $k$-유리점들의 조밀한 집합을 가지면
$X$는 기하적으로 기약이다.
\end{lemma}

\begin{proof}
$k'/k$를 유한 체 확대라 하고 $Z, Z' \subset X_{k'}$를 기약 성분들이라 하자.
$Z = Z'$임을 보이면 충분하다. 보조정리
\ref{varieties-lemma-characterize-geometrically-irreducible}을 보라.
보조정리 \ref{varieties-lemma-image-irreducible}에 의해
$p(Z) = p(Z') = X$이며, 여기서 $p : X_{k'} \to X$는 사영이다.
$Z \not = Z'$이면 $Z \cap Z'$는 $X_{k'}$에서 조밀한 곳이 없으므로
$p(Z \cap Z')$도 조밀하지 않다. Morphisms, Lemma
\ref{morphisms-lemma-image-nowhere-dense-finite}를 보라.
여기서 $p$가 유한 사상 $\Spec(k') \to \Spec(k)$의 기저 변환이므로
유한 사상이라는 사실도 사용한다. Morphisms, Lemma
\ref{morphisms-lemma-base-change-finite}를 보라.
따라서 $k$-유리점 $x \in X$를 택하여
$x \not \in p(Z \cap Z')$가 되게 할 수 있다.
$x$의 잉여체가 $k$이므로
$p^{-1}(\{x\}) = \{x'\}$이고, 여기서 $x' \in X_{k'}$의 잉여체는 $k'$이다.
$x \in p(Z) = p(Z')$이므로 $x' \in Z \cap Z'$를 얻는데,
이는 원하는 모순이다.
\end{proof}

\begin{lemma}
\label{varieties-lemma-galois-action-irreducible-components}
$k$를 체라 하고 그 분리 대수적 폐포를 $\overline{k}$라 하자.
$X$를 $k$ 위의 스킴이라 하자. 다음 작용이 존재한다.
$$
\text{Gal}(\overline{k}/k)^{opp} \times \text{IrredComp}(X_{\overline{k}})
\longrightarrow
\text{IrredComp}(X_{\overline{k}})
$$
이 작용은 다음 성질들을 갖는다.
\begin{enumerate}
\item 원소 $\overline{T} \in \text{IrredComp}(X_{\overline{k}})$가
이 작용에 의해 고정될 필요충분조건은 기약 성분 $T \subset X$가
$k$ 위에서 기하적으로 기약이고
$T_{\overline{k}} = \overline{T}$인 것이다.
\item 임의의 체 확대 $k'/k$와 그 분리 대수적 폐포
$\overline{k}'$에 대해 도식
$$
\xymatrix{
\text{Gal}(\overline{k}'/k') \times \text{IrredComp}(X_{\overline{k}'})
\ar[r] \ar[d] &
\text{IrredComp}(X_{\overline{k}'}) \ar[d] \\
\text{Gal}(\overline{k}/k) \times \text{IrredComp}(X_{\overline{k}})
\ar[r] &
\text{IrredComp}(X_{\overline{k}})
}
$$
은 가환이다(오른쪽 수직 화살표는 보조정리
\ref{varieties-lemma-separably-closed-field-irreducible-components}에 의해 전단사이다).
\end{enumerate}
\end{lemma}

\begin{proof}
$\text{Gal}(\overline{k}/k)$의 $X_{\overline{k}}$ 위 작용
(\ref{varieties-equation-galois-action-base-change-kbar})은 그 기약 성분들 위의 작용을 유도한다.
기약 성분은 항상 닫혀 있다
(Topology, Lemma \ref{topology-lemma-connected-components}).
따라서 $\overline{T}$가 (1)과 같으면 보조정리
\ref{varieties-lemma-closed-fixed-by-Galois}에 의해 닫힌 부분집합
$T \subset X$가 존재하여 $\overline{T} = T_{\overline{k}}$이다.
$T$는 $k$ 위에서 기하적으로 기약임에 유의하자. 보조정리
\ref{varieties-lemma-characterize-geometrically-irreducible}을 보라.
$T$가 $X$의 기약 성분임을 보이기 위해
$T \subset T'$, $T \not = T'$이고 $T'$가 $X$의 기약 성분이라고 가정하자.
$\overline{\eta}$를 $\overline{T}$의 일반점이라 하자.
이는 일반점 $\eta$로 사상되며, 그 점은 $T$의 일반점이다.
그러면 일반점 $\xi \in T'$는 $\eta$로 특수화된다.
$X_{\overline{k}} \to X$가 평탄이므로 점
$\overline{\xi} \in X_{\overline{k}}$가 존재하여 $\xi$로 사상되고
$\overline{\eta}$로 특수화된다.
따라서 한원소 집합 $\{\overline{\xi}\}$의 폐포는
$X_{\overline{\xi}}$의 기약 닫힌 부분집합이며
$\overline{T}$를 진부분집합으로 포함한다. 이것이 원하는 모순이다.

\medskip\noindent
(2)의 함자성에 대한 증명은 생략한다.
\end{proof}

\begin{lemma}
\label{varieties-lemma-orbit-irreducible-components}
$k$를 체라 하고 그 분리 대수적 폐포를 $\overline{k}$라 하자.
$X$를 $k$ 위의 스킴이라 하자. 보조정리
\ref{varieties-lemma-image-irreducible}의 사상
$$
\text{IrredComp}(X_{\overline{k}})
\longrightarrow
\text{IrredComp}(X)
$$
의 올들은 보조정리 \ref{varieties-lemma-galois-action-irreducible-components}의 작용 아래
$\text{Gal}(\overline{k}/k)$의 궤도들과 정확히 같다.
\end{lemma}

\begin{proof}
$T \subset X$를 $X$의 기약 성분이라 하고 $\eta \in T$를 그 일반점이라 하자.
보조정리 \ref{varieties-lemma-inverse-image-irreducible}과
\ref{varieties-lemma-image-irreducible}에 의해 $\overline{T}$의 기약 성분들의 일반점 중
$T$로 사상되는 것들은 $\eta$로 사상된다.
Algebra, Lemma \ref{algebra-lemma-Galois-orbit}에 의해 갈루아 군은
$X_{\overline{k}}$의 모든 점 가운데 $\eta$로 사상되는 것들 위에서 추이적으로 작용한다.
따라서 보조정리가 성립한다.
\end{proof}

\begin{lemma}
\label{varieties-lemma-galois-action-irreducible-components-locally-finite-type}
$k$를 체라 하자. $X \to \Spec(k)$가 국소 유한형이라고 가정하자.
이 경우
\begin{enumerate}
\item 작용
$$
\text{Gal}(\overline{k}/k)^{opp} \times \text{IrredComp}(X_{\overline{k}})
\longrightarrow
\text{IrredComp}(X_{\overline{k}})
$$
은 $\text{IrredComp}(X_{\overline{k}})$에 이산 위상을 주면 연속이고,
\item $X_{\overline{k}}$의 모든 기약 성분은 $k$의 유한 확대 위에서
정의될 수 있으며,
\item 임의의 기약 성분 $T \subset X$에 대해 스킴
$T_{\overline{k}}$는 $X_{\overline{k}}$의 기약 성분들의 유한 합집합이고,
이 성분들은 모두 같은 $\text{Gal}(\overline{k}/k)$-궤도에 속한다.
\end{enumerate}
\end{lemma}

\begin{proof}
$\overline{T}$를 $X_{\overline{k}}$의 기약 성분이라 하자.
아핀 열린집합 $U \subset X$를 택하여
$\overline{T} \cap U_{\overline{k}}$가 공집합이 아니게 할 수 있다.
$U = \Spec(A)$로 쓰자. 그러면 $A$는 유한형 $k$-대수이다.
Morphisms, Lemma \ref{morphisms-lemma-locally-finite-type-characterize}를 보라.
따라서 $A_{\overline{k}}$는 유한형 $\overline{k}$-대수이고 특히 뇌터이다.
$\mathfrak p = (f_1, \ldots, f_n)$를
$\overline{T} \cap U_{\overline{k}}$에 대응하는 소 아이디얼이라 하자.
$A_{\overline{k}} = A \otimes_k \overline{k}$이므로 유한 부분확대
$\overline{k}/k'/k$가 존재하여 각 $f_i \in A_{k'}$이다.
$\text{Gal}(\overline{k}/k')$가 $\overline{T}$를 고정함은 분명하고,
이것으로 (1)이 증명된다.

\medskip\noindent
(2)는 $k'$ 위의 상황에 보조정리
\ref{varieties-lemma-galois-action-irreducible-components}의 (1)을 적용하면 따른다.
그 결과 기약 성분 $\overline{T}$는 $T'_{\overline{k}}$의 꼴이며,
여기서 $T' \subset X_{k'}$는 어떤 기약 성분이다.

\medskip\noindent
(3)을 증명하기 위해 $T \subset X$를 기약 성분이라 하자.
기약 성분 $\overline{T} \subset X_{\overline{k}}$를 택하여 $T$로 사상되게 하자.
보조정리 \ref{varieties-lemma-image-irreducible}을 보라.
위에서 본 바에 의해 $\overline{T}$의 궤도는 유한하며, 이를
$\overline{T}_1, \ldots, \overline{T}_n$이라 하자. 그러면
$\overline{T}_1 \cup \ldots \cup \overline{T}_n$은
$\text{Gal}(\overline{k}/k)$-불변인 $X_{\overline{k}}$의 닫힌 부분집합이므로,
보조정리 \ref{varieties-lemma-closed-fixed-by-Galois}에 의해 $W_{\overline{k}}$의 꼴이며,
여기서 $W \subset X$는 어떤 닫힌 부분집합이다.
분명히 $W = T$이고 증명이 끝난다.
\end{proof}

\begin{lemma}
\label{varieties-lemma-finite-extension-geometrically-irreducible-components}
$k$를 체라 하고 $X \to \Spec(k)$가 국소 유한형이라고 하자.
$X$가 유한 개의 기약 성분을 갖는다고 가정하자.
그러면 유한 분리 확대 $k'/k$가 존재하여
$X_{k'}$의 모든 기약 성분이 $k'$ 위에서 기하적으로 기약이다.
\end{lemma}

\begin{proof}
$\overline{k}$를 $k$의 분리 대수적 폐포라 하자.
$X$가 유한 개의 기약 성분을 갖는다는 가정과 보조정리
\ref{varieties-lemma-galois-action-irreducible-components-locally-finite-type}의 (3)을 합치면
$X_{\overline{k}}$가 유한 개의 기약 성분
$\overline{T}_1, \ldots, \overline{T}_n$을 가짐을 알 수 있다.
보조정리 \ref{varieties-lemma-galois-action-irreducible-components-locally-finite-type}의 (2)에 의해
유한 확대 $\overline{k}/k'/k$와 기약 성분들
$T_i \subset X_{k'}$가 존재하여
$\overline{T}_i = T_{i, \overline{k}}$이고, 이것으로 증명이 끝난다.
\end{proof}

\begin{lemma}
\label{varieties-lemma-irreducible-components-geometrically-irreducible}
$X$를 체 $k$ 위의 스킴이라 하자.
$X$가 유한 개의 기약 성분을 가지며 이들이 모두 기하적으로 기약이라고 가정하자.
그러면 $X$는 유한 개의 연결 성분을 가지며 각각은 기하적으로 연결이다.
\end{lemma}

\begin{proof}
연결 성분은 기약 성분들의 합집합이므로 분명하다. 세부사항은 생략한다.
\end{proof}







\section{기하적으로 정역적인 스킴}
\label{varieties-section-geometrically-integral}

\noindent
$X$가 체 위의 정역적 스킴이어도 기초체를 확대하면 $X$가
비축약 또는 가약이 될 수 있다. 기하적으로 정역적인 스킴에서는
이런 일이 일어나지 않는다.

\begin{definition}
\label{varieties-definition-geometrically-integral}
$X$를 체 $k$ 위의 스킴이라 하자.
\begin{enumerate}
\item $x \in X$라 하자. $X$가 {\it $x$에서 기하적으로 점별 정역적}이라고 함은
모든 체 확대 $k'/k$와 모든 $x' \in X_{k'}$에 대해, 이것이 $x$ 위에 놓이면
국소환 $\mathcal{O}_{X_{k'}, x'}$가 정역인 경우이다.
\item $X$가 {\it 기하적으로 점별 정역적}이라고 함은
$X$가 모든 점에서 기하적으로 점별 정역적인 경우이다.
\item $X$가 $k$ 위에서 {\it 기하적으로 정역적}이라고 함은
스킴 $X_{k'}$가 모든 체 확대 $k'$가 $k$의 확대일 때 정역적인 경우이다.
\end{enumerate}
\end{definition}

\noindent
개념 (2)와 (3)은 구별할 필요가 있다. 예를 들어
$k = \mathbf{R}$이고 $X = \Spec(\mathbf{C}[x])$이면
$X$는 $\mathbf{R}$ 위에서 기하적으로 점별 정역적이지만,
물론 기하적으로 정역적이지는 않다.

\begin{lemma}
\label{varieties-lemma-geometrically-integral}
$k$를 체라 하고 $X$를 $k$ 위의 스킴이라 하자.
$X$가 $k$ 위에서 기하적으로 정역적일 필요충분조건은
$X$가 $k$ 위에서 기하적으로 축약이고 동시에 기하적으로 기약인 것이다.
\end{lemma}

\begin{proof}
Properties, Lemma \ref{properties-lemma-characterize-integral}을 보라.
\end{proof}

\begin{lemma}
\label{varieties-lemma-proper-geometrically-reduced-global-sections}
$k$를 체라 하고 $X$를 $k$ 위의 고유 스킴이라 하자.
\begin{enumerate}
\item $A = H^0(X, \mathcal{O}_X)$는 유한 차원 $k$-대수이다.
\item $A = \prod_{i = 1, \ldots, n} A_i$는 아르틴 국소
$k$-대수들의 곱이며, $X$의 각 연결 성분에 인자가 하나씩 대응한다.
\item $X$가 축약이면 $A = \prod_{i = 1, \ldots, n} k_i$는
체들의 곱이며 각 체는 $k$의 유한 확대이다.
\item $X$가 기하적으로 축약이면 $k_i$는 $k$의 유한 분리 확대이다.
\item $X$가 기하적으로 연결이면 $A$는 $k$ 위에서 기하적으로 기약이다.
\item $X$가 기하적으로 기약이면 $A$는 $k$ 위에서 기하적으로 기약이다.
\item $X$가 기하적으로 축약이고 기하적으로 연결이면 $A = k$이다.
\item $X$가 기하적으로 정역적이면 $A = k$이다.
\end{enumerate}
\end{lemma}

\begin{proof}
Cohomology of Schemes, Lemma
\ref{coherent-lemma-proper-over-affine-cohomology-finite}에 의해
$A = H^0(X, \mathcal{O}_X)$는 유한 차원 $k$-대수이다.
이것으로 (1)이 증명된다.

\medskip\noindent
그러면 Algebra, Lemma \ref{algebra-lemma-finite-dimensional-algebra}와
Proposition \ref{algebra-proposition-dimension-zero-ring}에 의해
$A$는 아르틴 국소 $k$-대수들의 곱이다.
$X = Y \amalg Z$이고 $Y$, $Z$가 $X$에서 열려 있으면, 멱등원
$e \in A$를 얻기 위해 $\mathcal{O}_X$의 절단 가운데 값이 $1$인 곳이
$Y$이고 값이 $0$인 곳이 $Z$인 것을 택한다. 역으로 $e \in A$가
멱등원이면 $X$의 대응하는 분해를 얻는다. 마지막으로 $X$의 바탕
위상공간은 뇌터이므로 그 연결 성분들은 열려 있다. 따라서 $X$의 연결
성분들은 $1$ 대 $1$로 $A$의 원시 멱등원들과 대응한다. 이것으로
(2)가 증명된다.

\medskip\noindent
$X$가 축약이면 $A$도 축약이다. 따라서 국소환 $A_i = k_i$들은
축약이고, 그러므로 체이다(예를 들어 Algebra, Lemma
\ref{algebra-lemma-minimal-prime-reduced-ring}에 의한다).
이것으로 (3)이 증명된다.

\medskip\noindent
$X$가 기하적으로 축약이면
$A \otimes_k \overline{k} =
H^0(X_{\overline{k}}, \mathcal{O}_{X_{\overline{k}}})$는 축약이다
(등식은 Cohomology of Schemes, Lemma
\ref{coherent-lemma-flat-base-change-cohomology}에 의한다).
이는 $k_i \otimes_k \overline{k}$가 체들의 곱임을 함의하므로
$k_i/k$는 분리 확대이다. 예를 들어 Algebra, Lemma
\ref{algebra-lemma-characterize-separable-field-extensions}와
\ref{algebra-lemma-geometrically-reduced-finite-purely-inseparable-extension}을 보라.
이것으로 (4)가 증명된다.

\medskip\noindent
$X$가 기하적으로 연결이면
$A \otimes_k \overline{k} =
H^0(X_{\overline{k}}, \mathcal{O}_{X_{\overline{k}}})$는 (2)에 의해
0차원 국소환이므로 그 스펙트럼은 점 하나를 갖고, 특히 기약이다.
따라서 $A$는 기하적으로 기약이다. 이것으로 (5)가 증명된다.
물론 (5)는 (6)을 함의한다.

\medskip\noindent
$X$가 기하적으로 축약이고 기하적으로 연결이면
$A = k_1$은 체이고 확대 $k_1/k$는 유한 분리이면서 기하적으로 기약이다.
그러나 이 경우 $k_1 \otimes_k \overline{k}$는
$[k_1 : k]$개의 $\overline{k}$ 복사본의 곱이므로
$k_1 = k$임을 얻는다. 이것으로 (7)이 증명된다.
물론 (7)은 (8)을 함의한다.
\end{proof}

\noindent
다음은 스타인 분해의 초보적 판본이다. 실제 스타인 분해는
More on Morphisms, Section
\ref{more-morphisms-section-stein-factorization}에서 다룰 것이다.

\begin{lemma}
\label{varieties-lemma-baby-stein}
$X$를 체 $k$ 위의 고유 스킴이라 하자.
$A = H^0(X, \mathcal{O}_X)$로 두자. 표준 사상
$X \to \Spec(A)$의 올들은 기하적으로 연결이다.
\end{lemma}

\begin{proof}
$S = \Spec(A)$로 두자. 표준 사상 $X \to S$는 Schemes, Lemma
\ref{schemes-lemma-morphism-into-affine}를 통해
$\Gamma(S, \mathcal{O}_S) = A = \Gamma(X, \mathcal{O}_X)$에 대응하는 사상이다.
$k$-대수 $A$는 국소 아르틴 대수들의 유한 곱 $A = \prod A_i$이며,
각 인자는 $k$-대수로서 $k$ 위에서 유한하다. 보조정리
\ref{varieties-lemma-proper-geometrically-reduced-global-sections}을 보라.
$s_i \in S$를 $A_i$의 극대 아이디얼에 대응하는 점이라 하자.
$\overline{k}$를 $k$의 대수적 폐포로 택하고
$\overline{A} = A \otimes_k \overline{k}$로 두자.
$\kappa(s_i) \to \overline{k}$를 $k$ 위의 매장으로 택하자. 이는
$\overline{k}$-대수 사상
$$
\sigma_i : \overline{A} = A \otimes_k \overline{k} \to
\kappa(s_i) \otimes_k \overline{k} \to \overline{k}
$$
을 결정한다. 다음 기저 변환을 생각하자.
$$
\xymatrix{
\overline{X} \ar[r] \ar[d] & X \ar[d] \\
\overline{S} \ar[r] & S
}
$$
이는 $X$를 $\overline{S} = \Spec(\overline{A})$로 기저 변환한 것이다.
Cohomology of Schemes, Lemma
\ref{coherent-lemma-flat-base-change-cohomology}에 의해
$\Gamma(\overline{X}, \mathcal{O}_{\overline{X}}) = \overline{A}$이다.
$\overline{s}_i \in \Spec(\overline{A})$가 $\overline{k}$-유리점이고
$\sigma_i$에 대응하면, $\overline{s}_i$는 $s_i \in S$로 사상되고
$\overline{X}_{\overline{s}_i}$는 $X_{s_i}$를 $\Spec(\sigma_i)$에 따라
기저 변환한 것이다. 따라서 $k$가 대수적으로 닫힌 경우를
증명하면 충분하다.

\medskip\noindent
$k$가 대수적으로 닫혀 있다고 가정하자. 이 경우 $\kappa(s_i)$는
대수적으로 닫혀 있고 $X_{s_i}$가 연결임을 보여야 한다.
곱 분해 $A = \prod A_i$는 서로소 합 분해
$\Spec(A) = \coprod \Spec(A_i)$에 대응한다. Algebra, Lemma
\ref{algebra-lemma-spec-product}를 보라.
$X_i$를 $\Spec(A_i)$의 역상이라 하자. 보조정리
\ref{varieties-lemma-proper-geometrically-reduced-global-sections}의 (2)에 의해
$A_i = \Gamma(X_i, \mathcal{O}_{X_i})$이다.
$X_{s_i} \to X_i$는 닫힌 몰입이고 바탕 위상공간들 사이의 동형을 유도한다
($\Spec(A_i)$가 한원소 집합이기 때문이다). 따라서 $X_{s_i}$가
연결이 아니면 $X_i$도 연결이 아니다. 그러므로 $X_i$가 공집합이고
$A_i = 0$이거나, $X_i$를 $U \amalg V$로 쓸 수 있다.
여기서 $U$와 $V$는 열려 있고 공집합이 아니며, 이 경우
$A_i$는 자명하지 않은 멱등원을 갖게 된다.
$A_i$는 국소환이므로 이는 모순이고, 증명이 끝난다.
\end{proof}

\begin{lemma}
\label{varieties-lemma-geometrically-reduced-stein}
$k$를 체라 하고 $X$를 $k$ 위의 고유하고 기하적으로 축약인 스킴이라 하자.
다음 조건들은 서로 동치이다.
\begin{enumerate}
\item $H^0(X, \mathcal{O}_X) = k$이다.
\item $X$는 기하적으로 연결이다.
\end{enumerate}
\end{lemma}

\begin{proof}
보조정리 \ref{varieties-lemma-baby-stein}에 의해 (1) $\Rightarrow$ (2)이다.
보조정리 \ref{varieties-lemma-proper-geometrically-reduced-global-sections}에 의해
(2) $\Rightarrow$ (1)이다.
\end{proof}




\section{기하적으로 정규인 스킴}
\label{varieties-section-geometrically-normal}

\noindent
Properties, Definition \ref{properties-definition-normal}에서
정규 스킴의 개념을 정의하였다. 이 개념은 뇌터가 아닌 스킴에 대해서도
정의된다. 따라서 ‘‘기하적으로 정칙인’’ 스킴을 논의할 때와 달리,
여기서는 기초체의 모든 체 확대를 고려한다.

\begin{definition}
\label{varieties-definition-geometrically-normal}
$X$를 체 $k$ 위의 스킴이라 하자.
\begin{enumerate}
\item $x \in X$라 하자. $X$가 {\it $x$에서 기하적으로 정규}라고 함은
모든 체 확대 $k'/k$와 모든 $x' \in X_{k'}$에 대해, 이것이
$x$ 위에 놓일 때 국소환 $\mathcal{O}_{X_{k'}, x'}$가 정규인 경우이다.
\item $X$가 $k$ 위에서 {\it 기하적으로 정규}라고 함은 $X$가
모든 $x \in X$에서 기하적으로 정규인 경우이다.
\end{enumerate}
\end{definition}

\begin{lemma}
\label{varieties-lemma-geometrically-normal-at-point}
$k$를 체라 하고 $X$를 $k$ 위의 스킴이라 하자.
$x \in X$라 하자. 다음 조건들은 서로 동치이다.
\begin{enumerate}
\item $X$는 $x$에서 기하적으로 정규이다.
\item 모든 유한 순수 비분리 체 확대 $k'$가 $k$의 확대이고
모든 $x' \in X_{k'}$가 $x$ 위에 놓일 때 국소환
$\mathcal{O}_{X_{k'}, x'}$가 정규이다.
\item 환 $\mathcal{O}_{X, x}$는 $k$ 위에서 기하적으로 정규이다
(Algebra, Definition \ref{algebra-definition-geometrically-normal} 참조).
\end{enumerate}
\end{lemma}

\begin{proof}
(1)이 (2)를 함의함은 분명하다. (2)를 가정하자. $k'/k$를
유한 순수 비분리 체 확대라 하자(예를 들어 $k = k'$일 수 있다).
환 $\mathcal{O}_{X, x} \otimes_k k'$를 생각하자.
Algebra, Lemma \ref{algebra-lemma-p-ring-map}에 의해 그 스펙트럼은
$\mathcal{O}_{X, x}$의 스펙트럼과 같다. 따라서 이 환도 국소환이다
(Algebra, Lemma \ref{algebra-lemma-characterize-local-ring}).
그러므로 유일한 점 $x' \in X_{k'}$가 $x$ 위에 놓이며,
$\mathcal{O}_{X_{k'}, x'} \cong \mathcal{O}_{X, x} \otimes_k k'$이다.
가정에 의해 이는 정규환이다. 따라서 Algebra, Lemma
\ref{algebra-lemma-geometrically-normal}에 의해 (3)을 얻는다.

\medskip\noindent
(3)을 가정하자. $k'/k$를 체 확대라 하자.
$\Spec(k') \to \Spec(k)$가 전사이므로 $X_{k'} \to X$도 전사이다
(Morphisms, Lemma \ref{morphisms-lemma-base-change-surjective}).
임의의 점 $x' \in X_{k'}$를 택하되 이것이 $x$ 위에 놓인다고 하자.
국소환 $\mathcal{O}_{X_{k'}, x'}$는
환 $\mathcal{O}_{X, x} \otimes_k k'$의 국소화이다.
따라서 가정에 의해 정규이고, (1)이 증명된다.
\end{proof}

\begin{lemma}
\label{varieties-lemma-geometrically-normal}
$k$를 체라 하고 $X$를 $k$ 위의 스킴이라 하자.
다음 조건들은 서로 동치이다.
\begin{enumerate}
\item $X$는 기하적으로 정규이다.
\item $X_{k'}$는 모든 체 확대 $k'/k$에 대해 정규 스킴이다.
\item $X_{k'}$는 모든 유한 생성 체 확대 $k'/k$에 대해 정규 스킴이다.
\item $X_{k'}$는 모든 유한 순수 비분리 체 확대 $k'/k$에 대해
정규 스킴이다.
\item 모든 아핀 열린집합 $U \subset X$에 대해 환 $\mathcal{O}_X(U)$는
기하적으로 정규이다
(Algebra, Definition \ref{algebra-definition-geometrically-normal} 참조).
\item $X_{k^{perf}}$는 정규 스킴이다.
\end{enumerate}
\end{lemma}

\begin{proof}
(1)을 가정하자. 그러면 모든 체 확대 $k'/k$와 모든 점
$x' \in X_{k'}$에 대해 $X_{k'}$의 $x'$에서의 국소환은 정규이다.
정의에 의해 이는 $X_{k'}$가 정규라는 뜻이다. 따라서 (2)가 성립한다.

\medskip\noindent
(2)가 (3)을 함의하고 (3)이 (4)를 함의함은 분명하다.

\medskip\noindent
(4)를 가정하고 $U \subset X$를 아핀 열린 부분스킴이라 하자.
$U_{k'}$는 모든 유한 순수 비분리 확대 $k'/k$에 대해
(여기에는 $k = k'$도 포함된다) 정규 스킴이다. 이는
$k' \otimes_k \mathcal{O}(U)$가 모든 유한 순수 비분리 확대
$k'/k$에 대해 정규환이라는 뜻이다.
따라서 정의에 의해 $\mathcal{O}(U)$는 기하적으로 정규인
$k$-대수이다. 그러므로 (4)는 (5)를 함의한다.

\medskip\noindent
(5)를 가정하자. 임의의 체 확대 $k'/k$에 대해 기저 변환
$X_{k'}$는 환 $\mathcal{O}_X(U) \otimes_k k'$들의 스펙트럼을
이어 붙여 얻어진다. 여기서 $U$는 $X$의 아핀 열린집합이다
(Schemes, Section \ref{schemes-section-fibre-products} 참조).
따라서 $X_{k'}$는 정규이다. 그러므로 (1)이 성립한다.

\medskip\noindent
(5)와 (6)의 동치는 기하적으로 정규인 대수의 정의와
방금 증명한 (3)과 (4)의 동치에서 따른다.
\end{proof}

\begin{lemma}
\label{varieties-lemma-geometrically-normal-upstairs}
$k$를 체라 하고 $X$를 $k$ 위의 스킴이라 하자.
$k'/k$를 체 확대라 하자. $x \in X$를 점이라 하고,
점 $x' \in X_{k'}$를 택하되 이것이 $x$ 위에 놓인다고 하자.
다음 조건들은 서로 동치이다.
\begin{enumerate}
\item $X$는 $x$에서 기하적으로 정규이다.
\item $X_{k'}$는 $x'$에서 기하적으로 정규이다.
\end{enumerate}
특히 $X$가 $k$ 위에서 기하적으로 정규일 필요충분조건은
$X_{k'}$가 $k'$ 위에서 기하적으로 정규인 것이다.
\end{lemma}

\begin{proof}
(1)이 (2)를 함의함은 분명하다. (2)를 가정하자.
$k''/k$를 유한 순수 비분리 체 확대라 하고,
점 $x'' \in X_{k''}$를 택하되 이것이 $x$ 위에 놓인다고 하자
(실제로 이 점은 유일하다).
Fields, Lemma \ref{fields-lemma-common-extension-field}에 따라
공통 체 확대 $k'''/k$를 택하되 이것이 $k'/k$와 $k''/k$의
공통 확대가 되게 하자. 점 $x''' \in X_{k'''}$를 택하되 이것이
$x'$와 $x''$ 둘 다의 위에 놓이게 하자.
다음 국소환 사상을 생각하자.
$$
\mathcal{O}_{X_{k''}, x''} \longrightarrow \mathcal{O}_{X_{k'''}, x''''}.
$$
이는 평탄 국소환 준동형이므로 충실히 평탄하다.
(2)에 의해 오른쪽 국소환은 정규이다. 따라서 Algebra, Lemma
\ref{algebra-lemma-descent-normal}에 의해
$\mathcal{O}_{X_{k''}, x''}$가 정규임을 얻는다.
Lemma \ref{varieties-lemma-geometrically-normal-at-point}에 의해 $X$는
$x$에서 기하적으로 정규이다.
\end{proof}

\begin{lemma}
\label{varieties-lemma-fibre-product-normal}
$k$를 체라 하자. $X$를 $k$ 위에서 기하적으로 정규인 스킴이라 하고,
$Y$를 $k$ 위의 정규 스킴이라 하자. 그러면 $X \times_k Y$는
정규 스킴이다.
\end{lemma}

\begin{proof}
이는 Algebra, Lemma \ref{algebra-lemma-geometrically-normal-tensor-normal}로
귀착된다. 여기에는 Lemma \ref{varieties-lemma-geometrically-normal}을 사용한다.
\end{proof}

\begin{lemma}
\label{varieties-lemma-base-change-normal-by-separable}
$k$를 체라 하고 $X$를 $k$ 위의 정규 스킴이라 하자.
$K/k$를 분리 체 확대라 하자. 그러면 $X_K$는 정규 스킴이다.
\end{lemma}

\begin{proof}
Lemma \ref{varieties-lemma-fibre-product-normal}와 Algebra, Lemma
\ref{algebra-lemma-separable-field-extension-geometrically-normal}에서 따른다.
\end{proof}

\begin{lemma}
\label{varieties-lemma-geometrically-normal-stein}
$k$를 체라 하고 $X$를 $k$ 위의 고유하고 기하적으로 정규인 스킴이라 하자.
다음 조건들은 서로 동치이다.
\begin{enumerate}
\item $H^0(X, \mathcal{O}_X) = k$이다.
\item $X$는 기하적으로 연결이다.
\item $X$는 기하적으로 기약이다.
\item $X$는 기하적으로 정역적이다.
\end{enumerate}
\end{lemma}

\begin{proof}
Lemma \ref{varieties-lemma-geometrically-reduced-stein}에 의해 (1)과 (2)는 동치이다.
국소 뇌터 정규 스킴(가령 $X_{\overline{k}}$)은 그 기약 성분들의
서로소 합이다(Properties, Lemma \ref{properties-lemma-normal-Noetherian}).
따라서 (2)와 (3)은 동치이다. $X_{\overline{k}}$는 축약이라고
가정되어 있으므로 (3)과 (4)도 동치이다.
\end{proof}







\section{체의 확대와 국소 뇌터 스킴}
\label{varieties-section-locally-Noetherian}

\noindent
$X$를 체 $k$ 위의 국소 뇌터 스킴이라 하자.
$X_{k'}$ 역시 언제나 국소 뇌터인 것은 아니다.
예를 들어 $X = \Spec(\overline{\mathbf{Q}})$이고
$k = \mathbf{Q}$이면, $X_{\overline{\mathbf{Q}}}$는
$\overline{\mathbf{Q}} \otimes_{\mathbf{Q}} \overline{\mathbf{Q}}$의
스펙트럼인데 이는 뇌터가 아니다(힌트: 멱등원이 너무 많다).
그러나 유한 생성 체 확대만을 사용하여 기저 변환하면 뇌터 성질은 보존된다.
(또는 $X$가 $k$ 위에서 국소 유한형이어도 된다. 이 성질은
기저 변환으로 보존되기 때문이다.)

\begin{lemma}
\label{varieties-lemma-locally-Noetherian-base-change}
$k$를 체라 하고 $X$를 $k$ 위의 스킴이라 하자.
$k'/k$를 유한 생성 체 확대라 하자. 그러면 $X$가 국소 뇌터일
필요충분조건은 $X_{k'}$가 국소 뇌터인 것이다.
\end{lemma}

\begin{proof}
Properties, Lemma \ref{properties-lemma-locally-Noetherian}을 사용하면
$X$가 아핀인 경우로 귀착된다. $X = \Spec(A)$라 하자.
이 경우 $A$가 뇌터일 필요충분조건이 $A_{k'}$가 뇌터임을 증명해야 한다.
$A \to A_{k'} = k' \otimes_k A$가 충실히 평탄하므로,
$A_{k'}$가 뇌터이면 Algebra, Lemma
\ref{algebra-lemma-descent-Noetherian}에 의해 $A$도 뇌터이다.
역으로 $A$가 뇌터이면 Algebra, Lemma
\ref{algebra-lemma-Noetherian-field-extension}에 의해 $A_{k'}$도 뇌터이다.
\end{proof}







\section{기하적으로 정칙인 스킴}
\label{varieties-section-geometrically-regular}

\noindent
체 $k$ 위의 기하적으로 정칙인 스킴이란, $k$ 위의 국소 뇌터
스킴으로서 적절히 기초체를 바꾸어도 정칙으로 남는 스킴이다.
$k$ 위의 유한형 스킴이 기하적으로 정칙일 필요충분조건은
$k$ 위에서 매끄러운 것이다
(Lemma \ref{varieties-lemma-geometrically-regular-smooth} 참조).
기하적 정칙성은 형식적 올처럼 매끄러움을 사용할 수 없는 상황에서
가장 흥미롭다(여기에 이후 참조를 삽입한다).

\medskip\noindent
아래 정의에서는 국소 뇌터 스킴으로 한정한다. 정칙 국소환이라는
성질이 뇌터 국소환에 대해서만 정의되기 때문이다. 위의 Lemma
\ref{varieties-lemma-locally-Noetherian-base-change}에 의해 유한 생성 체
확대만 고려하면 이 성질은 기초체를 바꾸어도 보존된다.
이 절에서는 이 사실을 별도의 인용 없이 사용한다.
특히 다음 정의는 잘 정의된다.

\begin{definition}
\label{varieties-definition-geometrically-regular}
$k$를 체라 하고 $X$를 $k$ 위의 국소 뇌터 스킴이라 하자.
\begin{enumerate}
\item $x \in X$라 하자. $X$가 {\it $x$에서 기하적으로 정칙}이라고
$k$ 위에서 함은, 모든 유한 생성 체 확대 $k'/k$와 모든
$x' \in X_{k'}$에 대해, 이것이 $x$ 위에 놓일 때 국소환
$\mathcal{O}_{X_{k'}, x'}$가 정칙인 경우이다.
\item $X$가 {\it $k$ 위에서 기하적으로 정칙}이라고 함은
$X$가 모든 점에서 기하적으로 정칙인 경우이다.
\end{enumerate}
\end{definition}

\noindent
비슷한 정의로 체 위에서 기하적으로 코언--매콜리인 스킴,
$(R_k)$ 스킴, $(S_k)$ 스킴을 정의할 수 있다.
필요해지면 이들을 위한 절을 따로 추가할 것이다.

\begin{lemma}
\label{varieties-lemma-geometrically-regular-at-point}
$k$를 체라 하고 $X$를 $k$ 위의 국소 뇌터 스킴이라 하자.
$x \in X$라 하자. 다음 조건들은 서로 동치이다.
\begin{enumerate}
\item $X$는 $x$에서 기하적으로 정칙이다.
\item 모든 유한 순수 비분리 체 확대 $k'$가 $k$의 확대이고
모든 $x' \in X_{k'}$가 $x$ 위에 놓일 때 국소환
$\mathcal{O}_{X_{k'}, x'}$가 정칙이다.
\item 환 $\mathcal{O}_{X, x}$는 $k$ 위에서 기하적으로 정칙이다
(Algebra, Definition \ref{algebra-definition-geometrically-regular} 참조).
\end{enumerate}
\end{lemma}

\begin{proof}
(1)이 (2)를 함의함은 분명하다. (2)를 가정하자. 특히
$\mathcal{O}_{X, x}$는 정칙 국소환이다. $k'/k$를 유한 순수 비분리
체 확대라 하자. 환 $\mathcal{O}_{X, x} \otimes_k k'$를 생각하자.
Algebra, Lemma \ref{algebra-lemma-p-ring-map}에 의해 그 스펙트럼은
$\mathcal{O}_{X, x}$의 스펙트럼과 같다. 따라서 이 환도 국소환이다
(Algebra, Lemma \ref{algebra-lemma-characterize-local-ring}).
그러므로 유일한 점 $x' \in X_{k'}$가 $x$ 위에 놓이며,
$\mathcal{O}_{X_{k'}, x'} \cong \mathcal{O}_{X, x} \otimes_k k'$이다.
가정에 의해 이는 정칙환이다. 따라서 기하적으로 정칙인 환의 정의로부터
(3)을 얻는다.

\medskip\noindent
(3)을 가정하자. $k'/k$를 체 확대라 하자.
$\Spec(k') \to \Spec(k)$가 전사이므로 $X_{k'} \to X$도 전사이다
(Morphisms, Lemma \ref{morphisms-lemma-base-change-surjective}).
임의의 점 $x' \in X_{k'}$를 택하되 이것이 $x$ 위에 놓인다고 하자.
국소환 $\mathcal{O}_{X_{k'}, x'}$는 환
$\mathcal{O}_{X, x} \otimes_k k'$의 국소화이다.
따라서 가정에 의해 정칙이고, (1)이 증명된다.
\end{proof}

\begin{lemma}
\label{varieties-lemma-geometrically-regular}
$k$를 체라 하고 $X$를 $k$ 위의 국소 뇌터 스킴이라 하자.
다음 조건들은 서로 동치이다.
\begin{enumerate}
\item $X$는 기하적으로 정칙이다.
\item $X_{k'}$는 모든 유한 생성 체 확대 $k'/k$에 대해 정칙 스킴이다.
\item $X_{k'}$는 모든 유한 순수 비분리 체 확대 $k'/k$에 대해
정칙 스킴이다.
\item 모든 아핀 열린집합 $U \subset X$에 대해 환 $\mathcal{O}_X(U)$는
기하적으로 정칙이다
(Algebra, Definition \ref{algebra-definition-geometrically-regular} 참조).
\item 아핀 열린 덮개 $X = \bigcup U_i$가 존재하여, 각
$\mathcal{O}_X(U_i)$는 $k$ 위에서 기하적으로 정칙이다.
\end{enumerate}
\end{lemma}

\begin{proof}
(1)을 가정하자. 그러면 모든 유한 생성 체 확대 $k'/k$와
모든 점 $x' \in X_{k'}$에 대해 $X_{k'}$의 $x'$에서의 국소환은
정칙이다. Properties, Lemma
\ref{properties-lemma-characterize-regular}에 의해 이는
$X_{k'}$가 정칙이라는 뜻이다. 따라서 (2)가 성립한다.

\medskip\noindent
(2)가 (3)을 함의함은 분명하다.

\medskip\noindent
(3)을 가정하고 $U \subset X$를 아핀 열린 부분스킴이라 하자.
$U_{k'}$는 모든 유한 순수 비분리 확대 $k'/k$에 대해
(여기에는 $k = k'$도 포함된다) 정칙 스킴이다. 이는
$k' \otimes_k \mathcal{O}(U)$가 모든 유한 순수 비분리 확대
$k'/k$에 대해 정칙환이라는 뜻이다. 따라서 $\mathcal{O}(U)$는
기하적으로 정칙인 $k$-대수이고, (4)가 성립한다.

\medskip\noindent
(4)가 (5)를 함의함은 분명하다. (5)와 같은 아핀 열린 덮개
$X = \bigcup U_i$를 택하자. 임의의 체 확대 $k'/k$에 대해
기저 변환 $X_{k'}$는 환
$\mathcal{O}_X(U_i) \otimes_k k'$들의 스펙트럼을 이어 붙여 얻어진다
(Schemes, Section \ref{schemes-section-fibre-products} 참조).
따라서 $X_{k'}$는 정칙이다. 그러므로 (1)이 성립한다.
\end{proof}

\begin{lemma}
\label{varieties-lemma-geometrically-regular-upstairs}
$k$를 체라 하고 $X$를 $k$ 위의 스킴이라 하자.
$k'/k$를 유한 생성 체 확대라 하자. $x \in X$를 점이라 하고,
점 $x' \in X_{k'}$를 택하되 이것이 $x$ 위에 놓인다고 하자.
다음 조건들은 서로 동치이다.
\begin{enumerate}
\item $X$는 $x$에서 기하적으로 정칙이다.
\item $X_{k'}$는 $x'$에서 기하적으로 정칙이다.
\end{enumerate}
특히 $X$가 $k$ 위에서 기하적으로 정칙일 필요충분조건은
$X_{k'}$가 $k'$ 위에서 기하적으로 정칙인 것이다.
\end{lemma}

\begin{proof}
(1)이 (2)를 함의함은 분명하다. (2)를 가정하자.
$k''/k$를 유한 순수 비분리 체 확대라 하고,
점 $x'' \in X_{k''}$를 택하되 이것이 $x$ 위에 놓인다고 하자
(실제로 이 점은 유일하다). 유한 생성 공통 체 확대
$k'''/k$를 택하되 이것이 $k'/k$와 $k''/k$의 공통 확대가 되게 하자.
Fields, Lemma \ref{fields-lemma-common-extension-field}와 그 증명을 보라.
그리고 점 $x''' \in X_{k'''}$를 택하되 이것이 $x'$와
$x''$ 둘 다의 위에 놓이게 하자. 다음 국소환 사상을 생각하자.
$$
\mathcal{O}_{X_{k''}, x''} \longrightarrow \mathcal{O}_{X_{k'''}, x''''}.
$$
이는 뇌터 국소환 사이의 평탄 국소환 준동형이므로 충실히 평탄하다.
(2)에 의해 오른쪽 국소환은 정칙이다. 따라서 Algebra, Lemma
\ref{algebra-lemma-flat-under-regular}에 의해
$\mathcal{O}_{X_{k''}, x''}$가 정칙임을 얻는다.
Lemma \ref{varieties-lemma-geometrically-regular-at-point}에 의해 $X$는
$x$에서 기하적으로 정칙이다.
\end{proof}

\noindent
다음 보조정리는 Algebra, Lemma
\ref{algebra-lemma-geometrically-regular-descent}의 기하학적 변형이다.

\begin{lemma}
\label{varieties-lemma-flat-under-geometrically-regular}
$k$를 체라 하자. $f : X \to Y$를 $k$ 위의 국소 뇌터 스킴 사이의
사상이라 하자. $x \in X$를 점이라 하고 $y = f(x)$로 두자.
$X$가 $x$에서 기하적으로 정칙이고 $f$가 $x$에서 평탄이면,
$Y$는 $y$에서 기하적으로 정칙이다. 특히 $X$가 $k$ 위에서
기하적으로 정칙이고 $f$가 평탄 전사이면 $Y$는 $k$ 위에서
기하적으로 정칙이다.
\end{lemma}

\begin{proof}
$k'$를 $k$의 유한 순수 비분리 확대라 하자.
$f' : X_{k'} \to Y_{k'}$를 $f$의 기저 변환이라 하자.
$x' \in X_{k'}$를 $x$ 위에 놓이는 유일한 점이라 하자.
$Y_{k'}$가 $y' = f'(x')$에서 정칙임을 보이면, Lemma
\ref{varieties-lemma-geometrically-regular}에 의해 $Y$는 $k$ 위에서
$y'$에서 기하적으로 정칙이다. Morphisms, Lemma
\ref{morphisms-lemma-base-change-module-flat}에 의해
사상 $X_{k'} \to Y_{k'}$는 $x'$에서 평탄하다.
따라서 환 사상
$$
\mathcal{O}_{Y_{k'}, y'}
\longrightarrow
\mathcal{O}_{X_{k'}, x'}
$$
은 뇌터 국소환 사이의 평탄 국소 준동형이고, 가정에 의해 오른쪽 환은
정칙이다. 따라서 Algebra, Lemma
\ref{algebra-lemma-flat-under-regular}에 의해 왼쪽 환도 정칙 국소환이다.
\end{proof}

\begin{lemma}
\label{varieties-lemma-geometrically-regular-smooth}
$k$를 체라 하자. $X$를 $k$ 위에서 국소 유한형인 스킴이라 하고
$x \in X$라 하자. $X$가 $x$에서 기하적으로 정칙일 필요충분조건은
$X \to \Spec(k)$가 $x$에서 매끄러운 것이다
(Morphisms, Definition \ref{morphisms-definition-smooth}).
\end{lemma}

\begin{proof}
문제는 $x$의 근방에서 국소적이므로 $X = \Spec(A)$라 가정해도 된다.
여기서 이는 어떤 유한형 $k$-대수에 대한 표현이다. $x$가 소 아이디얼
$\mathfrak p$에 대응한다고 하자.

\medskip\noindent
$A$가 $k$ 위에서 $\mathfrak p$에서 매끄러우면 $A$를 국소화하여
$A$가 $k$ 위에서 매끄럽다고 가정할 수 있다. 이 경우
$k' \otimes_k A$는 $k'$ 위에서 모든 확대체 $k'/k$에 대해 매끄럽고,
이 뇌터환들은 각각 Algebra, Lemma
\ref{algebra-lemma-characterize-smooth-over-field}에 의해 정칙이다.

\medskip\noindent
$X$가 $x$에서 기하적으로 정칙이라고 가정하자.
$K := \kappa(x) = \kappa(\mathfrak p)$를 생각하자. 이는 $x$의
잉여체이고 $k$의
유한 생성 확대이다. Algebra, Lemma
\ref{algebra-lemma-make-separable}에 의해 유한 순수 비분리 확대
$k'/k$가 존재하여 합성체 $k'K$는 $k'$의 분리 체 확대가 된다.
$\mathfrak p' \subset A' = k' \otimes_k A$를 $\mathfrak p$ 위에
놓이는 소 아이디얼이라 하자. Algebra, Lemma
\ref{algebra-lemma-p-ring-map}에 의해 이는 $\mathfrak p$ 위에
놓이는 유일한 소 아이디얼이다. 따라서 잉여체
$K' := \kappa(\mathfrak p')$는 합성체 $k'K$이다. 가정에 의해 국소환
$(A')_{\mathfrak p'}$는 정칙이다. 따라서 Algebra, Lemma
\ref{algebra-lemma-separable-smooth}에 의해
$k' \to A'$는 $\mathfrak p'$에서 매끄럽다. 그러면 다시 Algebra, Lemma
\ref{algebra-lemma-smooth-field-change-local}에 의해
$k \to A$가 $\mathfrak p$에서 매끄럽다. 보조정리가 증명되었다.
\end{proof}


\begin{example}
\label{varieties-example-geometrically-reduced-not-normal}
$k =\mathbf{F}_p(t)$라 하자. 정칙이지만 기하적으로 축약이지 않은
대수다양체 $V$의 예는 $k$ 위에서 쉽게 줄 수 있다. 예를 들어
$\Spec(k[x]/(x^p - t))$를 택할 수 있다. 사실 기하적으로 축약이지만
기하적으로 정규조차 아닌 정칙 대수다양체 $V$의 예도 존재한다.
즉 $p > 2$일 때 스킴
$V = \Spec(k[x, y]/(y^2 - x^p + t))$를 택하자. 다항식
$y^2 - x^p + t \in k[x, y]$가 기약이므로 이는 대수다양체이다.
사상 $V \to \Spec(k)$는 점 $v_0 \in V$를 제외한 모든 점에서
매끄럽다. 이 점은 극대 아이디얼 $(y, x^p - t)$에 대응한다
($2y$가 가역이기 때문이다). 특히 $V$는 $v_0$를 제외하면
모든 점에서 (기하적으로) 정칙이다. 국소환
$$
\mathcal{O}_{V, v_0} = \left(k[x, y]/(y^2 - x^p + t)\right)_{(y, x^p - t)}
$$
은 차원 $1$인 정역이다. 그 극대 아이디얼은 원소 $1$개, 즉
$y$로 생성된다. 따라서 이는 이산 값매김환이고 정칙이다.
$k' = k[t^{1/p}]$라 하자. $t' = t^{1/p} \in k'$,
$V' = V_{k'}$로 쓰고, $v'_0 \in V'$를 $v_0$ 위에 놓이는
유일한 점이라 하자. $k'$ 위에서
$x^p - t = (x - t')^p$로 쓸 수 있지만 다항식
$y^2 - (x - t')^p$는 여전히 기약이고 $V'$는 여전히 대수다양체이다.
그러나 원소
$$
\frac{y}{x - t'} \in (\text{fraction field of }\mathcal{O}_{V', v'_0})
$$
는 $\mathcal{O}_{V', v'_0}$ 위에서 정역적이지만
(그 제곱을 계산하면 된다) 그 안에 들어 있지 않다. 따라서
$V'$는 $v'_0$에서 정규가 아니다. 이것으로 예가 끝난다.
\end{example}







\section{체의 확대와 코언--매콜리 성질}
\label{varieties-section-CM}

\noindent
다음 보조정리는 기하적으로 코언--매콜리인 스킴을 정의할 필요가
없음을 말해 준다. 그렇게 정의한 스킴은 코언--매콜리 스킴과
같기 때문이다.

\begin{lemma}
\label{varieties-lemma-CM-base-change}
$X$를 체 $k$ 위의 국소 뇌터 스킴이라 하자.
$k'/k$를 유한 생성 체 확대라 하자. $x \in X$를 점이라 하고,
점 $x' \in X_{k'}$를 택하되 이것이 $x$ 위에 놓인다고 하자.
그러면
$$
\mathcal{O}_{X, x}\text{ is Cohen-Macaulay}
\Leftrightarrow
\mathcal{O}_{X_{k'}, x'}\text{ is Cohen-Macaulay}
$$
가 성립한다. $X$가 $k$ 위에서 국소 유한형이면, 모든 체 확대
$k'/k$에 대해서도 같은 명제가 성립한다.
\end{lemma}

\begin{proof}
첫 번째 경우는 Algebra, Lemma
\ref{algebra-lemma-CM-geometrically-CM}에서 따른다.
두 번째 경우는 Algebra, Lemma
\ref{algebra-lemma-extend-field-CM-locus}와 동치이다.
\end{proof}







\section{체의 확대와 야콥슨 성질}
\label{varieties-section-overfield}

\noindent
체 위에서 국소 유한형인 스킴에는 닫힌점이 충분히 많이 있다.
즉 이런 스킴은 야콥슨 스킴이다. 더욱이 그 잉여체들은
기초체의 유한 확대이다.

\begin{lemma}
\label{varieties-lemma-locally-finite-type-Jacobson}
$X$를 $k$ 위에서 국소 유한형인 스킴이라 하자. 그러면
\begin{enumerate}
\item 모든 닫힌점 $x \in X$에 대해 확대 $\kappa(x)/k$는 대수적이다.
\item $X$는 야콥슨 스킴이다
(Properties, Definition \ref{properties-definition-jacobson}).
\end{enumerate}
\end{lemma}

\begin{proof}
스킴이 야콥슨일 필요충분조건은 야콥슨 스킴들로 이루어진 아핀 열린
덮개를 갖는 것이다. Properties, Lemma
\ref{properties-lemma-locally-jacobson}을 보라. 닫힌점에서 잉여체에
관한 성질도 $X$ 위에서 국소적이다. 따라서 $X$가 아핀이라고 가정해도
된다. 이 경우 결과는 힐베르트 영점정리에서 따른다. Algebra, Theorem
\ref{algebra-theorem-nullstellensatz}을 보라. 또한 Morphisms, Lemmas
\ref{morphisms-lemma-jacobson-finite-type-points},
\ref{morphisms-lemma-Jacobson-universally-Jacobson},
\ref{morphisms-lemma-ubiquity-Jacobson-schemes}을 함께 사용해도 얻어진다.
\end{proof}

\noindent
$X$가 국소 유한형이 아니더라도 충분히 큰 기초체 확대를 택하면
같은 결과를 얻을 수 있음이 밝혀진다.

\begin{lemma}
\label{varieties-lemma-make-Jacobson}
$X$를 체 $k$ 위의 스킴이라 하자. 크기가 충분히 큰 모든 체 확대
$K/k$에 대해 다음이 성립한다.
\begin{enumerate}
\item 모든 닫힌점 $x \in X_K$에 대해 확대 $\kappa(x)/K$는 대수적이다.
\item $X_K$는 야콥슨 스킴이다
(Properties, Definition \ref{properties-definition-jacobson}).
\end{enumerate}
\end{lemma}

\begin{proof}
아핀 열린 덮개 $X = \bigcup U_i$를 택하자. Algebra, Lemma
\ref{algebra-lemma-base-change-Jacobson}와 Properties, Lemma
\ref{properties-lemma-affine-jacobson}에 의해 다음 성질을 갖는
기수들 $\kappa_i$가 존재하여 $U_{i, K}$가 $K$ 위에서 원하는
성질들을 갖게 된다. 여기에는 $\#(K) \geq \kappa_i$이면 충분하다.
$\kappa = \max\{\kappa_i\}$로 두자. 그러면 $K$의 기수가
$\kappa$보다 크면 각 $U_{i, K}$가 보조정리의 결론들을 만족한다.
따라서 Properties, Lemma \ref{properties-lemma-locally-jacobson}에 의해
$X_K$는 야콥슨이다. $X_K$의 닫힌점에서 잉여체에 관한 명제는
$U_{i, K}$의 닫힌점에서 잉여체에 관한 대응하는 명제들로부터 따른다.
\end{proof}





\section{체의 확대와 풍부한 가역층}
\label{varieties-section-change-fields-ample}

\noindent
다음 결과는 이 절의 결과들을 대표한다.

\begin{lemma}
\label{varieties-lemma-ample-after-field-extension}
$k$를 체라 하고 $X$를 $k$ 위의 스킴이라 하자. $X_K$ 위에
풍부한 가역층이 존재하게 하는 어떤 체 확대 $K/k$가 있으면,
$X$ 위에도 풍부한 가역층이 존재한다.
\end{lemma}

\begin{proof}
$K/k$를 $X_K$ 위에 풍부한 가역층 $\mathcal{L}$이 존재하는
체 확대라 하자. 사상 $X_K \to X$는 전사이다. 따라서 $X$는
준콤팩트 스킴의 상이므로 준콤팩트이다
(Properties, Definition \ref{properties-definition-ample}).
$X_K$가 준분리이므로(Properties, Lemma
\ref{properties-lemma-affine-s-opens-cover-quasi-separated})
$X$도 준분리이다. 실제로 $U, V \subset X$가 아핀 열린집합이면
$(U \cap V)_K = U_K \cap V_K$는 준콤팩트이고
$(U \cap V)_K \to U \cap V$는 전사이다. 따라서 Schemes, Lemma
\ref{schemes-lemma-characterize-quasi-separated}을 적용할 수 있다.

\medskip\noindent
$K = \colim A_i$로 쓰자. 여기서 이들은 $K$의 부분대수들이고
$k$ 위에서 유한형이다.
$X_i = X \times_{\Spec(k)} \Spec(A_i)$로 두자.
$X_K = \lim X_i$이므로 어떤 $i$와 가역층 $\mathcal{L}_i$를
$X_i$ 위에서 찾아, 이것을 $X_K$로 당긴 것이 $\mathcal{L}$이
되게 할 수 있다(Limits, Lemma
\ref{limits-lemma-descend-invertible-modules}; 여기와 아래에서는 방금 보인
$X$의 준콤팩트성과 준분리성을 사용한다). Limits, Lemma
\ref{limits-lemma-limit-ample}에 의해 $\mathcal{L}_i$가 풍부하다고
가정할 수 있으며, 필요하면 $i$를 키운다. 그런 $i$를 고정하고
극대 아이디얼 $\mathfrak m \subset A_i$를 택하자.
힐베르트 영점정리(Algebra, Theorem
\ref{algebra-theorem-nullstellensatz})에 의해 잉여체
$k' = A_i/\mathfrak m$는 $k$의 유한 확대이다. 따라서
$X_{k'} \subset X_i$는 닫힌 부분스킴이고, 그러므로 풍부한 가역층을
갖는다(Properties, Lemma \ref{properties-lemma-ample-on-closed}).
$X_{k'} \to X$가 유한 국소 자유 사상이므로 Divisors, Proposition
\ref{divisors-proposition-push-down-ample}에 의해 $X$에도
풍부한 가역층이 존재한다.
\end{proof}

\begin{lemma}
\label{varieties-lemma-quasi-affine-after-field-extension}
$k$를 체라 하고 $X$를 $k$ 위의 스킴이라 하자. $X_K$가
준아핀이 되게 하는 어떤 체 확대 $K/k$가 있으면 $X$도 준아핀이다.
\end{lemma}

\begin{proof}
$K/k$를 $X_K$가 준아핀이 되는 체 확대라 하자.
사상 $X_K \to X$는 전사이다. 따라서 $X$는 준콤팩트 스킴의
상이므로 준콤팩트이다
(Properties, Definition \ref{properties-definition-quasi-affine}).
$X_K$는 아핀 스킴의 열린 부분스킴이므로 준분리이고, 따라서
$X$도 준분리이다. 실제로 $U, V \subset X$가 아핀 열린집합이면
$(U \cap V)_K = U_K \cap V_K$는 준콤팩트이고
$(U \cap V)_K \to U \cap V$는 전사이다. 따라서 Schemes, Lemma
\ref{schemes-lemma-characterize-quasi-separated}을 적용할 수 있다.

\medskip\noindent
$K = \colim A_i$로 쓰자. 여기서 이들은 $K$의 부분대수들이고
$k$ 위에서 유한형이다.
$X_i = X \times_{\Spec(k)} \Spec(A_i)$로 두자.
$X_K = \lim X_i$이므로 어떤 $i$를 찾아 $X_i$가 준아핀이 되게 할 수 있다
(Limits, Lemma \ref{limits-lemma-limit-quasi-affine}; 여기서는 방금 보인
$X$의 준콤팩트성과 준분리성을 사용한다).
힐베르트 영점정리(Algebra, Theorem
\ref{algebra-theorem-nullstellensatz})에 의해 잉여체
$k' = A_i/\mathfrak m$는 $k$의 유한 확대이다. 따라서
$X_{k'} \subset X_i$는 닫힌 부분스킴이고, 그러므로 준아핀이다
(Properties, Lemma \ref{properties-lemma-quasi-affine-locally-closed}).
$X_{k'} \to X$가 유한 국소 자유 사상이므로 Divisors, Lemma
\ref{divisors-lemma-push-down-quasi-affine}에 의해 결론을 얻는다.
\end{proof}

\begin{lemma}
\label{varieties-lemma-quasi-projective-after-field-extension}
$k$를 체라 하고 $X$를 $k$ 위의 스킴이라 하자. $X_K$가
$K$ 위에서 준사영이 되게 하는 어떤 체 확대 $K/k$가 있으면,
$X$는 $k$ 위에서 준사영이다.
\end{lemma}

\begin{proof}
정의에 의해 스킴의 사상 $g : Y \to T$가 준사영이라고 함은
그 사상이 국소 유한형이고 준콤팩트이며, $g$-풍부한 가역층이
$Y$ 위에 존재하는 경우이다.
$K/k$를 $X_K$가 $K$ 위에서 준사영이 되는 체 확대라 하자.
아핀 열린집합 $\Spec(A) \subset X$를 택하자. 그러면 $U_K$는
$X_K$의 아핀 열린 부분스킴이고, 따라서 $A_K$는 유한형
$K$-대수이다. 그러면 Algebra, Lemma
\ref{algebra-lemma-finite-type-descends}에 의해 $A$는 유한형
$k$-대수이다. 따라서 $X \to \Spec(k)$는 국소 유한형이다.
$X_K \to \Spec(K)$가 준콤팩트이므로 $X_K$가 준콤팩트이고,
따라서 $X$가 준콤팩트이며, 따라서 $X \to \Spec(k)$는 유한형이다.
Morphisms, Lemma
\ref{morphisms-lemma-finite-type-over-affine-ample-very-ample}에 의해
$X_K$는 풍부한 가역층을 갖는다. 그러면 Lemma
\ref{varieties-lemma-ample-after-field-extension}에 의해 $X$도 풍부한 가역층을 갖는다.
따라서 Morphisms, Lemma
\ref{morphisms-lemma-finite-type-over-affine-ample-very-ample}에 의해
$X \to \Spec(k)$는 준사영이다.
\end{proof}

\noindent
다음 보조정리는 Descent, Lemma
\ref{descent-lemma-descending-property-proper}의 특수한 경우이다.

\begin{lemma}
\label{varieties-lemma-proper-after-field-extension}
$k$를 체라 하고 $X$를 $k$ 위의 스킴이라 하자. $X_K$가
$K$ 위에서 고유가 되게 하는 어떤 체 확대 $K/k$가 있으면,
$X$는 $k$ 위에서 고유이다.
\end{lemma}

\begin{proof}
$K/k$를 $X_K$가 $K$ 위에서 고유인 체 확대라 하자.
그러면 $X_K$는 분리이고 준콤팩트임을 상기하자
(Morphisms, Definition \ref{morphisms-definition-proper}).
사상 $X_K \to X$는 전사이다. 따라서 $X$는 준콤팩트 스킴의
상이므로 준콤팩트이다
(Properties, Definition \ref{properties-definition-ample}).
$X_K$가 분리이므로 $X$는 준분리이다. 실제로 $U, V \subset X$가
아핀 열린집합이면 $(U \cap V)_K = U_K \cap V_K$는 준콤팩트이고
$(U \cap V)_K \to U \cap V$는 전사이다. 따라서 Schemes, Lemma
\ref{schemes-lemma-characterize-quasi-separated}을 적용할 수 있다.

\medskip\noindent
$K = \colim A_i$로 쓰자. 여기서 이들은 $K$의 부분대수들이고
$k$ 위에서 유한형이다.
$X_i = X \times_{\Spec(k)} \Spec(A_i)$로 두자.
Limits, Lemma \ref{limits-lemma-eventually-proper}에 의해 어떤
$i$가 존재하여 $X_i \to \Spec(A_i)$가 고유이다. 여기서는
방금 보인 $X$의 준콤팩트성과 준분리성을 사용한다.
극대 아이디얼 $\mathfrak m \subset A_i$를 택하자.
힐베르트 영점정리(Algebra, Theorem
\ref{algebra-theorem-nullstellensatz})에 의해 잉여체
$k' = A_i/\mathfrak m$는 $k$의 유한 확대이다.
기저 변환 $X_{k'} \to \Spec(k')$는 고유이다
(Morphisms, Lemma \ref{morphisms-lemma-base-change-proper}).
$k'/k$가 유한이므로 $X_{k'} \to X$와 합성
$X_{k'} \to \Spec(k)$도 모두 고유이다
(Morphisms, Lemmas \ref{morphisms-lemma-finite-proper},
\ref{morphisms-lemma-base-change-proper},
\ref{morphisms-lemma-composition-proper}).
첫 번째 사실로부터 $X$가 $k$ 위에서 분리임을 얻는다. 여기에는
$X_{k'}$가 분리임을 사용한다(Morphisms, Lemma
\ref{morphisms-lemma-image-universally-closed-separated}).
두 번째 사실로부터 Morphisms, Lemma
\ref{morphisms-lemma-image-proper-is-proper}에 의해
$X \to \Spec(k)$가 고유임을 얻는다.
\end{proof}

\begin{lemma}
\label{varieties-lemma-projective-after-field-extension}
$k$를 체라 하고 $X$를 $k$ 위의 스킴이라 하자. $X_K$가
$K$ 위에서 사영이 되게 하는 어떤 체 확대 $K/k$가 있으면,
$X$는 $k$ 위에서 사영이다.
\end{lemma}

\begin{proof}
$k$ 위의 스킴이 $k$ 위에서 사영일 필요충분조건은
$k$ 위에서 준사영이고 고유인 것이다. Morphisms, Lemma
\ref{morphisms-lemma-projective-is-quasi-projective-proper}을 보라.
따라서 이 보조정리는 Lemmas
\ref{varieties-lemma-quasi-projective-after-field-extension}와
\ref{varieties-lemma-proper-after-field-extension}에서 따른다.
\end{proof}




\section{접공간}
\label{varieties-section-tangent-spaces}

\noindent
이 절에서는 이중수에 값을 갖는 점들을 사용하여, 스킴의 사상이
정의역의 한 점에서 갖는 접공간을 정의한다.

\begin{definition}
\label{varieties-definition-dual-numbers}
임의의 환 $R$에 대해 $R$ 위의 {\it 이중수}는
$R$-대수이며 $R[\epsilon]$으로 나타낸다. $R$-가군으로서는 기저
$1$, $\epsilon$을 갖는 자유 가군이며, $R$-대수 구조는
$\epsilon^2 = 0$으로 놓아 얻는다.
\end{definition}

\noindent
$f : X \to S$를 스킴의 사상이라 하자. $x \in X$를 점이라 하고,
그 상을 $s = f(x)$라 하자. 이 점은 $S$에 있다.
다음 실선 가환 그림을 생각하자.
\begin{equation}
\label{varieties-equation-tangent-space}
\vcenter{
\xymatrix{
\Spec(\kappa(x)) \ar[r] \ar[dr] \ar@/^1pc/[rr] &
\Spec(\kappa(x)[\epsilon]) \ar@{.>}[r] \ar[d]&
X \ar[d] \\
&
\Spec(\kappa(s)) \ar[r] &
S
}
}
\end{equation}
여기서 굽은 화살표는 $\Spec(\kappa(x))$에서 $X$로 가는 표준 사상이다.

\begin{lemma}
\label{varieties-lemma-tangent-space}
(\ref{varieties-equation-tangent-space})를 가환으로 만드는 점선 화살표들의 집합은
표준적인 $\kappa(x)$-벡터공간 구조를 갖는다.
\end{lemma}

\begin{proof}
$\kappa = \kappa(x)$로 두자. 스킴의 범주에서 다음 밂이 있음을 관찰하자.
$$
\Spec(\kappa[\epsilon]) \amalg_{\Spec(\kappa)} \Spec(\kappa[\epsilon])
= \Spec(\kappa[\epsilon_1, \epsilon_2])
$$
여기서 $\kappa[\epsilon_1, \epsilon_2]$는 $\kappa$-대수이며 기저
$1, \epsilon_1, \epsilon_2$를 갖고
$\epsilon_1^2 = \epsilon_1\epsilon_2 = \epsilon_2^2 = 0$이다.
이는 환에 대한 대응하는 결과와, 국소환의 스펙트럼에서 스킴으로 가는
사상에 대한 Schemes, Lemma
\ref{schemes-lemma-morphism-from-spec-local-ring}의 기술에서 바로 따른다.
두 화살표
$\theta_1, \theta_2 : \Spec(\kappa[\epsilon]) \to X$가 주어지면 사상
$$
\theta_1 + \theta_2 :
\Spec(\kappa[\epsilon]) \to
\Spec(\kappa[\epsilon_1, \epsilon_2])
\xrightarrow{\theta_1, \theta_2} X
$$
을 생각할 수 있다. 여기서 첫 번째 화살표는
$\epsilon_i \mapsto \epsilon$으로 주어진다. 한편
$\lambda \in \kappa$가 주어지면 $\Spec(\kappa[\epsilon])$의 자기 사상이
있는데, 이는 $\kappa$-대수 자기준동형, 구체적으로
$\kappa[\epsilon]$의 자기준동형 가운데
$\epsilon$을 $\lambda \epsilon$으로 보내는 것에 대응한다.
$\theta : \Spec(\kappa[\epsilon]) \to X$를 이 자기 사상과 먼저 합성하면
$\lambda \theta$를 얻는다. 독자는 적절한 스킴의 가환 그림이 존재함을
확인하여 벡터공간 공리들을 검증할 수 있다. 세부사항은 생략한다.
(또 다른 증명은 모든 것을 국소환으로 표현한 뒤 환 사상의 수준에서
벡터공간 공리들을 검증하는 것이다.)
\end{proof}

\begin{definition}
\label{varieties-definition-tangent-space}
$f : X \to S$를 스킴의 사상이라 하고 $x \in X$라 하자.
(\ref{varieties-equation-tangent-space})를 가환으로 만드는 점선 화살표들의 집합에
표준적인 $\kappa(x)$-벡터공간 구조를 준 것을
{\it $X$의 $S$ 위 $x$에서의 접공간}이라 하며 $T_{X/S, x}$로 나타낸다.
이 공간의 원소를 $X/S$의 $x$에서의 {\it 접벡터}라고 한다.
\end{definition}

\noindent
$x \in X$에서의 접벡터는 스킴론적 올 $X_s$ 안에 있다.
이 올은 $f : X \to S$의 $s = f(x)$ 위의 올이므로, 표준적인 동일시
\begin{equation}
\label{varieties-equation-tangent-space-fibre}
T_{X/S, x} = T_{X_s/s, x}
\end{equation}
를 얻는다. 점 함자를 사용하는 이 간결한 정의는 다음 대수적 기술을
갖는다. 이 기술은 {\it $X$의 $S$ 위 $x$에서의 여접공간}을
$\kappa(x)$-벡터공간
$$
T^*_{X/S, x} = \Omega_{X/S, x} \otimes_{\mathcal{O}_{X, x}} \kappa(x)
$$
으로 정의하게 한다. 이는 이 공간이 $\kappa(x)$에 관하여
$X$의 $S$ 위 $x$에서의 접공간과 표준적으로 쌍대이기 때문이다.

\begin{lemma}
\label{varieties-lemma-tangent-space-cotangent-space}
$f : X \to S$를 스킴의 사상이라 하고 $x \in X$라 하자.
표준 동형
$$
T_{X/S, x} = \Hom_{\mathcal{O}_{X, x}}(\Omega_{X/S, x}, \kappa(x))
$$
이 존재하며, 이는 $\kappa(x)$ 위의 벡터공간의 동형이다.
\end{lemma}

\begin{proof}
$\kappa = \kappa(x)$로 두자. $\theta \in T_{X/S, x}$가 주어지면 사상
$$
\theta^*\Omega_{X/S} \to
\Omega_{\Spec(\kappa[\epsilon])/\Spec(\kappa(s))} \to
\Omega_{\Spec(\kappa[\epsilon])/\Spec(\kappa)}
$$
을 얻는다. 절단을 취하면 $\mathcal{O}_{X, x}$-선형 사상
$\xi_\theta : \Omega_{X/S, x} \to \kappa \text{d}\epsilon$,
즉 보조정리 공식의 오른쪽 원소를 얻는다.
$\theta \mapsto \xi_\theta$가 동형임을 보이기 위해 $S$를 $s$로,
$X$를 스킴론적 올 $X_s$로 바꿀 수 있다. 실제로 공식의 양변은
스킴론적 올에만 의존한다. $T_{X/S, x}$에 대해서는 분명하고,
오른쪽에 대해서는 Morphisms, Lemma
\ref{morphisms-lemma-base-change-differentials}를 보라.
또한 $X$를 $\mathcal{O}_{X, x}$의 스펙트럼으로 바꿀 수 있다.
그러더라도 $T_{X/S, x}$는 바뀌지 않고
(Schemes, Lemma \ref{schemes-lemma-morphism-from-spec-local-ring}),
$\Omega_{X/S, x}$도 바뀌지 않는다
(Modules, Lemma \ref{modules-lemma-stalk-module-differentials}).

\medskip\noindent
$(A, \mathfrak m, \kappa)$를 체 $k$ 위의 국소환이라 하자.
증명을 끝내려면 임의의 $A$-선형 사상
$\xi : \Omega_{A/k} \to \kappa$가 유일한 $k$-대수 사상
$\varphi : A \to \kappa[\epsilon]$에서 생김을 보이면 된다.
이 사상은 표준 사상 $c : A \to \kappa$와 $\epsilon$을 법으로 일치한다.
$\varphi(a) = c(a) + D(a) \epsilon$으로 쓰면, 독자는
$a \mapsto D(a)$가 $k$-미분임을 알 수 있다.
$\Omega_{A/k}$의 보편 성질에 의해 각 $D$는 유일한 $\xi$에 대응하고
그 역도 성립한다. 이것으로 증명이 끝난다.
\end{proof}

\begin{lemma}
\label{varieties-lemma-tangent-space-rational-point}
$f : X \to S$를 스킴의 사상이라 하자. $x \in X$를 점이라 하고
$s = f(x) \in S$라 하자. $\kappa(x) = \kappa(s)$라고 가정하자.
그러면 표준 동형
$$
\mathfrak m_x/(\mathfrak m_x^2 + \mathfrak m_s\mathcal{O}_{X, x})
=
\Omega_{X/S, x} \otimes_{\mathcal{O}_{X, x}} \kappa(x)
$$
및
$$
T_{X/S, x} =
\Hom_{\kappa(x)}(
\mathfrak m_x/(\mathfrak m_x^2 + \mathfrak m_s\mathcal{O}_{X, x}),
\kappa(x))
$$
이 존재한다. 이는 더 일반적으로 $\kappa(x)/\kappa(s)$가
분리 대수 확대인 경우에도 성립한다.
\end{lemma}

\begin{proof}
두 번째 동형은 첫 번째 동형과 Lemma
\ref{varieties-lemma-tangent-space-cotangent-space}에서 따른다.
첫 번째 동형에 대해서는 $S$를 $s$로, $X$를 $X_s$로 바꿀 수 있다.
Morphisms, Lemma \ref{morphisms-lemma-base-change-differentials}를 보라.
또한 $X$를 $\mathcal{O}_{X, x}$의 스펙트럼으로 바꿀 수 있다.
Modules, Lemma \ref{modules-lemma-stalk-module-differentials}를 보라.
따라서 다음 대수적 사실을 보이면 된다. $(A, \mathfrak m, \kappa)$를
체 $k$ 위의 국소환이라 하고 $\kappa/k$가 분리 대수 확대라고 하자.
그러면 표준 사상
$$
\mathfrak m/\mathfrak m^2
\longrightarrow
\Omega_{A/k} \otimes \kappa
$$
은 동형이다. $\mathfrak m/\mathfrak m^2 = H_1(\NL_{\kappa/A})$임을
관찰하자. Algebra, Lemma \ref{algebra-lemma-exact-sequence-NL}에 의해
$\Omega_{\kappa/k} = 0$이고 $H_1(\NL_{\kappa/k}) = 0$임을 보이면 충분하다.
$\kappa$는 $k$ 안의 유한 분리 확대들의 합집합이므로,
$\kappa$가 $k$의 유한 분리 확대인 경우만 증명하면 충분하다
(Algebra, Lemma \ref{algebra-lemma-colimits-NL}).
이 경우 환 사상 $k \to \kappa$는 에탈이고, 따라서
$\NL_{\kappa/k} = 0$이다(거의 정의에서 바로 따르며,
Algebra, Section \ref{algebra-section-etale} 참조).
\end{proof}

\begin{lemma}
\label{varieties-lemma-map-tangent-spaces}
$f : X \to Y$를 기초 스킴 $S$ 위의 스킴 사상이라 하자.
$x \in X$를 점이라 하고 $y = f(x)$로 두자.
$\kappa(y) = \kappa(x)$이면 $f$는 자연스러운 선형 사상
$$
\text{d}f : T_{X/S, x} \longrightarrow T_{Y/S, y}
$$
을 유도한다. 이 사상은 선형 사상
$\Omega_{Y/S, y} \otimes \kappa(y) \to \Omega_{X/S, x} \otimes \kappa(x)$과 쌍대이며, 여기에는 Lemma
\ref{varieties-lemma-tangent-space-cotangent-space}의 동일시를 사용한다.
\end{lemma}

\begin{proof}
생략한다.
\end{proof}

\begin{lemma}
\label{varieties-lemma-tangent-space-product}
$X$, $Y$를 기초 $S$ 위의 스킴이라 하자. $x \in X$와 $y \in Y$가
같은 상점 $s \in S$를 가지며, $\kappa(s) = \kappa(x)$ 및
$\kappa(s) = \kappa(y)$가 성립한다고 하자. 표준 동형
$$
T_{X \times_S Y/S, (x, y)} = T_{X/S, x} \oplus T_{Y/S, y}
$$
이 존재한다. 왼쪽에서 오른쪽으로 가는 사상은 사영
$X \times_S Y \to X$와 $X \times_S Y \to Y$가 유도하는 접공간의
사상들로 유도된다. 오른쪽에서 왼쪽으로 가는 사상은 사상
$1 \times y : X_s \to X_s \times_s Y_s$와
$x \times 1 : Y_s \to X_s \times_s Y_s$가 유도한다.
여기서 접공간과 올의 접공간 사이의 동일시
(\ref{varieties-equation-tangent-space-fibre})를 사용한다.
\end{lemma}

\begin{proof}
직합 분해는 Morphisms, Lemma
\ref{morphisms-lemma-differential-product}에서 Lemma
\ref{varieties-lemma-tangent-space-rational-point}를 통해 따른다.
사상들과의 양립성은 Lemma \ref{varieties-lemma-map-tangent-spaces}에서 따른다.
\end{proof}

\begin{lemma}
\label{varieties-lemma-injective-tangent-spaces-unramified}
$f : X \to Y$를 기초 스킴 $S$ 위에서 국소 유한형인 스킴들의
사상이라 하자. $x \in X$를 점이라 하고 $y = f(x)$로 두며
$\kappa(y) = \kappa(x)$라고 가정하자. 다음 조건들은 서로 동치이다.
\begin{enumerate}
\item $\text{d}f : T_{X/S, x} \longrightarrow T_{Y/S, y}$는 단사이다.
\item $f$는 $x$에서 비분기이다.
\end{enumerate}
\end{lemma}

\begin{proof}
사상 $f$는 Morphisms, Lemma
\ref{morphisms-lemma-permanence-finite-type}에 의해 국소 유한형이다.
사상 $\text{d}f$가 단사일 필요충분조건은
$\Omega_{Y/S, y} \otimes \kappa(y) \to \Omega_{X/S, x} \otimes \kappa(x)$가 전사인 것이다
(Lemma \ref{varieties-lemma-map-tangent-spaces}).
완전열 $f^*\Omega_{Y/S} \to \Omega_{X/S} \to \Omega_{X/Y} \to 0$
(Morphisms, Lemma \ref{morphisms-lemma-triangle-differentials})에 의해
이는 $\Omega_{X/Y, x} \otimes \kappa(x) = 0$일 필요충분조건이다.
따라서 결과는 Morphisms, Lemma
\ref{morphisms-lemma-unramified-at-point}에서 따른다.
\end{proof}







\section{일반점에서 유한인 사상}
\label{varieties-section-generically-finite}

\noindent
이 절에서는 Morphisms, Section
\ref{morphisms-section-generically-finite}에서 연구한, 스킴 사이의
일반점에서 유한인 사상이라는 개념을 다시 살펴본다.

\begin{lemma}
\label{varieties-lemma-quasi-finite-in-codim-1}
$f : X \to Y$가 국소 유한형이라고 하자. $y \in Y$를 점이라 하고,
$\mathcal{O}_{Y, y}$가 차원 $\leq 1$인 뇌터환이라고 하자.
또한 다음 조건 가운데 하나가 성립한다고 가정하자.
\begin{enumerate}
\item 모든 일반점 $\eta$ 가운데 $X$의 기약 성분의 일반점인 것에 대해 체 확대
$\kappa(\eta)/\kappa(f(\eta))$가 유한이다(또는 대수적이다).
\item 모든 일반점 $\eta$ 가운데 $X$의 기약 성분의 일반점이고
$f(\eta) \leadsto y$를 만족하는 것에 대해 체 확대
$\kappa(\eta)/\kappa(f(\eta))$가 유한이다(또는 대수적이다).
\item $f$는 $X$의 기약 성분의 모든 일반점에서 준유한이다.
\item $Y$는 국소 뇌터이고 $f$는 $X$의 조밀한 점 집합에서 준유한이다.
\item 여기에 더 추가한다.
\end{enumerate}
그러면 $f$는 $X$의 모든 점 가운데 $y$ 위에 놓이는 점에서 준유한이다.
\end{lemma}

\begin{proof}
조건 (4)는 $X$가 국소 뇌터임을 함의한다
(Morphisms, Lemma \ref{morphisms-lemma-finite-type-noetherian}).
사상이 준유한인 점들의 집합은 열린집합이다
(Morphisms, Lemma \ref{morphisms-lemma-quasi-finite-points-open}).
국소 뇌터 스킴의 조밀한 열린집합은 기약 성분의 모든 일반점을 포함하므로
(4)는 (3)을 함의한다. 조건 (3)은 Morphisms, Lemma
\ref{morphisms-lemma-residue-field-quasi-finite}에 의해 조건 (1)을
함의한다. 조건 (1)은 조건 (2)를 함의한다. 따라서 조건 (2)가
성립하는 경우에 보조정리를 증명하면 충분하다.

\medskip\noindent
(2)가 성립한다고 가정하자. $\Spec(\mathcal{O}_{Y, y})$는
$Y$에서 $y$로 특수화되는 점들의 집합임을 상기하자.
Schemes, Lemma \ref{schemes-lemma-specialize-points}를 보라.
이를 Morphisms, Lemma
\ref{morphisms-lemma-base-change-quasi-finite}와 함께 사용하면
$Y$를 $\Spec(\mathcal{O}_{Y, y})$로 바꿀 수 있다.
따라서 $Y = \Spec(B)$라 가정할 수 있다. 여기서 $B$는 차원
$\leq 1$인 뇌터 국소환이고 $y$는 닫힌점이다.

\medskip\noindent
$X = \bigcup X_i$를 $X$의 기약 성분들을 축약 닫힌 부분스킴으로
본 것이라 하자. 각 올 $X_{i, y}$가 이산 공간임을 보일 수 있다면
$X_y = \bigcup X_{i, y}$도 이산이다. 그러면 Morphisms, Lemma
\ref{morphisms-lemma-quasi-finite-at-point-characterize}에 의해
$X \to Y$는 $X_y$의 모든 점에서 준유한이다.
따라서 $X$가 정역적 스킴이라고 가정할 수 있다.

\medskip\noindent
$X \to Y$가 일반점 $\eta$---곧 $X$의 일반점---를 $y$로 보내면, $X$는
$\kappa(y)$의 유한 확대의 스펙트럼이므로 결과가 성립한다.
$X$가 $\eta$를 극소 소 아이디얼 $\mathfrak q$에 대응하는 점으로
보낸다고 가정하자. 이 아이디얼은 $B$의 소 아이디얼이고
$\mathfrak m_B$와 다르다.
분해 $X \to \Spec(B/\mathfrak q) \to \Spec(B)$를 얻는다.
$x \in X$를 $y$ 위에 놓이는 점이라 하자. 차원 공식
(Morphisms, Lemma \ref{morphisms-lemma-dimension-formula})에 의해
$$
\dim(\mathcal{O}_{X, x}) \leq \dim(B/\mathfrak q) +
\text{trdeg}_{\kappa(\mathfrak q)}(R(X)) - \text{trdeg}_{\kappa(y)} \kappa(x)
$$
를 얻는다. $\dim(B/\mathfrak q) = 1$이고, $X$의 일반점은
$x$와 다르면서 $x$로 특수화되며, $R(X)$는
$\kappa(\mathfrak q)$ 위에서 대수적임을 안다. 따라서
$$
1 \leq 1 - \text{trdeg}_{\kappa(y)} \kappa(x)
$$
를 얻는다. 따라서 모든 점 $x$ 가운데 $X_y$에 속하는 것은 Morphisms, Lemma
\ref{morphisms-lemma-algebraic-residue-field-extension-closed-point-fibre}에
의해 $X_y$에서 닫혀 있다. 그러므로 Morphisms, Lemma
\ref{morphisms-lemma-quasi-finite-at-point-characterize}에 의해
$X \to Y$는 모든 점 $x$ 가운데 $X_y$에 속하는 점에서 준유한이다
(같은 보조정리는 $X_y$가 이산 위상공간임도 함의한다).
\end{proof}

\begin{lemma}
\label{varieties-lemma-finite-in-codim-1}
$f : X \to Y$를 고유 사상이라 하자. $y \in Y$를 점이라 하고
$\mathcal{O}_{Y, y}$가 차원 $\leq 1$인 뇌터환이라고 하자.
또한 다음 조건 가운데 하나가 성립한다고 가정하자.
\begin{enumerate}
\item 모든 일반점 $\eta$ 가운데 $X$의 기약 성분의 일반점인 것에 대해 체 확대
$\kappa(\eta)/\kappa(f(\eta))$가 유한이다(또는 대수적이다).
\item 모든 일반점 $\eta$ 가운데 $X$의 기약 성분의 일반점이고
$f(\eta) \leadsto y$를 만족하는 것에 대해 체 확대
$\kappa(\eta)/\kappa(f(\eta))$가 유한이다(또는 대수적이다).
\item $f$는 $X$의 모든 일반점에서 준유한이다.
\item $Y$는 국소 뇌터이고 $f$는 $X$의 조밀한 점 집합에서 준유한이다.
\item 여기에 더 추가한다.
\end{enumerate}
그러면 열린 근방 $V \subset Y$가 $y$에 대해 존재하여
$f^{-1}(V) \to V$는 유한이다.
\end{lemma}

\begin{proof}
Lemma \ref{varieties-lemma-quasi-finite-in-codim-1}에 의해 사상 $f$는
올 $X_y$의 모든 점에서 준유한이다. 따라서 $X_y$는 이산 위상공간이다
(Morphisms, Lemma \ref{morphisms-lemma-quasi-finite-at-point-characterize}).
$f$가 고유이므로 올 $X_y$는 준콤팩트, 즉 유한이다.
따라서 Cohomology of Schemes, Lemma
\ref{coherent-lemma-proper-finite-fibre-finite-in-neighbourhood}를
적용하여 결론을 얻는다.
\end{proof}

\begin{lemma}
\label{varieties-lemma-modification-normal-iso-over-codimension-1}
$X$를 뇌터 스킴이라 하자. $f : Y \to X$를 스킴의 쌍유리 고유
사상이라 하고 $Y$가 축약이라고 하자. $U \subset X$를 $f$가
동형이 되는 최대 열린집합이라 하자. 그러면 $U$는 다음을 포함한다.
\begin{enumerate}
\item 여차원 $0$을 $X$에서 갖는 모든 점.
\item 모든 $x \in X$ 가운데 여차원 $1$을 $X$에서 갖고
$\mathcal{O}_{X, x}$가 이산 값매김환인 점.
\item 모든 $x \in X$ 가운데 $Y \to X$의 $x$ 위의 올이 유한이고
$\mathcal{O}_{X, x}$가 정규인 점.
\item 모든 $x \in X$ 가운데 $f$가 어떤 $y \in Y$에서 준유한이고,
그 점이 $x$ 위에 놓이며 $\mathcal{O}_{X, x}$가 정규인 점.
\end{enumerate}
\end{lemma}

\begin{proof}
(1)은 Morphisms, Lemma
\ref{morphisms-lemma-birational-isomorphism-over-dense-open}에서 따른다.
(2)는 (3)과 Lemma \ref{varieties-lemma-finite-in-codim-1}에서 따른다
(유한 사상의 올이 유한하다는 사실도 사용한다).

\medskip\noindent
(3)은 (4)와 Morphisms, Lemma \ref{morphisms-lemma-finite-fibre}에서
따르지만, 직접 증명도 제시하겠다. $x \in X$가 (3)과 같다고 하자.
Cohomology of Schemes, Lemma
\ref{coherent-lemma-proper-finite-fibre-finite-in-neighbourhood}에 의해
$f$가 유한이라고 가정할 수 있다. $X$가 아핀이라고 가정할 수도 있다.
그러면 뇌터 아핀 스킴의 유한 쌍유리 사상 $Y \to X$와
$x \in X$를 다루면 된다. 여기서 $\mathcal{O}_{X, x}$는 정규 정역이다.
$\mathcal{O}_{X, x}$가 정역이고 $X$가 뇌터이므로, $X$를 $x$의
정역적 아핀 열린 근방으로 바꿀 수 있다. 이제 $Y \to X$가 쌍유리이고
$Y$가 축약이므로 $Y$도 정역적이다.
$X = \Spec(A)$ 및 $Y = \Spec(B)$로 쓰면, $A \subset B$는 같은
분수체를 갖는 정역들의 유한 포함이다. $\mathfrak p \subset A$가
$x$에 대응하는 소 아이디얼이면, $A_\mathfrak p$가 정규라는 사실로부터
$A_\mathfrak p \subset B_\mathfrak p$가 등식임을 얻는다.
$B$가 유한 $A$-가군이므로 $a \in A$, $a \not \in \mathfrak p$인
어떤 원소가 존재하여 $A_a \to B_a$가 동형이다.

\medskip\noindent
$x \in X$와 $y \in Y$가 (4)와 같다고 하자. $X$를 아핀 열린 근방으로
바꾸어 $X = \Spec(A)$이고 $A \subset \mathcal{O}_{X, x}$라고
가정할 수 있다. Properties, Lemma
\ref{properties-lemma-ring-affine-open-injective-local-ring}을 보라.
그러면 $A$는 정역이고, 따라서 $X$는 정역적이다.
$f$가 쌍유리이고 $Y$가 축약이므로 $Y$도 정역적이다.
환 사상 $\mathcal{O}_{X, x} \to \mathcal{O}_{Y, y}$를 생각하자.
이는 본질적 유한형인 환 사상이고, 잉여체 확대는 유한이며,
$\dim(\mathcal{O}_{Y, y}/\mathfrak m_x\mathcal{O}_{Y, y}) = 0$이다
(이를 확인하려면 Morphisms, Definition
\ref{morphisms-definition-quasi-finite}와 Algebra, Definition
\ref{algebra-definition-quasi-finite}에서 준유한 사상의 정의를 따라가라).
Algebra, Lemma
\ref{algebra-lemma-essentially-finite-type-fibre-dim-zero}에 의해
$\mathcal{O}_{Y, y}$는 유한 $\mathcal{O}_{X, x}$-대수 $B$의 국소화이다.
물론 $B$를 $B$의 $\mathcal{O}_{Y, y}$ 안에서의 상으로 바꾸고,
$B$가 $\mathcal{O}_{Y, y}$와 같은 분수체를 갖는 정역이라고
가정할 수 있다. 그러면 $\mathcal{O}_{X, x} \subset B$는
같은 분수체를 갖는다. 이는 $f$가 쌍유리이기 때문이다.
$\mathcal{O}_{X, x}$가 정규이므로 $\mathcal{O}_{X, x} = B$임을 얻는다
(유한이면 정수적이기 때문이다). 특히
$\mathcal{O}_{X, x} = \mathcal{O}_{Y, y}$이다. Morphisms, Lemma
\ref{morphisms-lemma-morphism-defined-local-ring}에 의해 $X$를
줄여서, $X \to Y$가 $f$의 절단이고 $x$를 $y$로 보내며 주어진
국소환의 동형을 유도한다고 가정할 수 있다.
$X \to Y$는 닫힌 사상이고(Schemes, Lemma
\ref{schemes-lemma-section-immersion}), 반드시 $X$의 일반점을
$Y$의 일반점으로 보내므로 $X \to Y$의 상은 $Y$이다.
따라서 $Y = X$이고, 원하는 바를 증명하였다.
\end{proof}






\section{뇌터 정규화의 변형들}
\label{varieties-section-noether-normalization}

\noindent
뇌터 정규화란 $k$가 체이고 $A$가 유한형 $k$-대수이며 차원이
$d$이면, 유한 단사 $k$-대수 준동형
$k[x_1, \ldots, x_d] \to A$가 존재한다는 명제이다.
Algebra, Lemma \ref{algebra-lemma-Noether-normalization}를 보라.
기하학적으로 이는 $\Spec(A) \to \mathbf{A}^d_k$라는 유한 전사
사상이 $\Spec(k)$ 위에서 존재한다는 뜻이다.

\begin{lemma}
\label{varieties-lemma-noether-normalization}
$f : X \to S$를 스킴의 사상이라 하자. $x \in X$의 상을
$s \in S$라 하자. $V \subset S$를 $s$의 아핀 열린 근방이라 하자.
$f$가 국소 유한형이고 $\dim_x(X_s) = d$이면, 아핀 열린집합
$U \subset X$로서 $x \in U$이고 $f(U) \subset V$인 것과 분해
$$
U \xrightarrow{\pi} \mathbf{A}^d_V \to V
$$
가 $f|_U : U \to V$에 대해 존재하며, 여기서 $\pi$는 준유한이다.
\end{lemma}

\begin{proof}
이는 Algebra, Lemma
\ref{algebra-lemma-quasi-finite-over-polynomial-algebra}에서 따른다.
\end{proof}

\begin{lemma}
\label{varieties-lemma-noether-normalization-affine}
$f : X \to S$를 아핀 스킴의 유한형 사상이라 하자. $s \in S$라 하자.
$\dim(X_s) = d$이면, 분해
$$
X \xrightarrow{\pi} \mathbf{A}^d_S \to S
$$
가 $f$에 대해 존재하며, 올 사이의 사상
$\pi_s : X_s \to \mathbf{A}^d_{\kappa(s)}$는 $s$ 위에서 유한이다.
\end{lemma}

\begin{proof}
$S = \Spec(A)$ 및 $X = \Spec(B)$로 쓰고, $A \to B$를 $f$에 대응하는
환 사상이라 하자. $\mathfrak p \subset A$를 $s$에 대응하는 소
아이디얼이라 하자. 전사 $A[x_1, \ldots, x_r] \to B$를 택할 수 있다.
Algebra, Lemma \ref{algebra-lemma-Noether-normalization}에 의해,
원소 $y_1, \ldots, y_d \in A$가 존재한다. 이들은
$\mathbf{Z}$-부분대수로서 $A$ 안에서 $x_1, \ldots, x_r$로 생성되는
것에 속하고, $A$-대수 준동형 $A[t_1, \ldots, t_d] \to B$ 가운데
$t_i$를 $y_i$로 보내는 것은 유한
$\kappa(\mathfrak p)$-대수 준동형
$\kappa(\mathfrak p)[t_1, \ldots, t_d] \to B \otimes_A \kappa(\mathfrak p)$를
유도한다. 이로써 보조정리가 증명된다.
\end{proof}

\begin{lemma}
\label{varieties-lemma-geometric-structure-unramified}
$f : X \to S$를 스킴의 사상이라 하자. $x \in X$라 하자.
$V = \Spec(A)$를 $f(x)$의 아핀 열린 근방이라 하되 $S$ 안에서
택하자. $f$가 $x$에서 비분기이면, 아핀 열린집합 $U \subset X$로서
$x \in U$이고 $f(U) \subset V$인 것이 존재하여 다음 가환 그림을 얻는다.
$$
\xymatrix{
X \ar[d] & U \ar[l] \ar[rd] \ar[r]^-j &
\Spec(A[t]_{g'}/(g)) \ar[d] \ar[r] &
\Spec(A[t]) = \mathbf{A}^1_V \ar[ld] \\
Y & & V \ar[ll]
}
$$
여기서 $j$는 몰입이고, $g \in A[t]$는 일계수 다항식이며,
$g'$는 $g$의 $t$에 관한 도함수이다. $f$가 $x$에서 에탈이면,
$j$가 열린 몰입이 되도록 그림을 택할 수 있다.
\end{lemma}

\begin{proof}
비분기인 경우는 Algebra, Proposition
\ref{algebra-proposition-unramified-locally-standard}을 기하학적으로
옮긴 것이다. 에탈인 경우에는 Algebra, Proposition
\ref{algebra-proposition-etale-locally-standard}을 옮긴 것이며,
진술들이 조금 다르기는 하지만 동치로 Morphisms, Lemma
\ref{morphisms-lemma-etale-locally-standard-etale}에서도 따른다.
\end{proof}

\begin{lemma}
\label{varieties-lemma-unramfied-over-affine}
$f : X \to S$를 아핀 스킴의 유한형 사상이라 하자.
$x \in X$의 상을 $s \in S$라 하자. 다음과 같이 놓자.
$$
r =
\dim_{\kappa(x)} \Omega_{X/S, x} \otimes_{\mathcal{O}_{X, x}} \kappa(x) =
\dim_{\kappa(x)} \Omega_{X_s/s, x} \otimes_{\mathcal{O}_{X_s, x}} \kappa(x) =
\dim_{\kappa(x)} T_{X/S, x}
$$
그러면 분해
$$
X \xrightarrow{\pi} \mathbf{A}^r_S \to S
$$
가 $f$에 대해 존재하며, $\pi$는 $x$에서 비분기이다.
\end{lemma}

\begin{proof}
Morphisms, Lemma \ref{morphisms-lemma-finite-type-differentials}에 의해
첫 번째 차원은 유한이다. 첫 번째 등식은 올에 제한한
$\Omega_{X/S}$가 미분 가군이라는 사실, 즉 Morphisms, Lemma
\ref{morphisms-lemma-base-change-differentials}에서 따른다.
마지막 등식은 Lemma \ref{varieties-lemma-tangent-space-cotangent-space}에서
따른다. 따라서 이 진술은 의미가 있다.

\medskip\noindent
보조정리를 증명하기 위해 $S = \Spec(A)$ 및 $X = \Spec(B)$로 쓰고
$A \to B$를 $f$에 대응하는 환 사상이라 하자.
$\mathfrak q \subset B$를 $x$에 대응하는 소 아이디얼이라 하자.
$A$-대수의 전사 $A[x_1, \ldots, x_t] \to B$를 택하자.
$\Omega_{B/A}$는 $\text{d}x_1, \ldots, \text{d}x_t$로 생성되므로,
그 상들은 $\Omega_{X/S, x} \otimes_{\mathcal{O}_{X, x}} \kappa(x)$를
$\kappa(x)$-벡터공간으로서 생성한다. 번호를 다시 매긴 뒤에는
$\text{d}x_1, \ldots, \text{d}x_r$의 상들이
$\Omega_{X/S, x} \otimes_{\mathcal{O}_{X, x}} \kappa(x)$의 기저라고
가정할 수 있다. $P = A[x_1, \ldots, x_r] \to B$가
$\mathfrak q$에서 비분기라고 주장한다. 이를 보이려면
$\Omega_{B/P, \mathfrak q} = 0$임을 보이면 충분하다
(Algebra, Lemma \ref{algebra-lemma-unramified}).
$\Omega_{B/P}$는 $\Omega_{B/A}$를
$\text{d}x_1, \ldots, \text{d}x_r$로 생성된 부분가군으로
나눈 몫임에 유의하자. 따라서
$\Omega_{B/P, \mathfrak q} \otimes_{B_\mathfrak q} \kappa(\mathfrak q) = 0$이다.
이는 $x_1, \ldots, x_r$의 선택에 따른다.
나카야마의 보조정리, 더 정확히는 $\Omega_{B/P}$가 유한이므로 적용되는
Algebra, Lemma \ref{algebra-lemma-NAK}의 (2)에 의해
$\Omega_{B/P, \mathfrak q} = 0$이라는 결론을 얻는다.
\end{proof}

\begin{lemma}
\label{varieties-lemma-immersion-into-affine}
$f : X \to S$를 스킴의 사상이라 하자. $x \in X$의 상을
$s \in S$라 하자. $V \subset S$를 $s$의 아핀 열린 근방이라 하자.
$f$가 국소 유한형이고
$$
r =
\dim_{\kappa(x)} \Omega_{X/S, x} \otimes_{\mathcal{O}_{X, x}} \kappa(x) =
\dim_{\kappa(x)} \Omega_{X_s/s, x} \otimes_{\mathcal{O}_{X_s, x}} \kappa(x) =
\dim_{\kappa(x)} T_{X/S, x}
$$
이면 다음이 존재한다.
\begin{enumerate}
\item 아핀 열린집합 $U \subset X$로서 $x \in U$이고
$f(U) \subset V$인 것 및 분해
$$
U \xrightarrow{j} \mathbf{A}^{r + 1}_V \to V
$$
로서 $f|_U$의 분해이고 $j$가 몰입인 것, 또는
\item 아핀 열린집합 $U \subset X$로서 $x \in U$이고
$f(U) \subset V$인 것 및 분해
$$
U \xrightarrow{j} D \to V
$$
로서 $f|_U$의 분해이고 $j$가 닫힌 몰입이며 $D \to V$가
상대차원 $r$인 매끄러운 사상인 것.
\end{enumerate}
\end{lemma}

\begin{proof}
임의의 아핀 열린집합 $U \subset X$로서 $x \in U$이고
$f(U) \subset V$인 것을 택하자. $U \to V$에 Lemma
\ref{varieties-lemma-unramfied-over-affine}를 적용하여 그 보조정리의 진술과 같은
$U \to \mathbf{A}^r_V \to V$를 얻는다.
Lemma \ref{varieties-lemma-geometric-structure-unramified}에 의해 분해
$$
U \xrightarrow{j} D \xrightarrow{j'} \mathbf{A}^{r + 1}_V
\xrightarrow{p} \mathbf{A}^r_V \to V
$$
를 얻는다. 여기서 $j$와 $j'$는 몰입이고, $p$는 사영이며,
$p \circ j'$는 표준 에탈이다. 따라서 특히 (1)과 (2)가 성립한다.
\end{proof}





\section{올의 차원}
\label{varieties-section-dimension-fibres}

\noindent
유한형 사상의 올의 차원은 보통 위로 뛰어오른다는 사실을 이미 보았다.
이 절에서는 여차원 $1$에서는 이 현상이 일어나지 않음을 논한다.
더 일반적으로는 올의 차원이 얼마나 뛰어오를 수 있는지를 논한다.
관련된 결과들은 다음과 같다.
\begin{enumerate}
\item 유한형 사상 $X \to S$에 대해, 조건
$x \in X$ 및 $\dim_x(X_{f(x)}) \leq d$를 만족하는 점들의 집합은
열려 있다. Algebra, Lemma
\ref{algebra-lemma-dimension-fibres-bounded-open-upstairs}와
Morphisms, Lemma
\ref{morphisms-lemma-openness-bounded-dimension-fibres}를 보라.
\item 차원 공식이 성립한다. Algebra, Lemma
\ref{algebra-lemma-dimension-formula}와 Morphisms, Lemma
\ref{morphisms-lemma-dimension-formula}를 보라.
\item 값매김환을 지배하는 정역적 유한형 스킴의 올의 차원은 일정하다.
Algebra, Lemma
\ref{algebra-lemma-finite-type-domain-over-valuation-ring-dim-fibres}를 보라.
\item $X \to S$가 유한형이고 $X$의 모든 일반점에서 준유한이면,
$X \to S$는 여차원 $1$에서 준유한이다. Algebra, Lemma
\ref{algebra-lemma-finite-in-codim-1}과 Lemma
\ref{varieties-lemma-quasi-finite-in-codim-1}을 보라.
\end{enumerate}
위에서 마지막으로 언급한 결과는 다음과 같이 일반화된다.

\begin{lemma}
\label{varieties-lemma-dimension-fibre-in-codim-1}
$f : X \to Y$가 국소 유한형이라고 하자. $x \in X$의 상을
$y \in Y$라 하고, $\mathcal{O}_{Y, y}$가 차원 $\leq 1$인
뇌터환이라고 하자. $d \geq 0$를 다음 조건을 만족하는 정수라 하자.
모든 일반점 $\eta$ 가운데 $X$의 기약 성분 중 $x$를 포함하는 것의
일반점인 것에 대해
$\dim_\eta(X_{f(\eta)}) = d$이다. 그러면 $\dim_x(X_y) = d$이다.
\end{lemma}

\begin{proof}
$\Spec(\mathcal{O}_{Y, y})$는 $Y$에서 $y$로 특수화되는 점들의
집합임을 상기하자. Schemes, Lemma
\ref{schemes-lemma-specialize-points}를 보라.
따라서 $Y$를 $\Spec(\mathcal{O}_{Y, y})$로 바꾸어
$Y = \Spec(B)$라고 가정할 수 있다. 여기서 $B$는 차원
$\leq 1$인 뇌터 국소환이고 $y$는 닫힌점이다.
$X$도 $x$의 아핀 근방으로 바꿀 수 있다.

\medskip\noindent
$X = \bigcup X_i$를 $X$의 기약 성분들을 축약 닫힌 부분스킴으로
본 것이라 하자. 각 올 $X_{i, y}$의 차원이 $d$임을 보일 수 있다면
$X_y = \bigcup X_{i, y}$의 차원도 $d$이다.
따라서 $X$가 정역적 스킴이라고 가정할 수 있다.

\medskip\noindent
$X \to Y$가 일반점 $\eta$ 가운데 $X$의 일반점을 $y$로 보내면, $X$는
$\kappa(y)$ 위의 스킴이므로 가정에 의해 결과가 성립한다.
$X$가 $\eta$를, 점 $\xi \in Y$ 가운데 극소 소 아이디얼
$\mathfrak q$에 대응하고 이 아이디얼이 $B$의 $\mathfrak m_B$와 다른 점으로 보낸다고
가정하자. 분해 $X \to \Spec(B/\mathfrak q) \to \Spec(B)$를 얻는다.
차원 공식(Morphisms, Lemma \ref{morphisms-lemma-dimension-formula})에 의해
$$
\dim(\mathcal{O}_{X, x}) + \text{trdeg}_{\kappa(y)} \kappa(x) \leq
\dim(B/\mathfrak q) + \text{trdeg}_{\kappa(\mathfrak q)}(R(X))
$$
를 얻는다. $\dim(B/\mathfrak q) = 1$이다. 또한
$\text{trdeg}_{\kappa(\mathfrak q)}(R(X)) = d$이다. 이는
$\dim_\eta(X_\xi) = d$라는 가정과 Morphisms, Lemma
\ref{morphisms-lemma-dimension-fibre-at-a-point}에서 따른다.
$\mathcal{O}_{X, x} \to \mathcal{O}_{X_s, x}$는 핵을 갖고
($\eta \mapsto \xi \not = y$이므로), $\mathcal{O}_{X, x}$는
뇌터 정역이므로
$\dim(\mathcal{O}_{X, x}) > \dim(\mathcal{O}_{X_y, x})$이다.
따라서
$$
\dim_x(X_s) =
\dim(\mathcal{O}_{X_s, x}) + \text{trdeg}_{\kappa(y)} \kappa(x) \leq d
$$
라는 결론을 얻는다(Morphisms, Lemma
\ref{morphisms-lemma-dimension-fibre-at-a-point}).
한편 Morphisms, Lemma
\ref{morphisms-lemma-openness-bounded-dimension-fibres}에 의해
$\dim_x(X_s) \geq \dim_\eta(X_{f(\eta)}) = d$이다.
\end{proof}

\begin{lemma}
\label{varieties-lemma-dominate-valuation-ring-dimension-fibres}
$f : X \to \Spec(R)$를 기약 스킴에서 값매김환의 스펙트럼으로 가는
사상이라 하자. $f$가 국소 유한형이고 전사이면, 특수 올은 순수차원이며
그 차원은 일반 올의 차원과 같다.
\end{lemma}

\begin{proof}
$X$를 그 축약으로 바꾸어도 $X$나 특수 올의 차원은 변하지 않으므로
그렇게 바꿀 수 있다. 그러면 $X$는 정역적이고, 보조정리는 Algebra, Lemma
\ref{algebra-lemma-finite-type-domain-over-valuation-ring-dim-fibres}에서
따른다.
\end{proof}

\noindent
다음 보조정리는 Lemma \ref{varieties-lemma-dimension-fibre-in-codim-1}을
일반화한다.

\begin{lemma}
\label{varieties-lemma-dimension-fibre-in-higher-codimension}
$f : X \to Y$가 국소 유한형이라고 하자. $x \in X$의 상을
$y \in Y$라 하고 $\mathcal{O}_{Y, y}$가 뇌터환이라고 하자.
$d \geq 0$를 다음 조건을 만족하는 정수라 하자. 모든 일반점
$\eta$ 가운데 $X$의 기약 성분 중 $x$를 포함하는 것의 일반점인 것에 대해
$f(\eta) \not = y$이고 $\dim_\eta(X_{f(\eta)}) = d$이다.
그러면 $\dim_x(X_y) \leq d + \dim(\mathcal{O}_{Y, y}) - 1$이다.
\end{lemma}

\begin{proof}
Lemma \ref{varieties-lemma-dimension-fibre-in-codim-1}의 증명과 정확히 같은
방식으로 다음 경우로 환원한다. 즉 $X = \Spec(A)$이고 $A$는 정역이며,
$Y = \Spec(B)$이고 $B$는 극대 아이디얼이 $y$에 대응하는 뇌터
국소환인 경우이다. $B$를 $B/\Ker(B \to A)$로 바꾸면
$B$가 정역이고 $B \subset A$라고 가정할 수 있다.
이제 차원 공식(Morphisms, Lemma
\ref{morphisms-lemma-dimension-formula})을 사용하여
$$
\dim(\mathcal{O}_{X, x}) + \text{trdeg}_{\kappa(y)} \kappa(x) \leq
\dim(B) + \text{trdeg}_B(A)
$$
를 얻는다. $\text{trdeg}_B(A) = d$이다. 이는
$\dim_\eta(X_\xi) = d$라는 가정과 Morphisms, Lemma
\ref{morphisms-lemma-dimension-fibre-at-a-point}에서 따른다.
$\mathcal{O}_{X, x} \to \mathcal{O}_{X_y, x}$는 핵을 갖고
($f(\eta) \not = y$이므로), $\mathcal{O}_{X, x}$는 뇌터 정역이므로
$\dim(\mathcal{O}_{X, x}) > \dim(\mathcal{O}_{X_y, x})$이다.
따라서
$$
\dim_x(X_y) =
\dim(\mathcal{O}_{X_y, x}) + \text{trdeg}_{\kappa(y)} \kappa(x)
< \dim(B) + d
$$
를 얻는다. 여기서 등식은 Morphisms, Lemma
\ref{morphisms-lemma-dimension-fibre-at-a-point}에 의한 것이며,
이는 원하는 바를 증명한다.
\end{proof}




\section{대수적 스킴}
\label{varieties-section-algebraic-schemes}

\noindent
다음 정의는 \cite[I Definition 6.4.1]{EGA}에서 가져왔다.

\begin{definition}
\label{varieties-definition-algebraic-scheme}
$k$를 체라 하자. {\it 대수적 $k$-스킴}이란 스킴 $X$ 가운데
$k$ 위에 놓이고 구조 사상 $X \to \Spec(k)$가 유한형인 것을 말한다.
{\it 국소 대수적 $k$-스킴}이란 스킴 $X$ 가운데 $k$ 위에 놓이고
구조 사상 $X \to \Spec(k)$가 국소 유한형인 것을 말한다.
\end{definition}

\noindent
모든 (국소) 대수적 $k$-스킴은 (국소) 뇌터임에 유의하자.
Morphisms, Lemma \ref{morphisms-lemma-finite-type-noetherian}을 보라.
대수적 $k$-스킴의 범주는 모든 곱과 올곱을 갖는다
($k$ 위의 다형체의 범주와는 달리). 국소 대수적 $k$-스킴의 범주도
마찬가지이다.

\begin{lemma}
\label{varieties-lemma-algebraic-scheme-dim-0}
$k$를 체라 하자. $X$를 국소 대수적 $k$-스킴 가운데 차원 $0$인
것이라 하자. 그러면 $X$는 국소 아르틴 $k$-대수 $A$ 가운데
$\dim_k(A) < \infty$인 것들의 스펙트럼의 서로소 합이다.
$X$가 대수적 $k$-스킴 가운데 차원 $0$인 것이면,
더욱이 $X$는 아핀이고 사상 $X \to \Spec(k)$는 유한이다.
\end{lemma}

\begin{proof}
$X$를 국소 대수적 $k$-스킴 가운데 차원 $0$인 것이라 하자.
$U = \Spec(A) \subset X$를 아핀 열린 부분스킴이라 하자.
$\dim(X) = 0$이므로 $\dim(A) = 0$이다.
뇌터 정규화, 즉 Algebra, Lemma
\ref{algebra-lemma-Noether-normalization}에 의해 유한 단사
$k \to A$가 존재한다. 다시 말해 $\dim_k(A) < \infty$이다.
따라서 $A$는 아르틴이다. Algebra, Lemma
\ref{algebra-lemma-finite-dimensional-algebra}를 보라.
그러므로 $A = A_1 \times \ldots \times A_r$는 유한 개의 아르틴
국소환의 곱이다. Algebra, Lemma
\ref{algebra-lemma-artinian-finite-length}를 보라.
$\dim_k(A_i) < \infty$가 각 $i$에 대해 성립함은 물론이다. 실제로 이
차원들의 합은 $\dim_k(A)$와 같다.

\medskip\noindent
위 논증은 $X$가 유한 이산 위상공간들로 이루어진 열린 덮개를 가짐을
보인다. 따라서 $X$는 이산 위상공간이다. 그러므로 $X$는 각각 한 점인
연결 성분들의 서로소 합과 동형이다. 한 점 스킴은 아핀이므로, 위 문단의
결과에 의해 각 한 점 성분은 국소 아르틴 $k$-대수 $A$ 가운데
$\dim_k(A) < \infty$인 것의 스펙트럼이다.

\medskip\noindent
마지막으로 $X$가 대수적 $k$-스킴 가운데 차원 $0$인 것이면
$X$는 준콤팩트이고,
따라서 유한 서로소 합
$X = \Spec(A_1) \amalg \ldots \amalg \Spec(A_r)$이다.
그러므로 아핀이다(Schemes, Lemma
\ref{schemes-lemma-disjoint-union-affines}를 보라).
또한 증명의 첫 문단에서 $X \to \Spec(k)$의 유한성을 이미 보았다.
\end{proof}

\noindent
다음 보조정리는 국소 대수적 스킴의 차원론에 관한 몇 가지 명제를
모아 둔다.

\begin{lemma}
\label{varieties-lemma-dimension-locally-algebraic}
$k$를 체라 하자. $X$를 국소 대수적 $k$-스킴이라 하자.
\begin{enumerate}
\item
\label{varieties-item-catenary}
$X$의 위상공간은 카테너리이다
(Topology, Definition \ref{topology-definition-catenary}).
\item
\label{varieties-item-dimension-at-closed-point}
닫힌점 $x \in X$에 대해
$\dim_x(X) = \dim(\mathcal{O}_{X, x})$이다.
\item
\label{varieties-item-dimension-irreducible}
$X$가 기약이면 $\dim(X) = \dim(U)$가 임의의 공집합이 아닌
열린집합 $U \subset X$에 대해 성립하고,
$\dim(X) = \dim_x(X)$가 임의의 $x \in X$에 대해 성립한다.
\item
\label{varieties-item-irreducible-maximal-chains}
$X$가 기약이면 기약 닫힌 부분집합들의 임의의 사슬은 극대 사슬로
연장되며, 기약 닫힌 부분집합들의 모든 극대 사슬의 길이는
$\dim(X)$와 같다.
\item
\label{varieties-item-dimension-irreducibles-passing-through}
$x \in X$에 대해
$\dim_x(X) = \max \dim(Z) = \min \dim(\mathcal{O}_{X, x'})$이다.
여기서 최댓값은 기약 성분 $Z \subset X$ 가운데 $x$를 포함하는 것
전체에 대해 취하고, 최솟값은 특수화 $x \leadsto x'$ 가운데
$x'$가 $X$에서 닫힌점인 것 전체에 대해 취한다.
\item
\label{varieties-item-dimension-irreducible-trdeg}
$X$가 기약이고 일반점이 $x$이면
$\dim(X) = \text{trdeg}_k(\kappa(x))$이다.
\item
\label{varieties-item-immediate-specialization}
$x \leadsto x'$가 $X$의 점들의 바로 다음 특수화이면
$\text{trdeg}_k(\kappa(x)) = \text{trdeg}_k(\kappa(x')) + 1$이다.
\item
\label{varieties-item-dimension-sup-trdeg}
$X$의 차원은 수들 $\text{trdeg}_k(\kappa(x))$의 상한이다.
여기서 $x$는 $X$의 기약 성분들의 일반점 전체를 달린다.
\item
\label{varieties-item-specialization}
$x \leadsto x'$가 $X$의 점들의 자명하지 않은 특수화이면 다음이 성립한다.
\begin{enumerate}
\item $\dim_x(X) \leq \dim_{x'}(X)$,
\item $\dim(\mathcal{O}_{X, x}) < \dim(\mathcal{O}_{X, x'})$,
\item $\text{trdeg}_k(\kappa(x)) > \text{trdeg}_k(\kappa(x'))$, 그리고
\item 자명하지 않은 특수화들로 이루어진 임의의 극대 사슬
$x = x_0 \leadsto x_1 \leadsto \ldots \leadsto x_n = x'$의 길이는
$n = \text{trdeg}_k(\kappa(x)) - \text{trdeg}_k(\kappa(x'))$이다.
\end{enumerate}
\item
\label{varieties-item-dimension-formula}
$x \in X$에 대해
$\dim_x(X) = \text{trdeg}_k(\kappa(x)) + \dim(\mathcal{O}_{X, x})$이다.
\item
\label{varieties-item-immediate-specialization-local-ring}
$x \leadsto x'$가 $X$의 점들의 바로 다음 특수화이고 $X$가 기약이거나
순수차원이면
$\dim(\mathcal{O}_{X, x'}) = \dim(\mathcal{O}_{X, x}) + 1$이다.
\end{enumerate}
\end{lemma}

\begin{proof}
앞에서 증명한 더 일반적인 결과에 의존하는 대신, 각 명제를 유한형
$k$-대수에 관한 대응 명제로 환원하고 가환대수 장의 결과를 인용하겠다.

\medskip\noindent
(\ref{varieties-item-catenary})의 증명. Topology, Lemma
\ref{topology-lemma-catenary}에 의해 이는 $X$ 위에서 국소적이다.
따라서 $X = \Spec(A)$이고 $A$가 유한형 $k$-대수라고 가정할 수 있다.
$A$가 카테너리임을 보여야 한다
(Algebra, Lemma \ref{algebra-lemma-catenary}).
Algebra, Lemma \ref{algebra-lemma-quotient-catenary}를 사용하여
$k[x_1, \ldots, x_n]$의 경우로 환원한 다음 Algebra, Lemma
\ref{algebra-lemma-dimension-height-polynomial-ring}을 적용할 수 있다.
또는 $k$가 (자명하게) 코언--매콜리이고 코언--매콜리 환은 보편
카테너리이므로 성립한다(Algebra, Lemma
\ref{algebra-lemma-CM-ring-catenary}).

\medskip\noindent
(\ref{varieties-item-dimension-at-closed-point})의 증명. 아핀 근방
$U = \Spec(A)$를 $x$에 대해 택하자. 그러면 $\dim_x(X) = \dim_x(U)$이다.
따라서 아핀인 경우로 환원되며, 이는 Algebra, Lemma
\ref{algebra-lemma-dimension-closed-point-finite-type-field}이다.

\medskip\noindent
(\ref{varieties-item-dimension-irreducible})의 증명. 공집합이 아닌 임의의 두
아핀 열린집합 $U, U' \subset X$가 같은 차원을 가짐을 보이면 충분하다
(기약 부분집합들의 임의의 유한 사슬은 한 아핀 열린집합과 만난다).
$x$를 $X$의 닫힌점 가운데 $x \in U \cap U'$인 것으로 택하자.
$X$가 기약이므로 $U \cap U'$는 공집합이 아니고, Lemma
\ref{varieties-lemma-locally-finite-type-Jacobson}에 의해 $X$가 야콥슨이므로
그러한 닫힌점이 존재한다. Algebra, Lemma
\ref{algebra-lemma-dimension-spell-it-out}에 의해
$\dim(U) = \dim(\mathcal{O}_{X, x}) = \dim(U')$이다
(엄밀하게 말하면 이 보조정리를 적용하기 전에 $X$를 그 축약으로
바꾸어야 한다).

\medskip\noindent
(\ref{varieties-item-irreducible-maximal-chains})의 증명. 기약 닫힌 부분집합들의
사슬이 주어졌을 때, 그 최소 원소와 만나는 아핀 열린집합
$U \subset X$를 찾을 수 있다. 따라서 이 명제는 Algebra, Lemma
\ref{algebra-lemma-dimension-spell-it-out} 및
(\ref{varieties-item-dimension-irreducible})에서 이미 본
$\dim(U) = \dim(X)$에서 따른다.

\medskip\noindent
(\ref{varieties-item-dimension-irreducibles-passing-through})의 증명.
아핀 근방 $U = \Spec(A)$를 $x$에 대해 택하자. 그러면
$\dim_x(X) = \dim_x(U)$이다. 대응 $Z \mapsto Z \cap U$는
$X$의 기약 성분 가운데 $x$를 지나는 것과 $U$의 기약 성분 가운데
$x$를 지나는 것 사이의 전단사이다. 또한
$\dim(Z \cap U) = \dim(Z)$가 그러한 $Z$에 대해 성립한다. 이는
(\ref{varieties-item-dimension-irreducible})에서 따른다.
따라서 명제는 Algebra, Lemma
\ref{algebra-lemma-dimension-at-a-point-finite-type-over-field}에서
따른다.

\medskip\noindent
(\ref{varieties-item-dimension-irreducible-trdeg})의 증명.
(\ref{varieties-item-dimension-irreducible})에 의해 $X = \Spec(A)$가
아핀인 경우로 환원된다. 이 경우 $A_{red}$에 Algebra, Lemma
\ref{algebra-lemma-dimension-prime-polynomial-ring}을 적용하면 된다.

\medskip\noindent
(\ref{varieties-item-immediate-specialization})의 증명.
$Z = \overline{\{x\}} \supset Z' = \overline{\{x'\}}$라 하자.
그러면 (\ref{varieties-item-irreducible-maximal-chains})에 의해
$Z \supset Z'$는 $Z$ 안의 기약 닫힌 부분스킴들의 극대 사슬의
시작이고, 따라서 $\dim(Z) = \dim(Z') + 1$이다.
(\ref{varieties-item-dimension-irreducible-trdeg})에 의해 결론을 얻는다.

\medskip\noindent
(\ref{varieties-item-dimension-sup-trdeg})의 증명. 간단한 위상 논증으로
$\dim(X) = \sup \dim(Z)$임을 알 수 있다. 여기서 상한은 $X$의
기약 성분 전체에 대해 취한다(힌트: Topology, Lemma
\ref{topology-lemma-irreducible}을 사용하라).
따라서 (\ref{varieties-item-dimension-irreducible-trdeg})에서 따른다.

\medskip\noindent
(\ref{varieties-item-specialization})의 증명. (a)는 임의의 열린집합
$U \subset X$ 가운데 $x'$를 포함하는 것이 $x$도 포함한다는 사실에서 따른다.
(b)는 $\mathcal{O}_{X, x}$가 $\mathcal{O}_{X, x'}$의 국소화이므로,
$\mathcal{O}_{X, x}$의 소 아이디얼 사슬은
$\mathcal{O}_{X, x'}$의 소 아이디얼 사슬에 대응하고, 끝에
$\mathfrak m_{x'}$를 붙여 이를 연장할 수 있다는 사실에서 따른다.
(c)와 (d)는 모두 (\ref{varieties-item-immediate-specialization})에서 형식적으로
따른다.

\medskip\noindent
(\ref{varieties-item-dimension-formula})의 증명. 아핀 근방
$U = \Spec(A)$를 $x$에 대해 택하자. 그러면 $\dim_x(X) = \dim_x(U)$이다.
따라서 아핀인 경우로 환원되며, 이는 Algebra, Lemma
\ref{algebra-lemma-dimension-at-a-point-finite-type-field}이다.

\medskip\noindent
(\ref{varieties-item-immediate-specialization-local-ring})의 증명.
$X$가 순수차원이면(Topology, Definition
\ref{topology-definition-equidimensional}), $\dim(X)$는 $X$의
각 기약 성분의 차원과 같다. 따라서
(\ref{varieties-item-dimension-irreducibles-passing-through})에 의해
$\dim_x(X) = \dim(X) = \dim_{x'}(X)$이다.
그러므로 (\ref{varieties-item-immediate-specialization})에서 따른다.
\end{proof}

\begin{lemma}
\label{varieties-lemma-dimension-fibres-locally-algebraic}
$k$를 체라 하자. $f : X \to Y$를 국소 대수적 $k$-스킴들의
사상이라 하자.
\begin{enumerate}
\item $y \in Y$에 대해 올 $X_y$는 $\kappa(y)$ 위의 국소 대수적
스킴이므로 Lemma \ref{varieties-lemma-dimension-locally-algebraic}의 모든
결과가 적용된다.
\item $X$가 기약이라고 가정하자.
$Z = \overline{f(X)}$ 및 $d = \dim(X) - \dim(Z)$로 놓자. 그러면
\begin{enumerate}
\item $\dim_x(X_{f(x)}) \geq d$가 모든 $x \in X$에 대해 성립하고,
\item 조건 $x \in X$ 및 $\dim_x(X_{f(x)}) = d$를 만족하는 점들의
집합은 조밀한 열린집합이고,
\item $\dim(\mathcal{O}_{Z, f(x)}) \geq 1$이면
$\dim_x(X_{f(x)}) \leq d + \dim(\mathcal{O}_{Z, f(x)}) - 1$이고,
\item $\dim(\mathcal{O}_{Z, f(x)}) = 1$이면
$\dim_x(X_{f(x)}) = d$이다.
\end{enumerate}
\item $x \in X$이고 $y = f(x)$이면
$\dim_x(X_y) \geq \dim_x(X) - \dim_y(Y)$이다.
\end{enumerate}
\end{lemma}

\begin{proof}
Morphisms, Lemma \ref{morphisms-lemma-permanence-finite-type}에 의해
사상 $f$는 국소 유한형이다. 따라서 밑변환
$X_y \to \Spec(\kappa(y))$는 국소 유한형이다. 이로써 (1)이 증명된다.
이하 증명에서는 $X$, $Y$, 그리고 $f$의 올에 대해 Lemma
\ref{varieties-lemma-dimension-locally-algebraic}의 결과들을 자유롭게 사용한다.

\medskip\noindent
(2)의 증명. $\eta \in X$를 일반점이라 하고 $\xi = f(\eta)$로 놓자.
그러면 $Z = \overline{\{\xi\}}$이다. 따라서
$$
d = \dim(X) - \dim(Z) =
\text{trdeg}_k \kappa(\eta) - \text{trdeg}_k \kappa(\xi) =
\text{trdeg}_{\kappa(\xi)} \kappa(\eta) = \dim_\eta(X_\xi)
$$
그러므로 (2)(a)와 (2)(b)는 Morphisms, Lemma
\ref{morphisms-lemma-openness-bounded-dimension-fibres}에서 따른다.
(2)(c)와 (2)(d)는 Lemmas
\ref{varieties-lemma-dimension-fibre-in-higher-codimension}와
\ref{varieties-lemma-dimension-fibre-in-codim-1}에서 따른다.

\medskip\noindent
(3)의 증명. $x \in X$라 하자. $X' \subset X$를 $X$의 기약 성분
가운데 $x$를 지나고 차원이 $\dim_x(X)$인 것이라 하자. 그러면 (2)에 의해
$\dim_x(X_y) \geq \dim(X') - \dim(Z')$이다. 여기서
$Z' \subset Y$는 $X'$의 상의 폐포이다. 이로써 (3)이 증명된다.
\end{proof}

\begin{lemma}
\label{varieties-lemma-dimension-product-locally-algebraic}
\begin{slogan}
곱의 차원은 차원들의 합이다.
\end{slogan}
$k$를 체라 하자. $X$, $Y$를 국소 대수적 $k$-스킴이라 하자.
\begin{enumerate}
\item $z \in X \times Y$ 가운데 $(x, y)$ 위에 놓이는 것에 대해
$\dim_z(X \times Y) = \dim_x(X) + \dim_y(Y)$이다.
\item $\dim(X \times Y) = \dim(X) + \dim(Y)$이다.
\end{enumerate}
\end{lemma}

\begin{proof}
(1)의 증명. 구조 사상의 분해
$$
X \times Y \longrightarrow Y \longrightarrow \Spec(k)
$$
를 생각하자. 첫 번째 사상 $p : X \times Y \to Y$는 평탄 사상
$X \to \Spec(k)$의 밑변환이므로 평탄하다. Morphisms, Lemma
\ref{morphisms-lemma-base-change-flat}을 보라. 더욱이 Morphisms, Lemma
\ref{morphisms-lemma-dimension-fibre-after-base-change}에 의해
$\dim_z(p^{-1}(y)) = \dim_x(X)$이다. 따라서 Morphisms, Lemma
\ref{morphisms-lemma-dimension-fibre-at-a-point-additive}에 의해
$\dim_z(X \times Y) = \dim_x(X) + \dim_y(Y)$이다.
(2)는 (1)의 직접적인 귀결이다.
\end{proof}





\section{완비 국소환}
\label{varieties-section-complete-local-rings}

\noindent
체 위의 스킴의 완비 국소환에 관한 몇 가지 결과를 제시한다.

\begin{lemma}
\label{varieties-lemma-complete-local-ring-structure-as-algebra}
$k$를 체라 하자. $X$를 $k$ 위의 국소 뇌터 스킴이라 하자.
$x \in X$를 잉여체가 $\kappa$인 점이라 하자. 잉여체 위에서 항등을
유도하는 동형
\begin{equation}
\label{varieties-equation-complete-local-ring}
\kappa[[x_1, \ldots, x_n]]/I \longrightarrow \mathcal{O}_{X, x}^\wedge
\end{equation}
이 존재한다. 일반적으로 (\ref{varieties-equation-complete-local-ring})을
$k$-대수의 동형으로 택할 수는 없다. 그러나 확대 $\kappa/k$가
분리적이면 (\ref{varieties-equation-complete-local-ring})을 $k$-대수의
동형으로 택할 수 있다.
\end{lemma}

\begin{proof}
동형의 존재는 코언 구조 정리의 직접적인 귀결이다.\footnote{$\kappa$의
표수가 $p$이면, 이 정리는 전사
$\Lambda[[x_1, \ldots, x_n]] \to \mathcal{O}_{X, x}^\wedge$를
얻는다는 것만을 말한다. 여기서 $\Lambda$는 $\kappa$의 코언환이다.
그러나 이 경우 사상은 물론
$\Lambda/p\Lambda[[x_1, \ldots, x_n]]$를 거쳐 분해되고
$\Lambda/p\Lambda = \kappa$이다.}
(Algebra, Theorem \ref{algebra-theorem-cohen-structure-theorem}).

\medskip\noindent
$p$를 홀수 소수라 하고, $k = \mathbf{F}_p(t)$ 및
$A = k[x, y]/(y^2 + x^p - t)$라 하자. 그러면 $A^\wedge$를
$A$의 극대 아이디얼 $\mathfrak m = (y)$에 관한 완비화라 할 때,
이는 환으로서는
$k(t^{1/p})[[z]]$와 동형이지만 $k$-대수로서는 그렇지 않다.
그 이유는 $A^\wedge$가 $p$제곱이 $t$인 원소를 포함하지 않기 때문이다
(독자는 $y^2$을 법으로 계산하여 이를 확인할 수 있다).
이는 또한 임의의 동형 (\ref{varieties-equation-complete-local-ring})이
$k$-대수의 동형일 수 없음을 보인다.

\medskip\noindent
$\kappa/k$가 분리적이면, 잉여체 위에서 항등을 유도하는 $k$-대수
준동형 $\kappa \to \mathcal{O}_{X, x}^\wedge$가 존재한다.
More on Algebra, Lemma \ref{more-algebra-lemma-lift-residue-field}를 보라.
$f_1, \ldots, f_n \in \mathfrak m_x$를 생성원들이라 하자.
사상
$$
\kappa[[x_1, \ldots, x_n]] \longrightarrow \mathcal{O}_{X, x}^\wedge,\quad
x_i \longmapsto f_i
$$
을 생각하자. 양변은 $(x_1, \ldots, x_n)$-진 완비이고
(오른쪽은 Algebra, Lemmas
\ref{algebra-lemma-hathat-finitely-generated}에 의해 그러하다),
구성상 $(x_1, \ldots, x_n)$을 법으로 이 사상이 전사이므로 Algebra,
Lemma \ref{algebra-lemma-completion-generalities}에 의해 이 사상은
전사이다.
\end{proof}

\begin{lemma}
\label{varieties-lemma-base-change-complete-local-ring}
$K/k$를 체의 확대라 하자. $X$를 국소 대수적 $k$-스킴이라 하자.
$Y = X_K$로 놓자. $y \in Y$의 상을 $x \in X$라 하자.
$\dim(\mathcal{O}_{X, x}) = \dim(\mathcal{O}_{Y, y})$이고
$\kappa(x)/k$가 분리적이라고 가정하자.
(\ref{varieties-equation-complete-local-ring})에서와 같은 $k$-대수의 동형
$$
\kappa(x)[[x_1, \ldots, x_n]]/(g_1, \ldots, g_m) \longrightarrow
\mathcal{O}_{X, x}^\wedge
$$
을 택하자. 그러면 (\ref{varieties-equation-complete-local-ring})에서와 같은
$K$-대수의 동형
$$
\kappa(y)[[x_1, \ldots, x_n]]/(g_1, \ldots, g_m) \longrightarrow
\mathcal{O}_{Y, y}^\wedge
$$
을 얻는다. 여기서는 $\kappa(x) \to \kappa(y)$를 사용하여 $g_j$를
$\kappa(y)$ 위의 형식적 멱급수로 본다.
\end{lemma}

\begin{proof}
국소환 사상 $\mathcal{O}_{X, x} \to \mathcal{O}_{Y, y}$는
국소환 사상 $\mathcal{O}_{X, x}^\wedge \to \mathcal{O}_{Y, y}^\wedge$를
유도한다. 유도된 사상
$$
\kappa(x) \to \kappa(x)[[x_1, \ldots, x_n]]/(g_1, \ldots, g_m)
\to \mathcal{O}_{X, x}^\wedge \to \mathcal{O}_{Y, y}^\wedge
$$
을 $\kappa(y)$로의 사영과 합성한 것은 표준 준동형
$\kappa(x) \to \kappa(y)$이다.
Lemma \ref{varieties-lemma-change-fields-flat}에 의해 잉여체 $\kappa(y)$는
$\kappa(x) \otimes_k K$를 핵 $\mathfrak p_0$에서 국소화한 것이다.
여기서 이는 사상 $\kappa(x) \otimes_k K \to \kappa(y)$의 핵이다.
한편 Lemma
\ref{varieties-lemma-change-fields-algebraic-unramified}에 의해 국소환
$(\kappa(x) \otimes_k K)_{\mathfrak p_0}$는 $\kappa(y)$와 같다.
따라서 사상
$$
\kappa(x) \otimes_k K \to \mathcal{O}_{Y, y}^\wedge
$$
은 표준적으로 $\kappa(y)$를 거쳐 분해된다. 다음 가환 그림을 얻는다.
$$
\xymatrix{
\kappa(y) \ar[rr] & & \mathcal{O}_{Y, y}^\wedge \\
\kappa(x) \ar[r] \ar[u] &
\kappa(x)[[x_1, \ldots, x_n]]/(g_1, \ldots, g_m) \ar[r] &
\mathcal{O}_{X, x}^\wedge \ar[u]
}
$$
$f_i \in \mathfrak m_x^\wedge \subset \mathcal{O}_{X, x}^\wedge$를
$x_i$의 상이라 하자. 사상이 전사이므로
$\mathfrak  m_x^\wedge = (f_1, \ldots, f_n)$임에 유의하자.
사상
$$
\kappa(y)[[x_1, \ldots, x_n]] \longrightarrow \mathcal{O}_{Y, y}^\wedge,\quad
x_i \longmapsto f_i
$$
을 생각하자. 여기서 $f_i$는 실제로는 $f_i$의
$\mathfrak m_y^\wedge$ 안에서의 상을 뜻한다. Lemma
\ref{varieties-lemma-change-fields-algebraic-unramified}에 의해
$\mathfrak m_x \mathcal{O}_{Y, y} = \mathfrak m_y$이므로, 오른쪽은
$(x_1, \ldots, x_n)$에 관해 완비이다(그것이 완비 국소환임을 보이려면
Algebra, Lemma \ref{algebra-lemma-hathat-finitely-generated}를 사용하라).
양변이 $(x_1, \ldots, x_n)$-진 완비이고 이 사상이
$(x_1, \ldots, x_n)$을 법으로 전사이므로, Algebra, Lemma
\ref{algebra-lemma-completion-generalities}에 의해 이 사상은 전사이다.
물론 형식적 멱급수 $g_1, \ldots, g_m$은 이 사상 아래에서 영으로
가는데, 이미 $\mathcal{O}_{X, x}^\wedge$에서 영으로 가기 때문이다.
따라서 다음 가환 그림을 얻는다.
$$
\xymatrix{
\kappa(y)[[x_1, \ldots, x_n]]/(g_1, \ldots, g_m) \ar[r] &
\mathcal{O}_{Y, y}^\wedge \\
\kappa(x)[[x_1, \ldots, x_n]]/(g_1, \ldots, g_m) \ar[r] \ar[u] &
\mathcal{O}_{X, x}^\wedge \ar[u]
}
$$
이제 위쪽 수평 화살표가 동형임을 보여야 한다. 이것이 전사임은 이미
안다. $\mathcal{O}_{X, x} \to \mathcal{O}_{Y, y}$가 평탄임도 안다
(Lemma \ref{varieties-lemma-change-fields-flat}). 이는
$\mathcal{O}_{X, x}^\wedge \to \mathcal{O}_{Y, y}^\wedge$가 평탄임을
함의한다(More on Algebra, Lemma
\ref{more-algebra-lemma-flat-completion}).
따라서 Algebra, Lemma \ref{algebra-lemma-mod-injective}를
$R = \kappa(x)[[x_1, \ldots, x_n]]/(g_1, \ldots, g_m)$,
$S = \kappa(y)[[x_1, \ldots, x_n]]/(g_1, \ldots, g_m)$,
$M = \mathcal{O}_{Y, y}^\wedge$, 그리고 $N = S$에 적용하면 사상이
단사라는 결론을 얻는다.
\end{proof}


\section{전역 생성}
\label{varieties-section-global-generation}

\noindent
준연접 가군의 전역 생성과 관련된 몇 가지 보조정리를 제시한다.

\begin{lemma}
\label{varieties-lemma-globally-generated-base-change}
$X \to \Spec(A)$를 스킴의 사상이라 하자. $A \subset A'$를 충실히
평탄한 환 사상이라 하자. $\mathcal{F}$를 준연접
$\mathcal{O}_X$-가군이라 하자. 그러면 $\mathcal{F}$가 전역 생성될
필요충분조건은 밑변환 $\mathcal{F}_{A'}$가 전역 생성되는 것이다.
\end{lemma}

\begin{proof}
더 정확히 $X_{A'} = X \times_{\Spec(A)} \Spec(A')$로 놓자.
$\mathcal{F}_{A'} = p^*\mathcal{F}$로 놓자. 여기서
$p : X_{A'} \to X$는 사영이다. Cohomology of Schemes, Lemma
\ref{coherent-lemma-flat-base-change-cohomology}에 의해
$H^0(X_{k'}, \mathcal{F}_{A'}) = H^0(X, \mathcal{F}) \otimes_A A'$이다.
따라서 $s_i$, $i \in I$가 $H^0(X, \mathcal{F})$의 $A$-가군으로서의
생성원들이면, 그 상들은 $H^0(X_{A'}, \mathcal{F}_{A'})$ 안에서
$H^0(X_{A'}, \mathcal{F}_{A'})$의 $A'$-가군으로서의 생성원들이다.
그러므로 사상
$\alpha : \bigoplus_{i \in I} \mathcal{O}_X \to \mathcal{F}$,
$(f_i) \mapsto \sum f_i s_i$가 전사일 필요충분조건이
$p^*\alpha$가 전사인 것임을 보이면 된다.
아핀 열린집합 $U = \Spec(B)$ 가운데 $X$에 속하는 것 위에서 이를 확인할 수 있다.
그러면 $\mathcal{F}|_U$는 $B$-가군 $M$에 대응하고, $s_i|_U$는
원소 $x_i \in M$에 대응한다. 따라서
$\bigoplus_{i \in I} B \to M$가 전사일 필요충분조건이 밑변환
$\bigoplus_{i \in I} B \otimes_A A' \to M \otimes_A A'$가 전사인
것임을 보이면 된다. 이는 $A \to A'$가 충실히 평탄하므로 참이다.
\end{proof}

\begin{lemma}
\label{varieties-lemma-very-ample-vanish-at-point}
$k$를 무한체라 하자. $X$를 $k$ 위의 유한형 스킴이라 하자.
$\mathcal{L}$을 $X$ 위의 매우 풍부한 가역층이라 하자.
$n \geq 0$이고 $x, x_1, \ldots, x_n \in X$를 점들이라 하되,
$x$는 $k$-유리점, 즉 $\kappa(x) = k$이고
$x \not = x_i$가 $i = 1, \ldots, n$에 대해 성립한다고 하자.
그러면 $s \in H^0(X, \mathcal{L})$가 존재하여 $x$에서는 영이고
$x_i$에서는 영이 아니다.
\end{lemma}

\begin{proof}
$n = 0$이면 결과는 자명하므로 $n > 0$이라고 가정한다.
매우 풍부한 가역층의 정의에 의해 보조정리는
$X = \mathbf{P}^r_k$이고 어떤 $r > 0$에 대해
$\mathcal{L} = \mathcal{O}_X(1)$인 경우로
곧바로 환원된다.
$\mathbf{P}^r_k = \text{Proj}(k[T_0, \ldots, T_r])$로 쓰고,
$V = H^0(X, \mathcal{L}) = kT_0 \oplus \ldots \oplus kT_r$로 놓자.
$x$가 $k$-유리점이므로 $s \in V$ 가운데 $x$에서 영이 되는 것들의 집합은
여차원 $1$인 부분공간 $W \subset V$이고, $W$는 $x$에 대응하는
동차 소 아이디얼을 생성한다. $x_i \not = x$이므로 이에 대응하는
동차 소 아이디얼 $\mathfrak p_i \subset k[T_0, \ldots, T_r]$는
$W$를 포함하지 않는다. $k$가 무한이므로
$W \not = \bigcup W \cap \mathfrak q_i$이고 증명이 끝난다.
\end{proof}

\begin{lemma}
\label{varieties-lemma-generated-by-dim-plus-1-sections}
$k$를 무한체라 하자. $X$를 대수적 $k$-스킴이라 하자.
$\mathcal{L}$을 가역 $\mathcal{O}_X$-가군이라 하자.
$V \to \Gamma(X, \mathcal{L})$를 $k$-벡터공간의 선형사상 가운데
그 상이 $\mathcal{L}$을 생성하는 것이라 하자.
그러면 부분공간 $W \subset V$ 가운데
$\dim_k(W) \leq \dim(X) + 1$을 만족하고 $\mathcal{L}$을
생성하는 것이 존재한다.
\end{lemma}

\begin{proof}
증명 전체에서 모든 $x \in X$에 대해 선형사상
$$
\psi_x : V \to \Gamma(X, \mathcal{L}) \to \mathcal{L}_x \to
\mathcal{L}_x \otimes_{\mathcal{O}_{X, x}} \kappa(x)
$$
이 영이 아니라는 사실을 사용한다. $\dim(X)$에 대한 귀납법으로
증명한다.

\medskip\noindent
기초 단계는 $\dim(X) = 0$인 경우이다. 이때 $X$는 유한 개의 점
$X = \{x_1, \ldots, x_n\}$을 갖는다(예를 들어 Lemma
\ref{varieties-lemma-algebraic-scheme-dim-0}을 보라). $k$가 무한이므로 모든
$v \in V$인 벡터 가운데 $\psi_{x_i}(v) \not = 0$이 모든
$i$에 대해 성립하는 것이 존재한다.
그러면 $W = k\cdot v$가 원하는 조건을 만족한다.

\medskip\noindent
$\dim(X) > 0$이라고 가정하자. $X_i \subset X$를 차원이
$\dim(X)$와 같은 기약 성분들이라 하자. $X$가 뇌터이므로 이러한
성분은 유한 개뿐이다. 각 $i$에 대해 점 $x_i \in X_i$를 택하자.
위와 같이 $v \in V$를 택하여 $\psi_{x_i}(v) \not = 0$이 모든
$i$에 대해 성립하게 하자. $Z \subset X$를 $v$의
$\Gamma(X, \mathcal{L})$ 안에서의 상의 영점 스킴이라 하자.
Divisors, Definition \ref{divisors-definition-zero-scheme-s}를 보라.
구성에 의해 $\dim(Z) < \dim(X)$이다. 귀납가정에 의해 부분공간
$W \subset V$ 가운데 $\dim(W) \leq \dim(X)$인 것을 찾을 수 있다.
이 $W$는 $\mathcal{L}|_Z$를 생성한다. 그러면 $W + k\cdot v$가
$\mathcal{L}$을 생성한다.
\end{proof}





\section{점과 접벡터의 분리}
\label{varieties-section-separating-points-tangent-vectors}

\noindent
이는 바로 다음 결과이다.

\begin{lemma}
\label{varieties-lemma-separate-points-tangent-vectors}
$k$를 대수적으로 닫힌 체라 하자.
$X$를 고유 $k$-스킴이라 하자.
$\mathcal{L}$을 가역 $\mathcal{O}_X$-가군이라 하자.
$V \subset H^0(X, \mathcal{L})$를 $k$-부분벡터공간이라 하자. 다음을 가정하자.
\begin{enumerate}
\item 서로 다른 임의의 두 닫힌점 $x, y \in X$에 대해 단면 $s \in V$가
존재하여, 그것은 $x$에서는 영이지만 $y$에서는 영이 아니고,
\item 임의의 닫힌점 $x \in X$와 영이 아닌 접벡터
$\theta \in T_{X/k, x}$에 대해 단면 $s \in V$가 존재하여,
그것은 $x$에서는 영이지만 $\theta$에 의한 당김은 영이 아니다.
\end{enumerate}
그러면 $\mathcal{L}$은 매우 풍부하며 정준 사상
$\varphi_{\mathcal{L}, V} : X \to \mathbf{P}(V)$는 닫힌 몰입이다.
\end{lemma}

\begin{proof}
조건 (1)은 특히 $V$의 원소들이 $\mathcal{L}$을 $X$ 위에서 생성함을 뜻한다.
따라서 Constructions, Example \ref{constructions-example-projective-space}에 의해
정준 사상
$$
\varphi = \varphi_{\mathcal{L}, V} :
X
\longrightarrow
\mathbf{P}(V)
$$
을 얻는다. 사상 $\varphi$는 Morphisms, Lemma
\ref{morphisms-lemma-image-proper-scheme-closed}에 의해 고유이다.
(1)에 의해 사상 $\varphi$는 닫힌점들 위에서 단사이다
(계산은 생략한다). 특히 $\mathbf{P}(V)$의 임의의 닫힌점 위의 올은
한 점으로 이루어진다(작은 세부사항은 생략한다). 따라서 예컨대
Cohomology of Schemes, Lemma
\ref{coherent-lemma-proper-finite-fibre-finite-in-neighbourhood}를 적용하면
$\varphi$가 유한임을 알 수 있다. 증명을 끝내려면 사상
$$
\varphi^\sharp :
\mathcal{O}_{\mathbf{P}(V)}
\longrightarrow
\varphi_*\mathcal{O}_X
$$
이 전사임을 보이면 충분하다. 이는 닫힌점에서의 줄기 위에서 확인할 수 있다.
$x \in X$를 닫힌점이라 하고 그 상인 닫힌점을
$p = \varphi(x) \in \mathbf{P}(V)$라 하자. (1)에 의해
$\varphi^{-1}(\{p\}) = \{x\}$이고 $\varphi$가 고유이므로
(따라서 닫힌 사상이므로), $\varphi^{-1}(U)$들이 $x$의 열린 근방들의
기본계를 이룬다. 여기서 $U$는 $p$의 열린 근방들의 기본계를 달린다.
따라서 $p$에서의 줄기 위에서 사상
$$
\varphi^\sharp_x :
\mathcal{O}_{\mathbf{P}(V), p}
\longrightarrow
\mathcal{O}_{X, x}
$$
을 얻는다. 특히 $\mathcal{O}_{X, x}$는
$\mathcal{O}_{\mathbf{P}(V), p}$-유한 가군이다.
또한 $x$와 $p$의 잉여체는 모두 $k$와 같다
($k$가 대수적으로 닫혀 있기 때문이다. 힐베르트 영점정리를 사용하라).
마지막으로 조건 (2)에 의해 사상
$$
T_{X/k, x} \longrightarrow T_{\mathbf{P}(V)/k, p}
$$
은 단사이다. 실제로 이 사상의 핵에 속하는 영이 아닌 임의의
$\theta$는 (2)의 결론을 만족할 수 없다. 위의 국소환 사상으로
표현하면 이는
$$
\mathfrak m_p/\mathfrak m_p^2 \longrightarrow
\mathfrak m_x/\mathfrak m_x^2
$$
이 전사라는 뜻이다. Lemma \ref{varieties-lemma-tangent-space-rational-point}를 보라.
이제 Algebra, Lemma \ref{algebra-lemma-when-surjective-local}를 적용하면
증명이 끝난다.
\end{proof}

\begin{lemma}
\label{varieties-lemma-variant-separate-points-tangent-vectors}
$k$를 대수적으로 닫힌 체라 하자.
$X$를 고유 $k$-스킴이라 하자.
$\mathcal{L}$을 가역 $\mathcal{O}_X$-가군이라 하자.
임의의 닫힌 부분스킴 $Z \subset X$ 가운데 차원이 $0$이고 차수가
$2$인 것을 $k$ 위에서 택하면, 사상
$$
H^0(X, \mathcal{L}) \longrightarrow H^0(Z, \mathcal{L}|_Z)
$$
이 전사라고 가정하자. 그러면 $\mathcal{L}$은 $X$에서 $k$ 위로 매우 풍부하다.
\end{lemma}

\begin{proof}
이는 Lemma \ref{varieties-lemma-separate-points-tangent-vectors}를 다시 쓴 것이다.
실제로 서로 다른 닫힌점 $x, y \in X$가 주어지면
$Z = x \cup y$로 두어(닫힌 부분스킴으로 본다) 그 보조정리의
조건 (1)을 얻는다. 또한 영이 아닌 접벡터
$\theta \in T_{X/k, x}$가 주어지면 사상
$\theta : \Spec(k[\epsilon]) \to X$는 닫힌 몰입이다.
$Z = \Im(\theta)$로 두면 그 보조정리의 조건 (2)를 얻는다.
\end{proof}





\section{곱의 폐포}
\label{varieties-section-closure-of-products}

\noindent
폐포와 곱 사이의 관계에 관한 몇 가지 결과를 제시한다.

\begin{lemma}
\label{varieties-lemma-closure-of-product}
$k$를 체라 하자.
$X$, $Y$를 $k$ 위의 스킴이라 하고,
$A \subset X$, $B \subset Y$를 부분집합이라 하자. 다음과 같이 두자.
$$
AB =
\{z \in X \times_k Y \mid \text{pr}_X(z) \in A, \ \text{pr}_Y(z) \in B\}
\subset X \times_k Y
$$
그러면 집합론적으로
$$
\overline{A} \times_k \overline{B} = \overline{AB}
$$
이다.
\end{lemma}

\begin{proof}
포함관계 $\overline{AB} \subset \overline{A} \times_k \overline{B}$는
자명하다. $X$와 $Y$를 각각 축소된 닫힌 부분스킴
$\overline{A}$와 $\overline{B}$로 바꾸어도 된다.
$W \subset X \times_k Y$를 공집합이 아닌 열린 부분집합이라 하자.
Morphisms, Lemma \ref{morphisms-lemma-scheme-over-field-universally-open}에 의해
부분집합 $U = \text{pr}_X(W)$는 $X$에서 공집합이 아닌 열린집합이다.
따라서 $A \cap U$는 공집합이 아니다. $a \in A \cap U$를 택하자.
$Y_a = \{a\} \times_k Y \cong Y_{\kappa(a)}$를
$\text{pr}_X : X \times_k Y \to X$의 $a$ 위의 올이라 쓰자.
다시 Morphisms, Lemma
\ref{morphisms-lemma-scheme-over-field-universally-open}에 의해 사상
$Y_a \to Y$는 열린 사상이다. 실제로
$\Spec(\kappa(a)) \to \Spec(k)$가 보편적으로 열린 사상이기 때문이다.
따라서 공집합이 아닌 열린 부분집합
$W_a = W \times_{X \times_k Y} Y_a$의 상은 $Y$에서 공집합이 아닌
열린 부분집합이다. 그러므로 그 상에 속하는 $b \in B$가 존재한다.
따라서 원하는 대로 $AB \cap W \not = \emptyset$이다.
\end{proof}

\begin{lemma}
\label{varieties-lemma-closure-image-product-map}
$k$를 체라 하자.
$f : A \to X$, $g : B \to Y$를 $k$ 위의 스킴 사상이라 하자.
그러면 집합론적으로
$$
\overline{f(A)} \times_k \overline{g(B)} =
\overline{(f \times g)(A \times_k B)}
$$
이다.
\end{lemma}

\begin{proof}
이는 Lemma \ref{varieties-lemma-closure-of-product}에서 따른다. 실제로 그 보조정리의
표기법으로 $f \times g$의 상은 $f(A)g(B)$이다.
\end{proof}

\begin{lemma}
\label{varieties-lemma-scheme-theoretic-image-product-map}
$k$를 체라 하자.
$f : A \to X$, $g : B \to Y$를 $k$ 위의 스킴 사이의 준콤팩트
사상이라 하자. $Z \subset X$를 $f$의 스킴론적 상이라 하자.
Morphisms, Definition \ref{morphisms-definition-scheme-theoretic-image}를 보라.
마찬가지로 $Z' \subset Y$를 $g$의 스킴론적 상이라 하자.
그러면 $Z \times_k Z'$는 $f \times g$의 스킴론적 상이다.
\end{lemma}

\begin{proof}
$Z$는 $X$의 닫힌 부분스킴 가운데 $f$가 이를 통해 분해되는 가장 작은 것임을
상기하자. $Z'$에 대해서도 마찬가지이다. $W \subset X \times_k Y$를
$f \times g$의 스킴론적 상이라 하자. $f \times g$가
$Z \times_k Z'$를 통해 분해되므로 $W \subset Z \times_k Z'$이다.

\medskip\noindent
반대쪽 포함관계를 보이기 위해 $U \subset X$와 $V \subset Y$를 아핀
열린집합이라 하자. Morphisms, Lemma
\ref{morphisms-lemma-quasi-compact-scheme-theoretic-image}에 의해 스킴
$Z \cap U$는 $f|_{f^{-1}(U)} : f^{-1}(U) \to U$의 스킴론적 상이며,
$Z' \cap V$와 $W \cap U \times_k V$에 대해서도 마찬가지이다.
따라서 $X$와 $Y$가 아핀이라고 가정해도 된다. $f$와 $g$가 준콤팩트이므로
이는 $A = \bigcup U_i$가 아핀들의 유한 합집합이고
$B = \bigcup V_j$가 아핀들의 유한 합집합임을 뜻한다.
그렇다면 $A$를 $\coprod U_i$로, $B$를 $\coprod V_j$로 바꾸어도 된다.
즉 $A$와 $B$도 아핀이라고 가정해도 된다. 이 경우 $Z$는
$\Ker(\Gamma(X, \mathcal{O}_X) \to \Gamma(A, \mathcal{O}_A))$가
정의하는 닫힌 부분스킴이고, $Z'$와 $W$도 마찬가지이다. 따라서 결과는 등식
$$
\Gamma(A \times_k B, \mathcal{O}_{A \times_k B})
=
\Gamma(A, \mathcal{O}_A) \otimes_k \Gamma(B, \mathcal{O}_B)
$$
에서 따른다. 이 등식은 $A$와 $B$가 아핀이므로 성립한다. 세부사항은 생략한다.
\end{proof}









\section{체 위에서 매끄러운 스킴}
\label{varieties-section-smooth}

\noindent
체 위에서 매끄러운 스킴을 특징짓는 두 보조정리를 제시한다.

\begin{lemma}
\label{varieties-lemma-char-zero-differentials-free-smooth}
$k$를 체라 하자. $X$를 $k$ 위의 스킴이라 하자. 다음을 가정하자.
\begin{enumerate}
\item $X$는 $k$ 위에서 국소 유한형이고,
\item $\Omega_{X/k}$는 국소 자유이며,
\item $k$의 표수는 영이다.
\end{enumerate}
그러면 구조 사상 $X \to \Spec(k)$는 매끄럽다.
\end{lemma}

\begin{proof}
이는 Algebra, Lemma
\ref{algebra-lemma-characteristic-zero-local-smooth}에서 따른다.
\end{proof}

\noindent
양의 표수에서는 미분층이 자유인 비축소 유한형 스킴이 존재한다. 예를 들어
$\Spec(\mathbf{F}_p[t]/(t^p))$는 $\Spec(\mathbf{F}_p)$ 위에서 그러하다.
바탕체 $k$가 표수 $p$인 비완전체라면, 축소 스킴 $X/k$ 가운데
$\Omega_{X/k}$가 자유이지만 매끄럽지 않은 것이 존재한다. 예를 들어
$\Spec(k[t]/(t^p-a)$가 그러한데, 여기서 $a \in k$는 $p$제곱이 아니다.

\begin{lemma}
\label{varieties-lemma-char-p-differentials-free-smooth}
$k$를 체라 하자. $X$를 $k$ 위의 스킴이라 하자. 다음을 가정하자.
\begin{enumerate}
\item $X$는 $k$ 위에서 국소 유한형이고,
\item $\Omega_{X/k}$는 국소 자유이며,
\item $X$는 축소이고,
\item $k$는 완전체이다.
\end{enumerate}
그러면 구조 사상 $X \to \Spec(k)$는 매끄럽다.
\end{lemma}

\begin{proof}
$x \in X$를 점이라 하자. $X$는 국소 뇌터 스킴이므로(Morphisms, Lemma
\ref{morphisms-lemma-finite-type-noetherian}을 보라), 유한 개의 기약 성분
$X_1, \ldots, X_n$만이 $x$를 지난다(Properties, Lemma
\ref{properties-lemma-Noetherian-topology} 및 Topology, Lemma
\ref{topology-lemma-Noetherian}을 보라). $\eta_i \in X_i$를 일반점이라 하자.
$X$가 축소이므로 $\mathcal{O}_{X, \eta_i} = \kappa(\eta_i)$이다.
Algebra, Lemma \ref{algebra-lemma-minimal-prime-reduced-ring}을 보라.
더욱이 $\kappa(\eta_i)$는 완전체 $k$의 유한 생성 체 확대이므로
$k$ 위에서 분리적으로 생성된다(Algebra, Section
\ref{algebra-section-separability}을 보라). 따라서
$\Omega_{X/k, \eta_i} = \Omega_{\kappa(\eta_i)/k}$는
$\kappa(\eta_i)$가 $k$ 위에서 갖는 초월차수를 계수로 갖는 자유 가군이다.
Morphisms, Lemma \ref{morphisms-lemma-dimension-fibre-at-a-point}에 의해
$\dim_{\eta_i}(X_i) = \text{rank}_{\eta_i}(\Omega_{X/k})$임을 얻는다.
$x \in X_1 \cap \ldots \cap X_n$이므로
$$
\text{rank}_x(\Omega_{X/k}) = \text{rank}_{\eta_i}(\Omega_{X/k}) = \dim(X_i).
$$
따라서 $\dim_x(X) = \text{rank}_x(\Omega_{X/k})$이다. Algebra, Lemma
\ref{algebra-lemma-dimension-at-a-point-finite-type-over-field}를 보라.
그러므로 예컨대 Algebra, Lemma
\ref{algebra-lemma-characterize-smooth-over-field}에 의해
$X \to \Spec(k)$는 $x$에서 매끄럽다.
\end{proof}

\begin{lemma}
\label{varieties-lemma-smooth-regular}
\begin{slogan}
체 위에서 매끄러우면 정칙이다
\end{slogan}
$X \to \Spec(k)$를 매끄러운 사상이라 하고 $k$를 체라 하자.
그러면 $X$는 정칙 스킴이다.
\end{lemma}

\begin{proof}
(Lemma \ref{varieties-lemma-geometrically-regular-smooth}도 보라.)
Algebra, Lemma \ref{algebra-lemma-characterize-smooth-over-field}에 의해
모든 국소환 $\mathcal{O}_{X, x}$는 정칙이다. 또한 $X$는 $k$ 위에서
국소 유한형이므로 국소 뇌터 스킴이다. 따라서 Properties, Lemma
\ref{properties-lemma-characterize-regular}에 의해 $X$는 정칙이다.
\end{proof}

\begin{lemma}
\label{varieties-lemma-smooth-geometrically-normal}
$X \to \Spec(k)$를 매끄러운 사상이라 하고 $k$를 체라 하자.
그러면 $X$는 $k$ 위에서 기하학적으로 정칙이고, 기하학적으로 정규이며,
기하학적으로 축소이다.
\end{lemma}

\begin{proof}
(Lemma \ref{varieties-lemma-geometrically-regular-smooth}도 보라.)
$k'$를 $k$의 유한 순수 비분리 확대라 하자. 다음 보조정리들에 의해
$X_{k'}$가 정칙이고 정규이며 축소임을 보이면 충분하다.
Lemmas \ref{varieties-lemma-geometrically-regular},
\ref{varieties-lemma-geometrically-normal}, and
\ref{varieties-lemma-check-only-finite-inseparable-extensions}.
Morphisms, Lemma \ref{morphisms-lemma-base-change-smooth}에 의해
사상 $X_{k'} \to \Spec(k')$도 매끄럽다. 따라서 체 위에서 매끄러운 스킴
$X$가 정칙이고 정규이며 축소임을 보이면 충분하다. Lemma
\ref{varieties-lemma-smooth-regular}에 의해 $X$는 정칙이다. 따라서 Properties, Lemma
\ref{properties-lemma-regular-normal}에 의해 $X$는 정규이다.
\end{proof}

\begin{lemma}
\label{varieties-lemma-affine-space-over-field}
$k$를 체라 하자. $d \geq 0$이라 하자. $W \subset \mathbf{A}^d_k$를
공집합이 아닌 열린집합이라 하자. 그러면 닫힌점 $w \in W$ 가운데
$k \subset \kappa(w)$가 유한 분리 확대인 것이 존재한다.
\end{lemma}

\begin{proof}
필요하다면 $W$를 줄여서 $W = \mathbf{A}^d_k \setminus V(f)$라고
가정할 수 있다. 여기서 $f \in k[x_1, \ldots, x_d]$이다. 이 보조정리가
거짓이라면 $f(a_1, \ldots, a_d) = 0$이 모든
$(a_1, \ldots, a_d) \in (k^{sep})^d$에 대해 성립한다. 그러나 $k^{sep}$는 무한체이므로 이는
모순이다.
\end{proof}

\begin{lemma}
\label{varieties-lemma-smooth-separable-closed-points-dense}
$k$를 체라 하자. $X$가 $\Spec(k)$ 위에서 매끄럽다면 집합
$$
\{x \in X\text{ closed such that }k \subset \kappa(x)
\text{ is finite separable}\}
$$
은 $X$에서 조밀하다.
\end{lemma}

\begin{proof}
공집합이 아닌 매끄러운 $X$가 $k$ 위에서 주어졌을 때 잉여체가
$k$의 유한 분리 확대인 닫힌점이 하나 이상 존재함을 보이면 충분하다.
이를 보이기 위해 다음 그림을 택하자.
$$
\xymatrix{
X & U \ar[l] \ar[r]^-\pi & \mathbf{A}^d_k
}
$$
여기서 $\pi$는 \'etale이다. Morphisms, Lemma
\ref{morphisms-lemma-smooth-etale-over-affine-space}를 보라.
사상 $\pi : U \to \mathbf{A}^d_k$는 열린 사상이다. Morphisms, Lemma
\ref{morphisms-lemma-etale-open}을 보라. Lemma
\ref{varieties-lemma-affine-space-over-field}에 의해 닫힌점 $w \in \pi(U)$를 택하여
그 잉여체가 $k$의 유한 분리 확대가 되게 할 수 있다. 임의의
$x \in U$ 가운데 $\pi(x) = w$인 것을 택하자. Morphisms, Lemma
\ref{morphisms-lemma-etale-over-field}에 의해 체 확대
$\kappa(x)/\kappa(w)$는 유한 분리 확대이다. 따라서
$\kappa(x)/k$는 유한 분리 확대이다. Morphisms, Lemma
\ref{morphisms-lemma-algebraic-residue-field-extension-closed-point-fibre}에 의해
점 $x$는 $X$의 닫힌점이다.
\end{proof}

\begin{lemma}
\label{varieties-lemma-geometrically-reduced-dense-smooth-open}
$X$를 체 $k$ 위에서 국소 유한형인 축소 스킴이라 하자.
그러면 $X$가 $k$ 위에서 기하학적으로 축소일 필요충분조건은
$X$가 $k$ 위에서 매끄러운 조밀 열린집합을 포함하는 것이다.
\end{lemma}

\begin{proof}
문제는 $X$ 위에서 국소적이므로 $X$가 준콤팩트라고 가정해도 된다.
$X = X_1 \cup \ldots \cup X_n$을 $X$의 기약 성분들이라 하자.
먼저 $X$가 $k$ 위에서 매끄러운 조밀 열린집합을 포함한다고 가정하자.
그러면 $X$는 각 $X_i$의 일반점에서 기하학적으로 축소이다
(예를 들어 Lemma \ref{varieties-lemma-smooth-geometrically-normal}에 의한다).
Lemma \ref{varieties-lemma-generic-points-geometrically-reduced}에 의해
$X$는 기하학적으로 축소이다.

\medskip\noindent
역으로 $X$가 기하학적으로 축소라고 가정하자.
$Z = \bigcup_{i \not = j} X_i \cap X_j$는 $X$에서 조밀하지 않은 닫힌집합이므로
$X$를 $X \setminus Z$로 바꾸어도 된다. $X \setminus Z$가 기약 스킴들의
분리합이므로 $X$가 기약인 경우로 귀착된다. $X$는 기약이고 축소이므로
정역 스킴이다. Properties, Lemma
\ref{properties-lemma-characterize-integral}을 보라.
$\eta \in X$를 그 일반점이라 하자. 그러면 함수체
$K = k(X) = \kappa(\eta) = \mathcal{O}_{X, \eta}$는 $k$ 위에서
기하학적으로 축소이므로 $k$ 위에서 분리적이다. Algebra, Lemma
\ref{algebra-lemma-characterize-separable-field-extensions}를 보라.
$U = \Spec(A) \subset X$를 공집합이 아닌 임의의 아핀 열린집합이라 하자.
그러면 $K = A_{(0)}$은 $A$의 분수체이다. Algebra, Lemma
\ref{algebra-lemma-separable-smooth}를 적용하면 $A$는 $(0)$에서
$k$ 위로 매끄럽다. 정의에 따라 이는 $A$의 어떤 주국소화가
$k$ 위에서 매끄럽다는 뜻이므로 결론을 얻는다.
\end{proof}

\begin{lemma}
\label{varieties-lemma-dense-smooth-open-variety-over-perfect-field}
$k$를 완전체라 하자. $X$를 국소 대수적인 축소 $k$-스킴, 예를 들어
$k$ 위의 대수다양체라 하자. 그러면
$$
\{x \in X \mid X \to \Spec(k)\text{ is smooth at }x\} =
\{x \in X \mid \mathcal{O}_{X, x}\text{ is regular}\}
$$
이고, 이는 $X$의 조밀 열린 부분스킴이다.
\end{lemma}

\begin{proof}
두 집합의 등식은 Algebra, Lemma
\ref{algebra-lemma-separable-smooth}와 정의에서 곧바로 따른다
(완전체의 정의는 Algebra, Definition \ref{algebra-definition-perfect}을 보라).
스킴 사상이 매끄러운 점들의 집합은 열려 있으므로 이 집합은 열린집합이다.
Morphisms, Definition \ref{morphisms-definition-smooth}를 보라.
마지막으로 이 집합이 조밀함을 보이는 두 가지 논증을 제시한다.
(1) $X$의 일반점에서의 국소환은 체이므로(Algebra, Lemma
\ref{algebra-lemma-minimal-prime-reduced-ring}) 정칙이다. 따라서 일반점들은
이 집합에 속한다. (2) Lemma \ref{varieties-lemma-perfect-reduced}에 의해 $X$가
기하학적으로 축소라는 사실을 사용하고, 따라서 Lemma
\ref{varieties-lemma-geometrically-reduced-dense-smooth-open}를 적용한다.
\end{proof}

\begin{lemma}
\label{varieties-lemma-flat-under-smooth}
$k$를 체라 하자. $f : X \to Y$를 $k$ 위에서 국소 유한형인 스킴들의
사상이라 하자. $x \in X$를 점이라 하고 $y = f(x)$로 두자.
$X \to \Spec(k)$가 $x$에서 매끄럽고 $f$가 $x$에서 평탄하면
$Y \to \Spec(k)$는 $y$에서 매끄럽다. 특히 $X$가 $k$ 위에서 매끄럽고
$f$가 평탄 전사이면 $Y$는 $k$ 위에서 매끄럽다.
\end{lemma}

\begin{proof}
$Y$가 $y$에서 기하학적으로 정칙임을 보이면 충분하다. Lemma
\ref{varieties-lemma-geometrically-regular-smooth}를 보라. 이는 Lemma
\ref{varieties-lemma-flat-under-geometrically-regular}에서 따른다
(그리고 Lemma \ref{varieties-lemma-geometrically-regular-smooth}를
$(X, x)$에 적용한다).
\end{proof}

\begin{lemma}
\label{varieties-lemma-variety-with-smooth-rational-point}
$k$를 체라 하자. $X$를 $k$ 위의 대수다양체라 하고, $k$-유리점 $x$를
가지며 $X$가 $x$에서 매끄럽다고 하자. 그러면 $X$는 $k$ 위에서
기하학적으로 정역이다.
\end{lemma}

\begin{proof}
$U \subset X$를 $X$의 매끄러운 자취라 하자. 가정에 의해 $U$는 공집합이
아니며, 따라서 조밀하고 스킴론적으로도 조밀하다. 그러면
$U_{\overline{k}} \subset X_{\overline{k}}$도 조밀하고 스킴론적으로 조밀하다
(몇몇 세부사항은 생략한다). 따라서 $U$가 기하학적으로 정역임을 보이면 충분하다.
$U$는 $k$-유리점을 가지므로 Lemma
\ref{varieties-lemma-geometrically-connected-if-connected-and-point}에 의해
기하학적으로 연결이다. 한편 $U_{\overline{k}}$는 축소이고 정규이다
(Lemma \ref{varieties-lemma-smooth-geometrically-normal}.
연결이고 정규인 뇌터 스킴은 정역이므로(Properties, Lemma
\ref{properties-lemma-normal-Noetherian}) 증명이 끝난다.
\end{proof}

\begin{lemma}
\label{varieties-lemma-smooth-stratification}
$X$를 체 $k$ 위의 유한형 스킴이라 하자. 유한 순수 비분리 확대
$k'/k$, 정수 $t \geq 0$, 그리고 닫힌 부분스킴들
$$
X_{k'} \supset Z_0 \supset Z_1 \supset \ldots \supset Z_t = \emptyset
$$
이 존재하여 $Z_0 = (X_{k'})_{red}$이고 $Z_i \setminus Z_{i + 1}$은
$k'$ 위에서 매끄럽다. 이는 모든 $i$에 대해 성립한다.
\end{lemma}

\begin{proof}
$\dim(X)$에 관한 귀납법을 사용할 수 있다. Lemma
\ref{varieties-lemma-finite-extension-geometrically-reduced}에 의해 유한 순수 비분리 확대
$k'/k$를 찾아서 $(X_{k'})_{red}$가 $k'$ 위에서 기하학적으로 축소가 되게
할 수 있다. Lemma \ref{varieties-lemma-geometrically-reduced-dense-smooth-open}에 의해
어디에서도 조밀하지 않은 닫힌 부분스킴 $X' \subset (X_{k'})_{red}$가
존재하여 $(X_{k'})_{red} \setminus X'$는 $k'$ 위에서 매끄럽다.
그러면 $\dim(X') < \dim(X)$이다. 귀납가정에 의해 유한 순수 비분리 확대
$k''/k'$, 정수 $t' \geq 0$, 그리고 닫힌 부분스킴들
$$
X'_{k''} \supset Y_0 \supset Y_1 \supset \ldots \supset Y_{t'} = \emptyset
$$
이 존재하여 $Y_0 = (X'_{k''})_{red}$이고 $Y_i \setminus Y_{i + 1}$은
$k''$ 위에서 매끄럽다. 이는 모든 $i$에 대해 성립한다. 이제 $t = t' + 1$로 두고
다음 닫힌 부분스킴들을 생각하자.
$$
X_{k''} \supset Z_0 \supset Z_1 \supset \ldots \supset Z_t = \emptyset
$$
여기서 $Z_0 = (X_{k''})_{red}$이고 $Z_i = Y_{i - 1}$로 둔다
($i > 0$). $X'_{k''}$가 $X_{k''}$의 닫힌 부분스킴이므로
이는 잘 정의된다. 명시한 모든 성질이 성립한다는 검증은 생략한다.
\end{proof}

\section{대수다양체의 유형}
\label{varieties-section-types}

\noindent
대수다양체의 몇 가지 초보적인 전역 성질을 논하는 짧은 절이다.

\begin{definition}
\label{varieties-definition-variety-type}
$k$를 체라 하자. $X$를 $k$ 위의 대수다양체라 하자.
\begin{enumerate}
\item $X$를 {\it 아핀 대수다양체}라 하는 것은 $X$가 아핀 스킴인
경우이다. 이는 $X$가 $\mathbf{A}^n_k$의 닫힌 부분스킴과 동형이라고
요구하는 것과 동치이며, 여기서 $n$은 어떤 정수이다.
\item $X$를 {\it 사영 대수다양체}라 하는 것은 구조 사상
$X \to \Spec(k)$가 사영적인 경우이다. Morphisms, Lemma
\ref{morphisms-lemma-characterize-locally-projective}에 의해 이는 $X$가
$\mathbf{P}^n_k$의 닫힌 부분스킴과 동형일 필요충분조건이며, 여기서
$n$은 어떤 정수이다.
\item $X$를 {\it 준사영 대수다양체}라 하는 것은 구조 사상
$X \to \Spec(k)$가 준사영적인 경우이다. Morphisms, Lemma
\ref{morphisms-lemma-characterize-locally-quasi-projective}에 의해 이는 $X$가
$\mathbf{P}^n_k$의 국소 닫힌 부분스킴과 동형일 필요충분조건이며, 여기서
$n$은 어떤 정수이다.
\item {\it 고유 대수다양체}란 사상 $X \to \Spec(k)$가 고유인
대수다양체이다.
\item {\it 매끄러운 대수다양체}란 사상 $X \to \Spec(k)$가 매끄러운
대수다양체이다.
\end{enumerate}
\end{definition}

\noindent
사영 대수다양체는 고유 대수다양체임에 주의하라. Morphisms, Lemma
\ref{morphisms-lemma-locally-projective-proper}를 보라. 또한
$\mathbf{A}^n_k$는 $\mathbf{P}^n_k$의 열린 부분스킴과 동형이므로 아핀
대수다양체는 준사영적이다. Constructions, Lemma
\ref{constructions-lemma-standard-covering-projective-space}를 보라.

\begin{lemma}
\label{varieties-lemma-regular-functions-proper-variety}
$X$를 $k$ 위의 고유 대수다양체라 하자. 그러면 다음이 성립한다.
\begin{enumerate}
\item $K = H^0(X, \mathcal{O}_X)$는 체이고 체 $k$의 유한 확대이다.
\item $X$가 기하학적으로 축소이면 $K/k$는 분리 확대이다.
\item $X$가 기하학적으로 기약이면 $K/k$는 순수 비분리 확대이다.
\item $X$가 기하학적으로 정역이면 $K = k$이다.
\end{enumerate}
\end{lemma}

\begin{proof}
이는 Lemma \ref{varieties-lemma-proper-geometrically-reduced-global-sections}의
특수한 경우이다.
\end{proof}

\section{정규화}
\label{varieties-section-normalization}

\noindent
정규화에 수반되는 몇 가지 문제를 다룬다.

\begin{lemma}
\label{varieties-lemma-normalization-locally-algebraic}
$k$를 체라 하자. $X$를 $k$ 위의 국소 대수적 스킴이라 하자.
$\nu : X^\nu \to X$를 정규화 사상이라 하자. Morphisms, Definition
\ref{morphisms-definition-normalization}을 보라. 그러면 다음이 성립한다.
\begin{enumerate}
\item $\nu$는 유한이고 우세이며, $X^\nu$는 $k$ 위의 정규이고 기약인
국소 대수적 스킴들의 분리합이다.
\item $\nu$는 $X^\nu \to X_{red} \to X$로 분해되고 첫 번째 사상은
$X_{red}$의 정규화 사상이다.
\item $X$가 축소 대수적 스킴이면 $\nu$는 쌍유리이다.
\item $X$가 대수다양체이면 $X^\nu$도 대수다양체이고, $\nu$는
대수다양체 사이의 유한 쌍유리 사상이다.
\end{enumerate}
\end{lemma}

\begin{proof}
$X$는 체 위에서 국소 유한형이므로 $X$는 국소 뇌터 스킴이다
(Morphisms, Lemma \ref{morphisms-lemma-finite-type-noetherian}).
따라서 각 준콤팩트 열린집합은 유한 개의 기약 성분만을 갖는다
(Properties, Lemma
\ref{properties-lemma-Noetherian-irreducible-components}).
그러므로 Morphisms, Definition
\ref{morphisms-definition-normalization}을 적용할 수 있다.
정규화 $X^\nu$는 언제나 정규 정역 스킴들의 분리합이고, 정규화 사상
$\nu$는 언제나 우세이다. Morphisms, Lemma
\ref{morphisms-lemma-normalization-normal}을 보라. $X$가 보편 나가타
스킴이므로(Morphisms, Lemma \ref{morphisms-lemma-ubiquity-nagata}),
$\nu$는 유한이다(Morphisms, Lemma
\ref{morphisms-lemma-nagata-normalization}). 따라서 $X^\nu$도 국소
대수적이다. 이로써 (1)을 증명했다.

\medskip\noindent
(2)는 Morphisms, Lemma \ref{morphisms-lemma-normalization-reduced}이다.

\medskip\noindent
(3)은 Morphisms, Lemma \ref{morphisms-lemma-normalization-birational}이다.

\medskip\noindent
(4)는 (1), (2), (3), 그리고 $X^\nu$가 분리 스킴 위의 유한 스킴이므로
분리라는 사실에서 따른다.
\end{proof}

\begin{lemma}
\label{varieties-lemma-normalize-projective}
$k$를 체라 하자. $X$를 $k$ 위의 고유 스킴이라 하자.
$\nu : X^\nu \to X$를 정규화 사상이라 하자. Morphisms, Definition
\ref{morphisms-definition-normalization}을 보라. 그러면 $X^\nu$는
$k$ 위에서 고유이다. $X$가 $k$ 위에서 사영적이면 $X^\nu$도
$k$ 위에서 사영적이다.
\end{lemma}

\begin{proof}
Lemma \ref{varieties-lemma-normalization-locally-algebraic}에 의해 사상 $\nu$는
유한이다. 따라서 Morphisms, Lemmas
\ref{morphisms-lemma-finite-proper} 및
\ref{morphisms-lemma-composition-proper}에 의해 $X^\nu$는 $k$ 위에서
고유이다. 사상 $\nu$는 Morphisms, Lemma
\ref{morphisms-lemma-finite-projective}에 의해 사영적이다. 그러므로
$X$가 $k$ 위에서 사영적이면 Morphisms, Lemma
\ref{morphisms-lemma-composition-projective}에 의해 $X^\nu$도 $k$ 위에서
사영적이다.
\end{proof}

\begin{lemma}
\label{varieties-lemma-relative-normalization-finite}
$k$를 체라 하자. $f : Y \to X$를 $k$ 위의 국소 대수적 스킴 사이의
준콤팩트 사상이라 하자. $X'$를 $X$의 $Y$ 안에서의 정규화라 하자.
$Y$가 축소이면 $X' \to X$는 유한이다.
\end{lemma}

\begin{proof}
$Y$는 준분리이므로(Properties, Lemma
\ref{properties-lemma-locally-Noetherian-quasi-separated} 및 Morphisms, Lemma
\ref{morphisms-lemma-finite-type-noetherian}에 의해), 사상 $f$는
준분리이다(Schemes, Lemma
\ref{schemes-lemma-compose-after-separated}). 따라서 Morphisms, Definition
\ref{morphisms-definition-normalization-X-in-Y}을 적용할 수 있다.
결과는 Morphisms, Lemma
\ref{morphisms-lemma-nagata-normalization-finite-general}에서 따른다.
여기서는 국소 대수적 스킴이 국소 뇌터 스킴이고(따라서 국소적으로 유한 개의
기약 성분을 가지며), 국소 대수적 스킴이 나가타 스킴이라는 사실을 사용한다
(Morphisms, Lemma \ref{morphisms-lemma-ubiquity-nagata}).
몇 가지 작은 세부사항은 생략한다.
\end{proof}

\begin{lemma}
\label{varieties-lemma-finite-extension-geometrically-normal}
$k$를 체라 하자. $X$를 대수적 $k$-스킴이라 하자. 유한 순수 비분리 확대
$k'/k$가 존재하여 정규화 $Y$, 곧 $X_{k'}$의 정규화는 $k'$ 위에서 기하학적으로
정규이다.
\end{lemma}

\begin{proof}
$K = k^{perf}$를 완전 폐포라 하자. $Y_K$를 $X_K$의 정규화라 하자.
Lemma \ref{varieties-lemma-normalization-locally-algebraic}를 보라. Limits, Lemma
\ref{limits-lemma-descend-finite-presentation}에 의해 유한 부분확대
$K/k'/k$와 유한 표시 사상 $\nu : Y \to X_{k'}$가 존재하여, 그 $K$로의
밑변환은 정규화 사상 $\nu_K : Y_K \to X_K$가 된다.
$Y$가 $k'$ 위에서 기하학적으로 정규임에 주의하라
(Lemma \ref{varieties-lemma-geometrically-normal}). $k'$를 더 크게 잡아서
$Y \to X_{k'}$가 유한이라고 가정해도 된다(Limits, Lemma
\ref{limits-lemma-descend-finite-finite-presentation}).
$\nu_K : Y_K \to X_K$가 정규화 사상이므로 쌍유리 사상
$Y_K \to (X_K)_{red}$를 유도한다. 따라서 조밀 열린집합
$V_K \subset X_K$가 존재하여 $\nu_K^{-1}(V_K) \to V_K$는 닫힌
몰입이다(이는 $\nu_K^{-1}(V_K)$와 $V_{K, red}$ 사이의 동형을 유도한다.
예를 들어 Morphisms, Lemma
\ref{morphisms-lemma-birational-isomorphism-over-dense-open}를 보라).
$k'$를 더 크게 잡으면 $V_K$가 조밀 열린집합 $V \subset Y$의 밑변환이고,
사상 $\nu^{-1}(V) \to V$가 닫힌 몰입이 되게 할 수 있다
(Limits, Lemmas \ref{limits-lemma-descend-opens} 및
\ref{limits-lemma-descend-closed-immersion-finite-presentation}).
이로부터 $\nu$가 정규화 사상임이 곧바로 따라 증명이 끝난다.
\end{proof}

\begin{lemma}
\label{varieties-lemma-normalization-and-change-of-fields}
$k$를 체라 하자. $X$를 국소 대수적 $k$-스킴이라 하자. $K/k$를 체의
확대라 하자. $\nu : X^\nu \to X$를 $X$의 정규화라 하고
$Y^\nu \to X_K$를 밑변환의 정규화라 하자. 그러면 정준 사상
$$
Y^\nu \longrightarrow X^\nu \times_{\Spec(k)} \Spec(K)
$$
은 $K/k$가 분리 확대이면 동형이고, 일반적으로는 보편 위상동형이다.
\end{lemma}

\begin{proof}
$Y = X_K$로 두자. $X^{(0)}$, 각각 $Y^{(0)}$를 $X$, 각각 $Y$의 기약
성분들의 일반점 집합이라 하자. 그러면 사영 사상
$\pi : Y \to X$는 $\pi(Y^{(0)}) = X^{(0)}$을 만족한다. 이는
$\pi$가 전사이고 열려 있으며 일반화 보존 사상이기 때문이다. Morphisms,
Lemmas \ref{morphisms-lemma-scheme-over-field-universally-open} 및
\ref{morphisms-lemma-open-generizing}을 보라.
$X^{(0)}$, 각각 $Y^{(0)}$를 (축소) 스킴으로 보면 $X^\nu$, 각각
$Y^\nu$는 $X$, 각각 $Y$의 $X^{(0)}$, 각각 $Y^{(0})$ 안에서의
정규화이다. 따라서 Morphisms, Lemma
\ref{morphisms-lemma-functoriality-normalization}는 정준 사상
$Y^\nu \to X^\nu$를 주며, 이는 $Y \to X$ 위에 있다. 다시 올곱의 보편 성질에 의해
보조정리의 정준 사상을 준다.

\medskip\noindent
이 사상이 보조정리에서 명시한 성질을 가짐을 보이기 위해
$X = \Spec(A)$가 아핀이라고 가정해도 된다. $Q(A_{red})$를
$A_{red}$의 전분수환이라 하자. 그러면 $X^\nu$는 정수적 폐포 $A'$, 곧
$A$의 $Q(A_{red})$ 안에서의 정수적 폐포의 스펙트럼이다. Morphisms,
Lemmas \ref{morphisms-lemma-normalization-reduced} 및
\ref{morphisms-lemma-description-normalization}을 보라.
마찬가지로 $Y^\nu$는 정수적 폐포 $B'$, 곧 $A \otimes_k K$의
$Q((A \otimes_k K)_{red})$ 안에서의 정수적 폐포의 스펙트럼이다.
정준 사상 $Q(A_{red}) \to Q((A \otimes_k K)_{red})$와 정준 사상
$A' \to B'$가 있으며, 보조정리의 사상은 유도된 $K$-대수 사상
$$
A' \otimes_k K \longrightarrow B'
$$
에 대응한다. 그 핵은 멱영 원소들로 이루어진다. 실제로
$Q(A_{red}) \otimes_k K \to Q((A \otimes_k K)_{red})$의 핵이
멱영 원소들의 집합이기 때문이다.

\medskip\noindent
$K/k$가 분리 확대이면 Lemma
\ref{varieties-lemma-base-change-normal-by-separable}에 의해 $A' \otimes_k K$는
정규이다. 특히 이는 축소이므로
$Q((A \otimes_k K)_{red}) = Q(A' \otimes_k K)$이고, Algebra, Lemma
\ref{algebra-lemma-characterize-reduced-ring-normal}에 의해
$B' = A' \otimes_k K$이다.

\medskip\noindent
$K/k$가 분리 확대가 아니라고 가정하자. 그러면 $k$의 표수는
$p > 0$이다. 모든 $b \in B'$에 대해 $q$가 존재하고, 이는 $p$의
거듭제곱이며 $b^q$가 $A' \otimes_k K$의 상에 속함을 보이겠다. 그러면 Algebra,
Lemma \ref{algebra-lemma-p-ring-map}에 의해 위에 표시한 사상이 보편
위상동형임이 증명된다. 주어진 $b$에 대해 부분체 $F \subset K$가 존재하여
$F/k$는 유한 생성이고, $b$는 $Q((A \otimes_k F)_{red})$에 속하며
$A \otimes_k F$ 위에서 정수적이다. 일계수 다항식
$P = T^d + \alpha_1 T^{d - 1} + \ldots + \alpha _d$를 택하여
$P(b) = 0$이고 $\alpha_i \in A \otimes_k F$가 되게 하자.
$t_1, \ldots, t_r$를 $F$의 $k$ 위의 초월기저로 택하자.
$F/F'/k(t_1, \ldots, t_r)$를 최대 분리 부분확대라 하자(Fields,
Lemma \ref{fields-lemma-separable-first}). $F/F'$는 유한 순수 비분리
확대이므로 어떤 $q$가 존재하여 $\lambda^q \in F'$가 모든
$\lambda \in F$에 대해 성립한다. 그러면 $b^q$는
$Q((A \otimes_k F')_{red})$에 속하고, 다항식
$T^d + \alpha_1^q T^{d - 1} + \ldots + \alpha _d^q$을 만족하며,
$\alpha_i^q \in A \otimes_k F'$이다. 분리 확대인 경우에 의해
$b^q \in A' \otimes_k F'$이고 증명이 끝난다.
\end{proof}

\begin{lemma}
\label{varieties-lemma-geometrically-normal-in-codim-1}
$k$를 체라 하자. $X$를 국소 대수적 $k$-스킴이라 하자.
$\nu : X^\nu \to X$를 $X$의 정규화라 하자. 점 $x \in X$가 다음을
만족한다고 하자. (a) $\mathcal{O}_{X, x}$는 축소이고,
(b) $\dim(\mathcal{O}_{X, x}) = 1$이며, (c) 모든
$x' \in X^\nu$ 가운데 $\nu(x') = x$인 것에 대해 확대
$\kappa(x')/k$는 분리 확대이다.
그러면 $X$는 $x$에서 기하학적으로 축소이고, $X^\nu$는
$x'$에서 기하학적으로 정칙이다. 여기서 $\nu(x') = x$이다.
\end{lemma}

\begin{proof}
Lemma \ref{varieties-lemma-normalization-locally-algebraic}의 결과들을 더 언급하지 않고
사용하겠다. $x' \in X^\nu$를 $x$ 위의 점이라 하자. 차원론
(Section \ref{varieties-section-algebraic-schemes})에 의해
$\dim(\mathcal{O}_{X^\nu, x'}) = 1$이다. $X^\nu$가 정규이므로
$\mathcal{O}_{X^\nu, x'}$는 이산 값매김환이다(Properties, Lemma
\ref{properties-lemma-criterion-normal}). 따라서
$\mathcal{O}_{X^\nu, x'}$는 잉여체가 $k$ 위에서 분리적인 정칙 국소
$k$-대수이다. 그러므로 $k \to \mathcal{O}_{X^\nu, x'}$는
$\mathfrak m_{x'}$-진 위상에서 형식적으로 매끄럽다. More on Algebra,
Lemma \ref{more-algebra-lemma-regular-implies-fs}를 보라. 이어서 More on
Algebra, Theorem \ref{more-algebra-theorem-regular-fs}에 의해
$\mathcal{O}_{X^\nu, x'}$는 $k$ 위에서 기하학적으로 정칙이다.
따라서 Lemma \ref{varieties-lemma-geometrically-regular-at-point}에 의해
$X^\nu$는 $x'$에서 기하학적으로 정칙이다.

\medskip\noindent
$\mathcal{O}_{X, x}$가 축소이므로 사상들의 족
$\mathcal{O}_{X, x} \to \mathcal{O}_{X^\nu, x'}$는 단사이다.
$\mathcal{O}_{X^\nu, x'}$가 기하학적으로 축소인 $k$-대수이므로
$\mathcal{O}_{X, x}$도 곧바로 기하학적으로 축소인 $k$-대수이다.
따라서 Lemma \ref{varieties-lemma-geometrically-reduced-at-point}에 의해
$X$는 $x$에서 기하학적으로 축소이다.
\end{proof}

\section{가역 함수들의 군}
\label{varieties-section-units}

\noindent
군 $\mathcal{O}^*(X)/k^*$는 $X$가 $k$ 위의 대수다양체일 때 유한 생성
아벨 군인 경우가 흔하지만 항상 그런 것은 아니다. 일련의 보조정리를 통해
이를 보인다. 모든 것은 다음 특수한 경우에 기초한다.

\begin{lemma}
\label{varieties-lemma-open-in-normal-proper}
$k$를 대수적으로 닫힌 체라 하자. $\overline{X}$를 $k$ 위의 고유
대수다양체라 하자. $X \subset \overline{X}$를 열린 부분스킴이라 하자.
$X$가 정규라고 가정하자. 그러면 $\mathcal{O}^*(X)/k^*$는 유한 생성
아벨 군이다.
\end{lemma}

\begin{proof}
명제는 $X$에만 관한 것이므로 $\overline{X}$를 $k$ 위의 다른 고유
대수다양체로 바꾸어도 된다. $\nu : \overline{X}^\nu \to
\overline{X}$를 정규화 사상이라 하자. Lemma
\ref{varieties-lemma-normalization-locally-algebraic}에 의해 $\nu$는 유한이고
$\overline{X}^\nu$는 대수다양체이다. $X$가 정규이므로
$\nu^{-1}(X) \to X$는 동형이다(아주 작은 세부사항은 생략한다).
마지막으로 $\overline{X}^\nu$는 $k$ 위에서 고유이다. 실제로 유한 사상은
고유이고(Morphisms, Lemma \ref{morphisms-lemma-finite-proper}), 고유
사상들의 합성은 고유이다(Morphisms, Lemma
\ref{morphisms-lemma-composition-proper}). 따라서
$\overline{X}$가 정규라고 가정하며, 이제부터 그렇게 하겠다.

\medskip\noindent
더 언급하지 않고 다음 사실을 사용하겠다. 임의의 아핀 열린집합 $U$가
$\overline{X}$에 들어 있을 때 환 $\mathcal{O}(U)$는 유한 생성 $k$-대수이고,
이는 뇌터 정규 정역이다. 다음을 보라: Algebra, Lemma
\ref{algebra-lemma-Noetherian-permanence}, Properties, Definition
\ref{properties-definition-integral}, Properties, Lemmas
\ref{properties-lemma-locally-Noetherian} 및
\ref{properties-lemma-locally-normal}, Morphisms, Lemma
\ref{morphisms-lemma-locally-finite-type-characterize}.

\medskip\noindent
$\xi_1, \ldots, \xi_r$를 $X$의 $\overline{X}$ 안에서의 여집합의 일반점들이라
하자. $\overline{X}$의 바탕 위상공간이 뇌터이므로 이들은 유한 개이다
(Morphisms, Lemma \ref{morphisms-lemma-finite-type-noetherian},
Properties, Lemma \ref{properties-lemma-Noetherian-topology}, 그리고
Topology, Lemma \ref{topology-lemma-Noetherian}을 보라).
각 $i$에 대해 국소환 $\mathcal{O}_i = \mathcal{O}_{X, \xi_i}$는
뇌터 정규 정역의 국소화이므로 뇌터 정규 국소 정역이다.
$J \subset \{1, \ldots, r\}$를 지표 $i$ 가운데
$\dim(\mathcal{O}_i) = 1$인 것들의 집합이라 하자. $j \in J$에 대해 국소환 $\mathcal{O}_j$는
이산 값매김환이다. Algebra, Lemma
\ref{algebra-lemma-characterize-dvr}을 보라. 따라서 값매김
$$
v_j : k(\overline{X})^* \longrightarrow \mathbf{Z}
$$
을 얻으며, 이는 $v_j(f) \geq 0 \Leftrightarrow f \in \mathcal{O}_j$라는
성질을 갖는다.

\medskip\noindent
$\mathcal{O}(X)$를 부분 $k$-대수로서 $k(X) = k(\overline{X})$ 안에 있다고 생각하자.
사상
$$
\mathcal{O}(X)^* \longrightarrow
\prod\nolimits_{j \in J} \mathbf{Z},
\quad
f \longmapsto \prod v_j(f)
$$
의 핵이 $k^*$라고 주장한다. 이 주장이 보조정리를 증명함은 분명하다.
실제로 $f \in \mathcal{O}(X)$를 그 핵의 원소라 하자. 임의의 아핀
열린집합 $U = \Spec(B) \subset \overline{X}$를 택하자. 그러면 $B$는
뇌터 정규 정역이다. 높이 일인 임의의 소아이디얼 $\mathfrak q \subset B$와
그에 대응하는 점 $\xi \in X$에 대해, $\xi = \xi_j$인 $j \in J$가
존재하거나 $\xi \in X$이다. 그 이유는 Properties, Lemma
\ref{properties-lemma-codimension-local-ring}에 의해
$\text{codim}(\overline{\{\xi\}}, \overline{X}) = 1$이고, 따라서
$\xi \in \overline{X} \setminus X$이면 이는
$\overline{X} \setminus X$의 일반점이어야 하므로 어떤 $\xi_j$와 같고,
$j \in J$이기 때문이다. 두 경우 모두
$f \in \mathcal{O}_{X, \xi} = B_{\mathfrak q}$이다. 실제로 $f$가
사상의 핵에 속하기 때문이다. 따라서
$f \in \bigcap_{\text{ht}(\mathfrak q) = 1} B_{\mathfrak q} = B$이다.
Algebra, Lemma
\ref{algebra-lemma-normal-domain-intersection-localizations-height-1}을 보라.
바꾸어 말하면 $f \in \Gamma(\overline{X}, \mathcal{O}_{\overline{X}})$이다.
그러나 $k$가 대수적으로 닫혀 있으므로 Lemma
\ref{varieties-lemma-regular-functions-proper-variety}에 의해 $f \in k$이다.
\end{proof}

\noindent
이제 여전히 체가 대수적으로 닫혀 있다고 두고, 몇 가지 초보적인 논증으로
위의 경우를 일반화한다.

\begin{lemma}
\label{varieties-lemma-units-integral-finite-type-algebraically-closed}
$k$를 대수적으로 닫힌 체라 하자. $X$를 $k$ 위에서 국소 유한형인 정역
스킴이라 하자. 그러면 $\mathcal{O}^*(X)/k^*$는 유한 생성 아벨 군이다.
\end{lemma}

\begin{proof}
$X$가 정역이므로 제한 사상 $\mathcal{O}(X) \to \mathcal{O}(U)$는
임의의 공집합이 아닌 열린 부분스킴 $U \subset X$에 대해
단사이다. 따라서 $X$가 아핀이라고 가정해도 된다. 닫힌 몰입
$X \to \mathbf{A}^n_k$를 택하고, $\overline{X}$를 $X$의
$\mathbf{P}^n_k$ 안에서의 폐포라 쓰자. 여기서는 보통의 몰입
$\mathbf{A}^n_k \to \mathbf{P}^n_k$를 사용한다. 따라서
$X$가 사영 대수다양체 $\overline{X}$의 아핀 열린집합이라고 가정해도 된다.

\medskip\noindent
$\nu : \overline{X}^\nu \to \overline{X}$를 정규화 사상이라 하자.
Morphisms, Definition \ref{morphisms-definition-normalization}을 보라.
$\nu$가 유한이고 우세이며, $\overline{X}^\nu$가 정규 기약 스킴임을
알고 있다. Morphisms, Lemmas
\ref{morphisms-lemma-normalization-normal},
\ref{morphisms-lemma-Japanese-normalization}, 및
\ref{morphisms-lemma-ubiquity-nagata}를 보라.
$\overline{X}^\nu$는 고유 대수다양체이다. 실제로
$\overline{X} \to \Spec(k)$는 유한 사상과 고유 사상의 합성이므로 고유이다
(Morphisms, Sections \ref{morphisms-section-proper} 및
\ref{morphisms-section-integral}의 결과를 보라).
또한 $\nu$는 전사이다. 실제로 $\nu$의 상은 닫혀 있고
$\overline{X}$의 일반점을 포함한다. 따라서 $X^\nu = \nu^{-1}(X)$로 두면 $X^\nu$에 대해
결과를 증명하는 것으로 충분하다. 바꾸어 말하면 $X$가 정규 고유
대수다양체 $\overline{X}$의 공집합이 아닌 열린집합이라고 가정해도 된다.
이 경우는 Lemma \ref{varieties-lemma-open-in-normal-proper}에서 다루었다.
\end{proof}

\noindent
앞의 보조정리에서 다음과 같은 작은 일반화가 나온다.

\begin{lemma}
\label{varieties-lemma-units-general-algebraically-closed}
$k$를 대수적으로 닫힌 체라 하자. $X$를 $k$ 위에서 국소 유한형이고,
연결이며 축소이고, 기약 성분이 유한 개인 스킴이라 하자.
그러면 $\mathcal{O}^*(X)/k^*$는 유한 생성 아벨 군이다.
\end{lemma}

\begin{proof}
$X = \bigcup X_i$를 기약 성분들로 분해하자. Lemma
\ref{varieties-lemma-units-integral-finite-type-algebraically-closed}에 의해
$\mathcal{O}(X_i)^*/k^*$는 유한 생성 아벨 군이다. $f \in \mathcal{O}(X)^*$를
사상
$$
\mathcal{O}(X)^* \longrightarrow \prod \mathcal{O}(X_i)^*/k^*.
$$
의 핵에 속하는 원소라 하자. 그러면 각 $i$에 대해 원소
$\lambda_i \in k$가 존재하여 $f|_{X_i} = \lambda_i$가 된다.
$X_i \cap X_j$에 제한하면 $\lambda_i = \lambda_j$임을 얻는다. 단,
$X_i \cap X_j \not = \emptyset$인 경우이다. $X$가 연결이므로 모든
$\lambda_i$가 일치하고, 따라서 $f \in k^*$이다. 이로써
$$
\mathcal{O}(X)^*/k^* \subset \prod \mathcal{O}(X_i)^*/k^*
$$
임을 증명했다. 오른쪽은 유한 개의 유한 생성 아벨 군들의 곱이므로
보조정리가 따른다.
\end{proof}

\begin{lemma}
\label{varieties-lemma-integral-closure-ground-field}
$k$를 체라 하자. $X$를 연결이고 축소인 $k$ 위의 스킴이라 하자.
그러면 $k$의 $\Gamma(X, \mathcal{O}_X)$ 안에서의 정수적 폐포는 체이다.
\end{lemma}

\begin{proof}
$k' \subset \Gamma(X, \mathcal{O}_X)$를 $k$의 정수적 폐포라 하자.
그러면 $X \to \Spec(k)$는 $\Spec(k')$를 통해 분해된다. Schemes, Lemma
\ref{schemes-lemma-morphism-into-affine}을 보라. $X$가 축소이므로 $k'$에는
영이 아닌 멱영 원소가 없다. $k \to k'$가 정수적이므로 $k'$의 모든
소아이디얼은 극대아이디얼이면서 극소소아이디얼이고, $\Spec(k')$는 완전
비연결이다. Algebra, Lemmas
\ref{algebra-lemma-integral-no-inclusion} 및
\ref{algebra-lemma-ring-with-only-minimal-primes}를 보라.
$X$가 연결이므로 사상 $X \to \Spec(k')$는 상수이며, 그 상을
$\mathfrak p \subset k'$에 대응하는 점이라 하자. 그러면 임의의
$f \in k'$, $f \not \in \mathfrak p$는 $\mathcal{O}_X$의 가역 원소로
보내진다. $k'$의 정의에 의해 이는 다시 $f$가 $k'$의 단원임을 강제한다.
따라서 $k'$는 극대아이디얼 $\mathfrak p$를 갖는 국소환이다. Algebra,
Lemma \ref{algebra-lemma-characterize-local-ring}을 보라. 이미 $k'$가
축소임을 보았으므로 $k'$는 체이다. Algebra, Lemma
\ref{algebra-lemma-minimal-prime-reduced-ring}을 보라.
\end{proof}

\begin{proposition}
\label{varieties-proposition-units-general}
$k$를 체라 하자. $X$를 $k$ 위의 스킴이라 하자. $X$가 $k$ 위에서 국소
유한형이고, 연결이며 축소이고, 기약 성분이 유한 개라고 가정하자.
위 조건들에 더하여 다음 조건 가운데 하나 이상이 성립하면
$\mathcal{O}(X)^*/k^*$는 유한 생성 아벨 군이다.
\begin{enumerate}
\item $k$의 $\Gamma(X, \mathcal{O}_X)$ 안에서의 정수적 폐포가 $k$이다.
\item $X$는 $k$-유리점을 갖는다.
\item $X$는 기하학적으로 정역이다.
\end{enumerate}
\end{proposition}

\begin{proof}
$\overline{k}$를 $k$의 대수적 폐포라 하자. $Y$를
$(X_{\overline{k}})_{red}$의 연결 성분이라 하자. 정준 사상
$p : Y \to X$는 열려 있고(Morphisms, Lemma
\ref{morphisms-lemma-scheme-over-field-universally-open}에 의해)
닫혀 있음에 주의하라(Morphisms, Lemma
\ref{morphisms-lemma-integral-universally-closed}에 의해).
따라서 $p(Y) = X$이다. 실제로 $X$는 연결이라고 가정했다. 특히 $X$가
축소이므로 이는 $\mathcal{O}(X) \subset \mathcal{O}(Y)$를 뜻한다.
Lemma \ref{varieties-lemma-galois-action-irreducible-components-locally-finite-type}에
의해 $Y$는 기약 성분을 유한 개만 갖는다(그 보조정리의 체
$\overline{k}$는 분리 대수적 폐포이지만, 두 밑변환
$X_{\overline{k}}$가 Lemma
\ref{varieties-lemma-separably-closed-field-irreducible-components}에 의해 같은 기약
성분들을 가지므로 이는 중요하지 않다). 따라서 Lemma
\ref{varieties-lemma-units-general-algebraically-closed}를 $Y$에 적용할 수 있다.
이에 따라 $\mathcal{O}(X)^*/k^*$가 유한 생성 아벨 군이 아니라면,
$f \in \mathcal{O}(X)$이고 $f \not \in k$인 원소가 존재하여
$\overline{k}$의 원소로 보내진다. 여기서 사상
$\mathcal{O}(X) \to \mathcal{O}(Y)$을 사용한다. 이 경우 $f$는
$k$ 위에서 대수적이므로 $k$ 위에서 정수적이다.
따라서 조건 (1)이 성립하면 이런 일은 일어날 수 없다. 증명을 끝내기 위해
조건 (2)와 (3)이 (1)을 함의함을 보인다.

\medskip\noindent
$k \subset k' \subset \Gamma(X, \mathcal{O}_X)$를
$k$의 $\Gamma(X, \mathcal{O}_X)$ 안에서의 정수적 폐포라 하자.
Lemma \ref{varieties-lemma-integral-closure-ground-field}에 의해 $k'$는 체이다.
$e : \Spec(k) \to X$가 $k$-유리점이면
$e^\sharp : \Gamma(X, \mathcal{O}_X) \to k$는 포함 사상
$k \to \Gamma(X, \mathcal{O}_X)$의 절단이다. 특히 $e^\sharp$의
$k'$로의 제한은 체 사상 $k' \to k$이며, 이는 $k$ 위의 사상이다.
따라서 (2)가 (1)을
함의함을 분명히 보여 준다.

\medskip\noindent
$k'$를 $k$의 $\Gamma(X, \mathcal{O}_X)$ 안에서의 정수적 폐포라 하자.
이것이 자명하지 않다면, $X$가 기하학적으로 연결이 아니거나
($k'/k$가 순수 비분리 확대가 아닌 경우), $X$가 기하학적으로 축소가
아니다($k'/k$가 자명하지 않은 순수 비분리 확대인 경우). 세부사항은
생략한다. 따라서 (3)은 (1)을 함의한다.
\end{proof}

\begin{lemma}
\label{varieties-lemma-units-variety}
$k$를 체라 하자. $X$를 $k$ 위의 대수다양체라 하자. 다음 조건 가운데
하나 이상이 성립하면 군 $\mathcal{O}(X)^*/k^*$는 유한 생성 아벨 군이다.
\begin{enumerate}
\item $k$는 $\Gamma(X, \mathcal{O}_X)$ 안에서 정수적으로 닫혀 있다.
\item $k$는 $k(X)$ 안에서 대수적으로 닫혀 있다.
\item $X$는 $k$ 위에서 기하학적으로 정역이다.
\item $k$는 체 확대 $\kappa(x)/k$들의 ‘교집합’이다. 여기서 $x$는
$x$의 닫힌점들을 달린다.
\end{enumerate}
\end{lemma}

\begin{proof}
Proposition \ref{varieties-proposition-units-general}에 의해 (1)이 충분함을 알 수 있다.
(2), (3), (4)가 각각 (1)을 함의한다는 검증은 생략한다.
\end{proof}

\section{퀴네트 공식, I}
\label{varieties-section-kunneth}

\noindent
이 절에서는 바탕이 체이고 준연접 가군의 코호몰로지를 다룰 때의 퀴네트
공식을 증명한다. 더 일반적인 판본은 Derived Categories of Schemes,
Section \ref{perfect-section-kunneth}를 보라.

\begin{lemma}
\label{varieties-lemma-kunneth}
$k$를 체라 하자. $X$와 $Y$를 $k$ 위의 스킴이라 하고,
$\mathcal{F}$, 각각 $\mathcal{G}$를 준연접 $\mathcal{O}_X$-가군,
각각 $\mathcal{O}_Y$-가군이라 하자. 그러면 정준 동형
$$
H^n(X \times_{\Spec(k)} Y, \text{pr}_1^*\mathcal{F}
\otimes_{\mathcal{O}_{X \times_{\Spec(k)} Y}} \text{pr}_2^*\mathcal{G}) =
\bigoplus\nolimits_{p + q = n}
H^p(X, \mathcal{F}) \otimes_k H^q(Y, \mathcal{G})
$$
이 존재한다. 단, $X$와 $Y$는 준콤팩트이고 아핀 대각을 갖는다고
가정한다\footnote{$X$와 $Y$가 준분리인 경우는 아래의 Lemma
\ref{varieties-lemma-kunneth-general}에서 논한다.}
(예를 들어 $X$와 $Y$가 분리인 경우).
\end{lemma}

\begin{proof}
이 증명에서 별도 표시가 없는 곱과 텐서곱은 $k$ 위에서 취한다. 사상
$$
H^p(X, \mathcal{F}) \otimes H^q(Y, \mathcal{G})
\longrightarrow
H^n(X \times Y, \text{pr}_1^*\mathcal{F}
\otimes_{\mathcal{O}_{X \times Y}} \text{pr}_2^*\mathcal{G})
$$
으로 다음 것을 사용한다. 먼저 코호몰로지의 함자성을 이용하여 사상
$H^p(X, \mathcal{F}) \to H^p(X \times Y, \text{pr}_1^*\mathcal{F})$와
$H^p(Y, \mathcal{G}) \to H^p(X \times Y, \text{pr}_2^*\mathcal{G})$를
얻고, 이어서 컵곱
$$
\cup :
H^p(X \times Y, \text{pr}_1^*\mathcal{F})
\otimes H^q(X \times Y, \text{pr}_2^*\mathcal{G})
\longrightarrow
H^n(X \times Y, \text{pr}_1^*\mathcal{F}
\otimes_{\mathcal{O}_{X \times Y}} \text{pr}_2^*\mathcal{G})
$$
을 사용한다. $X$와 $Y$가 아핀이면 아핀 위의 고차 코호몰로지 군이
사라진다는 사실(Cohomology of Schemes, Lemma
\ref{coherent-lemma-quasi-coherent-affine-cohomology-zero})과 정의
(준연접 가군의 당김 및 준연접 가군의 텐서곱의 정의)에 의해 결과가 성립한다.

\medskip\noindent
유한 아핀 열린 덮개
$\mathcal{U} : X = \bigcup_{i \in I} U_i$와
$\mathcal{V} : Y = \bigcup_{j \in J} V_j$를 택하자.
이는 아핀 열린 덮개
$\mathcal{W} : X \times Y = \bigcup_{(i, j) \in I \times J} U_i \times V_j$
를 정한다. $\mathcal{W}$는 $\text{pr}_1^{-1}\mathcal{U}$와
$\text{pr}_2^{-1}\mathcal{V}$ 각각의 세분임에 주의하라.
따라서 Cohomology, Lemma \ref{cohomology-lemma-functoriality-cech}에 의해
사상
$$
\check{\mathcal{C}}^\bullet(\mathcal{U}, \mathcal{F}) \to
\check{\mathcal{C}}^\bullet(\mathcal{W}, \text{pr}_1^*\mathcal{F})
\quad\text{and}\quad
\check{\mathcal{C}}^\bullet(\mathcal{V}, \mathcal{G}) \to
\check{\mathcal{C}}^\bullet(\mathcal{W}, \text{pr}_2^*\mathcal{G})
$$
을 얻으며, 이들은 코호몰로지 위의 당김 사상과 양립한다. Cohomology,
Equation (\ref{cohomology-equation-needs-signs})에서는 복합체 사상
$$
\text{Tot}(
\check{\mathcal{C}}^\bullet(\mathcal{W}, \text{pr}_1^*\mathcal{F})
\otimes
\check{\mathcal{C}}^\bullet(\mathcal{W}, \text{pr}_2^*\mathcal{G}))
\longrightarrow
\check{\mathcal{C}}^\bullet(\mathcal{W},
\text{pr}_1^*\mathcal{F} \otimes_{\mathcal{O}_{X \times Y}}
\text{pr}_2^*\mathcal{G})
$$
을 구성했으며, 이는 코호몰로지 위의 컵곱을 정의한다. 위의 것들을 합치면
복합체 사상
\begin{equation}
\label{varieties-equation-kunneth-on-cech}
\text{Tot}(
\check{\mathcal{C}}^\bullet(\mathcal{U}, \mathcal{F})
\otimes
\check{\mathcal{C}}^\bullet(\mathcal{V}, \mathcal{G}))
\longrightarrow
\check{\mathcal{C}}^\bullet(\mathcal{W},
\text{pr}_1^*\mathcal{F} \otimes_{\mathcal{O}_{X \times Y}}
\text{pr}_2^*\mathcal{G})
\end{equation}
을 얻는다. 이 사상이 복합체의 동형은 아님을 독자에게 경고한다. 우리의
준연접 층들의 코호몰로지는 이 덮개들을 사용해 계산할 수 있음을 상기하자
(Cohomology of Schemes, Lemmas
\ref{coherent-lemma-affine-diagonal} 및
\ref{coherent-lemma-cech-cohomology-quasi-coherent}).
따라서 코호몰로지 위에서 (\ref{varieties-equation-kunneth-on-cech})는 보조정리의
사상을 재현한다.

\medskip\noindent
준연접 가군들의 짧은 완전열
$0 \to \mathcal{F} \to \mathcal{F}' \to \mathcal{F}'' \to 0$을 생각하자.
(\ref{varieties-equation-kunneth-on-cech})의 구성이 $\mathcal{F}$에 대해 함자적이고,
관련 {\v C}ech 복합체의 형성이 변수 $\mathcal{F}$에 대해 완전하므로
(아핀 열린집합 위에서 단면을 취하기 때문이다), 짧은 완전 복합체 열들 사이의
사상
$$
\xymatrix{
\text{Tot}(
\check{\mathcal{C}}^\bullet(\mathcal{U}, \mathcal{F})
\otimes
\check{\mathcal{C}}^\bullet(\mathcal{V}, \mathcal{G})) \ar[r] \ar[d] &
\text{Tot}(
\check{\mathcal{C}}^\bullet(\mathcal{U}, \mathcal{F}')
\otimes
\check{\mathcal{C}}^\bullet(\mathcal{V}, \mathcal{G})) \ar[r] \ar[d] &
\text{Tot}(
\check{\mathcal{C}}^\bullet(\mathcal{U}, \mathcal{F}'')
\otimes
\check{\mathcal{C}}^\bullet(\mathcal{V}, \mathcal{G})) \ar[d] \\
\check{\mathcal{C}}^\bullet(\mathcal{W},
\text{pr}_1^*\mathcal{F} \otimes_{\mathcal{O}_{X \times Y}}
\text{pr}_2^*\mathcal{G}) \ar[r] &
\check{\mathcal{C}}^\bullet(\mathcal{W},
\text{pr}_1^*\mathcal{F}' \otimes_{\mathcal{O}_{X \times Y}}
\text{pr}_2^*\mathcal{G}) \ar[r] &
\check{\mathcal{C}}^\bullet(\mathcal{W},
\text{pr}_1^*\mathcal{F}'' \otimes_{\mathcal{O}_{X \times Y}}
\text{pr}_2^*\mathcal{G})
}
$$
을 얻는다(바깥쪽 영들은 생략했다). 긴 완전 코호몰로지 열을 살펴보고
$2$-out-of-$3$ 원리를 적용하면, $\mathcal{F}, \mathcal{F}', \mathcal{F}''$ 가운데 두 개에 대해 보조정리의 결과가 성립할 경우
나머지 하나에 대해서도 성립함을 알 수 있다.

\medskip\noindent
$X$는 준연접 가군에 대해 유한 코호몰로지 차원을 갖는다. Cohomology of
Schemes, Lemma \ref{coherent-lemma-vanishing-nr-affines}를 보라.
$d(\mathcal{F}) = \max \{d \mid H^d(X, \mathcal{F}) \not = 0\}$에
관한 귀납법을 사용하여 $d(\mathcal{F}) = 0$인 경우로 귀착시키겠다.
$d(\mathcal{F}) > 0$이라고 가정하자. Cohomology of Schemes, Lemma
\ref{coherent-lemma-affine-diagonal-universal-delta-functor}에서 보았듯,
매입 $\mathcal{F} \to \mathcal{F}'$가 존재하여
$H^p(X, \mathcal{F}') = 0$이 모든 $p \geq 1$에 대해 성립한다.
$\mathcal{F}'' = \Coker(\mathcal{F} \to \mathcal{F}')$로 두면
$d(\mathcal{F}'') < d(\mathcal{F})$이다. 이제 앞 단락의 결과를 적용하면
$\mathcal{F}'$와 $\mathcal{F}''$에 대해 보조정리를 증명하는 것으로
충분하며, 이로써 귀납 단계를 증명한다.

\medskip\noindent
$\mathcal{G}$에 대해서도 같은 방식으로 논증하면 $\mathcal{F}$와
$\mathcal{G}$가 모두 차수 $0$에서만 영이 아닌 코호몰로지를 갖는다고
가정해도 됨을 알 수 있다. $V \subset Y$를 아핀 열린집합이라 하자.
아핀 열린 덮개
$\mathcal{U}_V : X \times V = \bigcup_{i \in I} U_i \times V$를 생각하자.
다음은 자명하다.
$$
\check{\mathcal{C}}^\bullet(\mathcal{U}, \mathcal{F}) \otimes \mathcal{G}(V) =
\check{\mathcal{C}}^\bullet(\mathcal{U}_V,
\text{pr}_1^*\mathcal{F} \otimes_{\mathcal{O}_{X \times Y}}
\text{pr}_2^*\mathcal{G})
$$
(복합체의 등식이다). 따라서 $Y$ 위에서
$$
R\text{pr}_{2, *}(\text{pr}_1^*\mathcal{F} \otimes_{\mathcal{O}_{X \times Y}}
\text{pr}_2^*\mathcal{G})
\cong
\Gamma(X, \mathcal{F}) \otimes_k \mathcal{G} \cong
\bigoplus\nolimits_{\alpha \in A} \mathcal{G}
$$
임을 얻는다. 여기서 $A$는 $k$-벡터공간 $\Gamma(X, \mathcal{F})$의
기저이다. $Y$ 위의 코호몰로지는 직합과 교환한다(Cohomology, Lemma
\ref{cohomology-lemma-quasi-separated-cohomology-colimit}).
$\text{pr}_2$에 대한 르레 스펙트럼 열을 사용하면(Cohomology, Lemma
\ref{cohomology-lemma-apply-Leray}를 통해),
$H^n(X \times Y, \text{pr}_1^*\mathcal{F} \otimes_{\mathcal{O}_{X \times Y}}
\text{pr}_2^*\mathcal{G})$가 $n > 0$에서는 영이고
$H^0(X, \mathcal{F}) \otimes H^0(Y, \mathcal{G})$와 동형임을 얻는다.
후자는 $n = 0$인 경우이다.
이로써 증명이 끝난다(다만 이 동형이 차수 $0$의 컵곱으로 실제로 주어지는지
확인해야 하는데, 그 검증은 생략한다).
\end{proof}

\begin{lemma}
\label{varieties-lemma-kunneth-general}
$k$를 체라 하자. $X$와 $Y$를 $k$ 위의 스킴이라 하고,
$\mathcal{F}$, 각각 $\mathcal{G}$를 준연접 $\mathcal{O}_X$-가군,
각각 $\mathcal{O}_Y$-가군이라 하자. 그러면 정준 동형
$$
H^n(X \times_{\Spec(k)} Y, \text{pr}_1^*\mathcal{F}
\otimes_{\mathcal{O}_{X \times_{\Spec(k)} Y}} \text{pr}_2^*\mathcal{G}) =
\bigoplus\nolimits_{p + q = n}
H^p(X, \mathcal{F}) \otimes_k H^q(Y, \mathcal{G})
$$
이 존재한다. 단, $X$와 $Y$는 준콤팩트이고 준분리라고 가정한다.
\end{lemma}

\begin{proof}
$X$와 $Y$가 분리이거나 더 일반적으로 아핀 대각을 가지면, 더 ‘좋은’
증명은 Lemma \ref{varieties-lemma-kunneth}를 보라(그 증명이 이 증명보다 나은 점은
사상들을 당김 뒤의 컵곱으로 식별한다는 것이다).
$X'$, 각각 $Y'$를 $X$, 각각 $Y$의 무한소 두껍히기로 두고, 그 구조층을
$\mathcal{O}_{X'} = \mathcal{O}_X \oplus \mathcal{F}$,
각각 $\mathcal{O}_{Y'} = \mathcal{O}_Y \oplus \mathcal{G}$로 둔다.
여기서 $\mathcal{F}$, 각각 $\mathcal{G}$는 제곱이 영인 아이디얼이다.
그러면
$$
\mathcal{O}_{X' \times Y'} =
\mathcal{O}_{X \times Y}
\oplus \text{pr}_1^*\mathcal{F}
\oplus \text{pr}_2^*\mathcal{G}
\oplus \text{pr}_1^*\mathcal{F}
\otimes_{\mathcal{O}_{X \times Y}} \text{pr}_2^*\mathcal{G}
$$
이 $X \times Y$ 위의 층들의 등식으로 성립한다. 이로부터 임의의 준콤팩트
준분리 $k$-스킴 쌍에 대해
$$
H^n(X \times Y, \mathcal{O}_{X \times Y}) =
\bigoplus\nolimits_{p + q = n}
H^p(X, \mathcal{O}_X) \otimes_k H^q(Y, \mathcal{O}_Y)
$$
임을 증명하면 충분하다는 것을 알 수 있다. 몇 가지 세부사항은 생략한다.

\medskip\noindent
이 명제를 증명하기 위해 Cohomology of Schemes, Lemma
\ref{coherent-lemma-hypercoverings} 형태의 코호몰로지와 밑변환을 사용한다.
이 보조정리에 따르면 아래로 유계인 $k$-벡터공간의 복합체, 즉 복합체
$\mathcal{K}^\bullet$가 존재한다. 이는 $\Spec(k)$ 위의 준연접 가군들로 이루어져
$Y$의 $\Spec(k)$ 위의 코호몰로지를 보편적으로 계산한다. 특히
$$
R\text{pr}_{1, *}(\mathcal{O}_{X \times Y}) \cong
(X \to \Spec(k))^*\mathcal{K}^\bullet
$$
가 $D(\mathcal{O}_X)$에서 성립한다. 호모토피까지 고려하면 복합체
$\mathcal{K}^\bullet$는
$\bigoplus_{q \geq 0} H^q(Y, \mathcal{O}_Y)[-q]$와 동형이다.
이는 체 위의 벡터공간의 모든 복합체에 대해 참이기 때문이다. 따라서
$$
R\text{pr}_{1, *}(\mathcal{O}_{X \times Y}) \cong
\bigoplus\nolimits_{q \geq 0}
H^q(Y, \mathcal{O}_Y)[-q] \otimes_k \mathcal{O}_X
$$
가 $D(\mathcal{O}_X)$에서 성립한다. 이제 다음을 얻는다.
\begin{align*}
R\Gamma(X \times Y, \mathcal{O}_{X \times Y})
& =
R\Gamma(X, R\text{pr}_{1, *}(\mathcal{O}_{X \times Y})) \\
& =
R\Gamma(X, \bigoplus\nolimits_{q \geq 0}
H^q(Y, \mathcal{O}_Y)[-q] \otimes_k \mathcal{O}_X) \\
& =
\bigoplus\nolimits_{q \geq 0}
R\Gamma(X,  H^q(Y, \mathcal{O}_Y) \otimes \mathcal{O}_X)[-q] \\
& =
\bigoplus\nolimits_{q \geq 0}
R\Gamma(X, \mathcal{O}_X) \otimes_k H^q(Y, \mathcal{O}_Y)[-q] \\
& =
\bigoplus\nolimits_{p, q \geq 0}
H^p(X, \mathcal{O}_X)[-p] \otimes_k H^q(Y, \mathcal{O}_Y)[-q]
\end{align*}
원하는 등식이다. 첫 번째 등식은 $\text{pr}_1$에 대한 르레에서 따른다
(Cohomology, Lemma \ref{cohomology-lemma-before-Leray}).
두 번째 등식은 위에서 얻은 전체 직접상의 분해에서 따른다.
세 번째 등식은 코호몰로지가 언제나 유한 직합과 교환하기 때문이다
(그리고 $Y$의 코호몰로지는 충분히 큰 차수에서 사라진다. Cohomology of
Schemes, Lemma
\ref{coherent-lemma-vanishing-nr-affines-quasi-separated}를 보라).
네 번째 등식은 $X$ 위의 코호몰로지가 무한 직합과 교환하기 때문이다.
Cohomology, Lemma
\ref{cohomology-lemma-quasi-separated-cohomology-colimit}를 보라.
마지막 등식은 위에서 언급한 체의 유도범주에 관한 사실에서 따른다.
\end{proof}

\section{대수다양체의 피카르 군}
\label{varieties-section-picard-groups}

\noindent
이 절에서는 대수다양체의 피카르 군에 관한 몇 가지 초보적 결과를 모은다.

\begin{lemma}
\label{varieties-lemma-change-rings-pic-pre}
$A \to B$를 충실 평탄 환 준동형이라 하자. $X$를 $A$ 위의 준콤팩트
준분리 스킴이라 하자. $\mathcal{L}$을 가역 $\mathcal{O}_X$-가군이라
하고, 그 $X_B$로의 당김이 자명하다고 하자. 그러면
$H^0(X, \mathcal{L})$과 $H^0(X, \mathcal{L}^{\otimes -1})$은 가역
$H^0(X, \mathcal{O}_X)$-가군이고, 곱셈 사상은 동형
$$
H^0(X, \mathcal{L}) \otimes_{H^0(X, \mathcal{O}_X)}
H^0(X, \mathcal{L}^{\otimes -1}) \longrightarrow
H^0(X, \mathcal{O}_X)
$$
을 유도한다.
\end{lemma}

\begin{proof}
$\mathcal{L}_B$를 $\mathcal{L}$의 $X_B$로의 당김이라 쓰자. 동형
$\mathcal{L}_B \to \mathcal{O}_{X_B}$를 택하자.
$R = H^0(X, \mathcal{O}_X)$, $M = H^0(X, \mathcal{L})$로 두고
$M$을 $R$-가군으로 생각하자. 임의의 준연접 $\mathcal{O}_X$-가군
$\mathcal{F}$와 그 당김 $\mathcal{F}_B$를 생각하자. 이는 $X_B$ 위에 있다.
이때 정준 동형
$H^0(X_B, \mathcal{F}_B) = H^0(X, \mathcal{F}) \otimes_A B$가 있다.
Cohomology of Schemes, Lemma
\ref{coherent-lemma-flat-base-change-cohomology}를 보라. 따라서
$$
M \otimes_R (R \otimes_A B) =
M \otimes_A B = H^0(X_B, \mathcal{L}_B) \cong
H^0(X_B, \mathcal{O}_{X_B}) = R \otimes_A B
$$
이다. $R \to R \otimes_A B$는 충실 평탄 사상
$A \to B$의 밑변환이므로 충실 평탄이다. 따라서 Algebra, Proposition
\ref{algebra-proposition-ffdescent-finite-projectivity}에 의해 $M$은 가역
$R$-가군이다. 마찬가지로 $N = H^0(X, \mathcal{L}^{\otimes -1})$도 가역
$R$-가군이다. 텐서곱에 관한 명제가 참임을 보이려면, $X_B$로 당긴 뒤에는
그것이 참이라는 사실과 $R \to R \otimes_A B$의 충실 평탄성을 사용한다.
몇 가지 세부사항은 생략한다.
\end{proof}

\begin{lemma}
\label{varieties-lemma-change-rings-pic}
$A \to B$를 충실 평탄 환 준동형이라 하자. $X$를 $A$ 위의 스킴이라 하고
다음을 가정하자.
\begin{enumerate}
\item $X$는 준콤팩트이고 준분리이다.
\item $R = H^0(X, \mathcal{O}_X)$는 반국소환이다.
\end{enumerate}
그러면 당김 사상 $\Pic(X) \to \Pic(X_B)$는 단사이다.
\end{lemma}

\begin{proof}
$\mathcal{L}$을 가역 $\mathcal{O}_X$-가군이라 하고, 그 당김
$\mathcal{L}'$이 $X_B$ 위에서 자명하다고 하자.
$M = H^0(X, \mathcal{L})$과
$N = H^0(X, \mathcal{L}^{\otimes - 1})$로 두자. Lemma
\ref{varieties-lemma-change-rings-pic-pre}에 의해 $R$-가군 $M$과 $N$은 가역이다.
$R$이 반국소환이므로 $M \cong R$이고 $N \cong R$이다. Algebra, Lemma
\ref{algebra-lemma-locally-free-semi-local-free}를 보라. 생성원
$s \in M$과 $t \in N$을 택하자. 그러면
$st \in R = H^0(X, \mathcal{O}_X)$는 Lemma
\ref{varieties-lemma-change-rings-pic-pre}의 마지막 부분에 의해 단원이다.
따라서 $s$와 $t$는 $\mathcal{L}$과 $\mathcal{L}^{\otimes -1}$의
자명화를 각각 $X$ 위에서 정의한다.
\end{proof}

\begin{lemma}
\label{varieties-lemma-change-fields-pic}
$k'/k$를 체 확대라 하자. $X$를 $k$ 위의 스킴이라 하고 다음을 가정하자.
\begin{enumerate}
\item $X$는 준콤팩트이고 준분리이다.
\item $R = H^0(X, \mathcal{O}_X)$는 반국소환이다. 예를 들어
$\dim_k R < \infty$이면 그러하다.
\end{enumerate}
그러면 당김 사상 $\Pic(X) \to \Pic(X_{k'})$는 단사이다.
\end{lemma}

\begin{proof}
Lemma \ref{varieties-lemma-change-rings-pic}의 특수한 경우이다.
$\dim_k R < \infty$이면 $R$은 아르틴 환이고 따라서 반국소환이다
(Algebra, Lemmas \ref{algebra-lemma-finite-dimensional-algebra} 및
\ref{algebra-lemma-artinian-finite-nr-max}).
\end{proof}

\begin{example}
\label{varieties-example-need-condition}
Lemma \ref{varieties-lemma-change-fields-pic}는 스킴 $X$에 체 $k$ 위에서 아무 조건도
두지 않으면 참이 아니다. 다음이 그 예이다. $k$를 체라 하자.
$t \in \mathbf{P}^1_k$를 닫힌점이라 하자.
$X = \mathbf{P}^1 \setminus \{t\}$로 두자. 그러면 전사
$$
\mathbf{Z} = \Pic(\mathbf{P}^1_k) \longrightarrow \Pic(X)
$$
를 얻는다. 첫 번째 등식은 Divisors, Lemma
\ref{divisors-lemma-Pic-projective-space-UFD}에 의하고, 전사성은 Divisors,
Lemma \ref{divisors-lemma-extend-invertible-module}에 의한다
($\mathbf{P}^1_k$가 차원 $1$이고 $k$ 위에서 매끄러운 스킴이며, 따라서 그
모든 국소환이 이산 값매김환이기 때문이다). $\mathcal{L}$이 표시한 사상의
핵에 속하면 $\mathcal{L} \cong \mathcal{O}_{\mathbf{P}^1_k}(nt)$이다.
여기서 $n \in \mathbf{Z}$이다. 독자에게 다음을
보이는 일을 맡긴다.
$\mathcal{O}_{\mathbf{P}^1_k}(t) \cong \mathcal{O}_{\mathbf{P}^1_k}(d)$,
여기서 $d = [\kappa(t) : k]$이다. 따라서
$$
\Pic(X) = \mathbf{Z}/d\mathbf{Z}
$$
이다. 그러므로 $t$가 $k$-유리점이 아니면 $d > 1$이고 이 피카르 군은
영이 아니다. 한편 바탕체 $k$를 체 확대 $k'$로 확대했을 때
$k$-매입 $\kappa(t) \to k'$가 존재한다면
$\mathbf{P}^1_{k'} \setminus X_{k'}$는 $k'$-유리점 $t'$를 갖는다.
따라서 $\mathcal{O}_{\mathbf{P}^1_{k'}}(1) = \mathcal{O}_{\mathbf{P}^1_{k'}}(t')$는 사상
$\mathbf{Z} \to \Pic(X_{k'})$의 핵에 속하고, 위와 같은 방식으로
$\Pic(X_{k'}) = 0$임이 따른다.
\end{example}

\noindent
다음 보조정리는 ‘유리적으로 동치인 가역 가군들’이 정규 대수다양체 위에서
동형임을 말해 준다.

\begin{lemma}
\label{varieties-lemma-rational-equivalence-for-Pic}
$k$를 체라 하자. $X$를 $k$ 위의 정규 대수다양체라 하자.
$U \subset \mathbf{A}^n_k$를 $k$-유리점 $p, q \in U(k)$를 갖는 열린
부분스킴이라 하자. 모든 가역 가군 $\mathcal{L}$ 가운데
$X \times_{\Spec(k)} U$ 위에 있는 것에 대해 제한
$\mathcal{L}|_{X \times p}$와 $\mathcal{L}|_{X \times q}$는 동형이다.
\end{lemma}

\begin{proof}
$X \times_{\Spec(k)} U \to X$의 올들은 체 위의 아핀 $n$-공간의 열린
부분스킴이다. 따라서 Divisors, Lemma
\ref{divisors-lemma-open-subscheme-UFD}에 의해 이 올들의 피카르 군은
자명하다. Divisors, Lemma \ref{divisors-lemma-in-image-pullback}를
적용하면 $\mathcal{L}$은 가역 가군 $\mathcal{N}$의 당김이다. 이 가군은
$X$ 위에 있다.
\end{proof}

\section{바탕체의 유일성}
\label{varieties-section-base-field}

\noindent
‘$X$를 $k$ 위의 스킴이라 하자’라는 문구는 $X$가 사상
$X \to \Spec(k)$와 함께 주어진 스킴이라는 뜻이다. 이제 체 $k$가
스킴 $X$에 의해 유일하게 결정되는지 물을 수 있다. 물론 그렇지 않다.
예를 들어 $\mathbf{A}^1_{\mathbf{C}}$는 보통 복소수체 $\mathbf{C}$ 위의
스킴으로 생각하지만, $\mathbf{Q}$ 위의 스킴으로도 생각할 수 있다.
그러나 사상 $X \to \Spec(k)$가
$X \to \Spec(k') \to \Spec(k)$로 분해되지 않는다고 요구하면 어떨까?
여기서 $k'/k$는 임의의 자명하지 않은 체 확대이다. 바꾸어 말하면
$k$가 어떤 의미에서 극대여서 $X$가 $k$ 위에서 살 수 있다고 요구한다.

\medskip\noindent
이 조건도 $k$의 유일성을 보장하지 않음을 보여 주는 예는 스킴
$$
X =
\Spec\left(
\mathbf{Q}(x)[y]\left[\frac{1}{P(y)}, P \in \mathbf{Q}[y], P \not = 0\right]
\right)
$$
이다. 처음 보면 이는 $\mathbf{Q}(x)$ 위의 스킴인 듯하지만, 다시 보면
$\mathbf{Q}(y)$ 위의 스킴이기도 함이 분명하다. 더욱이 체
$\mathbf{Q}(x)$와 $\mathbf{Q}(y)$는
$R = \Gamma(X, \mathcal{O}_X)$의 부분체이고, $R$의 부분체 가운데
극대이다(세부사항은 생략한다). 특히 $\mathbf{Q}(x)$와
$\mathbf{Q}(y)$는 모두 위의 의미에서 극대이다. 두 사상
$X \to \Spec(\mathbf{Q}(x))$와
$X \to \Spec(\mathbf{Q}(y))$가 모두 ‘본질적으로 유한형’임에 주의하라
(즉 대응하는 환 준동형이 본질적으로 유한형이다). 따라서 $X$는 유한
차원의 뇌터 스킴이며, 즉 완전히 병리적인 예는 아니다.

\medskip\noindent
유일성을 방해할 수 있는 또 다른 문제는 스킴 $X$가 비축소일 수 있다는
것이다. 그 경우 $X$에서 주어진 체의 스펙트럼으로 가는 서로 다른 사상이
많이 존재할 수 있다. 구체적인 예로 쌍대수
$D = \mathbf{C}[y]/(y^2) = \mathbf{C} \oplus \epsilon \mathbf{C}$를
생각하자. 임의의 미분 $\theta : \mathbf{C} \to \mathbf{C}$가
$\mathbf{Q}$ 위에서 주어지면 환 준동형
$$
\mathbf{C} \longrightarrow D, \quad
c \longmapsto c + \epsilon \theta(c).
$$
을 얻는다. 이 사상들이 모두 일치하는 $\mathbf{C}$의 부분체는
$\mathbf{Q}$의 대수적 폐포이다. 이는 $X$가 살 수 있는 모든 체의
교집합을 취할 경우, $X$가 비축소이면 매우 작은 체를 얻게 될 수도
있다는 뜻이다.

\medskip\noindent
이와 관련된 한 가지 관찰은 다음과 같다. 체 $k$와 두 부분체
$k_1, k_2$, 즉 $k$의 부분체들이 주어지고, $k$가 $k_1$ 위에서도 $k_2$ 위에서도
유한이라고 하자. 그러면 일반적으로 $k$가
$k_1 \cap k_2$ 위에서 유한인 것은 {\it 아니다}. 예를 들어
체 $k = \mathbf{Q}(t)$와 그 부분체 $k_1 = \mathbf{Q}(t^2)$ 및
$\mathbf{Q}((t + 1)^2)$를 생각하자. 이 경우
$k_1 \cap k_2 = \mathbf{Q}$이다. 따라서 아래에서는 체들의 교집합을
취할 때 주의해야 한다.

\medskip\noindent
이제 이 모든 점을 감안하여, $X$가 체 위에서 국소 유한형이면 일정한
유일성이 성립함을 보인다. 정확한 결과는 다음과 같다.

\begin{proposition}
\label{varieties-proposition-unique-base-field}
$X$를 스킴이라 하자. $a : X \to \Spec(k_1)$과
$b : X \to \Spec(k_2)$를 $X$에서 체들의 스펙트럼으로 가는 사상이라 하자.
$a, b$가 국소 유한형이고 $X$가 축소이며 연결이라고 가정하자. 그러면
$k_1' = k_2'$이다. 여기서
$k_i' \subset \Gamma(X, \mathcal{O}_X)$는 $k_i$의
$\Gamma(X, \mathcal{O}_X)$ 안에서의 정수적 폐포이다.
\end{proposition}

\begin{proof}
먼저 $X$가 준콤팩트인 경우에 명제가 성립한다고 가정하자
(준콤팩트인 경우는 아래에서 증명한다).
$X$는 체 위에서 국소 유한형이므로 국소 뇌터 스킴이다. Morphisms,
Lemma \ref{morphisms-lemma-finite-type-noetherian}을 보라. 특히 이는
국소 연결이고, 열린 부분집합의 연결 성분들이 열려 있으며, 준콤팩트
열린집합들의 교집합이 준콤팩트이라는 뜻이다. 다음을 보라:
Properties, Lemma \ref{properties-lemma-Noetherian-topology},
Topology, Lemma \ref{topology-lemma-locally-connected},
Topology, Section \ref{topology-section-noetherian}, 및
Topology, Lemma \ref{topology-lemma-constructible-Noetherian-space}.
열린 덮개 $X = \bigcup_{i \in I} U_i$를 택하여 각 $U_i$가
준콤팩트이고 연결이 되게 하자. 각 $i$에 대해
$K_i \subset \mathcal{O}_X(U_i)$를 $k_1$과 $k_2$의 정수적 폐포라 하자.
각 쌍 $i, j \in I$에 대해
$$
U_i \cap U_j = \coprod U_{i, j, l}
$$
를 그 유한 개의 연결 성분들로 분해하자.
$K_{i, j, l} \subset \mathcal{O}(U_{i, j, l})$를 $k_1$과 $k_2$의
정수적 폐포라 쓰자. Lemma
\ref{varieties-lemma-integral-closure-ground-field}에 의해 환 $K_i$와
$K_{i, j, l}$은 체이다. 이제 $k_1'$과 $k_2'$가 모두 사상
$$
\prod K_i \longrightarrow \prod K_{i, j, l}, \quad
(x_i)_i \longmapsto x_i|_{U_{i, j, l}} - x_j|_{U_{i, j, l}}
$$
의 핵과 같다고 주장한다. 이것이 원하는 결론을 증명한다.
$k_1'$이 이 핵에 포함됨은 분명하다. 역으로 $(x_i)_i$가 핵에 속한다고
가정하자. 층 조건에 의해 $(x_i)_i$는
$f \in \mathcal{O}(X)$에 대응한다. 어떤 $i_0 \in I$를 택하고, 일계수 다항식
$P(T) \in k_1[T]$를 택하여 $P(x_{i_0}) = 0$이 되게 하자. 그러면
$P(f) = 0$이라고 주장하며, 이로부터 $f \in k_1$이 따라 결론을 얻는다.
이를 보이려면 $P(x_i) = 0$임을 모든 $i$에 대해 보이면 된다.
$i \in I$를 택하자. $X$가 연결이므로
수열 $i_0, i_1, \ldots, i_n = i \in I$가 존재하여
$U_{i_t} \cap U_{i_{t + 1}} \not = \emptyset$이다. 이는 각 $t$에 대해
어떤 $l_t$가 존재하여 $x_{i_t}$와 $x_{i_{t + 1}}$가 체
$K_{i_t, i_{t + 1}, l_t}$의 같은 원소로 보내진다는 뜻이다. 따라서
$P(x_{i_t}) = 0$이면 $P(x_{i_{t + 1}}) = 0$이다. 귀납법으로
$P(x_{i_0}) = 0$에서 시작하여 원하는 대로 $P(x_i) = 0$을 얻는다.

\medskip\noindent
명제의 증명을 끝내기 위해 $X$가 준콤팩트이라는 추가 가정 아래에서
명제를 증명한다. Lemma \ref{varieties-lemma-integral-closure-ground-field}에 의해
$k_i$를 $k_i'$로 바꾼 뒤, $k_i$가
$\Gamma(X, \mathcal{O}_X)$ 안에서 정수적으로 닫혀 있다고 가정해도 된다.
이는 $\mathcal{O}(X)^*/k_i^*$가 유한 생성 아벨 군임을 뜻한다.
Proposition \ref{varieties-proposition-units-general}을 보라.
$k_{12} = k_1 \cap k_2$를 $\mathcal{O}(X)$의 부분환으로 두자.
$k_{12}$는 체임에 주의하라. 사상
$$
k_1^*/k_{12}^* \longrightarrow \mathcal{O}(X)^*/k_2^*
$$
이 있으므로 $k_1^*/k_{12}^*$도 유한 생성 아벨 군이다. 따라서
$\alpha_1, \ldots, \alpha_n \in k_1^*$가 존재하여 모든
$\lambda \in k_1$는
$$
\lambda = c \alpha_1^{e_1} \ldots \alpha_n^{e_n}
$$
꼴로 쓸 수 있다. 여기서 어떤 $e_i \in \mathbf{Z}$와
$c \in k_{12}$를 취한다. 특히 환 준동형
$$
k_{12}[x_1, \ldots, x_n, \frac{1}{x_1 \ldots x_n}] \longrightarrow k_1, \quad
x_i \longmapsto \alpha_i
$$
은 전사이다. 힐베르트 영점정리 Algebra, Theorem
\ref{algebra-theorem-nullstellensatz}에 의해 $k_1$은 $k_{12}$의 유한
확대이다. 같은 방식으로 $k_2$도 $k_{12}$의 유한 확대임을 얻는다.
특히 $k_1$과 $k_2$는 모두 정수적 폐포 $k_{12}'$에 포함된다. 여기서 이는
$k_{12}$의 $\Gamma(X, \mathcal{O}_X)$ 안에서의 정수적 폐포이다. 그러나 Lemma
\ref{varieties-lemma-integral-closure-ground-field}에 의해 $k_{12}'$은 체이고,
$k_i$가 $\Gamma(X, \mathcal{O}_X)$ 안에서 정수적으로 닫히도록 택했으므로
원하는 대로 $k_1 = k_{12} = k_2$이다.
\end{proof}

\section{자기동형}
\label{varieties-section-automorphisms}

\noindent
체 위의 스킴의 자기동형을 다루는 절이다. 곡선의 (무한소) 자기동형에 관한
내용은 Algebraic Curves, Section
\ref{curves-section-vector-fields} 및 Moduli of Curves, Section
\ref{moduli-curves-section-finite-aut}를 보라.

\begin{lemma}
\label{varieties-lemma-automorphism}
$X$를 체 $k$ 위의 축소 유한형 스킴이라 하자. $f : X \to X$를
$k$ 위의 자기동형이라 하고, 이것이 $X$의 바탕 위상공간 위에서 항등
사상을 유도한다고 하자. 그러면 다음이 성립한다.
\begin{enumerate}
\item $f^*\mathcal{F} \cong \mathcal{F}$이다. 이는 모든 연접
$\mathcal{O}_X$-가군에 대해 성립한다.
\item $\dim(Z) > 0$이 모든 기약 성분 $Z \subset X$에 대해 성립하면
$f$는 항등 사상이다.
\end{enumerate}
\end{lemma}

\begin{proof}
(1)은 (2)와 $X$의 차원 $0$인 연결 성분들이 체의 스펙트럼이라는 사실에서
따른다.

\medskip\noindent
$Z \subset X$를 정역 닫힌 부분스킴으로 본 기약 성분이라 하자. 분명히
$f(Z) \subset Z$이고, $f|_Z : Z \to Z$는 $k$ 위의 자기동형으로서
$Z$의 바탕 위상공간 위에서 항등 사상을 유도한다. $X$가 축소이므로 사상들
$f|_Z : Z \to Z$가 항등임을 보이면 충분하다. 이는 다음 단락에서 다루는
경우로 귀착된다.

\medskip\noindent
$X$가 차원 $> 0$인 기약 스킴이라고 가정하자. 공집합이 아닌 아핀
열린집합 $U \subset X$를 택하자. $f(U) \subset U$이고
$U \subset X$가 스킴론적으로 조밀하므로
$f|_U : U \to U$가 항등임을 보이면 충분하다.

\medskip\noindent
$X = \Spec(A)$가 아핀이고 기약이며 차원 $> 0$이고, $k$가 무한체라고
가정하자. 상수가 아닌 $g \in A$를 택하자. 집합
$$
S = \bigcup\nolimits_{\lambda \in k} V(g - \lambda)
$$
은 $X$에서 조밀하다. 실제로 이는 조밀 부분집합
$\mathbf{A}^1_k(k)$의 비상수 사상 $g : X \to \mathbf{A}^1_k$에 의한
역상이다. $x \in S$이면 $g(x)$, 즉 $g$의 $\kappa(x)$에서의 상은
$k \to \kappa(x)$의 상에 속한다. 따라서
$f^\sharp : \kappa(x) \to \kappa(x)$는 $g(x)$를 고정한다.
그러므로 $f^\sharp(g)$의 $\kappa(x)$에서의 상은 $g(x)$와 같다.
따라서
$$
S \subset V(g - f^\sharp(g))
$$
이다. $X$가 축소이고 $S$가 조밀하므로 $g=f^\sharp(g)$임을 얻는다.
$f^\sharp = \text{id}_A$임이 증명된다. 실제로 $A$는 $k$-대수로서
위와 같은 원소 $g$들에 의해 생성된다(세부사항은 생략한다. 힌트:
상수 함수들의 집합은 유한 차원 $k$-부분벡터공간이며, 이는 $A$ 안에 있다).
따라서 $f = \text{id}_X$이다.

\medskip\noindent
$X = \Spec(A)$가 아핀이고 기약이며 차원 $> 0$이고, $k$가 유한체라고
가정하자. 모든 $1$차원 정역 닫힌 부분스킴 $C \subset X$에 대해 제한
$f|_C : C \to C$가 항등이면 $f$는 항등이다. 따라서 $X$가 곡선인
경우로 귀착된다. 유한체 위의 곡선은 유한 자기동형군을 갖는다
(세부사항은 생략한다). 따라서 $f$는 유한 위수, 이를테면 $n$을 갖는다.
그런 다음 위와 같이 비상수인 $g : X \to \mathbf{A}^1_k$를 택하고
집합
$$
S = \{x \in X\text{ closed such that }[\kappa(g(x)) : k]
\text{ is prime to }n\}
$$
을 생각하자. 앞과 같이 논증하면 $S$가 $X$에서 조밀함을 알 수 있다.
닫힌점 $x \in X$에 대해 사상
$f^\sharp : \kappa(x) \to \kappa(x)$는 위수가 $n$을 나누는
자기동형이다. 따라서 $x \in S$일 때 이 자기동형은 $g$의 상이
$\kappa(x)$ 안에서 생성하는 부분체 위에 자명하게 작용한다. 그러므로
$S \subset V(g - f^\sharp(g))$이고, 앞과 같이 결론을 얻는다.
\end{proof}

\section{오일러 표수}
\label{varieties-section-euler}

\noindent
이 절에서는 체 위에서 고유인 스킴 위의 연접 층의 오일러 표수에 관한
몇 가지 초보적인 성질을 증명한다.

\begin{definition}
\label{varieties-definition-euler-characteristic}
$k$를 체라 하자. $X$를 $k$ 위의 고유 스킴이라 하자. $\mathcal{F}$를
연접 $\mathcal{O}_X$-가군이라 하자. 이 상황에서
{\it $\mathcal{F}$의 오일러 표수}는 정수
$$
\chi(X, \mathcal{F}) = \sum\nolimits_i (-1)^i \dim_k H^i(X, \mathcal{F}).
$$
이다. 이 공식의 정당화는 아래를 보라.
\end{definition}

\noindent
정의의 상황에서는 벡터공간 $H^i(X, \mathcal{F})$ 가운데 유한 개만이
영이 아니다(Cohomology of Schemes, Lemma
\ref{coherent-lemma-quasi-coherence-higher-direct-images}). 또한 이 공간들은
각각 유한 차원이다(Cohomology of Schemes, Lemma
\ref{coherent-lemma-proper-over-affine-cohomology-finite}). 따라서
$\chi(X, \mathcal{F}) \in \mathbf{Z}$는 잘 정의된다. 이 정의는 체
$k$에 의존하며 쌍 $(X, \mathcal{F})$에만 의존하는 것이 아님에 주의하라.

\begin{lemma}
\label{varieties-lemma-euler-characteristic-additive}
$k$를 체라 하자. $X$를 $k$ 위의 고유 스킴이라 하자.
$0 \to \mathcal{F}_1 \to \mathcal{F}_2 \to \mathcal{F}_3 \to 0$을
$X$ 위의 연접 가군들의 짧은 완전열이라 하자. 그러면
$$
\chi(X, \mathcal{F}_2) = \chi(X, \mathcal{F}_1) + \chi(X, \mathcal{F}_3)
$$
이다.
\end{lemma}

\begin{proof}
보조정리의 짧은 완전열에 결부된 긴 완전 코호몰로지 열
$$
0 \to H^0(X, \mathcal{F}_1) \to H^0(X, \mathcal{F}_2) \to
H^0(X, \mathcal{F}_3) \to H^1(X, \mathcal{F}_1) \to \ldots
$$
을 생각하자. 선형대수의 계수-영차원 정리에 의해
$$
0 = \dim H^0(X, \mathcal{F}_1) - \dim H^0(X, \mathcal{F}_2)
+ \dim H^0(X, \mathcal{F}_3) - \dim H^1(X, \mathcal{F}_1) + \ldots
$$
이다. 이로부터 보조정리가 즉시 따른다.
\end{proof}

\begin{lemma}
\label{varieties-lemma-chi-tensor-finite}
$k$를 체라 하자. $X$를 $k$ 위의 고유 스킴이라 하자. $\mathcal{F}$를
$\dim(\text{Supp}(\mathcal{F})) \leq 0$인 연접 층이라 하자. 그러면
다음이 성립한다.
\begin{enumerate}
\item $\mathcal{F}$는 전역 단면들로 생성된다.
\item $H^0(X, \mathcal{F}) =
\bigoplus_{x \in \text{Supp}(\mathcal{F})} \mathcal{F}_x$이다.
\item $H^i(X, \mathcal{F}) = 0$은 $i > 0$에 대해 성립한다.
\item $\chi(X, \mathcal{F}) = \dim_k H^0(X, \mathcal{F})$이다.
\item $\chi(X, \mathcal{F} \otimes \mathcal{E}) = n\chi(X, \mathcal{F})$이다.
이는 모든 국소 자유 가군 $\mathcal{E}$ 가운데 계수가 $n$인 것에 대해 성립한다.
\end{enumerate}
\end{lemma}

\begin{proof}
Cohomology of Schemes, Lemma
\ref{coherent-lemma-coherent-support-closed}에 의해
$\mathcal{F} = i_*\mathcal{G}$이다. 여기서
$i : Z \to X$는 $\mathcal{F}$의 스킴론적 지지의 포함 사상이고,
$\mathcal{G}$는 연접 $\mathcal{O}_Z$-가군이다. 스킴론적 지지의 정의에
따라 $Z$의 바탕 위상공간은 $\text{Supp}(\mathcal{F})$이다.
$Z$의 차원이 $0$이므로 $Z$는 아핀이다(Properties, Lemma
\ref{properties-lemma-locally-Noetherian-dimension-0}). 따라서
$\mathcal{G}$는 전역 생성되고 $\mathcal{G}$의 고차 코호몰로지 군들은
영이다(Cohomology of Schemes, Lemma
\ref{coherent-lemma-quasi-coherent-affine-cohomology-zero}).
사실 Lemma \ref{varieties-lemma-algebraic-scheme-dim-0}에 의해 스킴 $Z$는 국소
아르틴 환들의 스펙트럼들의 유한 분리합이다. 따라서 그에 대응하여
$H^0(Z, \mathcal{G}) = \bigoplus_{z \in Z} \mathcal{G}_z$이다.
$\mathcal{F}$와 $\mathcal{G}$의 코호몰로지는 Cohomology of Schemes,
Lemma \ref{coherent-lemma-relative-affine-cohomology}에 의해 일치한다.
따라서 $H^i(X, \mathcal{F}) = 0$은 $i > 0$에 대해 성립하고
$H^0(X, \mathcal{F}) = H^0(Z, \mathcal{G})$이다. 특히 (3)이 참이다.
$z \in Z$가 $x \in \text{Supp}(\mathcal{F})$에 대응하면
$\mathcal{G}_z = (i_*\mathcal{G})_x = \mathcal{F}_x$이다.
따라서 (2)가 성립한다. 물론 (2)는 (1)을 함의한다. 오일러 표수
$\chi(X, \mathcal{F})$의 정의와 (3)에 의해 (4)를 얻는다.
사영 공식(Cohomology, Lemma
\ref{cohomology-lemma-projection-formula})에 의해
$$
i_*(\mathcal{G} \otimes i^*\mathcal{E}) = \mathcal{F} \otimes \mathcal{E}.
$$
이다. $Z$의 차원이 $0$이므로 국소 자유 층 $i^*\mathcal{E}$는
$\mathcal{O}_Z^{\oplus n}$과 동형이고, 위와 같이 논증하면 (5)를 얻는다.
\end{proof}

\begin{lemma}
\label{varieties-lemma-euler-characteristic-extend-base-field}
$k'/k$를 체 확대라 하자. $X$를 $k$ 위의 고유 스킴이라 하자.
$\mathcal{F}$를 $X$ 위의 연접 층이라 하자. $\mathcal{F}'$를
$\mathcal{F}$의 $X_{k'}$로의 당김이라 하자. 그러면
$\chi(X, \mathcal{F}) = \chi(X', \mathcal{F}')$이다.
\end{lemma}

\begin{proof}
이는 평탄 밑변환에 의한 등식
$$
H^i(X_{k'}, \mathcal{F}') = H^i(X, \mathcal{F}) \otimes_k k'
$$
때문에 참이다. Cohomology of Schemes, Lemma
\ref{coherent-lemma-flat-base-change-cohomology}를 보라.
\end{proof}

\begin{lemma}
\label{varieties-lemma-euler-characteristic-morphism}
$k$를 체라 하자. $f : Y \to X$를 $k$ 위의 고유 스킴 사이의 사상이라
하자. $\mathcal{G}$를 연접 $\mathcal{O}_Y$-가군이라 하자. 그러면
$$
\chi(Y, \mathcal{G}) = \sum (-1)^i \chi(X, R^if_*\mathcal{G})
$$
이다.
\end{lemma}

\begin{proof}
공식은 의미가 있다. 층 $R^if_*\mathcal{G}$는 연접이고 그중 유한 개만이
영이 아니다. Cohomology of Schemes, Proposition
\ref{coherent-proposition-proper-pushforward-coherent} 및 Lemma
\ref{coherent-lemma-quasi-coherence-higher-direct-images}를 보라.
Cohomology, Lemma \ref{cohomology-lemma-Leray}에 의해
$$
E_2^{p, q} = H^p(X, R^qf_*\mathcal{G})
$$
를 갖고 $H^{p + q}(Y, \mathcal{G})$로 수렴하는 스펙트럼 열이 있다.
$X$ 위의 코호몰로지의 유한성에 의해 $E_2^{p, q}$ 가운데 유한 개만
영이 아니고, 각 $E_2^{p, q}$는 유한 차원 벡터공간이다. 따라서
$E_r^{p, q}$도 $r \geq 2$에 대해 같은 성질을 가지며
$$
\sum (-1)^{p + q} \dim_k E_r^{p, q}
$$
는 $r$에 무관하다. 충분히 큰 $r$에 대해
$E_r^{p, q} = E_\infty^{p, q}$이고, 수렴한다는 말은
$H^n(Y, \mathcal{G})$ 위에 여과가 있어서 그 등급 조각들이
$E_\infty^{p, q}$이며 $p + q = n$이라는 뜻이므로(이것이 스펙트럼 열의
수렴의 의미이다), 결론을 얻는다. 더 일반적인 Homology, Lemma
\ref{homology-lemma-biregular-ss-relation-in-K0}와도 비교하라.
\end{proof}

\section{사영공간}
\label{varieties-section-projective-space}

\noindent
체 위의 사영공간에 관한 몇 가지 결과이다.

\begin{lemma}
\label{varieties-lemma-projective-space-smooth}
\begin{slogan}
사영공간은 매끄럽다.
\end{slogan}
$k$를 체라 하고 $n \geq 0$이라 하자. 그러면 $\mathbf{P}^n_k$는
차원 $n$이고 $k$ 위에서 매끄러운 사영 대수다양체이다.
\end{lemma}

\begin{proof}
생략한다.
\end{proof}

\begin{lemma}
\label{varieties-lemma-intersection-in-affine-space}
$k$를 체라 하고 $n \geq 0$이라 하자.
$X, Y \subset \mathbf{A}^n_k$를 닫힌 부분집합이라 하자. $X$와 $Y$가
순수차원이고 $\dim(X) = r$, $\dim(Y) = s$라고 가정하자. 그러면
$X \cap Y$의 모든 기약 성분의 차원은 $\geq r + s - n$이다.
\end{lemma}

\begin{proof}
닫힌 부분스킴 $X \times Y \subset \mathbf{A}^{2n}_k$을 생각하자.
여기서 좌표 $x_1, \ldots, x_n, y_1, \ldots, y_n$을 사용한다. 그러면
$X \cap Y = X \times Y \cap V(x_1 - y_1, \ldots, x_n - y_n)$이다.
닫힌점 $t \in X \cap Y \subset X \times Y$를 택하자. Lemma
\ref{varieties-lemma-dimension-product-locally-algebraic}에 의해
$\dim_t(X \times Y) = \dim(X) + \dim(Y)$이다. 따라서 Lemma
\ref{varieties-lemma-dimension-locally-algebraic}에 의해
$\dim(\mathcal{O}_{X \times Y, t}) = r + s$이다. Algebra, Lemma
\ref{algebra-lemma-one-equation}에 의해
$$
\dim(\mathcal{O}_{X \cap Y, t}) =
\dim(\mathcal{O}_{X \times Y, t}/(x_1 - y_1, \ldots, x_n - y_n)) \geq
r + s - n
$$
임을 얻는다. Lemma \ref{varieties-lemma-dimension-locally-algebraic}에 의해 이로부터
결과가 따른다.
\end{proof}

\begin{lemma}
\label{varieties-lemma-intersection-in-projective-space}
$k$를 체라 하고 $n \geq 0$이라 하자.
$X, Y \subset \mathbf{P}^n_k$를 공집합이 아닌 닫힌 부분집합이라 하자.
$\dim(X) = r$, $\dim(Y) = s$이고 $r + s \geq n$이면
$X \cap Y$는 공집합이 아니며
$\dim(X \cap Y) \geq r + s - n$이다.
\end{lemma}

\begin{proof}
$\mathbf{A}^n = \Spec(k[x_0, \ldots, x_n])$ 및
$\mathbf{P}^n = \text{Proj}(k[T_0, \ldots, T_n])$로 쓰자.
사상 $\pi : \mathbf{A}^{n + 1} \setminus \{0\} \to \mathbf{P}^n$을
생각하자. 이는 $(x_0, \ldots, x_n)$을 점
$[x_0 : \ldots : x_n]$으로 보낸다. 더 정확히 말하면 이는 쌍
$(\mathcal{O}_{\mathbf{A}^{n + 1} \setminus \{0\}}, (x_0, \ldots, x_n))$에
결부된 사상이다. Constructions, Lemma
\ref{constructions-lemma-projective-space}를 보라. 표준 아핀 열린집합
$D_+(T_i)$ 위에서는 환 준동형
$$
k\left[\frac{T_0}{T_i}, \ldots, \frac{T_n}{T_i}\right]
\longrightarrow
k\left[T_0, \ldots, T_n, \frac{1}{T_i}\right] \cong
k\left[\frac{T_0}{T_i}, \ldots, \frac{T_n}{T_i}\right]
\left[T_i, \frac{1}{T_i}\right]
$$
에 결부된 사상을 얻는다. 이는 전사이고, 상대차원 $1$로 매끄러우며,
기약 올들을 갖는다(세부사항은 생략한다). 따라서
$\pi^{-1}(X)$와 $\pi^{-1}(Y)$는 차원이 각각 $r + 1$과 $s + 1$인
공집합이 아닌 닫힌 부분집합이다. 기약 성분
$V \subset \pi^{-1}(X)$ 가운데 차원 $r + 1$인 것과 기약 성분
$W \subset \pi^{-1}(Y)$ 가운데 차원 $s + 1$인 것을 택하자.
이는 $V$와 $W$가 자신이 만나는
$\pi$의 모든 올을 포함함을 뜻한다. 실제로 $\pi$의 올들은 차원 $1$인
기약 스킴이고, Lemma
\ref{varieties-lemma-dimension-fibres-locally-algebraic}에 따르면
$V \to \pi(V)$와 $W \to \pi(W)$의 올들은 차원 $\geq 1$이다.
$\overline{V}$와 $\overline{W}$를 $V$와 $W$의
$\mathbf{A}^{n + 1}$ 안에서의 폐포라 하자.
$0 \in \mathbf{A}^{n + 1}$은 $\pi$의 모든 올의 폐포에 속하므로
$0 \in \overline{V} \cap \overline{W}$이다. Lemma
\ref{varieties-lemma-intersection-in-affine-space}에 의해
$\dim(\overline{V} \cap \overline{W}) \geq r + s - n + 1$이다.
다시 Lemma \ref{varieties-lemma-dimension-fibres-locally-algebraic}를 사용하여 위와
같이 논증하면 $\pi(V \cap W) \subset X \cap Y$의 차원이 적어도
$r + s - n$임을 얻는다.
\end{proof}

\begin{lemma}
\label{varieties-lemma-equation-codim-1-in-projective-space}
$k$를 체라 하자. $Z \subset \mathbf{P}^n_k$를 매장점이 없는 닫힌
부분스킴이라 하고, $Z$의 모든 기약 성분의 차원이 $n - 1$이라고 하자.
그러면 아이디얼 $I(Z) \subset k[T_0, \ldots, T_n]$, 즉 $Z$에
대응하는 아이디얼은 주아이디얼이다.
\end{lemma}

\begin{proof}
이는 Divisors, Lemma
\ref{divisors-lemma-equation-codim-1-in-projective-space}의 특수한 경우이다.
\end{proof}

\section{사영공간 위의 연접 층}
\label{varieties-section-coherent-Pn}

\noindent
이 절에서는 체 위의 $\mathbf{P}^n$ 위에 놓인 연접 층의 코호몰로지에
관하여 \cite{Mum}에서 찾을 수 있는 몇 가지 결과를 증명한다. 이 결과들은
나중에 Quot 스킴과 힐베르트 스킴을 논할 때 유용할 것이다.

\subsection{준비}
\label{varieties-subsection-preliminaries}

\noindent
$k$를 체, $n \geq 1$, $d \geq 1$이라 하고, 영이 아닌 단면
$s \in \Gamma(\mathbf{P}_k^n, \mathcal{O}(d))$를 택하자. 이 절에서는
$\mathcal{O}(d)$로 사영공간 구조층의 $d$번째 꼬임을 나타낸다
(Constructions, Definitions \ref{constructions-definition-twist} 및
\ref{constructions-definition-projective-space}).
$\mathbf{P}^n_k$가 대수다양체이므로 이 단면은 정칙이며, 따라서 $s$는
$\mathcal{O}(d)$의 정칙 단면이고 유효 카르티에 인자
$H = Z(s) \subset \mathbf{P}^n_k$를 정의한다. Divisors, Section
\ref{divisors-section-effective-Cartier-divisors}를 보라. 이러한 인자
$H$를 {\it 초곡면}이라 하고, $d = 1$이면 {\it 초평면}이라 한다.

\begin{lemma}
\label{varieties-lemma-hyperplane}
$k$를 체라 하자. $n \geq 1$이라 하자.
$i : H \to \mathbf{P}^n_k$를 초평면이라 하자. 그러면 동형
$$
\varphi : \mathbf{P}^{n - 1}_k \longrightarrow H
$$
이 존재하여 $i^*\mathcal{O}(1)$의 당김은 $\mathcal{O}(1)$이 된다.
\end{lemma}

\begin{proof}
$\mathbf{P}^n_k = \text{Proj}(k[T_0, \ldots, T_n])$이다.
단면 $s$는 $T_0, \ldots, T_n$에 관한 차수 $1$의 동차형식에 대응한다.
Cohomology of Schemes, Section
\ref{coherent-section-cohomology-projective-space}을 보라.
$s = \sum a_i T_i$라 하자. Constructions, Lemma
\ref{constructions-lemma-closed-in-projective-space}에 의해
$H = \text{Proj}(k[T_0, \ldots, T_n]/I)$이다. 여기서 등급 아이디얼
$I$는 $I_d$를 사상
$\Gamma(\mathbf{P}^n_k, \mathcal{O}(d)) \to \Gamma(H, i^*\mathcal{O}(d))$의
핵으로 둠으로써 정의한다.
Divisors, Definition \ref{divisors-definition-zero-scheme-s}에서 한
$Z(s)$의 구성에 따라, $D_{+}(T_j)$ 위에서 $H$의 아이디얼은
$\sum a_i T_i/T_j$로 생성되며, 이는 다항식환
$k[T_0/T_j, \ldots, T_n/T_j]$의 원소이다.
따라서 $I$가 $\sum a_i T_i$로 생성되는 아이디얼임은 분명하다.
다음 등급환의 동형에 주의하라.
$$
k[T_0, \ldots, T_n]/I = k[T_0, \ldots, T_n]/(\sum a_i T_i) \cong
k[S_0, \ldots, S_{n - 1}]
$$
예를 들어 $a_n \not = 0$이면 $S_i$를 $T_i$의 류로 보내면 된다.
$\text{Proj}$의 함자성에 의해 원하는 동형을 얻는다. 구조층 꼬임들의
동일성은 예를 들어 Constructions, Lemma
\ref{constructions-lemma-surjective-graded-rings-generated-degree-1-map-proj}에서
따른다.
\end{proof}

\begin{lemma}
\label{varieties-lemma-exact-sequence-induction}
$k$를 무한체라 하자. $n \geq 1$이라 하자. $\mathcal{F}$를
$\mathbf{P}^n_k$ 위의 연접 가군이라 하자. 그러면 영이 아닌 단면
$s \in \Gamma(\mathbf{P}^n_k, \mathcal{O}(1))$과 짧은 완전열
$$
0 \to \mathcal{F}(-1) \to \mathcal{F} \to i_*\mathcal{G} \to 0
$$
이 존재한다. 여기서 $i : H \to \mathbf{P}^n_k$는 초평면 $H$의 몰입이고,
이 초평면은 $s$에 결부되어 있으며, $\mathcal{G} = i^*\mathcal{F}$이다.
\end{lemma}

\begin{proof}
사상 $\mathcal{F}(-1) \to \mathcal{F}$는
Constructions, Equation
(\ref{constructions-equation-multiply-on-sheaf})에서 $n = 1$,
$m = -1$, 그리고 단면 $s$를 $\mathcal{O}(1)$에서 취하여 얻는다.
이 사상을 $D_{+}(T_0)$에 제한했을 때 어떤 모습인지 계산하자.
$D_{+}(T_0) = \Spec(k[x_1, \ldots, x_n])$로 쓰고
$x_i = T_i/T_0$로 두자. $\mathcal{O}(1)|_{D_{+}(T_0)}$를
$\mathcal{O}$와 동일시하되 단면 $T_0$를 사용한다. 따라서
$s = \sum a_iT_i$이면 위에서 택한 동일시 아래
$s|_{D_{+}(T_0)} = a_0 + \sum a_ix_i$이다. 더 나아가
$\mathcal{F}|_{D_{+}(T_0)}$가 유한
$k[x_1, \ldots, x_n]$-가군 $M$에 대응한다고 하자. 동일시
$\mathcal{F}(-1) = \mathcal{F} \otimes \mathcal{O}(-1)$과
$\mathcal{O}(1)$의 위에서 택한 자명화에 의해
$\mathcal{F}(-1)$도 $M$에 대응한다. 따라서
$\mathcal{F}(-1) \to \mathcal{F}$를 $D_{+}(T_0)$에 제한하면 사상
$$
M \xrightarrow{a_0 + \sum a_ix_i} M
$$
을 얻는다. 화살표가 단사임을 보이려면 $a_0 + \sum a_ix_i$를
$M$의 모든 결합소아이디얼 밖에서 택하면 충분하다. Algebra, Lemma
\ref{algebra-lemma-ass-zero-divisors}를 보라. Algebra, Lemma
\ref{algebra-lemma-finite-ass}에 의해 집합 $\text{Ass}(M)$은
$M$의 결합소아이디얼들로 이루어지며 유한하다. $\mathfrak p \in \text{Ass}(M)$에 대해
교집합 $\mathfrak p \cap \{a_0 + \sum a_i x_i\}$는 진
$k$-부분벡터공간임에 주의하라. 따라서 유한한 진 부분벡터공간들의 족
$V_1, \ldots, V_m \subset \Gamma(\mathbf{P}^n_k, \mathcal{O}(1))$가
존재하여, $s$를 $\bigcup V_i$ 밖에서 택하면 $s$에 의한 곱셈이
$D_{+}(T_0)$ 위에서 단사이다. $D_{+}(T_j)$로의 제한도 마찬가지이며,
여기서 $j = 1, \ldots, n$이다. $k$가 무한이므로 유한 개의 진
부분벡터공간들의 합집합은 전체 공간과 같을 수 없다. 따라서 사상이
단사가 되도록 $s$를 택할 수 있다.
$\mathcal{F}(-1) \to \mathcal{F}$의 여핵은
$\Im(s : \mathcal{O}(-1) \to \mathcal{O})$에 의해 소멸된다. 이것은
Divisors, Definition \ref{divisors-definition-zero-scheme-s}에 의해
$H$의 아이디얼층이다. 따라서 Cohomology of Schemes, Lemma
\ref{coherent-lemma-i-star-equivalence}를 이용하여 $\mathcal{G}$를
$H$ 위에서 얻는다.
\end{proof}

\begin{remark}
\label{varieties-remark-exact-sequence-induction}
$k$를 무한체라 하자. $n \geq 1$이라 하자. 유한 개의 연접 가군
$\mathcal{F}_i$가 $\mathbf{P}^n_k$ 위에 주어지면 단 하나의
$s \in \Gamma(\mathbf{P}^n_k, \mathcal{O}(1))$를 택하여 Lemma
\ref{varieties-lemma-exact-sequence-induction}의 명제가 각 가군에 대해 성립하게
할 수 있다. 이를 증명하려면 그 보조정리를
$\bigoplus \mathcal{F}_i$에 적용하면 된다.
\end{remark}

\begin{remark}
\label{varieties-remark-exact-sequence-induction-cohomology}
Lemmas \ref{varieties-lemma-hyperplane} 및
\ref{varieties-lemma-exact-sequence-induction}의 상황에서
$H \cong \mathbf{P}^{n - 1}_k$이고, 세르 꼬임은
$\mathcal{O}_H(d) = i^*\mathcal{O}_{\mathbf{P}^n_k}(d)$이다. 모든 $d \in \mathbf{Z}$에
대해 짧은 완전열
$$
0 \to \mathcal{F}(d - 1) \to \mathcal{F}(d) \to i_*(\mathcal{G}(d)) \to 0
$$
이 있다. 실제로 $\mathcal{O}_{\mathbf{P}^n_k}(d)$와의 텐서곱은 완전
함자이고, 사영 공식(Cohomology, Lemma
\ref{cohomology-lemma-projection-formula})에 의해
$i_*(\mathcal{G}(d)) = i_*\mathcal{G} \otimes \mathcal{O}_{\mathbf{P}^n_k}(d)$이다.
이에 대응하는 긴 완전열
$$
H^i(\mathbf{P}^n_k, \mathcal{F}(d - 1)) \to
H^i(\mathbf{P}^n_k, \mathcal{F}(d)) \to
H^i(H, \mathcal{G}(d)) \to
H^{i + 1}(\mathbf{P}^n_k, \mathcal{F}(d - 1))
$$
을 얻는다. 이는 위의 내용과 Cohomology of Schemes, Lemma
\ref{coherent-lemma-relative-affine-cohomology}에 의한 등식
$H^i(\mathbf{P}^n_k, i_*\mathcal{G}(d)) = H^i(H, \mathcal{G}(d))$에서
따른다(닫힌 몰입은 아핀이다).
\end{remark}

\subsection{정칙성}
\label{varieties-subsection-regularity}

\noindent
정의는 다음과 같다.

\begin{definition}
\label{varieties-definition-regularity}
$k$를 체라 하자. $n \geq 0$이라 하자. $\mathcal{F}$를
$\mathbf{P}^n_k$ 위의 연접 층이라 하자. 다음이 성립하면
$\mathcal{F}$를 {\it $m$-정칙}이라 한다.
$$
H^i(\mathbf{P}^n_k, \mathcal{F}(m - i)) = 0
$$
여기서 $i = 1, \ldots, n$이다.
\end{definition}

\noindent
$\mathcal{F} = \mathcal{O}(d)$가 $m$-정칙일 필요충분조건은
$d \geq -m$임에 주의하라. 이는 Cohomology of Schemes, Equation
(\ref{coherent-equation-identify})의 코호몰로지 군 계산에서 따른다.
실제로 $H^n(\mathbf{P}^n_k, \mathcal{O}(d)) = 0$일 필요충분조건은
$d \geq -n$이다.

\begin{lemma}
\label{varieties-lemma-m-regular-extend-base-field}
$k'/k$를 체의 확대라 하자. $n \geq 0$이라 하자.
$\mathcal{F}$를 $\mathbf{P}^n_k$ 위의 연접 층이라 하자.
$\mathcal{F}'$를 $\mathcal{F}$의 $\mathbf{P}^n_{k'}$로의 당김이라 하자.
그러면 $\mathcal{F}$가 $m$-정칙일 필요충분조건은 $\mathcal{F}'$가
$m$-정칙인 것이다.
\end{lemma}

\begin{proof}
평탄 기저변환에 의해
$$
H^i(\mathbf{P}^n_{k'}, \mathcal{F}') =
H^i(\mathbf{P}^n_k, \mathcal{F}) \otimes_k k'
$$
이기 때문이다. Cohomology of Schemes, Lemma
\ref{coherent-lemma-flat-base-change-cohomology}를 보라.
\end{proof}

\begin{lemma}
\label{varieties-lemma-m-regular}
Lemma \ref{varieties-lemma-exact-sequence-induction}의 상황에서
$\mathcal{F}$가 $m$-정칙이면 $\mathcal{G}$는 $m$-정칙이고,
그 바탕은 $H \cong \mathbf{P}^{n - 1}_k$이다.
\end{lemma}

\begin{proof}
Cohomology of Schemes, Lemma
\ref{coherent-lemma-relative-affine-cohomology}에 의해
$H^i(\mathbf{P}^n_k, i_*\mathcal{G}) = H^i(H, \mathcal{G})$임을 상기하라.
따라서 Remark \ref{varieties-remark-exact-sequence-induction-cohomology}에 의해
$i \geq 1$에 대해
$$
H^i(\mathbf{P}^n_k, \mathcal{F}(m - i)) \to
H^i(H, \mathcal{G}(m - i)) \to
H^{i + 1}(\mathbf{P}^n_k, \mathcal{F}(m - 1 - i))
$$
를 얻는다. 이제 보조정리가 따른다.
\end{proof}

\begin{lemma}
\label{varieties-lemma-m-regular-up}
$k$를 체라 하자. $n \geq 0$이라 하자.
$\mathcal{F}$를 $\mathbf{P}^n_k$ 위의 연접 층이라 하자.
$\mathcal{F}$가 $m$-정칙이면 $\mathcal{F}$는
$(m + 1)$-정칙이다.
\end{lemma}

\begin{proof}
$n$에 관한 귀납법으로 증명한다. $n = 0$이면 모든 층은 모든 $m$에
대해 $m$-정칙이므로 증명할 것이 없다. Lemma
\ref{varieties-lemma-m-regular-extend-base-field}에 의해 $k$를 무한 확대체로
바꾸어 $k$가 무한이라고 가정할 수 있다. 따라서 Lemma
\ref{varieties-lemma-exact-sequence-induction}을 적용할 수 있다. Lemma
\ref{varieties-lemma-m-regular}에 의해 $\mathcal{G}$가 $m$-정칙임을 안다.
$n$에 관한 귀납법에 의해 $\mathcal{G}$는 $(m + 1)$-정칙이다.
Remark \ref{varieties-remark-exact-sequence-induction-cohomology}에서와 같이 완전열
$$
0 \to \mathcal{F}(m - i) \to \mathcal{F}(m + 1 - i)
\to i_*\mathcal{G}(m + 1 - i) \to 0
$$
에 결부된 긴 코호몰로지 완전열을 살펴보면, 독자는 $i \geq 1$에 대해
$H^i(\mathbf{P}^n_k, \mathcal{F}(m + 1 - i))$의 소멸을 쉽게 얻는다.
이는 이미 알려진 $H^i(\mathbf{P}^n_k, \mathcal{F}(m - i))$의 소멸과
$H^i(\mathbf{P}^n_k, \mathcal{G}(m + 1 - i))$의 소멸에서 따른다.
\end{proof}

\begin{lemma}
\label{varieties-lemma-m-regular-multiply}
$k$를 체라 하자. $n \geq 0$이라 하자.
$\mathcal{F}$를 $\mathbf{P}^n_k$ 위의 연접 층이라 하자.
$\mathcal{F}$가 $m$-정칙이면 곱셈 사상
$$
H^0(\mathbf{P}^n_k, \mathcal{F}(m)) \otimes_k
H^0(\mathbf{P}^n_k, \mathcal{O}(1))
\longrightarrow
H^0(\mathbf{P}^n_k, \mathcal{F}(m + 1))
$$
은 전사이다.
\end{lemma}

\begin{proof}
$k'/k$를 체의 확대라 하자. $\mathcal{F}'$를 Lemma
\ref{varieties-lemma-m-regular-extend-base-field}에서와 같이 두자.
Cohomology of Schemes, Lemma
\ref{coherent-lemma-flat-base-change-cohomology}에 의해 보조정리의
선형사상을 $k'$로 기저변환하면 층 $\mathcal{F}'$에 대한 같은
선형사상을 얻는다. $k \to k'$가 충실히 평탄하므로 보조정리를
$k'$ 위에서 증명하면 충분하다. 즉, $k$가 무한이라고 가정할 수 있다.

\medskip\noindent
$k$가 무한이라고 가정하자. $n$에 관한 귀납법으로 보조정리를 증명한다.
$n = 0$인 경우에는 $\mathcal{O}(1) \cong \mathcal{O}$가 $T_0$로
생성되므로 자명하다. $n > 0$일 때 Lemma
\ref{varieties-lemma-exact-sequence-induction}을 적용하고 그 완전열에
$\mathcal{O}(m + 1)$을 텐서하여
$$
0 \to \mathcal{F}(m) \xrightarrow{s} \mathcal{F}(m + 1) \to
i_*\mathcal{G}(m + 1) \to 0
$$
을 얻는다. Remark \ref{varieties-remark-exact-sequence-induction-cohomology}을
보라. $t \in H^0(\mathbf{P}^n_k, \mathcal{F}(m + 1))$라 하자.
귀납가정에 의해 그 상
$\overline{t} \in H^0(H, \mathcal{G}(m + 1))$은
$\sum g_i \otimes \overline{s}_i$의 상이다. 여기서
$\overline{s}_i \in \Gamma(H, \mathcal{O}(1))$이고
$g_i \in H^0(H, \mathcal{G}(m))$이다. $\mathcal{F}$가 $m$-정칙이므로
$H^1(\mathbf{P}^n_k, \mathcal{F}(m - 1)) = 0$이다. 따라서 짧은 완전열
$$
0 \to \mathcal{F}(m - 1) \xrightarrow{s} \mathcal{F}(m) \to
i_*\mathcal{G}(m) \to 0
$$
에 결부된 긴 코호몰로지 완전열은 $g_i$를
$f_i \in H^0(\mathbf{P}^n_k, \mathcal{F}(m))$로 올릴 수 있음을 보인다.
또한 $\overline{s}_i$를
$s_i \in H^0(\mathbf{P}^n_k, \mathcal{O}(1))$로 올릴 수 있다
(예를 들어 Lemma \ref{varieties-lemma-hyperplane}의 증명을 보라).
$\sum f_i \otimes s_i$의 상을 $t$에서 빼면
$\overline{t} = 0$이라고 가정할 수 있음을 알 수 있다. 그러나 이는
정확히 $t$가 $f \otimes s$의 상이라는 뜻이다. 여기서
$f \in H^0(\mathbf{P}^n_k, \mathcal{F}(m))$이고, 이것이 원하는 바이다.
\end{proof}

\begin{lemma}
\label{varieties-lemma-m-regular-globally-generated}
$k$를 체라 하자. $n \geq 0$이라 하자.
$\mathcal{F}$를 $\mathbf{P}^n_k$ 위의 연접 층이라 하자.
$\mathcal{F}$가 $m$-정칙이면 $\mathcal{F}(m)$은 대역적으로 생성된다.
\end{lemma}

\begin{proof}
모든 $d \gg 0$에 대해 층 $\mathcal{F}(d)$는 대역적으로 생성된다.
이는 예를 들어 Cohomology of Schemes, Lemma
\ref{coherent-lemma-coherent-projective}의 첫째 부분에서 따른다.
$d \geq m$이고 $\mathcal{F}(d)$가 대역적으로 생성되도록 값을 택하자.
기저 $f_1, \ldots, f_r \in H^0(\mathbf{P}^n_k, \mathcal{F})$를 택하자.
Lemma \ref{varieties-lemma-m-regular-multiply}에 의해 모든 원소
$f \in H^0(\mathbf{P}^n_k, \mathcal{F}(d))$는
$f = \sum P_if_i$로 쓸 수 있다. 여기서
$P_i \in k[T_0, \ldots, T_n]$는 차수 $d - m$인 동차다항식이다.
단면들 $f$가 $\mathcal{F}(d)$를 생성하므로 단면들 $f_i$가
$\mathcal{F}(m)$을 생성한다.
\end{proof}

\subsection{힐베르트 다항식}
\label{varieties-subsection-hilbert}

\noindent
다음 보조정리는 더 일반적인 Lemma
\ref{varieties-lemma-numerical-polynomial-from-euler}에 의해 불필요해질 것이다.

\begin{lemma}
\label{varieties-lemma-hilbert-polynomial}
$k$를 체라 하자. $n \geq 0$이라 하자. $\mathcal{F}$를
$\mathbf{P}^n_k$ 위의 연접 층이라 하자. 함수
$$
d \longmapsto \chi(\mathbf{P}^n_k, \mathcal{F}(d))
$$
는 다항식이다.
\end{lemma}

\begin{proof}
$n$에 관한 귀납법으로 증명한다. $n = 0$이면
$\mathbf{P}^n_k = \Spec(k)$이고 $\mathcal{F}(d) = \mathcal{F}$이다.
따라서 이 경우 함수는 상수, 즉 차수 $0$의 다항식이다. $n > 0$이라고
가정하자. Lemma \ref{varieties-lemma-euler-characteristic-extend-base-field}에
의해 $k$가 무한이라고 가정할 수 있다. Lemma
\ref{varieties-lemma-exact-sequence-induction}을 적용하자. 꼬인 완전열
$0 \to \mathcal{F}(d - 1) \to \mathcal{F}(d) \to i_*\mathcal{G}(d) \to 0$에
Lemma \ref{varieties-lemma-euler-characteristic-additive}를 적용하면
$$
\chi(\mathbf{P}^n_k, \mathcal{F}(d)) -
\chi(\mathbf{P}^n_k, \mathcal{F}(d - 1)) =
\chi(H, \mathcal{G}(d))
$$
을 얻는다. Remark
\ref{varieties-remark-exact-sequence-induction-cohomology}을 보라.
$H \cong \mathbf{P}^{n - 1}_k$이므로 귀납가정에 의해 우변은
다항식이다. 이제 함수 $f : \mathbf{Z} \to \mathbf{Z}$가 사상
$d \mapsto f(d) - f(d - 1)$이 다항식이라는 성질을 가지면
그 함수 자체도 다항식이라는 사실로부터 보조정리가 끝난다. 이 사실의
증명은 생략한다(힌트: Algebra, Lemma
\ref{algebra-lemma-numerical-polynomial}과 비교하라).
\end{proof}

\begin{definition}
\label{varieties-definition-hilbert-polynomial}
$k$를 체라 하자. $n \geq 0$이라 하자. $\mathcal{F}$를
$\mathbf{P}^n_k$ 위의 연접 층이라 하자. 함수
$d \mapsto \chi(\mathbf{P}^n_k, \mathcal{F}(d))$를
$\mathcal{F}$의 {\it 힐베르트 다항식}이라 한다.
\end{definition}

\noindent
힐베르트 다항식의 계수는 $\mathbf{Q}$에 속하지만 일반적으로
$\mathbf{Z}$에는 속하지 않는다. 예를 들어 $\mathcal{O}_{\mathbf{P}^n_k}$의
힐베르트 다항식은
$$
d \longmapsto {d + n \choose n} = \frac{d^n}{n!} + \ldots
$$
이다. 이는 다음 보조정리와
$$
H^0(\mathbf{P}^n_k, \mathcal{O}_{\mathbf{P}^n_k}(d)) = k[T_0, \ldots, T_n]_d
$$
라는 사실에서 따른다. 이 공간은 차수 $d$ 부분이며, $k$ 위의 차원은
${d + n \choose n}$이다.

\begin{lemma}
\label{varieties-lemma-hilbert-polynomial-H0}
$k$를 체라 하자. $n \geq 0$이라 하자. $\mathcal{F}$를
$\mathbf{P}^n_k$ 위의 연접 층이라 하고, 그 힐베르트 다항식을
$P \in \mathbf{Q}[t]$라 하자. 그러면
$$
P(d) = \dim_k H^0(\mathbf{P}^n_k, \mathcal{F}(d))
$$
가 모든 $d \gg 0$에 대해 성립한다.
\end{lemma}

\begin{proof}
이는 $\mathcal{F}$의 충분히 높은 꼬임의 코호몰로지가 소멸한다는
사실에서 따른다. Cohomology of Schemes, Lemma
\ref{coherent-lemma-coherent-projective}를 보라.
\end{proof}



\subsection{몫들의 유계성}
\label{varieties-subsection-boundedness}

\noindent
이 소절에서는 $\mathbf{P}^n$ 위에 주어진 연접 층의 몫들의 정칙성을
힐베르트 다항식으로 한정한다.

\begin{lemma}
\label{varieties-lemma-bound-quotients-free}
$k$를 체라 하자. $n \geq 0$이라 하자. $r \geq 1$이라 하자.
$P \in \mathbf{Q}[t]$라 하자. 정수 $m$이 존재하며 이는
$n$, $r$, $P$에 의존하고 다음 성질을 가진다. $\mathbf{P}^n_k$ 위의 연접 층들의 짧은
완전열
$$
0 \to \mathcal{K} \to \mathcal{O}^{\oplus r} \to \mathcal{F} \to 0
$$
이 주어지고 $\mathcal{F}$의 힐베르트 다항식이 $P$이면,
$\mathcal{K}$는 $m$-정칙이다.
\end{lemma}

\begin{proof}
$n$에 관한 귀납법으로 증명한다. $n = 0$이면
$\mathbf{P}^n_k = \Spec(k)$이고 모든 연접 가군은 $0$-정칙이며,
모든 전사 사상은 대역단면 위에서도 전사이다. $n > 0$이라고 가정하자.
보조정리에 나오는 완전열을 생각하자. $P' \in \mathbf{Q}[t]$를
$P'(t) = P(t) - P(t - 1)$인 다항식이라 하자. $m'$를
$n - 1$, $r$, $P'$에 대해 성립하는 정수라 하자. Lemmas
\ref{varieties-lemma-m-regular-extend-base-field} 및
\ref{varieties-lemma-euler-characteristic-extend-base-field}에 의해 $k$를
체의 확대로 바꿀 수 있으므로 $k$가 무한이라고 가정할 수 있다.
연접 층 $\mathcal{F}$에 Lemma
\ref{varieties-lemma-exact-sequence-induction}을 적용하자.
$\mathcal{F}' = i^*\mathcal{F}$의 힐베르트 다항식은 $P'$이다
(Lemma \ref{varieties-lemma-hilbert-polynomial}의 증명을 보라).
$i^*$가 오른쪽 완전이므로 $\mathcal{F}'$는
$\mathcal{O}_H^{\oplus r} = i^*\mathcal{O}^{\oplus r}$의 몫이다.
따라서 귀납가정은 $\mathcal{F}'$에 적용된다. 그 바탕은
$H \cong \mathbf{P}^{n - 1}_k$이다(Lemma \ref{varieties-lemma-hyperplane}).
$\mathcal{K}(-1) \to \mathcal{K}$는 단사임에 주의하라. 실제로
$\mathcal{K} \subset \mathcal{O}^{\oplus r}$이다. 그 여핵은
$i_*\mathcal{H}$이고, 여기서 $\mathcal{H} = i^*\mathcal{K}$이다.
뱀 보조정리(Homology, Lemma \ref{homology-lemma-snake})에 의해
열과 행이 완전인 다음 가환 그림을 얻는다.
$$
\xymatrix{
& 0 \ar[d] & 0 \ar[d] & 0 \ar[d] \\
0 \ar[r] &
\mathcal{K}(-1) \ar[r] \ar[d] &
\mathcal{O}^{\oplus r}(-1) \ar[r] \ar[d] &
\mathcal{F}(-1) \ar[d] \ar[r] & 0 \\
0 \ar[r] &
\mathcal{K} \ar[r] \ar[d] &
\mathcal{O}^{\oplus r} \ar[r] \ar[d] &
\mathcal{F} \ar[d] \ar[r] & 0\\
0 \ar[r] &
i_*\mathcal{H} \ar[r] \ar[d] &
i_*\mathcal{O}_H^{\oplus r} \ar[r] \ar[d] &
i_*\mathcal{F}' \ar[r] \ar[d] & 0 \\
& 0 & 0 & 0
}
$$
따라서 귀납가정은 완전열
$0 \to \mathcal{H} \to \mathcal{O}_H^{\oplus r} \to \mathcal{F}' \to 0$에
적용된다. 이는 $H \cong \mathbf{P}^{n - 1}_k$ 위의 열이다
(Lemma \ref{varieties-lemma-hyperplane}). 그러므로 $\mathcal{H}$는
$m'$-정칙이다. 이는 $\mathcal{H}$가 $d$-정칙임을 함의하며,
모든 $d \geq m'$에 대해 그러하다는 것을 상기하라
(Lemma \ref{varieties-lemma-m-regular-up}).

\medskip\noindent
$i \geq 2$이고 $d \geq m'$이라 하자. 위 그림의 왼쪽 열에 결부된
긴 코호몰로지 완전열과 $H^{i - 1}(H, \mathcal{H}(d))$의 소멸로부터
사상
$$
H^i(\mathbf{P}^n_k, \mathcal{K}(d - 1))
\longrightarrow
H^i(\mathbf{P}^n_k, \mathcal{K}(d))
$$
이 단사임을 얻는다. 이 군들은 $d \gg 0$이면 영이므로
(Cohomology of Schemes, Lemma
\ref{coherent-lemma-coherent-projective}),
$H^i(\mathbf{P}^n_k, \mathcal{K}(d))$가 모든 $d \geq m'$ 및
$i \geq 2$에 대해 영이라고 결론짓는다.

\medskip\noindent
아직 $H^1$을 제어해야 한다. 먼저 모든 사상
$$
H^1(\mathbf{P}^n_k, \mathcal{K}(m' - 1)) \to
H^1(\mathbf{P}^n_k, \mathcal{K}(m')) \to
H^1(\mathbf{P}^n_k, \mathcal{K}(m' + 1)) \to \ldots
$$
이 전사임을 관찰한다. 이는 $H^1(H, \mathcal{H}(d))$가
$d \geq m'$에 대해 소멸하기 때문이다. $d > m'$이고 사상
$$
H^1(\mathbf{P}^n_k, \mathcal{K}(d - 1))
\longrightarrow
H^1(\mathbf{P}^n_k, \mathcal{K}(d))
$$
이 단사라고 하자. 그러면
$H^0(\mathbf{P}^n_k, \mathcal{K}(d)) \to H^0(H, \mathcal{H}(d))$는
전사이다. 가환 그림
$$
\xymatrix{
H^0(\mathbf{P}^n_k, \mathcal{K}(d)) \otimes_k
H^0(\mathbf{P}^n_k, \mathcal{O}(1))
\ar[r] \ar[d] &
H^0(\mathbf{P}^n_k, \mathcal{K}(d + 1)) \ar[d] \\
H^0(H, \mathcal{H}(d)) \otimes_k
H^0(H, \mathcal{O}_H(1))
\ar[r] &
H^0(H, \mathcal{H}(d + 1))
}
$$
을 생각하자. Lemma \ref{varieties-lemma-m-regular-multiply}에 의해 아래쪽
수평 화살표는 전사이다. 따라서 오른쪽 수직 화살표도 전사이다.
그러므로
$$
H^1(\mathbf{P}^n_k, \mathcal{K}(d))
\longrightarrow
H^1(\mathbf{P}^n_k, \mathcal{K}(d + 1))
$$
이 단사이다. 귀납법으로
$$
H^1(\mathbf{P}^n_k, \mathcal{K}(d - 1)) \to
H^1(\mathbf{P}^n_k, \mathcal{K}(d)) \to
H^1(\mathbf{P}^n_k, \mathcal{K}(d + 1)) \to \ldots
$$
가 모두 단사임을 안다. 이 군들이 결국 소멸하므로
$H^1(\mathbf{P}^n_k, \mathcal{K}(d - 1)) = 0$이라고 결론짓는다.
따라서 군들 $H^1(\mathbf{P}^n_k, \mathcal{K}(d))$의 차원은
$d \geq m'$에 대해 영이 될 때까지 엄밀히 감소한다. 그러므로
$\mathcal{K}$의 정칙성은
$m' + \dim_k H^1(\mathbf{P}^n_k, \mathcal{K}(m'))$로 한정된다.
한편 고차 코호몰로지 군들의 소멸에 의해
$$
\dim_k H^1(\mathbf{P}^n_k, \mathcal{K}(m')) = 
- \chi(\mathbf{P}^n_k, \mathcal{K}(m')) +
\dim_k H^0(\mathbf{P}^n_k, \mathcal{K}(m'))
$$
이다. 여기서 $H^0$의 차원은
$H^0(\mathbf{P}^n_k, \mathcal{O}^{\oplus r}(m'))$의 차원으로
한정되며, 이는 기껏해야 $r{n + m' \choose n}$이다. 여기서는
$m' > 0$인 경우를 말하며, 그렇지 않으면 영이다. 마지막으로
$\chi(\mathbf{P}^n_k, \mathcal{K}(m'))$ 항은
$r{n + m' \choose n} - P(m')$과 같다. 이로써 원하는 유형의
상계를 얻고 보조정리의 증명이 끝난다.
\end{proof}






\section{프로베니우스 사상}
\label{varieties-section-frobenius}

\noindent
$p$를 소수라 하자. 스킴 $X$에 대하여 “$X$의 표수는 $p$이다”,
“$X$는 표수 $p$이다”, 또는 “$X$는 표수 $p$ 안에 있다”라고 말한다.
이는 $p$가 $\mathcal{O}_X$에서 영일 때이다.

\begin{definition}
\label{varieties-definition-absolute-frobenius}
$p$를 소수라 하자. $X$를 표수 $p$인 스킴이라 하자.
{\it $X$의 절대 프로베니우스}는 사상 $F_X : X \to X$이다.
이 사상은 바탕 위상공간에서는 항등사상이고,
$F_X^\sharp : \mathcal{O}_X \to \mathcal{O}_X$는 $g \mapsto g^p$로 주어진다.
\end{definition}

\noindent
이 정의가 가능한 까닭은 임의의 환 $A$가 표수 $p$일 때 사상
$F_A : A \to A$, $a \mapsto a^p$가 $\Spec(A)$ 위에 항등사상을
유도하는 환 자기준동형이기 때문이다. 더구나 $A$가 국소환이면
$F_A$는 국소 준동형이다. 이로부터 $X$의 절대 프로베니우스가 스킴의
범주에서 $X$의 자기준동형임을 알 수 있다. 실제로 절대 프로베니우스는
표수 $p$인 스킴의 범주 위 항등 함자에서 자기 자신으로 가는 사상을
정의한다.

\begin{lemma}
\label{varieties-lemma-frobenius-endomorphism-identity}
$p > 0$을 소수라 하자. $f : X \to Y$를 표수 $p$인 스킴들의
사상이라 하자. 그러면 그림
$$
\xymatrix{
X \ar[d]_f \ar[r]_{F_X} & X \ar[d]^f \\
Y \ar[r]^{F_Y} & Y
}
$$
은 가환이다.
\end{lemma}

\begin{proof}
다음 자명한 대수적 사실에서 따른다. $\varphi : A \to B$가
표수 $p$인 환들의 준동형이면
$\varphi(a^p) = \varphi(a)^p$이다.
\end{proof}

\begin{lemma}
\label{varieties-lemma-frobenius}
$p > 0$을 소수라 하자. $X$를 표수 $p$인 스킴이라 하자.
그러면 절대 프로베니우스 $F_X : X \to X$는 보편 위상동형사상이고
정수적이며, 순비분리 잉여체 확대를 유도한다.
\end{lemma}

\begin{proof}
이는 프로베니우스 자기준동형 $F_A : A \to A$에 대한 대응하는
결과들에서 따른다. 여기서 환 $A$의 표수는 $p > 0$이다. 예를 들어
Algebra, Section \ref{algebra-section-universal-homeomorphism}의
논의와 Lemma \ref{algebra-lemma-p-ring-map}을 보라.
\end{proof}

\noindent
고정된 바탕 위의 스킴을 다룰 때에는 프로베니우스 사상의 상대적
형태가 있다.

\begin{definition}
\label{varieties-definition-relative-frobenius}
$p > 0$을 소수라 하자. $S$를 표수 $p$인 스킴이라 하자.
$X$를 $S$ 위의 스킴이라 하자. 다음과 같이 정의한다.
$$
X^{(p)} = X^{(p/S)} = X \times_{S, F_S} S
$$
이를 $S$ 위의 스킴으로 본다. Lemma
\ref{varieties-lemma-frobenius-endomorphism-identity}을 적용하면 다음 가환
그림에 들어맞는 유일한 사상 $F_{X/S} : X \longrightarrow X^{(p)}$가
존재하고, 이 사상은 $S$ 위의 사상임을 알 수 있다.
$$
\xymatrix{
X \ar[rr]_{F_{X/S}} \ar[rrd] \ar@/^1em/[rrrr]^{F_X}
& & X^{(p)} \ar[rr] \ar[d] & & X \ar[d] \\
& & S \ar[rr]^{F_S} & & S
}
$$
여기서 오른쪽 사각형은 데카르트 사각형이다. 사상 $F_{X/S}$를
{\it $X/S$의 상대 프로베니우스 사상}이라 한다.
\end{definition}

\noindent
$X \mapsto X^{(p)}$는 함자임에 주의하라. 이는 절대 프로베니우스
사상 $F_S : S \to S$에 대한 기저변환 함자이다. 상대 프로베니우스
사상에 관해서도 앞과 같은 보조정리들이 성립한다.

\begin{lemma}
\label{varieties-lemma-relative-frobenius-endomorphism-identity}
$p > 0$을 소수라 하자. $S$를 표수 $p$인 스킴이라 하자.
$f : X \to Y$를 $S$ 위 스킴들의 사상이라 하자. 그러면 그림
$$
\xymatrix{
X \ar[d]_f \ar[r]_{F_{X/S}} & X^{(p)} \ar[d]^{f^{(p)}} \\
Y \ar[r]^{F_{Y/S}} & Y^{(p)}
}
$$
은 가환이다.
\end{lemma}

\begin{proof}
Lemma \ref{varieties-lemma-frobenius-endomorphism-identity}과 정의들에서 따른다.
\end{proof}

\begin{lemma}
\label{varieties-lemma-relative-frobenius}
$p > 0$을 소수라 하자. $S$를 표수 $p$인 스킴이라 하자.
$X$를 $S$ 위의 스킴이라 하자. 그러면 상대 프로베니우스
$F_{X/S} : X \to X^{(p)}$는 보편 위상동형사상이고 정수적이며,
순비분리 잉여체 확대를 유도한다.
\end{lemma}

\begin{proof}
Lemma \ref{varieties-lemma-frobenius}에 의해 사상 $F_X : X \to X$와 기저변환
$h : X^{(p)} \to X$는 보편 위상동형사상이다. 여기서 후자는
$F_S$의 기저변환이다.
$h \circ F_{X/S} = F_X$이므로 $F_{X/S}$가 보편 위상동형사상이라고
결론짓는다(Morphisms, Lemma
\ref{morphisms-lemma-permanence-universal-homeomorphism}).
Morphisms, Lemmas \ref{morphisms-lemma-universal-homeomorphism} 및
\ref{morphisms-lemma-universally-injective}에 의해 $F_{X/S}$가
나머지 성질들도 가짐을 알 수 있다.
\end{proof}

\begin{lemma}
\label{varieties-lemma-relative-frobenius-omega}
$p > 0$을 소수라 하자. $S$를 표수 $p$인 스킴이라 하자.
$X$를 $S$ 위의 스킴이라 하자. 그러면
$\Omega_{X/S} = \Omega_{X/X^{(p)}}$이다.
\end{lemma}

\begin{proof}
이는 다음 대수적 사실로 옮겨진다. $A \to B$를 표수 $p$인 환들의
준동형이라 하자. $B' = B \otimes_{A, F_A} A$로 두고 환 준동형
$F_{B/A} : B' \to B$, $b \otimes a \mapsto b^pa$를 생각하자.
그러면 우리의 주장은 $\Omega_{B/A} = \Omega_{B/B'}$라는 것이다.
$\text{d}(b^pa) = 0$이다. 여기서
$\text{d} : B \to \Omega_{B/A}$는 보편 미분이며, 따라서 $\text{d}$는
$B'$-미분이므로 이 주장은 참이다.
\end{proof}

\begin{lemma}
\label{varieties-lemma-relative-frobenius-finite}
$p > 0$을 소수라 하자. $S$를 표수 $p$인 스킴이라 하자.
$X$를 $S$ 위의 스킴이라 하자. $X \to S$가 국소 유한형이면
$F_{X/S}$는 유한이다.
\end{lemma}

\begin{proof}
이는 다음 대수적 사실로 옮겨진다. $A \to B$를 표수 $p$인 환들의
유한형 준동형이라 하자. $B' = B \otimes_{A, F_A} A$로 두고 환 준동형
$F_{B/A} : B' \to B$, $b \otimes a \mapsto b^pa$를 생각하자.
그러면 우리의 주장은 $F_{B/A}$가 유한이라는 것이다. 실제로
$x_1, \ldots, x_n \in B$가 $A$ 위의 생성원들이면
$x_i$는 $B'$ 위에서 정수적이다. 실제로
$x_i^p = F_{B/A}(x_i \otimes 1)$이다.
따라서 Algebra, Lemma
\ref{algebra-lemma-characterize-finite-in-terms-of-integral}에 의해
$F_{B/A} : B' \to B$는 유한이다.
\end{proof}

\begin{lemma}
\label{varieties-lemma-geometrically-reduced-p}
$k$를 표수 $p > 0$인 체라 하자. $X$를 $k$ 위의 스킴이라 하자.
그러면 $X$가 기하적으로 기약인 것과 $X^{(p)}$가 기약인 것은 동치이다.
\end{lemma}

\begin{proof}
절대 프로베니우스 $F_k : k \to k$를 생각하자. 그러면
$F_k(k) = k^p$이다. 다시 말해 $F_k : k \to k$는
$k$를 $k^{1/p}$에 넣는 매입과 동형이다. 따라서 보조정리는
Lemma \ref{varieties-lemma-geometrically-reduced}에서 따른다.
\end{proof}

\begin{lemma}
\label{varieties-lemma-inseparable-deg-p-smooth}
$k$를 표수 $p > 0$인 체라 하자. $X$를 $k$ 위의 대수다양체라 하자.
다음 조건들은 동치이다.
\begin{enumerate}
\item $X^{(p)}$는 기약이다.
\item $X$는 기하적으로 기약이다.
\item 공집합이 아닌 열린부분집합 $U \subset X$가 존재하고, 이는
$k$ 위에서 매끄럽다.
\end{enumerate}
이 경우 $X^{(p)}$는 $k$ 위의 대수다양체이고,
$F_{X/k} : X \to X^{(p)}$는 차수 $p^{\dim(X)}$의 유한 우세 사상이다.
\end{lemma}

\begin{proof}
Lemma \ref{varieties-lemma-geometrically-reduced-p}에서 (1)과 (2)가 동치임을
보았다. Lemma \ref{varieties-lemma-geometrically-reduced-dense-smooth-open}에서
(2)가 (3)을 함의함을 보았다. (3)이 성립하면 $U$는 기하적으로
기약이고(예를 들어 Lemma
\ref{varieties-lemma-geometrically-regular-smooth}를 보라), 따라서 Lemma
\ref{varieties-lemma-generic-points-geometrically-reduced}에 의해 $X$는
기하적으로 기약이다. 이로써 (1), (2), (3)이 동치임을 알 수 있다.

\medskip\noindent
(1), (2), (3)이 성립한다고 가정하자. $F_{X/k}$는 위상동형사상이므로
(Lemma \ref{varieties-lemma-relative-frobenius}), $X^{(p)}$가 대수다양체임을
알 수 있다. 이어서 Lemma \ref{varieties-lemma-relative-frobenius-finite}에
의해 $F_{X/k}$는 유한이다. 이 사상은 전사이므로 우세하다. 차수
(Morphisms, Definition \ref{morphisms-definition-degree})를
계산하려면 $F_{U/k} : U \to U^{(p)}$의 차수를 계산하면 충분하다.
실제로 Lemma \ref{varieties-lemma-relative-frobenius-endomorphism-identity}에
의해 $F_{U/k} = F_{X/k}|_U$이다. $U$를 조금 줄여서 에탈 사상
$h : U \to \mathbf{A}^n_k$가 존재한다고 가정할 수 있다. Morphisms,
Lemma \ref{morphisms-lemma-smooth-etale-over-affine-space}를 보라.
물론 $n = \dim(U)$이다. 실제로
$\mathbf{A}^n_k \to \Spec(k)$는 상대차원 $n$의 매끄러운 사상이고,
에탈 사상 $h$는 상대차원 $0$의 매끄러운 사상이며,
$U \to \Spec(k)$는 상대차원 $\dim(U)$의 매끄러운 사상이고,
상대차원들은 올바르게 더해진다(Morphisms, Lemma
\ref{morphisms-lemma-composition-relative-dimension-d}).
$h$는 대수다양체들의 생성적으로 유한인 우세 사상이고, 따라서
$\deg(h)$가 정의됨에 주의하라. Lemma
\ref{varieties-lemma-relative-frobenius-endomorphism-identity}에 의해 가환 그림
$$
\xymatrix{
U \ar[rr]_{F_{X/k}} \ar[d]_h & & U^{(p)} \ar[d]^{h^{(p)}} \\
\mathbf{A}^n_k \ar[rr]^{F_{\mathbf{A}^n_k/k}} & &
(\mathbf{A}^n_k)^{(p)}
}
$$
을 얻는다. $h^{(p)}$는 $h$의 기저변환이므로 역시 에탈이고, 따라서
$h^{(p)}$ 역시 대수다양체들의 생성적으로 유한인 우세 사상이다.
$h^{(p)}$의 차수는 확대
$k(X^{(p)})/k((\mathbf{A}^n_k)^{(p)})$의 차수이고, 이는 확대
$k(X)/k(\mathbf{A}^n_k)$의 차수와 같다. 왜냐하면 $h^{(p)}$는
$h$의 기저변환이기 때문이다(작은 세부사항은 생략한다).
차수의 곱셈성(Morphisms, Lemma
\ref{morphisms-lemma-degree-composition})에 의해
$F_{\mathbf{A}^n_k/k}$의 차수가 $p^n$임을 보이면 충분하다.
이를 보려면 $(\mathbf{A}^n_k)^{(p)} = \mathbf{A}^n_k$이고
$F_{\mathbf{A}^n_k/k}$가 각 좌표를 그 $p$제곱으로 보내는
사상으로 주어짐을 관찰하면 된다.
\end{proof}

\begin{remark}
\label{varieties-remark-n-fold-relative-frobenius}
$p > 0$을 소수라 하자. $S$를 표수 $p$인 스킴이라 하자.
$X$를 $S$ 위의 스킴이라 하자. $n \geq 1$에 대해
$$
X^{(p^n)} = X^{(p^n/S)} = X \times_{S, F_S^n} S
$$
로 두고 이를 $S$ 위의 스킴으로 본다. $X \mapsto X^{(p^n)}$는
함자임에 주의하라. Lemma
\ref{varieties-lemma-frobenius-endomorphism-identity}을 적용하면
$F_{X/S, n} = (F_X^n, \text{id}_S) : X \longrightarrow X^{(p^n)}$가
다음 가환 그림에 들어맞는 $S$ 위의 사상임을 알 수 있다.
$$
\xymatrix{
X \ar[rr]_{F_{X/S, n}} \ar[rrd] \ar@/^1em/[rrrr]^{F_X^n}
& & X^{(p^n)} \ar[rr] \ar[d] & & X \ar[d] \\
& & S \ar[rr]^{F_S^n} & & S
}
$$
여기서 오른쪽 사각형은 데카르트 사각형이다. 사상 $F_{X/S, n}$를
때때로 {\it $n$중 상대 프로베니우스 사상, 즉 $X/S$의 그러한 사상}이라 한다.
이는 다음 공식 때문에 자연스럽다.
$$
F_{X/S, n} =
F_{X^{(p^{n - 1})}/S} \circ \ldots \circ F_{X^{(p)}/S} \circ F_{X/S}
$$
이 공식은 $F_{X/S, n}$이 $n$개의 상대 프로베니우스의 합성임을 보인다.
또한
$$
F_{X^{(p^m)}/S} = F_{X^{(p^{m - 1})}/S}^{(p)} = \ldots = F_{X/S}^{(p^m)}
$$
이므로(세부사항은 생략한다), 다음 식도 얻는다.
$$
F_{X/S, n} =
F_{X/S}^{(p^{n - 1})} \circ \ldots \circ F_{X/S}^{(p)} \circ F_{X/S}
$$
\end{remark}




\section{1차원 환의 붙이기}
\label{varieties-section-glueing-dim-1}

\noindent
이 절에서는 분리 스킴의 여차원 $1$인 유한 개의 점들이 하나의 아핀
열린부분집합에 포함된다는 것을 증명하기 위한 몇 가지 대수적 준비를
한다.

\begin{situation}
\label{varieties-situation-glue}
여기서는 다음 환들의 가환 그림이 주어져 있다.
$$
\xymatrix{
A \ar[r] & K \\
R \ar[u] \ar[r] & B \ar[u]
}
$$
여기서 $K$는 체이고 $A$, $B$는 $K$의 부분환이며 분수체가
$K$이다. 마지막으로 $R = A \times_K B = A \cap B$이다.
\end{situation}

\begin{lemma}
\label{varieties-lemma-glue-valuation-ring}
Situation \ref{varieties-situation-glue}에서 $B$가 값매김환이라고 가정하자.
그러면 모든 단원 $u$에 대해, 그것이 $A$의 단원이면 $u \in R$이거나
$u^{-1} \in R$이다.
\end{lemma}

\begin{proof}
실제로 상 $c$를 생각하자. 이것은 $u$의 $K$에서의 상이다.
이 상이 $B$에 속하면 $u \in R$이다.
그렇지 않으면 $c^{-1} \in B$이고(Algebra, Lemma
\ref{algebra-lemma-valuation-ring-x-or-x-inverse}),
$u^{-1} \in R$이다.
\end{proof}

\noindent
다음 보조정리는 이후 자주 나타나는 조건
“$A \otimes B \to K$가 전사이다”의 의미를 설명한다.

\begin{lemma}
\label{varieties-lemma-glue-separated}
Situation \ref{varieties-situation-glue}에서 $A$가 차원 $1$의 뇌터 환이라고
가정하자. 다음 조건들은 동치이다.
\begin{enumerate}
\item $A \otimes B \to K$는 전사가 아니다.
\item 이산 값매김환 $\mathcal{O} \subset K$가 존재하고, 이는
$A$와 $B$를 모두 포함한다.
\end{enumerate}
\end{lemma}

\begin{proof}
(2)가 (1)을 함의함은 분명하다. 반대로 $A \otimes B \to K$가
전사가 아니면 그 상 $C \subset K$는 체가 아니므로 $C$는 영이 아닌
극대아이디얼 $\mathfrak m$을 가진다. 값매김환
$\mathcal{O} \subset K$를 택하자. 이는 $C_\mathfrak m$을 지배한다.
$A \subset \mathcal{O}$에
Algebra, Lemma \ref{algebra-lemma-krull-akizuki}를 적용하면 환
$\mathcal{O}$는 뇌터 환이다. 따라서 Algebra, Lemma
\ref{algebra-lemma-valuation-ring-Noetherian-discrete}에 의해
$\mathcal{O}$는 이산 값매김환이다.
\end{proof}

\begin{lemma}
\label{varieties-lemma-semi-local}
Situation \ref{varieties-situation-glue}에서 다음을 가정하자.
\begin{enumerate}
\item $A$는 차원 $1$의 뇌터 반국소 정역이다.
\item $B$는 이산 값매김환이다.
\end{enumerate}
그러면 다음 두 가능성이 있다.
\begin{enumerate}
\item[(a)] $A^*$가 $R$에 포함되지 않으면
$\Spec(A) \to \Spec(R)$와 $\Spec(B) \to \Spec(R)$는
$\Spec(R)$을 덮는 열린 몰입들이고 $K = A \otimes_R B$이다.
\item[(b)] $A^*$가 $R$에 포함되면 $B$는 $A$의 한 극대아이디얼에서의
국소환을 지배하고, $A \otimes B \to K$는 전사가 아니다.
\end{enumerate}
\end{lemma}

\begin{proof}
가정 (a)에 의해 단원 $u$가 존재한다. 이는 $A$의 단원이고 $K$에서의 그 상이
$B$의 극대아이디얼에 속한다. 그러면 $u$는 $R$의 영인자가 아니고,
모든 $a \in A$에 대해 어떤 $n$이 존재하여 $u^n a \in R$이다.
따라서 $A = R_u$이다.

\medskip\noindent
$\mathfrak m_A$를 $A$의 제이콥슨 근기라 하자.
$x \in \mathfrak m_A$를 영이 아닌 원소라 하자. $\dim(A) = 1$이므로
$K = A_x$임을 알 수 있다. $x$를 $x^n u^m$으로 바꾸되, 어떤
$n \geq 1$ 및 $m \in \mathbf{Z}$를 택하면 $x$의
$B$에서의 상이 단원이라고 가정할 수 있다. 모든 $b \in B$에 대해
$x^nb$가 $R$의 상에 속하게 하는 어떤 $n$이 존재한다. 따라서
$B = R_x$이다.

\medskip\noindent
$z \in R$이라 하자. $z \not \in \mathfrak m_A$이고 $z$의 상이
$\mathfrak m_B$에 속하지 않으면 $z$는 가역이다. 따라서
$x + u$는 $R$에서 가역이다. 그러므로
$\Spec(R) = D(x) \cup D(u)$이다. 위에서
$D(u) = \Spec(A)$이고 $D(x) = \Spec(B)$임을 보았다.

\medskip\noindent
경우 (b)를 생각하자. $x \in \mathfrak m_A$이면 $1 + x$는 단원이고,
따라서 $1 + x \in R$, 즉 $x \in R$이다. 그러므로
$\mathfrak m_A \subset R \subset A$임을 알 수 있다. 실제로 이 경우
$A$는 $R$ 위에서 정수적이다. 구체적으로 체들의 곱
$A/\mathfrak m_A = \kappa_1 \times \ldots \times \kappa_n$으로 쓰자.
$x = (c_1, \ldots, c_r, 0, \ldots, 0)$가 원소이고
$c_i \not = 0$이라고 하자. 그러면
$$
x^2 - x(c_1, \ldots, c_r, 1, \ldots, 1)  = 0
$$
이다. $R$이 모든 단원을 포함하므로 $A/\mathfrak m_A$는 그 안에서
$R$의 상 위에 정수적이고, 따라서 $A$는 $R$ 위에 정수적이다.
$R \subset A \subset B$이다. 이는 $B$가 정수적으로 닫혀 있기 때문이다.
더구나 영이 아닌 $x \in \mathfrak m_A$에 대해
$K = A_x = \bigcup x^{-n}A = \bigcup x^{-n}R$이다. 따라서
$x^{-1} \not \in B$, 즉 $x \in \mathfrak m_B$이다. 그러므로
$\mathfrak m_A \subset \mathfrak m_B$이다. 이에 따라
$A \cap \mathfrak m_B$는 $A$의 극대아이디얼이며, 증명이 끝난다.
\end{proof}

\begin{lemma}
\label{varieties-lemma-semi-local-dimension-one-conductor}
$B$를 차원 $1$의 뇌터 반국소 정역이라 하자. $B'$를 그 분수체
안에서 $B$의 정수적 폐포라 하자. 그러면 $B'$는 반국소 데데킨트
정역이다. $x$를 $B'$의 제이콥슨 근기에 속하는 영이 아닌 원소라
하자. 그러면 모든 $y \in B'$에 대해 어떤 $n$이 존재하여
$x^n y \in B$이다.
\end{lemma}

\begin{proof}
$\mathfrak m_B$를 $B$의 제이콥슨 근기라 하자. $B'$의 구조는
Algebra, Lemma \ref{algebra-lemma-integral-closure-Dedekind}에서
따른다. 보조정리의 명제에서와 같은 $x, y \in B'$가 주어졌다고 하고,
부분환 $B \subset A \subset B'$를 생각하자. 이는 $x$와 $y$로 생성된다.
그러면 $A$는 $B$ 위에서 유한이다(Algebra, Lemma
\ref{algebra-lemma-characterize-finite-in-terms-of-integral}).
$B$와 $A$의 분수체가 같으므로 유한 가군 $A/B$의 지지는
$B$의 닫힌점들의 집합에 놓인다. 따라서
$\mathfrak m_B^n A \subset B$이며, 여기서 $n$은 알맞게 택한다. 또한
$\Spec(B') \to \Spec(A)$는 전사이므로(Algebra, Lemma
\ref{algebra-lemma-integral-overring-surjective}), $A$도 반국소이다.
또한 $x$가 제이콥슨 근기 $\mathfrak m_A$에 속함도 따른다. 이는
$A$의 제이콥슨 근기이다.
$\mathfrak m_A = \sqrt{\mathfrak m_B A}$임에 주의하라.
그러므로 $x^m y \in \mathfrak m_B A$가 되는 어떤 $m$이 존재한다.
따라서 $x^{nm} y \in B$이다.
\end{proof}

\begin{lemma}
\label{varieties-lemma-semi-local-both-side}
Situation \ref{varieties-situation-glue}에서 다음을 가정하자.
\begin{enumerate}
\item $A$는 차원 $1$의 뇌터 반국소 정역이다.
\item $B$는 차원 $1$의 뇌터 반국소 정역이다.
\item $A \otimes B \to K$는 전사이다.
\end{enumerate}
그러면 $\Spec(A) \to \Spec(R)$와 $\Spec(B) \to \Spec(R)$는
$\Spec(R)$을 덮는 열린 몰입들이고 $K = A \otimes_R B$이다.
\end{lemma}

\begin{proof}
먼저 특수한 경우로 $B$가 $K$ 안에서 정수적으로 닫혀 있다고 하자.
이는 $B$가 데데킨트 정역이라는 뜻이다(Algebra, Lemma
\ref{algebra-lemma-characterize-Dedekind}). 따라서 극대아이디얼에서의
모든 국소화는 이산 값매김환이다. $\mathfrak m_1, \ldots, \mathfrak m_r$을
$B$의 극대아이디얼들이라 하자. 다음과 같이 둔다.
$$
R_1 = A \times_K B_{\mathfrak m_1}
$$
$A \otimes_{R_1} B_{\mathfrak m_1} \to K$가 전사임을 관찰하면,
Lemma \ref{varieties-lemma-semi-local}에 의해 $A$와
$B_{\mathfrak m_1}$이 $\Spec(R_1)$을 덮는 열린 부분스킴들을 정의하고
$K = A \otimes_{R_1} B_{\mathfrak m_1}$임을 알 수 있다. 특히
$R_1$은 차원 $1$의 뇌터 반국소 환이다. 귀납적으로
$$
R_{i + 1} = R_i \times_K B_{\mathfrak m_{i + 1}}
$$
로 정의한다. 여기서 $i = 1, \ldots, r - 1$이다.
$R = R_r$임에 주의하라. 이는
$B = B_{\mathfrak m_1} \cap \ldots \cap B_{\mathfrak m_r}$이기 때문이다
(Algebra, Lemma
\ref{algebra-lemma-normal-domain-intersection-localizations-height-1}를
보라). 귀납적 절차로부터 $R \to A$가 열린 몰입
$\Spec(A) \to \Spec(R)$을 정의함이 따른다. 한편 극대아이디얼
$\mathfrak n_i$들은 $R$의 이 열린집합에 속하지 않으며, 극대아이디얼
$\mathfrak m_i$들과 대응한다. 후자들은 $B$의 극대아이디얼들이다. 실제로 환 준동형
$R \to B$는 동형
$R_{\mathfrak n_i} \to B_{\mathfrak m_i}$를 정의한다
(세부사항은 생략한다. 힌트: 각 단계에서 정확히 하나의 극대아이디얼을
$\Spec(R_i)$에 더했다). 따라서 원하는 대로
$\Spec(B) \to \Spec(R)$는 열린 몰입이다.

\medskip\noindent
이제 일반적인 경우를 보자. $B' \subset K$를 $B$의 정수적 폐포라
하자. Lemma \ref{varieties-lemma-semi-local-dimension-one-conductor}를 보라.
그러면 특수한 경우를 $R' = A \times_K B'$에 적용할 수 있다.
$x \in R'$를 택하자. 이는 $A$의 극대아이디얼들에는 들어가지 않고
$B'$의 극대아이디얼들에는 들어간다(Algebra, Lemma
\ref{algebra-lemma-chinese-remainder}를 보라). Lemma
\ref{varieties-lemma-semi-local-dimension-one-conductor}에 의해 어떤 정수
$n$이 존재하여 $x^n \in R = A \times_K B$이다. $x$를 $x^n$으로
바꾸어 $x \in R$이라고 하자. 모든 $y \in R'$에 대해 어떤 정수
$n$이 존재하여 $x^n y \in R$이다. 한편 $R'_x = A$임은 분명하다.
따라서 $R_x = A$이다. $A$와 $B$의 역할을 바꾸어도 어떤
$y \in R$가 존재하여 $B = R_y$임을 얻는다. $x$와 $y$를 모두
가역화하면 $(0)$ 이외의 소아이디얼은 남지 않음에 주의하라. 따라서
$K = R_{xy} = R_x \otimes_R R_y$이다. 이것으로 증명이 끝난다.
\end{proof}

\begin{lemma}
\label{varieties-lemma-glue-a-bunch-of-local-rings}
$K$를 체라 하자. $A_1, \ldots, A_r \subset K$를 차원 $1$인 뇌터
반국소 환들이라 하고, 그 분수체가 $K$라고 하자.
$A_i \otimes A_j \to K$가 모든 $i \not = j$에 대해 전사이면,
뇌터 반국소 정역 $A \subset K$가 존재하고, 그 차원은 $1$이다.
이는 $A_1, \ldots, A_r$에 포함되고 다음 성질들을 가진다.
\begin{enumerate}
\item $A \to A_i$는 열린 몰입
$j_i : \Spec(A_i) \to \Spec(A)$를 유도한다.
\item $\Spec(A)$는 열린집합들 $j_i(\Spec(A_i))$의 합집합이다.
\item $\Spec(A)$의 각 닫힌점은 이 열린집합들 중 정확히 하나에 속한다.
\end{enumerate}
\end{lemma}

\begin{proof}
실제로 $A = A_1 \cap \ldots \cap A_r$로 둘 수 있다. 먼저 (1)과 (2)가
증명되었다고 하면 (3)은 가정 $A_i \otimes A_j \to K$가 전사라는
것에서 따른다. 왜냐하면
$\mathfrak m \in j_i(\Spec(A_i)) \cap j_j(\Spec(A_j))$이면
$A_i \otimes A_j \to K$의 상은 $A_\mathfrak m$에 들어가기 때문이다.
(1)과 (2)를 증명하기 위해 $r$에 관한 귀납법을 쓴다. $r > 1$이면
귀납가정에 의해 $B = A_2 \cap \ldots \cap A_r$에 대해 (1)과 (2)를
얻는다. 이어서 Lemma \ref{varieties-lemma-semi-local-both-side}를 적용하면
$A = A_1 \cap B$에 대해서도 이들이 성립한다.
\end{proof}

\begin{lemma}
\label{varieties-lemma-create-globally-generated}
$A$를 분수체 $K$를 가지는 정역이라 하자.
$B_1, \ldots, B_r \subset K$를 차원 $1$의 뇌터 반국소 정역들이라
하고, 그 분수체들이 $K$라고 하자. $A \otimes B_i \to K$가
$i = 1, \ldots, r$에 대해 전사이면 $x \in A$가 존재하여
$x^{-1}$은 $B_i$의 제이콥슨 근기에 속한다. 이는
$i = 1, \ldots, r$에 대해 성립한다.
\end{lemma}

\begin{proof}
$B_i'$를 $B_i$의 정수적 폐포라 하되, 이는 $K$ 안에서 취한다. 영이 아닌
$x \in A$를 찾아서 $x^{-1}$이 $B'_i$의 제이콥슨 근기에 속한다고 하자.
이는 $i = 1, \ldots, r$에 대해 성립한다고 가정한다. 그러면 Lemma
\ref{varieties-lemma-semi-local-dimension-one-conductor}에 의해 $x$를
거듭제곱으로 바꾼 뒤 $x^{-1} \in B_i$를 얻는다.
$\Spec(B'_i) \to \Spec(B_i)$가 전사이므로, 이때 $x^{-1}$은
$B_i$의 제이콥슨 근기에도 속함을 알 수 있다. 따라서 각 $B_i$가
반국소 데데킨트 정역이라고 가정할 수 있다.

\medskip\noindent
$B_i$가 국소환이 아니면 목록에서 $B_i$를 제거하고 유한한 국소환들의
모음 $(B_i)_\mathfrak m$을 다시 넣는다. 따라서 $B_i$가 이산
값매김환이라고 가정할 수 있고, 이는 $i = 1, \ldots, r$에 대해 성립한다.

\medskip\noindent
$v_i : K \to \mathbf{Z}$, $i = 1, \ldots, r$을 대응하는 이산
값매김들이라 하자(Algebra, Lemma
\ref{algebra-lemma-characterize-Dedekind}를 보라). 우리는 영이 아닌
$x \in A$를 찾아서 $v_i(x) < 0$이 되게 하려 한다. 이는
$i = 1, \ldots, r$에 대해 요구된다. 이를 $r$에 관한 귀납법으로 증명한다.

\medskip\noindent
$r = 1$이고 결론이 거짓이면 $A \subset B$이며, 사상
$A \otimes B \to K$는 전사가 아니다. 이는 모순이다.

\medskip\noindent
$r > 1$이면 귀납가정에 의해 영이 아닌 $x \in A$를 찾아서
$v_i(x) < 0$이 되게 할 수 있다. 이는 $i = 1, \ldots, r - 1$에 대해 성립한다.
$v_r(x) < 0$이면 끝났으므로 $v_r(x) \geq 0$이라고 가정할 수 있다.
기초 단계에 의해 영이 아닌 $y \in A$를 찾아서
$v_r(y) < 0$이 되게 할 수 있다. $x$를 거듭제곱으로 바꾸면
$v_i(x) < v_i(y)$라고 가정할 수 있고, 이는
$i = 1, \ldots, r - 1$에 대해 성립한다.
그러면 $x + y$가 우리가 찾는 원소이다.
\end{proof}

\begin{lemma}
\label{varieties-lemma-power-equal}
$A$를 차원 $1$의 뇌터 국소환이라 하자. $L = \prod A_\mathfrak p$로
두되, 곱은 $A$의 극소소아이디얼들에 걸쳐 취한다.
$a_1, a_2 \in \mathfrak m_A$의 $L$에서의 상이 같다고 하자.
그러면 $a_1^n = a_2^n$이고, 여기서 어떤 $n > 0$을 택할 수 있다.
\end{lemma}

\begin{proof}
$a_1 = a_2 + x$로 쓰자. 그러면 $x$의 $L$에서의 상은 영이다.
따라서 $x$는 $A$의 멱영원소이다. 왜냐하면
$\bigcap \mathfrak p$는 $(0)$의 근기이고, 소멸아이디얼 $I$는
$x$를 소멸시키며,
극대아이디얼의 한 거듭제곱을 포함하기 때문이다. 실제로 모든
극소소아이디얼에 대해 $\mathfrak p \not \in V(I)$이다.
$x^k = 0$이고 $\mathfrak m^n \subset I$라고 하자. 그러면
$$
a_1^{k + n} = a_2^{k + n} + {n + k \choose 1} a_2^{n + k - 1} x +
{n + k \choose 2} a_2^{n + k - 2} x^2 + \ldots +
{n + k \choose k - 1} a_2^{n + 1} x^{k - 1} = a_2^{n + k}
$$
이다. 이는 $a_2 \in \mathfrak m_A$이기 때문이다.
\end{proof}

\begin{lemma}
\label{varieties-lemma-power-works}
$A$를 차원 $1$의 뇌터 국소환이라 하자.
$L = \prod A_\mathfrak p$ 및 $I = \bigcap \mathfrak p$로 두되,
곱과 교집합은 $A$의 극소소아이디얼들에 걸쳐 취한다.
$f \in L$이 $f = i + a$ 꼴이라고 하자. 여기서
$a \in \mathfrak m_A$이고 $i \in IL$이다. 그러면 $f$의 어떤
거듭제곱은 $A \to L$의 상에 속한다.
\end{lemma}

\begin{proof}
$A$가 뇌터 환이므로 $I^t = 0$이며, 여기서 어떤 $t > 0$을 택할 수 있다.
$f = a + i$이고 $i \in I^kL$인 표현을 안다고 하자. 그러면
$f^n = a^n + na^{n - 1}i \bmod I^{k + 1}L$이다. 따라서
$na^{n - 1}i$가 $I^k \to I^kL$의 상에 속함을 보이면 충분하며,
이는 어떤 $n \gg 0$에 대해 보이면 된다. 이를 위해 $g \in A$를 택하여
$\mathfrak m_A = \sqrt{(g)}$가 되게 하자(Algebra, Lemma
\ref{algebra-lemma-height-1}). 그러면 예를 들어 Algebra,
Proposition \ref{algebra-proposition-dimension-zero-ring}에 의해
$L = A_g$이다. 한편 어떤 $n$에 대해 $a^n \in (g)$이다.
따라서 $L$의 원소들의 분모는 $a$의 높은 거듭제곱을 곱하여 없앨 수 있다.
\end{proof}

\begin{lemma}
\label{varieties-lemma-good-intersection}
$A$를 차원 $1$의 뇌터 국소환이라 하자. $L = \prod A_\mathfrak p$로
두되, 곱은 $A$의 극소소아이디얼들에 걸쳐 취한다.
$K \to L$을 정수적 환 준동형이라 하자. 그러면
$a \in \mathfrak m_A$와 $x \in K$가 존재하여 $L$에서 같은 원소로
가고, $\mathfrak m_A = \sqrt{(a)}$이다.
\end{lemma}

\begin{proof}
Lemma \ref{varieties-lemma-power-works}에 의해 $A$를
$A/(\bigcap \mathfrak p)$로 바꾸어 $A$가 기약환이라고 가정할 수 있다
(몇 가지 세부사항은 생략한다). 또한 $K$를 $K \to L$의 상으로
바꿀 수 있다. 그러면 $K$는 기약환이다. 사상
$\Spec(L) \to \Spec(K)$는 전사이고 닫힌 사상이다(세부사항은 생략한다).
따라서 $\Spec(K)$는 유한 이산공간이다. 그러므로 $K$는 체들의 유한곱이다.

\medskip\noindent
$\mathfrak p_j$, $j = 1, \ldots, m$을 $A$의 극소소아이디얼들이라
하자. $L_j$를 분수체라 하되 그 정역을 $A_j$라 두어서
$L = \prod_{j = 1, \ldots, m} L_j$가 되게 하자.
$A_j$를 $A/\mathfrak p_j$의 정규화라 하자. 그러면 $A_j$는 적어도
하나의 극대아이디얼을 가지는 반국소 데데킨트 정역이다. Algebra,
Lemma \ref{algebra-lemma-integral-closure-Dedekind}를 보라.
$n$을 $A_1, \ldots, A_m$에 있는 극대아이디얼들의 개수의 합이라 하자.
그러한 극대아이디얼 $\mathfrak m \subset A_j$에 대해 함수
$$
v_{\mathfrak m} : L \to L_j \to \mathbf{Z} \cup \{\infty\}
$$
를 생각하자. 여기서 두 번째 화살표는 이산 값매김환
$(A_j)_{\mathfrak m}$에 대응하는 이산 값매김을 $0$을 $\infty$로
보내도록 확장한 것이다. 이와 같이 $n$개의 함수
$v_1, \ldots, v_n : L \to \mathbf{Z} \cup \{\infty\}$를 얻는다.
원소 $x \in K$를 찾아서 $v_i(x) < 0$이 되게 할 것이다. 이는 모든
$i = 1, \ldots, n$에 대해 요구된다.

\medskip\noindent
먼저 각 $i$에 대해 원소 $x \in K$가 존재하여 $v_i(x) < 0$이라고
주장한다. 실제로 $v_i$가 $\mathfrak m \subset A_j$에 대응한다고
하자. $v_i(x) \geq 0$이 모든 $x \in K$에 대해 성립하면 $K$는
$(A_j)_{\mathfrak m}$으로 사상된다. 이는 $L_j$, 즉 $A_j$의 분수체
안에서의 진술이다. $K$의 $L_j$에서의 상은 체이고, 그 위에서
$L_j$는 대수적이다.
이는 Algebra, Lemma \ref{algebra-lemma-integral-under-field}에서
따른다. 이 두 사실을 합치면 Algebra, Lemma
\ref{algebra-lemma-valuation-ring-cap-field-finite}와 모순된다.

\medskip\noindent
원소 $x \in K$를 찾아서 $v_1(x) < 0, \ldots, v_r(x) < 0$이
되게 했다고 하자. 여기서 어떤 $r < n$을 택한다.
$v_{r + 1}(x) < 0$이면 $x$는
$r + 1$에 대해서도 조건을 만족한다. 그렇지 않으면 앞 문단의
결과에 따라 가능한 대로 어떤 $y \in K$를 택하여
$v_{r + 1}(y) < 0$이 되게 하자. $x$를 $x^n$으로 바꾸되
어떤 $n > 0$을 택하면, $v_i(x) < v_i(y)$라고 가정할 수 있다.
이는 $i = 1, \ldots, r$에 대해 성립한다. 그러면 값매김의 성질에 의해
$v_j(x + y) = v_j(x) < 0$이고, 이는 $j = 1, \ldots, r$에 대해
성립한다. 마찬가지로
$v_{r + 1}(x + y) = v_{r + 1}(y) < 0$이다. 귀납적으로 논증하여
$x \in K$를 얻으며, $v_i(x) < 0$은 $i = 1, \ldots, n$에 대해 성립한다.

\medskip\noindent
특히 원소 $x \in K$는 $K$의 각 인자에서 영이 아닌 사영을 가진다
($K$가 체들의 유한곱이고, $x$의 어떤 성분이 영이면 값들 중 하나
$v_i(x)$가 $\infty$가 됨을 상기하라). 따라서 $x$는 가역이고
$x^{-1} \in K$는 $\infty > v_i(x^{-1}) > 0$인 원소이다. 이는 모든
$i$에 대해 성립한다. Lemma
\ref{varieties-lemma-semi-local-dimension-one-conductor}에 의해 어떤
$e < 0$에 대해 원소 $x^e \in K$는
$\mathfrak m_A/\mathfrak p_j \subset A/\mathfrak p_j$의 원소로
사상된다. 이는 모든 $j = 1, \ldots, m$에 대해 성립한다. 사상
$\mathfrak m_A \to \prod \mathfrak m_A/\mathfrak p_j$의 여핵은
$\mathfrak m_A$의 한 거듭제곱에 의해 소멸됨에 주의하라. 따라서
$e$를 더 작은 $e$로 바꾸면 $a \in \mathfrak m_A$를 찾아서
$\mathfrak m_A/\mathfrak p_j$에서의 그 상이 $x^e$의 상과 같게 할 수
있다. 쌍 $(a, x^e)$은 보조정리의 결론을 만족한다.
\end{proof}

\begin{lemma}
\label{varieties-lemma-localization-semi-local}
$A$를 환이라 하자. $\mathfrak p_1, \ldots, \mathfrak p_r$을
$A$의 유한한 소아이디얼 집합이라 하자.
$S = A \setminus \bigcup \mathfrak p_i$로 두자. 그러면 $S$는
곱셈계이고 $S^{-1}A$는 반국소환이며, 그 극대아이디얼들은 집합
$\{\mathfrak p_i\}$의 극대원소들과 대응한다.
\end{lemma}

\begin{proof}
$a, b \in A$이고 $a, b \in S$이면
$a, b \not \in \mathfrak p_i$이므로 $ab \not \in \mathfrak p_i$이고,
따라서 $ab \in S$이다. 또한 $1 \in S$이다. 그러므로 $S$는
$A$의 곱셈적 부분집합이다. Algebra, Lemma
\ref{algebra-lemma-spec-localization}의 $\Spec(S^{-1}A)$에 대한
설명과 Algebra, Lemma \ref{algebra-lemma-silly}에 의해
$S^{-1}A$의 소아이디얼들은 $A$의 소아이디얼들 가운데 어떤
$\mathfrak p_i$에 포함되는 것들과 대응함을 알 수 있다. 따라서
$S^{-1}A$의 극대아이디얼들은 집합
$\{\mathfrak p_1, \ldots, \mathfrak p_r\}$의 포함관계에 관한
극대원소들과 일대일 대응한다.
\end{proof}




\section{1차원 뇌터 스킴}
\label{varieties-section-dimension-one}

\noindent
이 절의 주된 결과는 차원 $1$의 뇌터 분리 스킴이 풍부한 가역층을
가진다는 것이다. Proposition
\ref{varieties-proposition-dim-1-noetherian-separated-has-ample}을 보라.

\begin{lemma}
\label{varieties-lemma-affine}
$X$를 모든 국소환이 뇌터이고 차원이 $\leq 1$인 스킴이라 하자.
$U \subset X$를 레트로콤팩트 열린집합이라 하자. 포함사상을
$j : U \to X$로 나타내자. 그러면
$R^pj_*\mathcal{F} = 0$, $p > 0$이다. 이는 모든 준연접
$\mathcal{O}_U$-가군 $\mathcal{F}$에 대해 성립한다.
\end{lemma}

\begin{proof}
$R^pj_*\mathcal{F}$의 소멸은 줄기에서 확인할 수 있다.
$R^qj_*$의 형성은 평탄 기저변환과 가환한다. Cohomology of Schemes,
Lemma \ref{coherent-lemma-flat-base-change-cohomology}를 보라.
따라서 $X$가 차원 $\leq 1$인 뇌터 국소환의 스펙트럼이라고 가정할
수 있다. 이 경우 $X$는 닫힌점 $x$와 유한 개의 다른 점
$x_1, \ldots, x_n$을 가진다. 후자들은 $x$로 특수화되지만 서로에게는
특수화되지 않는다
(Algebra, Lemma
\ref{algebra-lemma-Noetherian-irreducible-components}를 보라).
$x \in U$이면 $U = X$이고 결과는 분명하다. 그렇지 않으면 점들의
번호를 바꾸어 $U = \{x_1, \ldots, x_r\}$라고 할 수 있으며, 여기서
어떤 $r$을 택한다. 그러면 $U$는 아핀이다(Schemes, Lemma
\ref{schemes-lemma-scheme-finite-discrete-affine}). 따라서 결과는
Cohomology of Schemes, Lemma
\ref{coherent-lemma-relative-affine-vanishing}에서 따른다.
\end{proof}

\begin{lemma}
\label{varieties-lemma-open-in-affine-curve-affine}
$X$를 모든 국소환이 뇌터이고 차원이 $\leq 1$인 아핀 스킴이라 하자.
그러면 임의의 준콤팩트 열린집합 $U \subset X$는 아핀이다.
\end{lemma}

\begin{proof}
포함사상을 $j : U \to X$로 나타내자. $\mathcal{F}$를 준연접
$\mathcal{O}_U$-가군이라 하자. Lemma \ref{varieties-lemma-affine}에 의해
고차 직상들 $R^pj_*\mathcal{F}$는 영이다.
$\mathcal{O}_X$-가군 $j_*\mathcal{F}$는 준연접이다
(Schemes, Lemma \ref{schemes-lemma-push-forward-quasi-coherent}).
따라서 Cohomology of Schemes, Lemma
\ref{coherent-lemma-quasi-coherent-affine-cohomology-zero}에 의해
그 고차 코호몰로지 군들은 소멸한다. Leray 스펙트럼 열
Cohomology, Lemma \ref{cohomology-lemma-apply-Leray}에 의해
$H^p(U, \mathcal{F}) = 0$이며, 이는 모든 $p > 0$에 대해 성립한다.
따라서 예를 들어
Cohomology of Schemes, Lemma
\ref{coherent-lemma-quasi-compact-h1-zero-covering}에 의해 $U$는 아핀이다.
\end{proof}

\begin{lemma}
\label{varieties-lemma-complement-codim-1-closed-points}
$X$를 스킴이라 하자. $U \subset X$를 열린집합이라 하자. 다음을
가정하자.
\begin{enumerate}
\item $U$는 $X$의 레트로콤팩트 열린집합이다.
\item $X \setminus U$는 이산집합이다.
\item $x \in X \setminus U$에 대해 국소환
$\mathcal{O}_{X, x}$는 차원 $\leq 1$인 뇌터 환이다.
\end{enumerate}
그러면 (1) 가역 $\mathcal{O}_X$-가군 $\mathcal{L}$과 단면 $s$가
존재하여 $U = X_s$이고, (2) 사상 $\Pic(X) \to \Pic(U)$는 전사이다.
\end{lemma}

\begin{proof}
$X \setminus U = \{x_i; i \in I\}$로 두자. 아핀 열린집합
$U_i \subset X$를 택하자. 이때 $x_i \in U_i$이고
$x_j \not \in U_i$이며, 후자는 $j \not = i$에 대해 요구한다.
조건 (2)에 의해 이것이 가능하다.
$U_i = \Spec(A_i)$라고 하자. $\mathfrak m_i \subset A_i$를
$x_i$에 대응하는 극대아이디얼이라 하자. 국소환들에 대한 가정에 의해
$\mathfrak q \subset \mathfrak m_i$이고
$\mathfrak q \not = \mathfrak m_i$인 소아이디얼은 유한 개뿐이다
(Algebra, Lemma
\ref{algebra-lemma-Noetherian-irreducible-components}를 보라).
따라서 소아이디얼 회피(Algebra, Lemma
\ref{algebra-lemma-silly})에 의해 $f_i \in \mathfrak m_i$를 택하여
그 어느 소아이디얼에도 들어가지 않게 할 수 있다. 그러면
$V(f_i) = \{\mathfrak m_i\} \amalg Z_i$이며, 여기서
$Z_i \subset U_i$는 어떤 닫힌부분집합이다. 실제로 $Z_i$는
$V(f_i)$의 레트로콤팩트 열린부분집합이고 특수화에 대해 닫혀 있다.
Algebra, Lemma
\ref{algebra-lemma-constructible-stable-specialization-closed}를 보라.
$U_i$를 줄여서 $V(f_i) = \{x_i\}$라고 가정할 수 있다. 그러면
$$
\mathcal{U} : X = U \cup \bigcup U_i
$$
는 $X$의 열린 덮개이다. $2$-코사이클을 생각하자. 그 값은
$\mathcal{O}_X^*$에 속하며, 이는 $f_i$로 $U \cap U_i$에서 주어지고,
$f_i/f_j$로 $U_i \cap U_j$에서 주어진다. 이는 선다발
$\mathcal{L}$을 정의한다. 또한 단면 $s$를 정의하되, $1$로
$U$에서, $f_i$로 $U_i$에서 정의하면 보조정리의 성질을 가진다.

\medskip\noindent
$\mathcal{N}$을 가역 $\mathcal{O}_U$-가군이라 하자.
$N_i$를 가역 $(A_i)_{f_i}$-가군으로서
$\mathcal{N}|_{U \cap U_i}$가 $\tilde N_i$와 같게 하는 것이라 하자.
$(A_{\mathfrak m_i})_{f_i}$는 아르틴 환임에 주의하라
(차원 영의 뇌터 환이므로, Algebra, Lemma
\ref{algebra-lemma-Noetherian-dimension-0}을 보라). 따라서 이는
국소환들의 곱이고(Algebra, Lemma
\ref{algebra-lemma-artinian-finite-length}), 그 피카르 군은 자명하다.
그러므로 $U_i$를 줄인 뒤, 즉 $A_i$를 $(A_i)_g$로 바꾼 뒤
(어떤 $g \in A_i$, $g \not \in \mathfrak m_i$에 대해),
$N_i = (A_i)_{f_i}$라고 가정할 수 있다. 즉,
$\mathcal{N}|_{U \cap U_i}$는 자명하다. 이 경우
$\mathcal{N}$을 $X$ 위의 가역층으로 연장하는 방법은 분명하다
(각 $U_i$ 위에서 자명한 가역 가군으로 연장하면 된다).
\end{proof}

\begin{lemma}
\label{varieties-lemma-find-globally-generated}
$X$를 정수적 분리 스킴이라 하자. $U \subset X$를 공집합이 아닌
아핀 열린집합이라 하고, $X \setminus U$가 유한한 점들의 집합
$x_1, \ldots, x_r$이라고 하자. 또한 $\mathcal{O}_{X, x_i}$는
차원 $1$의 뇌터 환이라고 하자. 그러면 대역적으로 생성되는 가역
$\mathcal{O}_X$-가군 $\mathcal{L}$과 단면 $s$가 존재하여
$U = X_s$이다.
\end{lemma}

\begin{proof}
$U = \Spec(A)$라고 하고 $K$를 $X$의 함수체라 하자.
$B_i = \mathcal{O}_{X, x_i}$ 및
$\mathfrak m_i = \mathfrak m_{x_i}$로 쓰자.
$x_i \not \in U$이므로 열린부분집합
$U \times_X \Spec(B_i)$는 $\Spec(B_i)$ 안에서 점 하나만을 가진다. 즉,
$U \times_X \Spec(B_i) = \Spec(K)$이다. $X$가 분리이므로
$\Spec(K)$는 $U \times \Spec(B_i)$의 닫힌 부분스킴이다. 즉, 사상
$A \otimes B_i \to K$는 전사이다. Lemma
\ref{varieties-lemma-create-globally-generated}에 의해 영이 아닌
$f \in A$를 찾아서 $f^{-1} \in \mathfrak m_i$가 되게 할 수 있다.
이는 $i = 1, \ldots, r$에 대해 성립한다. 열린집합
$x_i \in U_i \subset X$를 택하여
$f^{-1} \in \mathcal{O}(U_i)$가 되게 하자. 그러면
$$
\mathcal{U} : X = U \cup \bigcup U_i
$$
는 $X$의 열린 덮개이다. $2$-코사이클을 생각하자. 그 값은
$\mathcal{O}_X^*$에 속하며, 이는 $f$로 $U \cap U_i$에서 주어지고,
$1$로 $U_i \cap U_j$에서 주어진다. 이는 다음 두 단면을 가지는 선다발
$\mathcal{L}$을 정의한다.
\begin{enumerate}
\item 단면 $s$를 $1$로 $U$에서, $f^{-1}$로 $U_i$에서 정의하면
보조정리의 명제에서 말한 성질을 가진다.
\item 단면 $t$를 $f$로 $U$에서, $1$로 $U_i$에서 정의한다.
\end{enumerate}
$X_t \supset U_1 \cup \ldots \cup U_r$임에 주의하라. 따라서
$s, t$는 $\mathcal{L}$을 생성하고 보조정리가 증명되었다.
\end{proof}

\begin{lemma}
\label{varieties-lemma-enough-globally-generated-ample}
$X$를 준콤팩트 스킴이라 하자. 모든 $x \in X$에 대해 쌍
$(\mathcal{L}, s)$가 존재하되, 이는 대역적으로 생성되는 가역층
$\mathcal{L}$과 대역단면 $s$로 이루어지고 $x \in X_s$이며 $X_s$가 아핀이면,
$X$는 풍부한 가역층을 가진다.
\end{lemma}

\begin{proof}
$X$가 준콤팩트이므로 유한한 쌍들의 모음
$(\mathcal{L}_i, s_i)$, $i = 1, \ldots, n$을 택하여
$\mathcal{L}_i$가 대역적으로 생성되고 $X_{s_i}$가 아핀이며
$X = \bigcup X_{s_i}$가 되게 할 수 있다. 다시 $X$의 준콤팩트성에
의해 각 $i$에 대해 유한한 단면들의 모음
$t_{i, j}$를 $\mathcal{L}_i$에서 택하되, $j = 1, \ldots, m_i$로 두어
$X = \bigcup X_{t_{i, j}}$가 되게 할 수 있다.
$t_{i, 0} = s_i$로 두자. 가역층
$$
\mathcal{L} = \mathcal{L}_1 
\otimes_{\mathcal{O}_X} \ldots
\otimes_{\mathcal{O}_X} \mathcal{L}_n
$$
과 대역단면들
$$
\tau_J = t_{1, j_1} \otimes \ldots \otimes t_{n, j_n}
$$
을 생각하자. Properties, Lemma
\ref{properties-lemma-affine-cap-s-open}에 의해 열린집합
$X_{\tau_J}$는 $j_i = 0$이면 아핀이다. 여기서 조건은 어떤 $i$에
대해 성립하면 된다. 이 열린집합들이
$X$를 덮는다는 것은 쉽게 알 수 있다. 따라서 정의에 의해
$\mathcal{L}$은 풍부하다.
\end{proof}

\begin{lemma}
\label{varieties-lemma-dim-1-noetherian-integral-separated-has-ample}
$X$를 차원 $1$의 뇌터 정수적 분리 스킴이라 하자. 그러면
$X$는 풍부한 가역층을 가진다.
\end{lemma}

\begin{proof}
아핀 열린 덮개 $X = U_1 \cup \ldots \cup U_n$을 택하자.
$X$가 뇌터이므로 각 집합 $X \setminus U_i$는 유한하다. 따라서
Lemma \ref{varieties-lemma-find-globally-generated}에 의해 쌍
$(\mathcal{L}_i, s_i)$를 찾을 수 있다. 이는 대역적으로 생성되는
가역층 $\mathcal{L}_i$와 대역단면 $s_i$로 이루어지고
$U_i = X_{s_i}$를 만족한다.
Lemma \ref{varieties-lemma-enough-globally-generated-ample}에 의해
$X$가 풍부한 가역층을 가진다고 결론짓는다.
\end{proof}

\begin{lemma}
\label{varieties-lemma-surjective-pic-birational-finite}
$f : X \to Y$를 스킴들의 유한 사상이라 하자. 열린집합
$V \subset Y$가 존재하여 $f^{-1}(V) \to V$가 동형이고
$Y \setminus V$가 이산공간이라고 가정하자. 그러면 모든 가역
$\mathcal{O}_X$-가군은 어떤 가역 $\mathcal{O}_Y$-가군의 당김이다.
\end{lemma}

\begin{proof}
$\Pic(X) = H^1(X, \mathcal{O}_X^*)$라는 사실을 사용할 것이다.
Cohomology, Lemma \ref{cohomology-lemma-h1-invertible}를 보라.
아벨 층 $\mathcal{O}_X^*$와 $f$에 대한 Leray 스펙트럼 열을 생각하자.
Cohomology, Lemma \ref{cohomology-lemma-Leray}를 보라. 유도된 사상
$$
H^1(X, \mathcal{O}_X^*) \longrightarrow H^0(Y, R^1f_*\mathcal{O}_X^*)
$$
을 생각하자. Divisors, Lemma
\ref{divisors-lemma-finite-trivialize-invertible-upstairs}는 정확히 이
사상이 영임을 말한다. 따라서 Leray로부터
$H^1(X, \mathcal{O}_X^*) = H^1(Y, f_*\mathcal{O}_X^*)$를 얻는다.
이제 사상
$$
f^\sharp : \mathcal{O}_Y^* \longrightarrow f_*\mathcal{O}_X^*
$$
을 생각하자. 가정에 의해 이 사상의 핵과 여핵은 닫힌부분집합
$T = Y \setminus V$에 지지되며, 이는 $Y$의 닫힌부분집합이다.
가정에 의해 $T$는 이산 위상공간이므로
$Y$ 위에서 $T$에 지지되는 임의의 아벨 층의 고차 코호몰로지 군은
영이다. 이는 Cohomology, Lemma
\ref{cohomology-lemma-cohomology-and-closed-immersions}, Modules,
Lemma \ref{modules-lemma-i-star-exact}, 그리고
$H^i(T, \mathcal{F}) = 0$이라는 사실에서 따른다. 이는 임의의
$i > 0$과 임의의 아벨 층 $\mathcal{F}$에 대해 성립하며, 층은
$T$ 위에 있다.
위에 표시한 사상을 짧은 완전열들
$$
0 \to \Ker(f^\sharp) \to \mathcal{O}_Y^* \to \Im(f^\sharp) \to 0,\quad
0 \to \Im(f^\sharp) \to f_*\mathcal{O}_X^* \to \Coker(f^\sharp) \to 0
$$
로 나누면, 먼저
$H^1(Y, \mathcal{O}_Y^*) \to H^1(Y, \Im(f^\sharp))$가 전사임을
얻고, 이어서
$H^1(Y, \Im(f^\sharp)) \to H^1(Y, f_*\mathcal{O}_X^*)$가 전사임을
얻는다. 위의 모든 내용을 합치면 원하는 대로
$H^1(Y, \mathcal{O}_Y^*) \to
H^1(X, \mathcal{O}_X^*)$가 전사이다.
\end{proof}

\begin{lemma}
\label{varieties-lemma-glue-invertible-sheaves}
$X$를 스킴이라 하자. $Z_1, \ldots, Z_n \subset X$를 닫힌
부분스킴들이라 하자. $\mathcal{L}_i$를 $Z_i$ 위의 가역층이라 하자.
다음을 가정하자.
\begin{enumerate}
\item $X$는 기약이다.
\item 집합론적으로 $X = \bigcup Z_i$이다.
\item $Z_i \cap Z_j$는 이산 위상공간이며, 이는 $i \not = j$에 대해 요구한다.
\end{enumerate}
그러면 가역층 $\mathcal{L}$이 $X$ 위에 존재하여 $Z_i$로의 제한이
$\mathcal{L}_i$이다. 더 나아가 단면들
$s_i \in \Gamma(Z_i, \mathcal{L}_i)$가 주어지고 이들이
$Z_i \cap Z_j$의 점들에서 소멸하지 않으면, $\mathcal{L}$을 택하여
$s \in \Gamma(X, \mathcal{L})$가 존재하고 $s|_{Z_i} = s_i$가 모든
$i$에 대해 성립하게 할 수 있다.
\end{lemma}

\begin{proof}
$\mathcal{L}$의 존재는 Lemma
\ref{varieties-lemma-surjective-pic-birational-finite}에서 얻을 수 있다.
그러나 직접적인 증명도 제시하고, 그 증명으로 단면에 관한 명제가
참임을 보겠다. $T = \bigcup_{i \not = j} Z_i \cap Z_j$로 두자.
$X$가 기약이므로 스킴으로서
$$
X \setminus T = \bigcup (Z_i \setminus T)
$$
이다. 가정 (3)에 의해 $T$는 $X$의 이산 부분집합이다. 따라서 각
$t \in T$에 대해 열린집합 $U_t \subset X$를 찾을 수 있다. 이때
$t \in U_t$이지만 $t' \not \in U_t$이며, 후자는
$t' \in T$, $t' \not = t$일 때 요구한다. 필요하다면 $U_t$를 줄여서 동형들
$\varphi_{t, i} : \mathcal{L}_i|_{U_t \cap Z_i} \to
\mathcal{O}_{U_t \cap Z_i}$가 존재한다고 가정할 수 있다. 더 나아가
각 $i$에 대해 열린 덮개
$$
Z_i \setminus T = \bigcup\nolimits_j U_{ij}
$$
를 택하여 동형들
$\varphi_{i, j} : \mathcal{L}_i|_{U_{ij}} \cong \mathcal{O}_{U_{ij}}$가
존재하게 하자.
$$
\mathcal{U} : X = \bigcup U_t \cup \bigcup U_{ij}
$$
가 $X$의 열린 덮개임에 주의하라. 동형들 $\varphi_{t, i}$와
$\varphi_{i, j}$를 이용하여 이 덮개의 $2$-코사이클을 정의할 수
있다고 주장한다. 그 값은 $\mathcal{O}_X^*$에 있다. 이 코사이클은 보조정리의
명제에서와 같은 $\mathcal{L}$을 정의한다.

\medskip\noindent
실제로 $i \not = i'$이면
$U_{i, j} \cap U_{i', j'} = \emptyset$이므로 할 일이 없다.
$U_{i, j} \cap U_{i, j'}$에 대해서는 증명의 첫째 관찰에 의해
$\mathcal{O}_X(U_{i, j} \cap U_{i, j'}) =
\mathcal{O}_{Z_i}(U_{i, j} \cap U_{i, j'})$이다. 따라서
$\mathcal{L}_i$의 전이함수, 더 정확히는
$\varphi_{i, j} \circ \varphi_{i, j'}^{-1}$가 이 교집합에서 우리
코사이클의 값을 정의한다. $U_t \cap U_{i, j}$에 대해서도 같은
방법을 쓸 수 있다. 마지막으로 $t \not = t'$이면
$$
U_t \cap U_{t'} = \coprod (U_t \cap U_{t'}) \cap Z_i
$$
이고, 더구나 교집합 $U_t \cap U_{t'} \cap Z_i$는
$Z_i \setminus T$에 포함된다. 그러므로 앞과 같은 논증으로
$$
\mathcal{O}_X(U_t \cap U_{t'}) =
\prod \mathcal{O}_{Z_i}(U_t \cap U_{t'} \cap Z_i)
$$
임을 안다. 이제 $\mathcal{L}_i$의 전이함수, 더 정확히는
$\varphi_{t, i} \circ \varphi_{t', i}^{-1}$를 이용하여
$U_t \cap U_{t'}$에서 우리 코사이클의 값을 정의할 수 있다.
이로써 $\mathcal{L}$의 존재 증명이 끝난다.

\medskip\noindent
보조정리의 마지막 주장에 나오는 단면들 $s_i$가 주어졌다고 하자.
위 논증에서 $U_t$를 택하여 $s_i|_{U_t \cap Z_i}$가 소멸하지 않게 하고,
$\varphi_{t, i}$를 택하여
$\varphi_{t, i}(s_i|_{U_t \cap Z_i}) = 1$이 되게 한다. 그러면 $1$을
$U_t$ 위에서, $\varphi_{i, j}(s_i|_{U_{i, j}})$를
$U_{i, j}$ 위에서 사용하면 $\mathcal{L}$의 한 단면을 정의하며,
이는 $s_i$로 $Z_i$ 위에서 제한된다.
\end{proof}

\begin{remark}
\label{varieties-remark-conductor}
$A$를 기약환이라 하자. $I, J$를 $A$의 아이디얼들이라 하고
$V(I) \cup V(J) = \Spec(A)$라고 하자. $B = A/J$로 두자.
그러면 $I \to IB$는 $A$-가군의 동형이다. 실제로
$IB = I + J/J = I/(I \cap J)$이고 $I \cap J$는 영이다. 이는 $A$가
기약이고 $\Spec(A) = V(I) \cup V(J) = V(I \cap J)$이기 때문이다.
따라서 임의의 사영 $A$-가군 $P$에 대해서도
$IP = I(P/JP)$이다.
\end{remark}

\begin{lemma}
\label{varieties-lemma-dim-1-noetherian-reduced-separated-has-ample}
$X$를 차원 $1$의 뇌터 기약 분리 스킴이라 하자. 그러면 $X$는
풍부한 가역층을 가진다.
\end{lemma}

\begin{proof}
$Z_i$, $i = 1, \ldots, n$을 $X$의 기약성분들이라 하자. 이들을
$X$의 기약 닫힌 부분스킴들로 본다. Lemma
\ref{varieties-lemma-dim-1-noetherian-integral-separated-has-ample}에 의해
풍부한 가역층들 $\mathcal{L}_i$가 $Z_i$ 위에 존재한다.
$T = \bigcup_{i \not = j} Z_i \cap Z_j$로 두자. $X$는 차원 $1$의
뇌터 스킴이므로 집합 $T$는 유한하고 $X$의 닫힌점들로 이루어진다.
각 $i$에 대해 필요하면 $\mathcal{L}_i$를 거듭제곱으로 바꾼 뒤
$s_i \in \Gamma(Z_i, \mathcal{L}_i)$를 택하여
$(Z_i)_{s_i}$가 아핀이고 $T \cap Z_i$를 포함하게 할 수 있다.
Properties, Lemma
\ref{properties-lemma-ample-finite-set-in-principal-affine}를 보라.

\medskip\noindent
Lemma \ref{varieties-lemma-glue-invertible-sheaves}에 의해 가역층
$\mathcal{L}$을 $X$ 위에서 찾고 $s \in \Gamma(X, \mathcal{L})$를 택하여
$(\mathcal{L}, s)|_{Z_i} = (\mathcal{L}_i, s_i)$가 되게 할 수 있다.
$X_s$는 $T$를 포함하고, 집합론적으로 아핀 닫힌 부분스킴들
$(Z_i)_{s_i}$와 같음에 주의하라. 따라서 Limits, Lemma
\ref{limits-lemma-affines-glued-in-closed-affine}에 의해 이는 아핀이다.
증명을 끝내려면 모든 $x \in X$, $x \not \in T$에 대해 정수
$m > 0$과 단면 $t \in \Gamma(X, \mathcal{L}^{\otimes m})$을 찾아서
$X_t$가 아핀이고 $x \in X_t$가 되게 하면 충분하다.
$x \not \in T$이므로 $x \in Z_i$이고, 이러한 $i$는 유일하다.
$i = 1$이라고 하자. $Z \subset X$를 그 바탕 위상공간이
$Z_2 \cup \ldots \cup Z_n$인 기약 닫힌 부분스킴이라 하자.
$\mathcal{I} \subset \mathcal{O}_X$를 $Z$의 아이디얼층이라 하자.
$\mathcal{I}_1 \subset \mathcal{O}_{Z_1}$로 $Z_1 \to X$ 아래 이
아이디얼층의 역상을 나타내자.
$$
\Gamma(X, \mathcal{I}\mathcal{L}^{\otimes m}) =
\Gamma(Z_1, \mathcal{I}_1 \mathcal{L}_1^{\otimes m})
$$
임에 주의하라. Remark \ref{varieties-remark-conductor}를 보라. 따라서
$m > 0$과 $t \in \Gamma(Z_1, \mathcal{I}_1 \mathcal{L}_1^{\otimes m})$을
찾아서 $x \in (Z_1)_t$가 아핀이 되게 하면 충분하다.
$\mathcal{L}_1$은 풍부하고 $x$는
$Z_1 \cap T = V(\mathcal{I}_1)$에 속하지 않으므로 단면
$t_1 \in \Gamma(Z_1, \mathcal{I}_1 \mathcal{L}_1^{\otimes m_1})$을
찾아서 $x \in (Z_1)_{t_1}$이 되게 할 수 있다. Properties,
Proposition \ref{properties-proposition-characterize-ample}을 보라.
$\mathcal{L}_1$이 풍부하므로 단면
$t_2 \in \Gamma(Z_1, \mathcal{L}_1^{\otimes m_2})$를 찾아서
$x \in (Z_1)_{t_2}$이고 $(Z_1)_{t_2}$가 아핀이 되게 할 수 있다.
Properties, Definition \ref{properties-definition-ample}을 보라.
$m = m_1 + m_2$ 및 $t = t_1 t_2$로 두자. 그러면
$t \in \Gamma(Z_1, \mathcal{I}_1 \mathcal{L}_1^{\otimes m})$이고
구성에 의해 $x \in (Z_1)_t$이며, Properties, Lemma
\ref{properties-lemma-affine-cap-s-open}에 의해 $(Z_1)_t$는 아핀이다.
\end{proof}

\begin{lemma}
\label{varieties-lemma-lift-line-bundle-from-reduction-dimension-1}
$i : Z \to X$를 스킴들의 닫힌 몰입이라 하자. $X$의 바탕 위상공간이
뇌터이고 $\dim(X) \leq 1$이면
$\Pic(X) \to \Pic(Z)$는 전사이다.
\end{lemma}

\begin{proof}
$X$ 위의 아벨 군층들의 짧은 완전열
$$
0 \to (1 + \mathcal{I}) \cap \mathcal{O}_X^* \to
\mathcal{O}^*_X \to i_*\mathcal{O}^*_Z \to 0
$$
을 생각하자. 여기서 $\mathcal{I}$는 $Z$에 대응하는 준연접
아이디얼층이다. $\dim(X) \leq 1$이므로
$H^2(X, \mathcal{F}) = 0$이고, 이는 임의의 아벨 층
$\mathcal{F}$에 대해 성립한다. Cohomology,
Proposition \ref{cohomology-proposition-vanishing-Noetherian}을 보라.
따라서 사상
$H^1(X, \mathcal{O}^*_X) \to H^1(X, i_*\mathcal{O}_Z^*)$는 전사이다.
Cohomology, Lemma
\ref{cohomology-lemma-cohomology-and-closed-immersions}에 의해
$H^1(X, i_*\mathcal{O}_Z^*) = H^1(Z, \mathcal{O}_Z^*)$이다.
Cohomology, Lemma \ref{cohomology-lemma-h1-invertible}에 의해
이것이 보조정리를 증명한다.
\end{proof}

\begin{proposition}
\label{varieties-proposition-dim-1-noetherian-separated-has-ample}
$X$를 차원 $1$의 뇌터 분리 스킴이라 하자. 그러면 $X$는 풍부한
가역층을 가진다.
\end{proposition}

\begin{proof}
$Z \subset X$를 $X$의 기약화라 하자. Lemma
\ref{varieties-lemma-dim-1-noetherian-reduced-separated-has-ample}에 의해 스킴
$Z$는 풍부한 가역층을 가진다. 따라서 Lemma
\ref{varieties-lemma-lift-line-bundle-from-reduction-dimension-1}에 의해
가역 $\mathcal{O}_X$-가군 $\mathcal{L}$이 존재하고, 이는
$X$ 위에 있으며
$Z$로의 제한이 풍부하다. 그러면 Cohomology of Schemes, Lemma
\ref{coherent-lemma-ample-on-reduction}를 적용하면
$\mathcal{L}$은 풍부하다.
\end{proof}

\begin{remark}
\label{varieties-remark-useless-generalization}
실제로 $X$가 그 기약화가 차원 $1$의 뇌터 분리 스킴인 스킴이면
$X$는 풍부한 가역층을 가진다. 이를 증명하는 논증은 Proposition
\ref{varieties-proposition-dim-1-noetherian-separated-has-ample}의 증명과 같다.
다만 Limits, Lemma
\ref{limits-lemma-ample-on-reduction}을 사용하고,
Cohomology of Schemes, Lemma
\ref{coherent-lemma-ample-on-reduction}은 사용하지 않는다.
\end{remark}

\noindent
다음 보조정리는 실제로 준유한 분리 사상들에 대해 성립한다.
독자는 Zariski 주정리(More on Morphisms, Lemma
\ref{more-morphisms-lemma-quasi-finite-separated-pass-through-finite})와
Lemma \ref{varieties-lemma-complement-codim-1-closed-points}를 사용하여 이를
확인할 수 있다.

\begin{lemma}
\label{varieties-lemma-surjection-on-pic-quasi-finite}
$f : X \to Y$를 스킴들의 사상이라 하자. $Y$가 차원 $\leq 1$인
뇌터 스킴이고, $f$가 유한이며, 조밀한 열린집합 $V \subset Y$가
존재하여 $f^{-1}(V) \to V$가 닫힌 몰입이라고 가정하자. 그러면 모든
가역 $\mathcal{O}_X$-가군은 어떤 가역 $\mathcal{O}_Y$-가군의 당김이다.
\end{lemma}

\begin{proof}
$f$를 $X \to Z \to Y$로 분해하자. 여기서 $Z$는 $f$의 스킴론적
상이다. 그러면 $X \to Z$는 $V \cap Z$ 위에서 동형이고 Lemma
\ref{varieties-lemma-surjective-pic-birational-finite}를 적용할 수 있다.
한편 Lemma \ref{varieties-lemma-lift-line-bundle-from-reduction-dimension-1}을
$Z \to Y$에 적용할 수 있다. 몇 가지 세부사항은 생략한다.
\end{proof}








\section{델타 불변량}
\label{varieties-section-delta-invariant}

\noindent
이 절에서는 $\delta$-불변량을 정의한다. 이는 기약인 $1$차원
나가타 스킴의 특이점에 대한 불변량이다.

\begin{lemma}
\label{varieties-lemma-pre-pre-delta-invariant}
$(A, \mathfrak m)$을 뇌터 $1$차원 국소환이라 하자.
$f \in \mathfrak m$이라 하자. 다음 조건들은 동치이다.
\begin{enumerate}
\item $\mathfrak m = \sqrt{(f)}$.
\item $f$는 $A$의 어느 극소소아이디얼에도 속하지 않는다.
\item $A_f = \prod_{\mathfrak p\text{ minimal}} A_\mathfrak p$이며,
이는 $A$-대수로서의 등식이다.
\end{enumerate}
그러한 $f \in \mathfrak m$은 존재한다. $\text{depth}(A) = 1$이면
(예를 들어 $A$가 기약이면), (1)--(3)은 다음 조건들과도 동치이다.
\begin{enumerate}
\item[(4)] $f$는 영인자가 아니다.
\item[(5)] $A_f$는 $A$의 전체 분수환이다.
\end{enumerate}
$A$가 기약이면 (1)--(5)는 다음 조건과도 동치이다.
\begin{enumerate}
\item[(6)] $A_f$는 $A$의 극소소아이디얼들에서의 잉여체들의 곱이다.
\end{enumerate}
\end{lemma}

\begin{proof}
$A$의 스펙트럼에는 유한 개의 소아이디얼
$\mathfrak p_1, \ldots, \mathfrak p_n$이 $\mathfrak m$ 이외에 있고,
이들은 모두 극소이다.
Algebra, Lemma
\ref{algebra-lemma-Noetherian-irreducible-components}를 보라. 그러면
(1)과 (2)의 동치는 Algebra, Lemma
\ref{algebra-lemma-Zariski-topology}에서 따른다. (3)이 (2)를
함의함은 분명하다. 반대로 (2)가 참이면 $A_f$의 스펙트럼은 부분집합
$\{\mathfrak p_1, \ldots, \mathfrak p_n\}$에 유도된 위상을 준 것이다.
이 부분집합은 $\Spec(A)$ 안에 있다.
Algebra, Lemma \ref{algebra-lemma-spec-localization}을 보라. 이는 유한
이산 위상공간이다. 따라서 Algebra, Proposition
\ref{algebra-proposition-dimension-zero-ring}에 의해
$A_f = \prod_{\mathfrak p\text{ minimal}} A_\mathfrak p$이다.
$f$의 존재는 Algebra, Lemma \ref{algebra-lemma-height-1}에서 주장한다.

\medskip\noindent
$A$의 깊이가 $1$이라고 가정하자. 이는 Algebra, Lemma
\ref{algebra-lemma-bound-depth}에 의해 가능한 최댓값이고, $A$가
기약이면 Algebra, Lemma \ref{algebra-lemma-criterion-reduced}에 의해
성립한다. 그러면 $\mathfrak m$은 $A$의 결합소아이디얼이 아니다.
$A$의 모든 극소소아이디얼은 결합소아이디얼이다(Algebra, Proposition
\ref{algebra-proposition-minimal-primes-associated-primes}). 따라서
Algebra, Lemma \ref{algebra-lemma-ass-zero-divisors}에 의해 $A$의
영인자가 아닌 원소들의 집합은 정확히 어느 극소소아이디얼에도 속하지
않는 원소들의 집합이다. 그러므로 (4)는 (2)와 동치이다. (5)는
Algebra, Lemma
\ref{algebra-lemma-total-ring-fractions-no-embedded-points}에 의해
(3)과 동치이다.

\medskip\noindent
그러면 $A_\mathfrak p$는 체이며, 여기서 $\mathfrak p \subset A$는 극소이다.
Algebra, Lemma \ref{algebra-lemma-minimal-prime-reduced-ring}을 보라.
따라서 (3)은 (6)과 동치이다.
\end{proof}

\begin{lemma}
\label{varieties-lemma-pre-delta-invariant}
$(A, \mathfrak m)$을 기약인 나가타 $1$차원 국소환이라 하자.
$A'$를 $A$의 전체 분수환 안에서 $A$의 정수적 폐포라 하자.
그러면 $A'$는 정규 나가타 환이고, $A \to A'$는 유한이며,
$A'/A$는 $A$-가군으로서 유한 길이를 가진다.
\end{lemma}

\begin{proof}
전체 분수환은 $A$ 위에서 본질적 유한형이므로, $A \to A'$는
유한이다. 이는 $A$가 나가타라는 것에 의한다. Algebra, Lemma
\ref{algebra-lemma-nagata-in-reduced-finite-type-finite-integral-closure}를
보라. 환 $A'$는 예를 들어 Algebra, Lemmas
\ref{algebra-lemma-characterize-reduced-ring-normal} 및
\ref{algebra-lemma-Noetherian-irreducible-components}에 의해
정규이다. 환 $A'$는 예를 들어 Algebra, Lemma
\ref{algebra-lemma-quasi-finite-over-nagata}에 의해 나가타이다.
Lemma \ref{varieties-lemma-pre-pre-delta-invariant}에서와 같은
$f \in \mathfrak m$을 택하자. $A' \subset A_f$이므로
$A_f = A'_f$임은 분명하다. 따라서 유한 $A$-가군 $A'/A$의 지지는
$\{\mathfrak m\}$에 포함된다. 그러므로 Algebra, Lemma
\ref{algebra-lemma-support-point}에 의해 이는 유한 길이를 가진다.
\end{proof}

\begin{definition}
\label{varieties-definition-delta-invariant-algebra}
$A$를 차원 $1$인 기약 나가타 국소환이라 하자.
{\it $\delta$-불변량, 즉 $A$의 델타 불변량}은
$\text{length}_A(A'/A)$이다. 여기서
$A'$는 Lemma \ref{varieties-lemma-pre-delta-invariant}에서와 같다.
\end{definition}

\noindent
이 불변량의 거동에 관한 몇 가지 보조정리를 증명한다.

\begin{lemma}
\label{varieties-lemma-delta-invariant-is-zero}
$A$를 차원 $1$인 기약 나가타 국소환이라 하자.
$\delta$-불변량이 $A$에 대해 $0$일 필요충분조건은
$A$가 이산 값매김환인 것이다.
\end{lemma}

\begin{proof}
$A$가 이산 값매김환이면 $A$는 정규이고 환 $A'$는 $A$와 같다.
반대로 $\delta$-불변량이 $A$에 대해 $0$이면 $A$는 그 전체 분수환
안에서 정수적으로 닫혀 있다. 이는 $A$가 정규임을 함의하고
(Algebra, Lemma
\ref{algebra-lemma-characterize-reduced-ring-normal}), Algebra,
Lemma \ref{algebra-lemma-characterize-dvr}에 의해 $A$가 이산
값매김환일 수밖에 없다.
\end{proof}

\begin{lemma}
\label{varieties-lemma-normalization-same-after-completion}
$A$를 차원 $1$인 기약 나가타 국소환이라 하자.
$A \to A'$를 Lemma \ref{varieties-lemma-pre-delta-invariant}에서와 같이 두자.
$A^h$, $A^{sh}$, 각각 $A^\wedge$를 $A$의 헨젤화, 엄밀 헨젤화,
각각 완비화라 하자. 그러면 $A^h$, $A^{sh}$, 각각 $A^\wedge$는
차원 $1$인 기약 나가타 국소환이고,
$A' \otimes_A A^h$, $A' \otimes_A A^{sh}$, 각각
$A' \otimes_A A^\wedge$는 $A^h$, $A^{sh}$, 각각 $A^\wedge$의
전체 분수환 안에서의 정수적 폐포이다.
\end{lemma}

\begin{proof}
$A^\wedge$가 기약임에 주의하라. More on Algebra, Lemma
\ref{more-algebra-lemma-completion-reduced}를 보라. 환들 $A^h$와
$A^{sh}$는 More on Algebra, Lemma
\ref{more-algebra-lemma-henselization-reduced}에 의해 기약이다.
$A$, $A^h$, $A^{sh}$, $A^\wedge$의 차원들은 More on Algebra,
Lemmas \ref{more-algebra-lemma-completion-dimension} 및
\ref{more-algebra-lemma-henselization-dimension}에 의해 같다.

\medskip\noindent
뇌터 국소환이 나가타일 필요충분조건은 $A$의 형식적 올들이
기하적으로 기약인 것임을 상기하라. More on Algebra, Lemma
\ref{more-algebra-lemma-Nagata-local-ring}을 보라. 이 성질은
$A^h$와 $A^{sh}$에 유전된다. More on Algebra, Section
\ref{more-algebra-section-properties-formal-fibres}의 내용, 특히
Lemma \ref{more-algebra-lemma-henselization-P-ring}을 보라.
완비화는 Algebra, Lemma
\ref{algebra-lemma-Noetherian-complete-local-Nagata}에 의해 나가타이다.

\medskip\noindent
이제 정수적 폐포에 관한 명제를 다룬다. 계속하기 전에 Lemma
\ref{varieties-lemma-pre-pre-delta-invariant}에서와 같은
$f \in \mathfrak m$을 택하자. 그러면 $f$의 상은 $A^h$, $A^{sh}$,
$A^\wedge$에서 분명히 그 환에서 성질 (1)--(6)을
만족하는 원소이다.

\medskip\noindent
$A \to A'$가 유한이므로 $A' \otimes_A A^h$와
$A' \otimes_A A^{sh}$는 각각 $A^h$와 $A^{sh}$ 위에서 유한인
헨젤 국소환들의 곱이다. Algebra, Lemma
\ref{algebra-lemma-finite-over-henselian}을 보라. 이 국소환들 각각은
$A'$의 헨젤화이다. 이는 극대아이디얼 $\mathfrak m' \subset A'$에서
취하며, 그 아이디얼은 $\mathfrak m$ 위에 놓인다. Algebra, Lemma
\ref{algebra-lemma-quasi-finite-henselization} 또는
\ref{algebra-lemma-quasi-finite-strict-henselization}을 보라.
따라서 이 국소환들은 More on Algebra, Lemma
\ref{more-algebra-lemma-henselization-normal}에 의해 정규 정역이다.
그러므로 $A' \otimes_A A^h$와 $A' \otimes_A A^{sh}$는 정규환이다.
$A^h \to A' \otimes_A A^h$와
$A^{sh} \to A' \otimes_A A^{sh}$는 유한이고, 따라서 정수적이다.
또한
$A' \otimes_A A^h \subset (A^h)_f = Q(A^h)$ 및
$A' \otimes_A A^{sh} \subset (A^{sh})_f = Q(A^{sh})$이므로
$A' \otimes_A A^h$와 $A' \otimes_A A^{sh}$가 원하는 정수적
폐포라고 결론짓는다.

\medskip\noindent
완비화에 대해서도 완전히 같은 방식으로 논증한다. 먼저 Algebra,
Lemma \ref{algebra-lemma-completion-finite-extension}에 의해
$$
A' \otimes_A A^\wedge = (A')^\wedge =
\prod\nolimits (A'_{\mathfrak m'})^\wedge
$$
이다. 국소환들 $A'_{\mathfrak m'}$은 정규이고 차원이 $1$이다
(예를 들어 Algebra, Lemma
\ref{algebra-lemma-finite-in-codim-1} 또는 Algebra, Section
\ref{algebra-section-homomorphism-dimension}의 논의에 의한다).
따라서 $A'_{\mathfrak m'}$은 이산 값매김환이다. Algebra, Lemma
\ref{algebra-lemma-characterize-dvr}를 보라. 그러므로 More on Algebra,
Lemma \ref{more-algebra-lemma-completion-dvr}에 의해
$(A'_{\mathfrak m'})^\wedge$는 이산 값매김환이다. 따라서
$A' \otimes_A A^\wedge$는 정규환이고, 앞과 정확히 같은 방식으로
결론지을 수 있다.
\end{proof}

\begin{lemma}
\label{varieties-lemma-delta-same-after-completion}
$A$를 차원 $1$인 기약 나가타 국소환이라 하자.
$\delta$-불변량은 $A$에 대해 계산하든, $\delta$-불변량을 헨젤화,
엄밀 헨젤화, 또는 $A$의 완비화에 대해 계산하든 같다.
\end{lemma}

\begin{proof}
완비화 $B = A^\wedge$의 경우를 보자. 다른 경우들도 정확히 같은
방식으로 증명된다. $A'$, 각각 $B'$를 $A$, 각각 $B$의 전체 분수환
안에서의 정수적 폐포라 하자. 그러면 Lemma
\ref{varieties-lemma-normalization-same-after-completion}에 의해
$B' = A' \otimes_A B$이다. 따라서
$B'/B = (A'/A) \otimes_A B$이다. 이제 Algebra, Lemma
\ref{algebra-lemma-pullback-module}와
$B \otimes_A \kappa_A = \kappa_B$라는 사실에서 등식이 따른다.
\end{proof}

\begin{definition}
\label{varieties-definition-delta-invariant}
$k$를 체라 하자. $X$를 국소 대수적 $k$-스킴이라 하자.
$x \in X$를 $\mathcal{O}_{X, x}$가 기약이고
$\dim(\mathcal{O}_{X, x}) = 1$인 점이라 하자.
{\it $\delta$-불변량, 즉 $X$의 $x$에서의 불변량}은 Definition
\ref{varieties-definition-delta-invariant-algebra}에서 정의한
$\delta$-불변량, 즉 $\mathcal{O}_{X, x}$의 불변량이다.
\end{definition}

\noindent
국소 대수적 스킴의 국소환은 Algebra, Proposition
\ref{algebra-proposition-ubiquity-nagata}에 의해 나가타이므로 이
정의는 의미가 있다. 물론 더 일반적으로, $x \in X$가 어떤 스킴의
점이고 국소환 $\mathcal{O}_{X, x}$가 기약인 차원 $1$의 나가타
환일 때마다 이 정의를 내릴 수 있다. Lemma
\ref{varieties-lemma-delta-same-after-completion}에 의해 $\delta$-불변량, 즉
$X$의 $x$에서의 불변량은
$$
\delta\text{-invariant of }X\text{ at }x =
\delta\text{-invariant of }\mathcal{O}_{X, x}^h =
\delta\text{-invariant of }\mathcal{O}_{X, x}^\wedge
$$
이다. 따라서 $\delta$-불변량은 그 점의 완비 국소환의 불변량이다.

\begin{lemma}
\label{varieties-lemma-delta-invariant-and-change-of-fields}
$k$를 체라 하자. $X$를 국소 대수적 $k$-스킴이라 하자.
$K/k$를 체의 확대라 하고 $Y = X_K$로 두자. $y \in Y$의 상을
$x \in X$라 하자. $X$가 $x$에서 기하적으로 기약이고
$\dim(\mathcal{O}_{X, x}) = \dim(\mathcal{O}_{Y, y}) = 1$이라고
가정하자. 그러면
$$
\delta\text{-invariant of }X\text{ at }x \leq
\delta\text{-invariant of }Y\text{ at }y
$$
이다.
\end{lemma}

\begin{proof}
$A = \mathcal{O}_{X, x}$ 및 $B = \mathcal{O}_{Y, y}$로 두자.
Lemma \ref{varieties-lemma-geometrically-reduced-at-point}에 의해 $A$는
기하적으로 기약이다. 따라서 $B$는 $A \otimes_k K$의 국소화이다.
$A \to A'$를 Lemma \ref{varieties-lemma-pre-delta-invariant}에서와 같이 두자.
그러면
$$
B' = B \otimes_{(A \otimes_k K)} (A' \otimes_k K)
$$
은 $B$ 위에서 유한이고, $B \to B'$는 전체 분수환 위에 동형을
유도한다. 실제로 Lemma \ref{varieties-lemma-pre-pre-delta-invariant}의
(1)--(6)을 만족하는 $f \in \mathfrak m_A$를 택하자.
$\dim(B) = 1$이므로 $f \in \mathfrak m_B$는 $B$에 대해 같은 역할을
하고, $B_f = B'_f$임을 알 수 있다. 이는 $A_f = A'_f$이기 때문이다.
$B''$를 Lemma \ref{varieties-lemma-pre-delta-invariant}에서와 같이 $B$의
전체 분수환 안에서의 정수적 폐포라 하자. 그러면
$B' \subset B''$이다. 따라서 $\delta$-불변량, 즉 $Y$의 $y$에서의 불변량은
$\text{length}_B(B''/B)$이고
\begin{align*}
\text{length}_B(B''/B)
& \geq
\text{length}_B(B'/B) \\
& =
\text{length}_B((A'/A) \otimes_A B) \\
& =
\text{length}_B(B/\mathfrak m_A B) \text{length}_A(A'/A)
\end{align*}
이다. 이는 Algebra, Lemma \ref{algebra-lemma-pullback-module}에
의하며, $A \to B$가 평탄하기 때문이다(이는
$A \to A \otimes_k K$의 국소화이다).
$\text{length}_A(A'/A)$는 $\delta$-불변량, 즉 $X$의 $x$에서의 불변량이고,
$\text{length}_B(B/\mathfrak m_A B) \geq 1$이므로 보조정리가 증명되었다.
\end{proof}

\begin{lemma}
\label{varieties-lemma-delta-invariant-and-change-of-fields-better}
$k$를 체라 하자. $X$를 국소 대수적 $k$-스킴이라 하자.
$K/k$를 체의 확대라 하고 $Y = X_K$로 두자. $y \in Y$의 상을
$x \in X$라 하자. Lemma
\ref{varieties-lemma-geometrically-normal-in-codim-1}의 가정 (a), (b), (c)가
$x \in X$에 대해 성립하고 $\dim(\mathcal{O}_{Y, y}) = 1$이라고
가정하자. 그러면 $\delta$-불변량, 즉 $X$의 $x$에서의 불변량은
$\delta$-불변량, 즉 $Y$의 $y$에서의 불변량과 같다.
\end{lemma}

\begin{proof}
$A = \mathcal{O}_{X, x}$ 및 $B = \mathcal{O}_{Y, y}$로 두자.
Lemma \ref{varieties-lemma-geometrically-normal-in-codim-1}에 의해 $A$는
기하적으로 기약이다. 따라서 $B$는 $A \otimes_k K$의 국소화이다.
$A \to A'$를 Lemma \ref{varieties-lemma-pre-delta-invariant}에서와 같이 두자.
Lemma \ref{varieties-lemma-geometrically-normal-in-codim-1}에 의해
$A' \otimes_k K$는 정규이다. 따라서
$$
B' = B \otimes_{(A \otimes_k K)} (A' \otimes_k K)
$$
은 정규이고 $B$ 위에서 유한이며, $B \to B'$는 전체 분수환 위에
동형을 유도한다. 실제로 Lemma
\ref{varieties-lemma-pre-pre-delta-invariant}의 (1)--(6)을 만족하는
$f \in \mathfrak m_A$를 택하자. $\dim(B) = 1$이므로
$f \in \mathfrak m_B$는 $B$에 대해 같은 역할을 하고,
$B_f = B'_f$임을 알 수 있다. 이는 $A_f = A'_f$이기 때문이다. 따라서
$B \to B'$는 $B$에 대해 Lemma
\ref{varieties-lemma-pre-delta-invariant}에서와 같은 사상이다. 그러므로
$\text{length}_A(A'/A) =
\text{length}_B(B'/B) = \text{length}_B((A'/A) \otimes_A B)$를
보이면 된다. $A \to B$는 평탄하고(이는
$A \to A \otimes_k K$의 국소화이다),
$\mathfrak m_B = \mathfrak m_A B$이다. 실제로 위의 관찰에 의해
$B/\mathfrak m_A B$는 차원 영이고
$K \otimes_k \kappa(x)$의 국소화이며, $\kappa(x)$가 $k$ 위에서
분리적이므로 이 환은 기약이다. 따라서 Algebra, Lemma
\ref{algebra-lemma-pullback-module}에 의해 결론이 따른다.
\end{proof}





\section{가지의 수}
\label{varieties-section-number-of-branches}

\noindent
한 점에서 스킴의 가지 수를 Properties, Section
\ref{properties-section-unibranch}에서 정의하였다.

\begin{lemma}
\label{varieties-lemma-number-of-branches}
$X$를 스킴이라 하자. $X$의 각 준콤팩트 열린집합이 유한 개의
기약성분을 가진다고 가정하자. $\nu : X^\nu \to X$를 $X$의
정규화라 하자. $x \in X$라 하자.
\begin{enumerate}
\item $X$의 가지 수를 $x$에서 세면, 이는 $x$의 역상들을
$X^\nu$에서 센 수이다.
\item $X$의 기하적 가지 수를 $x$에서 세면, 이는
$\sum_{\nu(x^\nu) = x} [\kappa(x^\nu) : \kappa(x)]_s$이다.
\end{enumerate}
\end{lemma}

\begin{proof}
먼저 $X$에 대한 가정은 정확히 정규화가 정의된다는 뜻임에 주의하라.
Morphisms, Definition \ref{morphisms-definition-normalization}을 보라.
그러면 줄기 $A' = (\nu_*\mathcal{O}_{X^\nu})_x$는
$A = \mathcal{O}_{X, x}$의 전체 분수환 안에서 $A_{red}$의 정수적
폐포이다. Morphisms, Lemma
\ref{morphisms-lemma-stalk-normalization}을 보라.
$\nu$가 정수적 사상이므로 $X^\nu$에서 $x$ 위에 놓인 점들은
$A'$의 소아이디얼들과 대응한다. 이들은 극대아이디얼 $\mathfrak m$
위에 놓이며, 이는 $A$의 극대아이디얼이다.
$A \to A'$가 정수적이므로 이는 $A'$의 극대아이디얼들과
같은 것이다(Algebra, Lemmas
\ref{algebra-lemma-integral-no-inclusion} 및
\ref{algebra-lemma-integral-going-up}). 이제 보조정리는 그 대수적
대응물인 More on Algebra, Lemma
\ref{more-algebra-lemma-number-of-branches-1}에서 따른다.
\end{proof}

\begin{lemma}
\label{varieties-lemma-geometric-branches-and-change-of-fields}
$k$를 체라 하자. $X$를 국소 대수적 $k$-스킴이라 하자.
$K/k$를 체의 확대라 하자. $y \in X_K$를 점이라 하고, 그 상을
$x$라 하되 이는 $X$의 점이다. 그러면 $X$의 기하적 가지 수를
$x$에서 센 값은 $X_K$의 기하적 가지 수를 $y$에서 센 값과 같다.
\end{lemma}

\begin{proof}
$Y = X_K$로 쓰고 $X^\nu$, 각각 $Y^\nu$를 $X$, 각각 $Y$의
정규화라 하자. 다음 가환 그림을 생각하자.
$$
\xymatrix{
Y^\nu \ar[r] \ar[d] & X^\nu_K \ar[r] \ar[d]_{\nu_K} & X^\nu \ar[d]_\nu \\
Y \ar@{=}[r] & Y \ar[r] & X
}
$$
Lemma \ref{varieties-lemma-normalization-and-change-of-fields}에 의해 왼쪽 위의
수평 화살표는 보편 위상동형사상이다. 따라서 순비분리 잉여체 확대를
유도한다. Morphisms, Lemmas
\ref{morphisms-lemma-universal-homeomorphism} 및
\ref{morphisms-lemma-universally-injective}를 보라. 따라서
$Y$의 기하적 가지 수를 $y$에서 센 값은 Lemma
\ref{varieties-lemma-number-of-branches}에 의해
$\sum_{\nu_K(y') = y} [\kappa(y') : \kappa(y)]_s$이다. 마찬가지로
$\sum_{\nu(x') = x} [\kappa(x') : \kappa(x)]_s$는 $X$의
기하적 가지 수를 $x$에서 센 값이다. Schemes, Lemma
\ref{schemes-lemma-points-fibre-product}를 사용하면 우리의 명제는
다음 대수적 사실에서 따른다. 체의 확대 $l/\kappa$와 대수적 체의 확대
$m/\kappa$가 주어지면
$$
\sum\nolimits_{m \otimes_\kappa l \to m'} [m' : l']_s = [m : \kappa]_s
$$
이다. 여기서 합은 $m \otimes_\kappa l$의 몫체들에 걸쳐 취한다.
이를 초등적인 방법으로 증명할 수도 있고, Lemma
\ref{varieties-lemma-separably-closed-field-connected-components}를
$$
\Spec(m \otimes_\kappa l) \times_{\Spec(l)} \Spec(\overline{l}) =
\Spec(m) \otimes_{\Spec(\kappa)} \Spec(\overline{l})
\longrightarrow
\Spec(m) \times_{\Spec(\kappa)} \Spec(\overline{\kappa})
$$
에 적용할 수도 있다. 실제로 $[m : \kappa]_s$를 우변의 연결성분 수로,
합 $\sum_{m \otimes_\kappa l \to m'} [m' : l']_s$를 좌변의
연결성분 수로 해석할 수 있기 때문이다.
\end{proof}

\begin{lemma}
\label{varieties-lemma-geometrically-unibranch-and-change-of-fields}
$k$를 체라 하자. $X$를 국소 대수적 $k$-스킴이라 하자.
$K/k$를 체의 확대라 하자. $y \in X_K$를 점이라 하고, 그 상을
$x$라 하되 이는 $X$의 점이다. 그러면 $X$가 $x$에서 기하적으로 단일가지일
필요충분조건은 $X_K$가 $y$에서 기하적으로 단일가지인 것이다.
\end{lemma}

\begin{proof}
Lemma \ref{varieties-lemma-geometric-branches-and-change-of-fields}와 More on
Algebra, Lemma \ref{more-algebra-lemma-number-of-branches-1}에서 즉시
따른다.
\end{proof}

\begin{definition}
\label{varieties-definition-wedge}
$A$와 $A_i$, $1 \leq i \leq n$을 국소환들이라 하자.
{\it $A$가 $A_1, \ldots, A_n$의 쐐기}라는 것은 동형들
$$
\kappa_{A_1} \to \kappa_{A_2} \to \ldots \to \kappa_{A_n}
$$
이 존재하고, $A$가 $n$-튜플
$(a_1, \ldots, a_n) \in A_1 \times \ldots \times A_n$ 가운데
$\kappa_{A_n}$의 같은 원소로 가는 것들로 이루어진 환과 동형이면
성립한다고 정의한다.
\end{definition}

\noindent
바탕환 $\Lambda$가 주어지고 $A$와 $A_i$가 $\Lambda$-대수들이면,
$\kappa_{A_i} \to \kappa_{A_{i + 1}}$가 $\Lambda$-대수 동형일 것을
요구하고, $A$가 $\Lambda$-대수로서 다음 $\Lambda$-대수와 동형일
것을 요구한다. 이는 $n$-튜플
$(a_1, \ldots, a_n) \in A_1 \times \ldots \times A_n$ 가운데
$\kappa_{A_n}$의 같은 원소로 가는 것들로 이루어진다. 특히
$\Lambda = k$가 체이고 사상들 $k \to \kappa_{A_i}$가 동형이면,
동형들 $\kappa_{A_i} \to \kappa_{A_{i + 1}}$의 선택은 유일하며,
흔히 {\it $A_1, \ldots, A_n$의 쐐기}라고 말한다.

\begin{lemma}
\label{varieties-lemma-delta-number-branches-inequality-sh}
$(A, \mathfrak m)$을 엄밀 헨젤인 기약 나가타 $1$차원 국소환이라 하자.
그러면
$$
\delta\text{-invariant of }A \geq \text{number of geometric branches of }A - 1
$$
이다. 등호가 성립하면 $A$는 $n \geq 1$개의 엄밀 헨젤 이산
값매김환들의 쐐기이다.
\end{lemma}

\begin{proof}
기하적 가지 수는 $A$의 가지 수와 같다(More on Algebra, Definition
\ref{more-algebra-definition-number-of-branches}에서 즉시 따른다).
$A \to A'$를 Lemma \ref{varieties-lemma-pre-delta-invariant}에서와 같이 두자.
$A$의 가지 수는 $A'$의 극대아이디얼 수임에 주의하라. More on
Algebra, Lemma \ref{more-algebra-lemma-number-of-branches-1}을 보라.
전사
$$
A'/A \longrightarrow
\left(\prod\nolimits_{\mathfrak m'} \kappa(\mathfrak m')\right)/
\kappa(\mathfrak m)
$$
가 존재한다. $\dim_{\kappa(\mathfrak m)} \prod \kappa(\mathfrak m')$은
$\geq$ 가지 수이므로 부등식은 분명하다.

\medskip\noindent
등호가 성립하면 $\kappa(\mathfrak m') = \kappa(\mathfrak m)$이고,
이는 모든 $\mathfrak m' \subset A'$에 대해 성립한다. 위에 표시한
화살표는 동형이다. $A$가 헨젤이고 $A \to A'$가 유한이므로
$A'$는 국소 헨젤 환들의 곱이다. Algebra, Lemma
\ref{algebra-lemma-finite-over-henselian}을 보라. 인자들은 국소환
$A'_{\mathfrak m'}$들이고, $A'$가 정규이므로 이 인자들은 이산
값매김환이다(Algebra, Lemma
\ref{algebra-lemma-characterize-dvr}). 표시한 화살표가 동형이므로
$A$는 실제로 이 국소환들의 쐐기이다.
\end{proof}

\begin{lemma}
\label{varieties-lemma-delta-number-branches-inequality}
$(A, \mathfrak m)$을 기약 나가타 $1$차원 국소환이라 하자. 그러면
$$
\delta\text{-invariant of }A \geq \text{number of geometric branches of }A - 1
$$
\end{lemma}

\begin{proof}
$A$를 $A$의 엄밀 헨젤화로 바꿀 수 있으며, 이때 $\delta$-불변량은
바뀌지 않는다(Lemma \ref{varieties-lemma-delta-same-after-completion}). 또한 $A$의 기하적
가지 수도 바뀌지 않는다(이는 정의에서 즉시 따른다. More on Algebra,
Definition \ref{more-algebra-definition-number-of-branches}를 보라).
따라서 $A$가 엄밀 헨젤이라고 가정하고 Lemma
\ref{varieties-lemma-delta-number-branches-inequality-sh}를 적용할 수 있다.
\end{proof}






\section{1차원 스킴의 정규화}
\label{varieties-section-normalization-one-dimensional}

\noindent
Krull--Akizuki 결과에 의해 차원 $1$인 뇌터 스킴의 정규화 사상은
예상보다 좋은 성질들을 가진다.

\begin{lemma}
\label{varieties-lemma-normalize-noetherian-dim-1}
$X$를 차원 $1$의 국소 뇌터 스킴이라 하자.
$\nu : X^\nu \to X$를 정규화라 하자. 그러면
\begin{enumerate}
\item $\nu$는 정수적이고 전사이며 기약성분들 사이의 전단사를 유도한다.
\item 분해 $X^\nu \to X_{red} \to X$가 존재하고,
$X^\nu \to X_{red}$는 $X_{red}$의 정규화이다.
\item $X^\nu \to X_{red}$는 쌍유리이다.
\item 모든 닫힌점 $x \in X$에 대해 줄기
$(\nu_*\mathcal{O}_{X^\nu})_x$는 $\mathcal{O}_{X, x}$의 정수적
폐포이며, 이 폐포는 $(\mathcal{O}_{X, x})_{red} = \mathcal{O}_{X_{red}, x}$의
전체 분수환 안에서 취한다.
\item $\nu$의 올들은 유한하고 잉여체 확대들도 유한하다.
\item $X^\nu$는 정수적 정규 국소 뇌터 스킴들의 서로소 합집합이고,
각 아핀 열린집합은 데데킨트 정역들의 유한곱의 스펙트럼이다.
\end{enumerate}
\end{lemma}

\begin{proof}
많은 결과는 사실 정규화 사상의 일반적인 성질이다. Morphisms,
Lemmas \ref{morphisms-lemma-normalization-reduced},
\ref{morphisms-lemma-stalk-normalization},
\ref{morphisms-lemma-normalization-normal}, 및
\ref{morphisms-lemma-normalization-birational}을 보라. 분명하지 않은
것은 올들이 유한하다는 것, 유도된 잉여체 확대들이 유한하다는 것,
그리고 $X^\nu$가 국소적으로 데데킨트 정역의 스펙트럼처럼 보인다는
것이다(따라서 뇌터이다). 이를 보기 위해 $X = \Spec(A)$가 아핀이고
뇌터이며 차원이 $1$이고, $A$가 축약이라고 가정할 수 있다.
그러면 Morphisms, Lemma
\ref{morphisms-lemma-description-normalization}의 설명을 이용하여
$A$가 차원 $1$의 뇌터 정역인 경우로 환원할 수 있다. 이 경우 원하는
성질들은 Algebra, Lemma
\ref{algebra-lemma-integral-closure-Dedekind}에 진술된 형태의
Krull--Akizuki에서 따른다.
\end{proof}

\noindent
물론 다음 보조정리에는 $X$가 축약이 아닌 경우의 변형도 있다.

\begin{lemma}
\label{varieties-lemma-prepare-delta-invariant}
$X$를 차원 $1$인 축약 나가타 스킴이라 하자.
$\nu : X^\nu \to X$를 정규화라 하자. $x \in X$를 닫힌점이라 하자.
그러면
\begin{enumerate}
\item $\nu : X^\nu \to X$는 유한이고 전사이며 쌍유리이다.
\item $\mathcal{O}_X \subset \nu_*\mathcal{O}_{X^\nu}$이고,
$\nu_*\mathcal{O}_{X^\nu}/\mathcal{O}_X$는 값 $\mathcal{Q}_x$를
$x$에서 가지는 마천루층들의 직합이다. 이 값이 영이 아닐
필요충분조건은 $x$가 $X$의 특이점인 것이다.
\item $A' = (\nu_*\mathcal{O}_{X^\nu})_x$는
$A = \mathcal{O}_{X, x}$의 전체 분수환 안에서의 정수적 폐포이다.
\item $\mathcal{Q}_x = A'/A$는 유한 길이를 가지며 그 길이는
$\delta$-불변량, 곧 $X$의 $x$에서의 불변량과 같다.
\item $A'$는 데데킨트 정역들의 유한곱인 반국소환이다.
\item $A^\wedge$는 차원 $1$인 축약 뇌터 완비 국소환이다.
\item $(A')^\wedge$는 $A^\wedge$의 전체 분수환 안에서의 정수적 폐포이다.
\item $(A')^\wedge$는 완비 이산 값매김환들의 유한곱이다.
\item $A'/A \cong (A')^\wedge/A^\wedge$이다.
\end{enumerate}
\end{lemma}

\begin{proof}
Lemma \ref{varieties-lemma-normalize-noetherian-dim-1}의 모든 결과를 사용할 수
있고 실제로 사용한다. $\nu$의 유한성은 Morphisms, Lemma
\ref{morphisms-lemma-nagata-normalization}에서 따른다.
$X$가 차원 $1$인 축약 나가타 스킴이므로 More on Algebra,
Proposition \ref{more-algebra-proposition-ubiquity-J-2}에 의해 정칙
자취는 조밀한 열린집합 $U \subset X$이다. 정칙 스킴은 정규이므로
이는 $\nu$가 $U$ 위에서 동형임을 보인다. $\dim(X) \leq 1$이므로
이는 $\nu$가 닫힌점들의 이산집합 $x \in X$ 위에서만 동형이 아님을
뜻한다. 특히 짧은 완전열
$$
0 \to \mathcal{O}_X \to \nu_*\mathcal{O}_{X^\nu} \to
\bigoplus\nolimits_{x \in X \setminus U} \mathcal{Q}_x \to 0
$$
을 얻는다. Lemma \ref{varieties-lemma-normalize-noetherian-dim-1}에 의해
$\nu_*\mathcal{O}_{X^\nu}$의 줄기들에 대한 설명을 가지므로,
$Q_x = A'/A$의 길이는 실제로 $\delta$-불변량, 곧 $X$의 $x$에서의
불변량과 같다고 결론짓는다. 예를 들어 Lemma
\ref{varieties-lemma-delta-invariant-is-zero}에 의해 $Q_x \not = 0$일
필요충분조건은 $x$가 특이점인 것이다. $A'$가 반국소 데데킨트
정역들의 곱이라는 설명도 Lemma
\ref{varieties-lemma-normalize-noetherian-dim-1}에서 따른다. $A$, $A'$,
$(A')^\wedge$ 사이의 관계는 Lemma
\ref{varieties-lemma-normalization-same-after-completion}와 그 증명에서 보았다.
\end{proof}




\section{아핀 열린집합 찾기}
\label{varieties-section-finding-affine-opens}

\noindent
Properties, Section \ref{properties-section-finding-affine-opens}에서 시작한
논의를 계속한다. 분리 스킴 위에서 여차원 $1$인 주어진 유한 개의 점을
포함하는 아핀 열린집합을 찾을 수 있음이 밝혀진다. Proposition
\ref{varieties-proposition-finite-set-of-points-of-codim-1-in-affine}을 보라.

\medskip\noindent
다음 보조정리는 Descent, Lemma
\ref{descent-lemma-characterize-open-immersion}에서 개선할 것이다.

\begin{lemma}
\label{varieties-lemma-characterize-open-immersion}
$f : X \to Y$를 스킴의 사상이라 하자. $X^0$로 $X$의 기약 성분들의
일반점 집합을 나타내자. 다음을 가정하자.
\begin{enumerate}
\item $f$는 분리이다.
\item 열린 덮개 $X = \bigcup U_i$가 존재하여
$f|_{U_i} : U_i \to Y$가 열린 몰입이다.
\item $\xi, \xi' \in X^0$이고 $\xi \not = \xi'$이면
$f(\xi) \not = f(\xi')$이다.
\end{enumerate}
그러면 $f$는 열린 몰입이다.
\end{lemma}

\begin{proof}
$y = f(x) = f(x')$라 하자. 특수화 $y_0 \leadsto y$를 택하되,
$y_0$는 $Y$의 한 기약 성분의 일반점이라 하자. $f$는 정의역에서
국소적으로 동형이므로 특수화 $x_0 \leadsto x$와
$x'_0 \leadsto x'$를 택하여 $y_0 \leadsto y$로 사상되게 할 수 있다.
$x_0, x'_0 \in X^0$임에 유의하자. 따라서 가정 (3)에 의해
$x_0 = x'_0$이다. $f$가 분리이므로 $x = x'$라고 결론짓는다.
그러므로 $f$는 열린 몰입이다.
\end{proof}

\begin{lemma}
\label{varieties-lemma-local-isomorphism}
$X \to S$를 스킴의 사상이라 하자. $x \in X$를 상이 $s \in S$인
점이라 하자. 다음을 가정하자.
\begin{enumerate}
\item $\mathcal{O}_{X, x} = \mathcal{O}_{S, s}$이다.
\item $S$는 축약이다.
\item $X \to S$는 유한형이다.
\item $S$의 기약 성분은 유한 개이다.
\end{enumerate}
그러면 열린 근방 $U$가 존재하여 $x$를 포함하고 $f|_U$는 열린 몰입이다.
\end{lemma}

\begin{proof}
$S$의 유한 개 기약 성분 중 $s$를 포함하지 않는 것들을 제거할 수 있다.
$S$를 $s$의 아핀 열린 근방으로 바꿀 수 있다. $X$를 $x$의 아핀 열린
근방으로 바꿀 수 있다. $S = \Spec(A)$이고 $X = \Spec(B)$라 쓰자.
$\mathfrak q \subset B$, resp.\ $\mathfrak p \subset A$를 $x$,
resp.\ $s$에 대응하는 소 아이디얼이라 하자. $A$가 축약이고 $A$의
모든 극소 소 아이디얼이 $\mathfrak p$에 포함되므로
$A \subset A_\mathfrak p$임을 알 수 있다. $X \to S$가 유한형이므로
$B$는 $A$ 위 유한형이다. $b_1, \ldots, b_n \in B$를 택하되, 이들이
$B$를 $A$ 위에서 생성하게 하자. $A_\mathfrak p = B_\mathfrak q$이므로
$f \in A$, $f \not \in \mathfrak p$ 및 $a_i \in A$를 택하여
$b_i$와 $a_i/f$가 $B_\mathfrak q$에서 같은 상을 가지게 할 수 있다.
따라서 $g \in B$, $g \not \in \mathfrak q$를 택하여
$g(fb_i - a_i) = 0$이 $B$에서 성립하게 할 수 있다. 그러므로
$A_f \to B_{fg}$의 상은 $b_1, \ldots, b_n$의 상들을 포함하고,
특히 $g$의 상도 포함한다. $n \geq 0$과 $f' \in A$를 택하여
$f'/f^n$이 $g$의 $B_{fg}$에서의 상으로 사상되게 하자.
$A_\mathfrak p = B_\mathfrak q$이므로 $f' \not \in \mathfrak p$임을
알 수 있다. 따라서 $A_{ff'} \to B_{fg}$는 전사이다. 마지막으로
$A_{ff'} \subset A_\mathfrak p = B_\mathfrak q$이므로(위 참조)
사상 $A_{ff'} \to B_{fg}$는 단사이고, 따라서 동형이다.
\end{proof}

\begin{lemma}
\label{varieties-lemma-points-in-affine}
$f : T \to X$를 스킴의 사상이라 하자. $X^0$, resp.\ $T^0$로
기약 성분들의 일반점 집합들을 나타내자.
$t_1, \ldots, t_m \in T$를 상들이 $x_j = f(t_j)$인 유한 점 집합이라
하자. 다음을 가정하자.
\begin{enumerate}
\item $T$는 아핀이다.
\item $X$는 준분리이다.
\item $X^0$는 유한하다.
\item $f(T^0) \subset X^0$이고 $f : T^0 \to X^0$는 단사이다.
\item $\mathcal{O}_{X, x_j} = \mathcal{O}_{T, t_j}$이다.
\end{enumerate}
그러면 $X$에는 $x_1, \ldots, x_r$을 포함하는 아핀 열린집합이 존재한다.
\end{lemma}

\begin{proof}
Limits, Proposition \ref{limits-proposition-affine}을 사용하면 $X$와
$T$가 축약인 경우로 즉시 환원된다. 세부사항은 생략한다.

\medskip\noindent
$X$와 $T$가 축약이라고 가정하자. Limits, Lemma
\ref{limits-lemma-relative-approximation}에 의해 $T = \lim_{i \in I} T_i$를
$X$ 위 유한 표시 스킴들의 유향 극한으로 쓸 수 있으며 전이 사상들은
아핀이다. Limits, Lemma \ref{limits-lemma-limit-affine}에 의해
$i \in I$를 택하되 $T_i$가 아핀이 되게 하자.
$T_i = \Spec(R_i)$이고 $T = \Spec(R)$라 쓰자. $R' \subset R$를
$R_i \to R$의 상이라 하자. 그러면 $T' = \Spec(R')$는 아핀이고
축약이며 $X$ 위 유한형이고, $T \to T'$는 우세이다.
$j = 1, \ldots, r$에 대해 $t'_j \in T'$를 $t_j$의 상이라 하자.
국소환 사상들을 생각하자.
$$
\mathcal{O}_{X, x_j} \to
\mathcal{O}_{T', t'_j} \to
\mathcal{O}_{T, t_j}
$$
$(T')^0$로 $T'$의 기약 성분들의 일반점 집합을 나타내자.
특수화 $\xi \leadsto t'_j$를 택하되 $\xi \in (T')^0$라 하자.
$T \to T'$가 우세이므로 $\eta \in T^0$를 택하여 $\xi$로 사상되게
할 수 있다(주의: 선험적으로 $\eta$가 $t_j$로 특수화되는지는 알지 못한다).
가정 (3)을 $\eta$에 적용하면 상 $\theta$, 즉 $\xi$의 $X$에서의 상이
$\mathcal{O}_{X, x_j}$의 극소 소 아이디얼에 대응함을 알 수 있다.
(5)의 동형을 통해 $\xi$를 올리면 특수화 $\eta' \leadsto t_j$를 얻으며,
여기서 $\eta' \in T^0$는 $\theta \leadsto x_j$로 사상된다.
(4)의 단사성으로부터 $\eta = \eta'$임을 안다. 따라서
$\mathcal{O}_{T', t'_j}$의 모든 극소 소 아이디얼은
$\mathcal{O}_{T, t_j}$의 한 극소 소 아이디얼 아래에 놓인다.
그러므로 $\mathcal{O}_{T', t'_j} \to \mathcal{O}_{T, t_j}$는
단사이며, 위의 두 사상은 모두 동형이다.

\medskip\noindent
Lemma \ref{varieties-lemma-local-isomorphism}에 의해 열린집합 $U \subset T'$가
존재하여 모든 점 $t'_j$를 포함하고 $U \to X$는 Lemma
\ref{varieties-lemma-characterize-open-immersion}에서와 같은 국소 동형이다.
그 보조정리에 의해 $U \to X$는 열린 몰입이다. 마지막으로 Properties,
Lemma \ref{properties-lemma-ample-finite-set-in-affine}에 의해 열린집합
$W \subset U \subset T'$를 찾아 모든 $t'_j$를 포함하게 할 수 있다.
$W$의 $X$에서의 상이 원하는 아핀 열린집합이다.
\end{proof}

\begin{lemma}
\label{varieties-lemma-finite-set-codim-1-points-in-affine}
$X$를 정수적 분리 스킴이라 하자. $x_1, \ldots, x_r \in X$를
$\mathcal{O}_{X, x_i}$가 뇌터이고 차원이 $\leq 1$인 유한 점 집합이라
하자. 그러면 $X$에는 모든 $x_1, \ldots, x_r$을 포함하는 아핀 열린
부분스킴이 존재한다.
\end{lemma}

\begin{proof}
$K$를 $X$의 유리함수체라 하자. $A_i = \mathcal{O}_{X, x_i}$로 놓자.
그러면 $A_i \subset K$이고 $K$는 $A_i$의 분수체이다. $X$가
분리이고 $x_i \not = x_j$이므로 값매김환 $\mathcal{O} \subset K$가
$A_i$와 $A_j$를 모두 지배하는 일은 있을 수 없다. 실제로 그림
$$
\xymatrix{
\Spec(\mathcal{O}) \ar[r] \ar[d] & \Spec(A_1) \ar[d] \\
\Spec(A_2) \ar[r] & X
}
$$
을 생각하고 분리성의 값매김 판정법(Schemes, Lemma
\ref{schemes-lemma-separated-implies-valuative})을 적용하면
$x_i = x_j$를 얻게 된다. 따라서 Lemma \ref{varieties-lemma-glue-separated}에
의해 $A_i \otimes A_j \to K$는 모든 $i \not = j$에 대해 전사이다.
Lemma \ref{varieties-lemma-glue-a-bunch-of-local-rings}에 의해
$A = A_1 \cap \ldots \cap A_r$은 정확히 $r$개의 극대 아이디얼
$\mathfrak m_1, \ldots, \mathfrak m_r$을 가지며
$A_i = A_{\mathfrak m_i}$인 뇌터 반국소환이다. 또한
$$
\Spec(A) = \Spec(A_1) \cup \ldots \cup \Spec(A_r)
$$
는 열린 덮개이고, 이 덮개의 서로 다른 두 조각의 교집합은
$\Spec(K)$이다. 따라서 주어진 사상들 $\Spec(A_i) \to X$는 스킴의
사상
$$
\Spec(A) \longrightarrow X
$$
으로 이어 붙으며, 이는 $\mathfrak m_i$를 $x_i$로 보내고 국소환의
동형들을 유도한다. 이제 결과는 Lemma \ref{varieties-lemma-points-in-affine}에서
따른다.
\end{proof}

\begin{lemma}
\label{varieties-lemma-extra-silly}
$A$를 환, $I \subset A$를 아이디얼,
$\mathfrak p_1, \ldots, \mathfrak p_r$을 $A$의 소 아이디얼들,
$\overline{f} \in A/I$를 원소라 하자. $I \not \subset \mathfrak p_i$가
모든 $i$에 대해 성립하면, $f \in A$, $f \not \in \mathfrak p_i$가
존재하여 $\overline{f}$로 $A/I$에서 사상된다.
\end{lemma}

\begin{proof}
$\mathfrak p_i$들 사이에 포함관계가 없다고 가정할 수 있다
(더 작은 소 아이디얼들을 제거한다). 먼저 임의의 $f \in A$를 택하여
$\overline{f}$를 올리게 하자. $S$를 $s \in \{1, \ldots, r\}$ 중
$f \in \mathfrak p_s$인 것들의 집합이라 하자. $S$가 공집합이면 끝이다.
그렇지 않으면 아이디얼 $J = I \prod_{i \not \in S} \mathfrak p_i$를
생각하자. $J$는 $\mathfrak p_s$에 포함되지 않는데, 여기서
$s \in S$이다. 실제로 $\mathfrak p_i$들 사이에 포함관계가 없고
$I$는 어느 $\mathfrak p_i$에도 포함되지 않는다. 따라서 Algebra,
Lemma \ref{algebra-lemma-silly}에 의해 $g \in J$를 택하여
$g \not \in \mathfrak p_s$가 모든 $s \in S$에 대해 성립하게 할 수 있다.
그러면 $f + g$가 보조정리에서 제기한 문제의 해이다.
\end{proof}

\begin{lemma}
\label{varieties-lemma-finite-set-codim-1-points-in-affine-per-component}
$X$를 스킴이라 하자. $T \subset X$를 유한 점 집합이라 하자.
다음을 가정하자.
\begin{enumerate}
\item $X$는 유한 개의 기약 성분 $Z_1, \ldots, Z_t$를 가진다.
\item $Z_i \cap T$는 $Z_i$에 대응하는 유도된 축약 부분스킴의 한
아핀 열린집합에 포함된다.
\end{enumerate}
그러면 $X$에는 $T$를 포함하는 아핀 열린 부분스킴이 존재한다.
\end{lemma}

\begin{proof}
Limits, Proposition \ref{limits-proposition-affine}을 사용하면 $X$가
축약인 경우로 즉시 환원된다. 세부사항은 생략한다. 증명의 나머지에서는
$X$의 모든 닫힌 부분집합에 유도된 축약 닫힌 부분스킴 구조를 준다.

\medskip\noindent
아핀 열린집합 $U \subset Z_1 \cup \ldots \cup Z_r$을 찾아
$T \cap (Z_1 \cup \ldots \cup Z_r)$을 포함하게 할 수 있음을 귀납법으로
보이자. $r = 1$인 경우에는 가정으로 성립한다. $r > 1$이라 하고,
$U \subset Z_1 \cup \ldots \cup Z_{r -  1}$을 아핀 열린집합으로 택하되
$T \cap (Z_1 \cup \ldots \cup Z_{r - 1})$을 포함하게 하자.
$V \subset X_r$을 $T \cap Z_r$을 포함하는 아핀 열린집합이라 하자
(가정에 의해 존재한다). 그러면 $U \cap V$는
$T \cap ( Z_1 \cup \ldots \cup Z_{r - 1} ) \cap Z_r$을 포함한다.
따라서
$$
\Delta = (U \cap Z_r) \setminus (U \cap V)
$$
는 $T$의 어떤 원소도 포함하지 않는다. $\Delta$는 $U$의 닫힌
부분집합임에 유의하자. 소 아이디얼 회피(Algebra, Lemma
\ref{algebra-lemma-silly})에 의해 $U'$를 $U$의 표준 열린집합으로
택하여 $T \cap U$를 포함하고 $\Delta$와 만나지 않게 할 수 있다. 즉,
$U' \cap Z_r \subset U \cap V$이다. $U$를 $U'$로 바꾸면
$U \cap V$가 $U$에서 닫혀 있다고 가정할 수 있다.

\medskip\noindent
위와 같은 논증에 의해 집합
$\Delta' = (U \cap (Z_1 \cup \ldots \cup Z_{r - 1})) \setminus (U \cap V)$
도 $T$의 어떤 원소도 포함하지 않는다. 이를 사용하여
$h \in \mathcal{O}(V)$를 택해서 $D(h) \subset V$가 $T \cap V$를
포함하고 $U \cap D(h) \subset U \cap V$가 되게 할 수 있다.
$U \cap V$가 $U$에서 닫혀 있다는 사실과 Lemma
\ref{varieties-lemma-extra-silly}를 사용하여 원소 $g \in \mathcal{O}(U)$를
찾을 수 있고, 그 제한은 $U \cap V$ 위에서 $h$의 $U \cap V$로의
제한과 같으며, 또한 $T \cap U \subset D(g)$이다.
그런 다음 $U$를 $D(g)$로, $V$를 $D(h)$로 바꾸어 $U \cap V$가
$U$와 $V$ 모두에서 닫힌 경우에 도달한다. 이 경우 스킴 $U \cup V$는
Limits, Lemma \ref{limits-lemma-affines-glued-in-closed-affine}에 의해
아핀이다. 이로써 귀납 단계를 증명하고 따라서 보조정리를 증명한다.
\end{proof}

\noindent
위의 내용에서 다음 결론을 얻을 수 있다.

\begin{proposition}
\label{varieties-proposition-finite-set-of-points-of-codim-1-in-affine}
$X$를 모든 준콤팩트 열린집합이 유한 개의 기약 성분을 가지는 분리
스킴이라 하자. $x_1, \ldots, x_r \in X$를
$\mathcal{O}_{X, x_i}$가 뇌터이고 차원이 $\leq 1$인 점들이라 하자.
그러면 $X$에는 모든 $x_1, \ldots, x_r$을 포함하는 아핀 열린
부분스킴이 존재한다.
\end{proposition}

\begin{proof}
$X$를 $x_1, \ldots, x_r$을 포함하는 준콤팩트 열린집합으로 바꿀 수
있으므로 $X$의 기약 성분이 유한 개라고 가정할 수 있다. Lemma
\ref{varieties-lemma-finite-set-codim-1-points-in-affine-per-component}에 의해
$X$가 정수적인 경우로 환원된다. 이 경우는 Lemma
\ref{varieties-lemma-finite-set-codim-1-points-in-affine}이다.
\end{proof}





\section{곡선}
\label{varieties-section-curves}

\noindent
Stacks 프로젝트에서는 다음을 곡선의 정의로 사용한다.

\begin{definition}
\label{varieties-definition-curve}
$k$를 체라 하자. {\it 곡선}은 차원 $1$인 $k$ 위의 대수다양체이다.
\end{definition}

\noindent
$k$ 위 곡선의 표준적인 두 예는 아핀 직선 $\mathbf{A}^1_k$와 사영 직선
$\mathbf{P}^1_k$이다. 스킴 $X = \Spec(k[x, y]/(f))$가 곡선일
필요충분조건은 $f \in k[x, y]$가 기약인 것이다.

\medskip\noindent
우리의 곡선 정의에는 대수다양체 정의와 같은 문제가 있다. Definition
\ref{varieties-definition-variety} 뒤의 논의를 보라. 또한 이 정의에 따르면 모든
곡선에는 정의체가 지정되어 있다. 예를 들어
$X = \Spec(\mathbf{C}[x])$는 $\mathbf{C}$ 위 곡선이지만
$\mathbf{R}$ 위 곡선으로도 볼 수 있다. 스킴 $\Spec(\mathbf{Z})$는
곡선이 아니지만, 스킴 $\Spec(\mathbf{Z})$와
$\mathbf{A}^1_{\mathbf{F}_p}$는 여러 면에서 비슷하게 행동한다.

\begin{lemma}
\label{varieties-lemma-proper-minus-point}
$X$를 차원 $> 0$인 체 $k$ 위의 분리 기약 스킴이라 하자.
$x \in X$를 닫힌점이라 하자. 열린 부분스킴 $X \setminus \{x\}$는
$k$ 위에서 고유가 아니다.
\end{lemma}

\begin{proof}
$X$가 기약이므로 $U = X \setminus \{x\}$는 $X$에서 닫혀 있지 않다.
특히 몰입 $U \to X$는 고유가 아니다. Morphisms, Lemma
\ref{morphisms-lemma-image-proper-scheme-closed}에 의해(여기서는
$X$가 분리임을 사용한다) $U \to \Spec(k)$도 고유가 아니다.
\end{proof}

\begin{lemma}
\label{varieties-lemma-dim-1-quasi-projective}
$X$를 체 $k$ 위의 분리 유한형 스킴이라 하자.
$\dim(X) \leq 1$이면 $X$는 $k$ 위에서 H-준사영이다.
\end{lemma}

\begin{proof}
Proposition \ref{varieties-proposition-dim-1-noetherian-separated-has-ample}에 의해
스킴 $X$는 풍부한 가역층 $\mathcal{L}$을 가진다. Morphisms, Lemma
\ref{morphisms-lemma-quasi-projective-finite-type-over-S}에 의해
$X$는 $\mathbf{P}^n_k$의 국소 닫힌 부분스킴과 동형이며, 이 동형은
$\Spec(k)$ 위의 것이다. 이는 $k$ 위 H-준사영의 정의이다. Morphisms, Definition
\ref{morphisms-definition-quasi-projective}을 보라.
\end{proof}

\begin{lemma}
\label{varieties-lemma-dim-1-proper-projective}
$X$를 체 $k$ 위의 고유 스킴이라 하자.
$\dim(X) \leq 1$이면 $X$는 $k$ 위에서 H-사영이다.
\end{lemma}

\begin{proof}
Lemma \ref{varieties-lemma-dim-1-quasi-projective}에 의해 $X$는
$\mathbf{P}^n_k$의 국소 닫힌 부분스킴이며, 여기서 $k$는 어떤 체이다.
$X$가 $k$ 위에서 고유이므로 $X$는 $\mathbf{P}^n_k$의 닫힌 부분스킴이다
(Morphisms, Lemma \ref{morphisms-lemma-image-proper-scheme-closed}).
\end{proof}

\begin{lemma}
\label{varieties-lemma-dim-1-projective-completion}
$X$를 $k$ 위의 분리 유한형 스킴이라 하자.
$\dim(X) \leq 1$이면 다음 성질들을 가지는 열린 몰입
$j : X \to \overline{X}$가 존재한다.
\begin{enumerate}
\item $\overline{X}$는 $k$ 위에서 H-사영이다. 즉, $\overline{X}$는
$\mathbf{P}^d_k$의 닫힌 부분스킴이며, 여기서 $d$는 어떤 정수이다.
\item $j(X) \subset \overline{X}$는 조밀하고 스킴론적으로 조밀하다.
\item $\overline{X} \setminus X = \{x_1, \ldots, x_n\}$이고,
여기서 $x_i \in \overline{X}$는 닫힌점들이다.
\end{enumerate}
\end{lemma}

\begin{proof}
Lemma \ref{varieties-lemma-dim-1-quasi-projective}에 의해 $X$가
$\mathbf{P}^d_k$의 국소 닫힌 부분스킴이라고 가정할 수 있으며,
여기서 $d$는 어떤 정수이다.
$\overline{X} \subset \mathbf{P}^d_k$를 $X \to \mathbf{P}^d_k$의
스킴론적 상이라 하자. Morphisms, Definition
\ref{morphisms-definition-scheme-theoretic-image}을 보라.
Morphisms, Lemma \ref{morphisms-lemma-quasi-compact-immersion}의 설명으로
성질 (1)과 (2)를 얻는다. 일반점들을 살펴보면
$\dim(X) = 1 \Rightarrow \dim(\overline{X}) = 1$이다. 예를 들어 Lemma
\ref{varieties-lemma-dimension-locally-algebraic}을 보라. $\overline{X}$가
뇌터이므로 $\overline{X} \setminus X = \{x_1, \ldots, x_n\}$는
유한한 닫힌점 집합이다.
\end{proof}

\begin{lemma}
\label{varieties-lemma-reduced-dim-1-projective-completion}
$X$를 $k$ 위의 분리 유한형 스킴이라 하자. $X$가 축약이고
$\dim(X) \leq 1$이면 다음 성질들을 가지는 열린 몰입
$j : X \to \overline{X}$가 존재한다.
\begin{enumerate}
\item $\overline{X}$는 $k$ 위에서 H-사영이다. 즉, $\overline{X}$는
$\mathbf{P}^d_k$의 닫힌 부분스킴이며, 여기서 $d$는 어떤 정수이다.
\item $j(X) \subset \overline{X}$는 조밀하고 스킴론적으로 조밀하다.
\item $\overline{X} \setminus X = \{x_1, \ldots, x_n\}$이고,
여기서 $x_i \in \overline{X}$는 닫힌점들이다.
\item 국소환 $\mathcal{O}_{\overline{X}, x_i}$들은
$i = 1, \ldots, n$에 대해 이산 값매김환이다.
\end{enumerate}
\end{lemma}

\begin{proof}
$j : X \to \overline{X}$를 Lemma
\ref{varieties-lemma-dim-1-projective-completion}에서와 같이 택하자.
$X'$를 $\overline{X}$의 $X$ 안에서의 정규화라 하자. Lemma
\ref{varieties-lemma-relative-normalization-finite}에 의해 사상
$X' \to \overline{X}$는 유한이다. Morphisms, Lemma
\ref{morphisms-lemma-finite-projective}에 의해
$X' \to \overline{X}$는 사영적이다. Morphisms, Lemma
\ref{morphisms-lemma-projective-over-quasi-projective-is-H-projective}에
의해 $X' \to \overline{X}$는 H-사영이다. Morphisms, Lemma
\ref{morphisms-lemma-H-projective-composition}에 의해
$X' \to \Spec(k)$는 H-사영이다.
$\{x'_1, \ldots, x'_m\} \subset X'$를
$\{x_1, \ldots, x_n\} = \overline{X} \setminus X$의 역상이라 하자.
$\dim(\mathcal{O}_{X', x'_i}) = 1$은 모든 $1 \leq i \leq m$에 대해
성립한다. 따라서 Morphisms, Lemma
\ref{morphisms-lemma-relative-normalization-normal-codim-1}에 의해
국소환 $\mathcal{O}_{X', x'}$들은 이산 값매김환이다.
그러면 $X \to X'$와 $\{x'_1, \ldots, x'_m\}$가 원하는 것이다.
\end{proof}

\begin{lemma}
\label{varieties-lemma-dim-1-projective-completion-after-insep}
$X$를 $k$ 위의 분리 유한형 스킴이라 하고 $\dim(X) \leq 1$이라 하자.
그러면 다음 가환 그림이 존재한다.
$$
\xymatrix{
\overline{Y}_1 \amalg \ldots \amalg \overline{Y}_n \ar[rd] &
Y_1 \amalg \ldots \amalg Y_n \ar[r]_-\nu \ar[d] \ar[l]^j &
X_{k'} \ar[r] \ar[d] &
X \ar[d]^f \\
& \Spec(k'_1) \amalg \ldots \amalg \Spec(k'_n) \ar[r] &
\Spec(k') \ar[r] &
\Spec(k)
}
$$
이 그림은 다음 성질들을 가진다.
\begin{enumerate}
\item $k'/k$는 체의 유한 순비분리 확대이다.
\item $\nu$는 $X_{k'}$의 정규화이다.
\item $j$는 조밀한 상을 가지는 열린 몰입이다.
\item $k'_i/k'$는 $i = 1, \ldots, n$에 대해 유한 분리 확대이다.
\item $\overline{Y}_i$는 차원이 $\leq 1$인 $k'_i$ 위의 매끄럽고
사영적이며 기하적으로 기약인 스킴이다.
\end{enumerate}
\end{lemma}

\begin{proof}
$X$를 그 축약으로 바꿀 수 있으므로 $X$가 축약이라고 가정하고 실제로
그렇게 한다. $X \to \overline{X}$를 Lemma
\ref{varieties-lemma-reduced-dim-1-projective-completion}에서와 같이 택하자.
$\overline{X}$에 대해 보조정리를 보일 수 있으면 $X$에 대해서도
보조정리가 따른다(세부사항은 생략한다). 따라서 $X$가 사영적이라고
가정하고 실제로 그렇게 한다.

\medskip\noindent
$k'/k$를 유한 순비분리 확대 중 $X_{k'}$의 정규화가 $k'$ 위에서
기하적으로 정규가 되게 하는 것으로 택하자. Lemma
\ref{varieties-lemma-finite-extension-geometrically-normal}을 보라.
$Y = (X_{k'})^\nu$로 그 정규화를 나타내자. 정규화의 성질은 Section
\ref{varieties-section-normalization}을 보라. $Y$는 차원 $\leq 1$에서 정규와
정칙이 같으므로 기하적으로 정칙이다. Properties, Lemma
\ref{properties-lemma-normal-dimension-1-regular}을 보라.
따라서 Lemma \ref{varieties-lemma-geometrically-regular-smooth}에 의해
$Y$는 $k'$ 위에서 매끄럽다. $Y = Y_1 \amalg \ldots \amalg Y_n$을
$Y$의 기약 성분들로의 분해라 하자.
$k'_i = \Gamma(Y_i, \mathcal{O}_{Y_i})$로 놓자. 이들은 Lemma
\ref{varieties-lemma-proper-geometrically-reduced-global-sections}에 의해
$k'$의 유한 분리 확대이다. Lemma \ref{varieties-lemma-baby-stein}으로
증명이 끝난다.
\end{proof}

\begin{lemma}
\label{varieties-lemma-regular-point-on-curve}
$k$를 체라 하자. $X$를 $k$ 위의 곡선이라 하자.
$x \in X$를 닫힌점이라 하자. $x$를 $X$의 (축약) 닫힌 부분스킴으로
생각하고 그 아이디얼층을 $\mathcal{I}$라 하자. 다음은 동치이다.
\begin{enumerate}
\item $\mathcal{O}_{X, x}$는 정칙이다.
\item $\mathcal{O}_{X, x}$는 정규이다.
\item $\mathcal{O}_{X, x}$는 이산 값매김환이다.
\item $\mathcal{I}$는 가역 $\mathcal{O}_X$-가군이다.
\item $x$는 $X$ 위의 유효 카르티에 인자이다.
\end{enumerate}
$k$가 완전체이거나 $\kappa(x)$가 $k$ 위에서 분리이면, 이들은 다음과도
동치이다.
\begin{enumerate}
\item[(6)] $X \to \Spec(k)$는 $x$에서 매끄럽다.
\end{enumerate}
\end{lemma}

\begin{proof}
$X$가 곡선이므로 국소환 $\mathcal{O}_{X, x}$는 차원 $1$인 뇌터
국소 정역이다(Lemma \ref{varieties-lemma-dimension-locally-algebraic}).
(4)와 (5)는 정의에 의해 동치이고, 이것들은
$\mathcal{I}_x = \mathfrak m_x \subset \mathcal{O}_{X, x}$가 하나의
생성원을 가진다는 것과 동치이다(Divisors, Lemma
\ref{divisors-lemma-effective-Cartier-in-points}). 따라서
(1), (2), (3), (4), (5)의 동치는 Algebra, Lemma
\ref{algebra-lemma-characterize-dvr}에서 따른다. 마지막 명제는 $k$가
완전체인 경우 Lemma
\ref{varieties-lemma-dense-smooth-open-variety-over-perfect-field}에서 따른다.
$\kappa(x)/k$가 분리이면 동치는 Algebra, Lemma
\ref{algebra-lemma-separable-smooth}에서 따른다.
\end{proof}

\begin{remark}
\label{varieties-remark-smooth-projective-compactification}
$k$를 체라 하자. $X$를 $k$ 위의 정칙 곡선이라 하자. Lemmas
\ref{varieties-lemma-regular-point-on-curve}와
\ref{varieties-lemma-reduced-dim-1-projective-completion}에 의해 비특이 사영 곡선
$\overline{X}$가 존재하며, 이는 $X$의 콤팩트화이다. 즉, 여집합이 유한 개의
닫힌점으로 이루어진 열린 몰입 $j : X \to \overline{X}$가 존재한다.
$k$가 완전체이면 $X$와 $\overline{X}$는 $k$ 위에서 매끄럽고,
$\overline{X}$는 $X$의 매끄러운 사영 콤팩트화이다.
\end{remark}

\noindent
$X$가 $k$ 위의 아핀 스킴이고 $k$ 위에서 고유이면 $X$는 $k$ 위에서
유한임에 유의하자(Morphisms, Lemma
\ref{morphisms-lemma-finite-proper}). 따라서 차원은 $0$이다
(Algebra, Lemma \ref{algebra-lemma-finite-dimensional-algebra} 및
Proposition \ref{algebra-proposition-dimension-zero-ring}).
그러므로 차원 $> 0$인 $k$ 위의 스킴은 아핀이면서 $k$ 위에서 고유일 수 없다.
따라서 다음 보조정리의 두 가능성은 서로 배타적이다.

\begin{lemma}
\label{varieties-lemma-curve-affine-projective}
$X$를 $k$ 위의 곡선이라 하자. 그러면 $X$는 아핀 스킴이거나,
$X$는 $k$ 위에서 H-사영이다.
\end{lemma}

\begin{proof}
$X \to \overline{X}$를
$\overline{X} \setminus X = \{x_1, \ldots, x_r\}$가 되도록 Lemma
\ref{varieties-lemma-reduced-dim-1-projective-completion}에서와 같이 택하자.
그러면 $\overline{X}$도 곡선이다. $r = 0$이면
$X = \overline{X}$는 $k$ 위에서 H-사영이다. 따라서 $r \geq 1$이라
가정할 수 있고 목표는 $X$가 아핀임을 보이는 것이다. Lemma
\ref{varieties-lemma-open-in-affine-curve-affine}에 의해
$\overline{X} \setminus \{x_1\}$가 아핀임을 보이면 충분하다.
이는 다음 문단의 주장으로 환원한다.

\medskip\noindent
$X$를 $k$ 위의 H-사영 곡선이라 하자. $x \in X$를
$\mathcal{O}_{X, x}$가 이산 값매김환인 닫힌점이라 하자. 주장:
$U = X \setminus \{x\}$는 아핀이다. Lemma
\ref{varieties-lemma-regular-point-on-curve}에 의해 점 $x$는 $X$의 유효 카르티에
인자를 정의한다. $n \geq 1$에 대해 $nx = x + \ldots + x$로 $n$중 합을
나타내자. Divisors, Definition
\ref{divisors-definition-sum-effective-Cartier-divisors}를 보라.
$\mathcal{O}_{nx}$로 $nx$의 구조층을 나타내되 이를 $X$ 위의 연접
가군으로 보자. 국소 스킴 $nx$ 위의 모든 가역 가군은 자명하므로 Divisors,
Remark \ref{divisors-remark-ses-regular-section}의 첫 번째 짧은 완전열은
이 경우
$$
0 \to \mathcal{O}_X
\xrightarrow{1} \mathcal{O}_X(nx) \to \mathcal{O}_{nx} \to 0
$$
이다. $\dim_k H^0(X, \mathcal{O}_{nx}) \geq n$임에 유의하자. 실제로
Lemma \ref{varieties-lemma-chi-tensor-finite}에 의해
$H^0(X, \mathcal{O}_{nx}) = \mathcal{O}_{X, x}/(\pi^n)$이고, 여기서
$\pi$는 $\mathcal{O}_{X, x}$의 균일화원이며 거듭제곱 $\pi^i$들은
$k$-선형 독립인 원소들로 $\mathcal{O}_{X, x}/(\pi^n)$에 사상된다.
여기서 $i = 0, 1, \ldots, n - 1$이다.
$\dim_k H^1(X, \mathcal{O}_X) < \infty$는 Cohomology of Schemes,
Lemma \ref{coherent-lemma-proper-over-affine-cohomology-finite}에서
따른다. $n > \dim_k H^1(X, \mathcal{O}_X)$이면 긴 완전 코호몰로지
열로부터 $s \in \Gamma(X, \mathcal{O}_X(nx))$가 존재하고
$\mathcal{O}_X$의 단면은 아니라고 결론짓는다.
이 성질을 가지는 $n$을 최소로 택하면 $s$는 줄기
$\left(\mathcal{O}_X(nx)\right)_x$의 생성원으로 사상된다. 그렇지 않으면
$\mathcal{O}_X((n - 1)x) \subset \mathcal{O}_X(nx)$의 단면을 정의할
것이기 때문이다. 이 $n$에 대해 $s_0 = 1$과 $s_1 = s$가 가역 가군
$\mathcal{L} = \mathcal{O}_X(nx)$을 생성한다고 결론짓는다.

\medskip\noindent
Constructions, Section \ref{constructions-section-projective-space}의
대응하는 사상
$f = \varphi_{\mathcal{L}, (s_0, s_1)} : X \to \mathbf{P}^1_k$를
생각하자. $D_{+}(T_0)$의 역상은 $U = X \setminus \{x\}$이다.
단면 $s_0$가 $\mathcal{L}$의 단면으로서 $x$에서만 영이 되기 때문이다.
특히 $f$는 상수가 아니다. 즉, $\Im(f)$는 둘 이상의 점을 가진다.
따라서 $f$는 일반점 $\eta$, 즉 $X$의 일반점을 $\mathbf{P}^1_k$의 일반점으로
사상해야 한다. 그러므로 $y \in \mathbf{P}^1_k$가 닫힌점이면
$f^{-1}(\{y\})$는 $X$의 닫힌집합이고 $\eta$를 포함하지 않으며, 따라서
유한하다. 마지막으로 $f$는 고유이다\footnote{실제로 H-사영 대수다양체는
Morphisms, Lemma
\ref{morphisms-lemma-projective-is-quasi-projective-proper}에 의해 고유
대수다양체이다. 정의역이 고유 대수다양체인 대수다양체의 사상은
Morphisms, Lemma \ref{morphisms-lemma-image-proper-scheme-closed}에 의해
고유 사상이다.}. Cohomology of Schemes, Lemma
\ref{coherent-lemma-proper-finite-fibre-finite-in-neighbourhood}\footnote{형식
함수 정리에 의존하는 이 보조정리의 사용을 피할 수 있다. 실제로 $X$는
사영적이므로, 유한한 올들을 가지는 고유 사상 $f : X \to Y$가
$k$ 위의 준사영 스킴들 사이에 있을 때 유한임을 보이면 충분하다. 이를
위해 $X$의 아핀 열린집합을 택하여 $f$의 한 점 $y$ 위의 올을 포함하게 한다.
$k$ 위 준사영 스킴의 유한 점 집합은 아핀에 포함된다는 사실을 사용한다.
$Y$를 $y$의 작은 아핀 근방으로 줄이면 아핀들 사이의 고유 사상의
경우로 환원된다. 그러한 사상은 Morphisms, Lemma
\ref{morphisms-lemma-integral-universally-closed}에 의해 유한이다.}에
의해 $f$는 유한이다. 따라서 $U = f^{-1}(D_{+}(T_0))$는 아핀이다.
\end{proof}

\noindent
다음 보조정리를 Lemma \ref{varieties-lemma-proper-minus-point}와 결합하면 다음을
알 수 있다. $X$가 차원 $1$이고 $k$ 위 유한형인 분리 스킴일 때,
$X \setminus Z$는 아핀이다. 단, 닫힌 부분집합 $Z$가 $X$의 모든
기약 성분과 만난다고 가정한다.

\begin{lemma}
\label{varieties-lemma-dim-1-nonproper-affine}
$X$를 $k$ 위의 분리 유한형 스킴이라 하자. $\dim(X) \leq 1$이고
$X$의 어떤 기약 성분도 차원 $1$인 고유 스킴이 아니면 $X$는 아핀이다.
\end{lemma}

\begin{proof}
$X = \bigcup X_i$를 $X$의 기약 성분들로의 분해라 하자. $X_i$를
정수적 스킴으로 생각한다(유도된 축약 스킴 구조를 사용한다. Schemes,
Definition \ref{schemes-definition-reduced-induced-scheme}을 보라).
특히 $X_i$는 한원소 집합(따라서 아핀)이거나 곡선이고, 후자의 경우
Lemma \ref{varieties-lemma-curve-affine-projective}에 의해 아핀이다.
그러면 $\coprod X_i \to X$는 유한 전사이고 $\coprod X_i$는 아핀이다.
따라서 Cohomology of Schemes, Lemma
\ref{coherent-lemma-image-affine-finite-morphism-affine-Noetherian}에 의해
$X$는 아핀이다.
\end{proof}





\section{곡선 위의 차수}
\label{varieties-section-divisors-curves}

\noindent
체 위의 차원 $1$인 고유 스킴에서 가역층, 더 일반적으로 국소 자유층의
차수를 정의하기 시작한다. Section \ref{varieties-section-euler}에서 연접층
$\mathcal{F}$의 오일러 표수를 정의했다. 이 층은 고유 스킴 $X$ 위에
놓이고 기초체는 $k$이며, 공식은 다음과 같다.
$$
\chi(X, \mathcal{F}) = \sum (-1)^i \dim_k H^i(X, \mathcal{F}).
$$

\begin{definition}
\label{varieties-definition-degree-invertible-sheaf}
$k$를 체, $X$를 차원 $\leq 1$인 $k$ 위의 고유 스킴, $\mathcal{L}$을
가역 $\mathcal{O}_X$-가군이라 하자. $\mathcal{L}$의 {\it 차수}를
다음과 같이 정의한다.
$$
\deg(\mathcal{L}) = \chi(X, \mathcal{L}) - \chi(X, \mathcal{O}_X)
$$
더 일반적으로 $\mathcal{E}$가 계수 $n$인 국소 자유층이면
$\mathcal{E}$의 {\it 차수}를 다음과 같이 정의한다.
$$
\deg(\mathcal{E}) = \chi(X, \mathcal{E}) - n\chi(X, \mathcal{O}_X)
$$
\end{definition}

\noindent
이는 삼중항 $\mathcal{E}/X/k$에 의존함에 유의하자. $X$가 비연결이고
$\mathcal{E}$가 유한 국소 자유이지만 상수 계수가 아니면, $\mathcal{E}$를
$X$의 연결 성분들로 제한하여 얻는 차수들을 합하는 식으로 정의를 수정할 수 있다.
$\mathcal{E}$가 단지 연접층이면 정의를 확장하는 몇 가지 서로 다른
방법이 있다\footnote{$X$가 고유 곡선이고 $\mathcal{F}$가 $X$ 위의
연접층이면, 흔히 차수를
$\chi(X, \mathcal{F}) - r\chi(X, \mathcal{O}_X)$로 정의한다. 여기서
$r = \dim_{\kappa(\xi)} \mathcal{F}_\xi$는 $\mathcal{F}$의 계수이고,
$\xi$는 $X$의 일반점이다.}. 일련의 보조정리에서 이 정의가 차수에
기대되는 모든 성질을 가짐을 보인다.

\begin{lemma}
\label{varieties-lemma-degree-base-change}
$k'/k$를 체의 확대라 하자. $X$를 차원 $\leq 1$인 $k$ 위의 고유
스킴이라 하자. $\mathcal{E}$를 국소 자유 $\mathcal{O}_X$-가군 중
상수 계수 $n$인 것으로 하자. 그러면 $\mathcal{E}/X/k$의 차수는
$\mathcal{E}_{k'}/X_{k'}/k'$의 차수와 같다.
\end{lemma}

\begin{proof}
더 정확히 $X_{k'} = X \times_{\Spec(k)} \Spec(k')$로 놓자.
$\mathcal{E}_{k'} = p^*\mathcal{E}$라 하되
$p : X_{k'} \to X$는 사영이라 하자. Cohomology of Schemes, Lemma
\ref{coherent-lemma-flat-base-change-cohomology}에 의해
$H^i(X_{k'}, \mathcal{E}_{k'}) = H^i(X, \mathcal{E}) \otimes_k k'$이고
$H^i(X_{k'}, \mathcal{O}_{X_{k'}}) = H^i(X, \mathcal{O}_X) \otimes_k k'$이다.
따라서 오일러 표수들이 변하지 않고, 그러므로 차수도 변하지 않는다.
\end{proof}

\begin{lemma}
\label{varieties-lemma-degree-additive}
$k$를 체라 하자. $X$를 차원 $\leq 1$인 $k$ 위의 고유 스킴이라 하자.
$0 \to \mathcal{E}_1 \to \mathcal{E}_2 \to \mathcal{E}_3 \to 0$을 각각
유한 상수 계수인 국소 자유 $\mathcal{O}_X$-가군들의 짧은 완전열이라
하자. 그러면
$$
\deg(\mathcal{E}_2) = \deg(\mathcal{E}_1) + \deg(\mathcal{E}_3)
$$
\end{lemma}

\begin{proof}
오일러 표수의 가법성(Lemma
\ref{varieties-lemma-euler-characteristic-additive})과 계수의 가법성에서 즉시
따른다.
\end{proof}

\begin{lemma}
\label{varieties-lemma-degree-birational-pullback}
$k$를 체라 하자. $f : X' \to X$를 차원 $\leq 1$인 $k$ 위의 고유
스킴들 사이의 쌍유리 사상이라 하자. 그러면 상수 계수인 모든 유한
국소 자유층에 대해
$$
\deg(f^*\mathcal{E}) = \deg(\mathcal{E})
$$
이다. 더 일반적으로 $f$가 차원 $1$인 기약 성분들 사이의 전단사를
유도하고 대응하는 일반점들의 국소환 사이에 동형을 유도하면 충분하다.
\end{lemma}

\begin{proof}
사상 $f$는 고유이고(Morphisms, Lemma
\ref{morphisms-lemma-image-proper-scheme-closed}) 올들의 차원은
$\leq 0$이다. 따라서 $f$는 유한이다(Cohomology of Schemes, Lemma
\ref{coherent-lemma-proper-finite-fibre-finite-in-neighbourhood}).
그러므로
$$
Rf_*f^*\mathcal{E} = f_*f^*\mathcal{E} =
\mathcal{E} \otimes_{\mathcal{O}_X} f_*\mathcal{O}_{X'}
$$
이다. $f$가 모든 차원 $1$의 기약 성분의 일반점에서 국소환의 동형을
유도하므로 핵과 여핵
$$
0 \to \mathcal{K} \to \mathcal{O}_X \to f_*\mathcal{O}_{X'}
\to \mathcal{Q} \to 0
$$
의 지지는 차원이 $\leq 0$이다. $\mathcal{E}$가 국소 자유이므로 이를
$\mathcal{E}$와 텐서해도 여전히 완전열임에 유의하자. 다음을 얻는다.
\begin{align*}
\chi(X, \mathcal{E}) - \chi(X', f^*\mathcal{E})
& =
\chi(X, \mathcal{E}) - \chi(X, f_*f^*\mathcal{E}) \\
& =
\chi(X, \mathcal{E}) - \chi(X, \mathcal{E} \otimes f_*\mathcal{O}_{X'}) \\
& =
\chi(X, \mathcal{K} \otimes \mathcal{E}) -
\chi(X, \mathcal{Q} \otimes \mathcal{E}) \\
& =
n\chi(X, \mathcal{K}) -
n\chi(X, \mathcal{Q}) \\
& =
n\chi(X, \mathcal{O}_X) - n\chi(X, f_*\mathcal{O}_{X'}) \\
& =
n\chi(X, \mathcal{O}_X) - n\chi(X', \mathcal{O}_{X'})
\end{align*}
이는 원하는 것을 증명한다. 첫 번째 등식은 $f$가 유한이기 때문이다.
Cohomology of Schemes, Lemma
\ref{coherent-lemma-relative-affine-cohomology}를 보라. 두 번째 등식은
사영 공식에서 따른다. Cohomology, Lemma
\ref{cohomology-lemma-projection-formula}를 보라. 세 번째 등식은
오일러 표수의 가법성에서 따른다. Lemma
\ref{varieties-lemma-euler-characteristic-additive}를 보라. 네 번째 등식은
Lemma \ref{varieties-lemma-chi-tensor-finite}에서 따른다.
\end{proof}

\begin{lemma}
\label{varieties-lemma-degree-on-proper-curve}
$k$를 체라 하자. $X$를 $k$ 위의 고유 곡선이라 하고 그 일반점을
$\xi$라 하자. $\mathcal{E}$를 국소 자유 $\mathcal{O}_X$-가군 중
계수 $n$인 것으로 하고,
$\mathcal{F}$를 연접 $\mathcal{O}_X$-가군이라 하자. 그러면
$$
\chi(X, \mathcal{E} \otimes \mathcal{F}) =
r \deg(\mathcal{E}) + n \chi(X, \mathcal{F})
$$
이다. 여기서 $r = \dim_{\kappa(\xi)} \mathcal{F}_\xi$는
$\mathcal{F}$의 계수이다.
\end{lemma}

\begin{proof}
$\mathcal{P}$를 연접층 $\mathcal{F}$가 $X$ 위에 놓이고 보조정리의 식이
성립한다는 성질이라 하자. Cohomology of Schemes, Lemma
\ref{coherent-lemma-property}의 가정 (1)과 (2)가 $\mathcal{P}$에 대해
성립한다고 주장한다. 실제로 오일러 표수와 계수 $r$은 연접층들의 짧은
완전열에서 가법적이므로 (1)이 성립한다. (2)도 성립한다.
$Z = X$이면 $\mathcal{G} = \mathcal{O}_X$로 택할 수 있고
$\mathcal{P}(\mathcal{O}_X)$는 차수의 정의에 의해 참이다.
$i : Z \to X$가 닫힌점의 포함이면 $\mathcal{G} = i_*\mathcal{O}_Z$로
택할 수 있고, $\mathcal{P}$가 Lemma \ref{varieties-lemma-chi-tensor-finite}와
이 경우 $r = 0$이라는 사실에 의해 성립한다.
\end{proof}

\noindent
$k$를 체라 하자. $X$를 $k$ 위의 차원 $\leq 1$인 유한형 스킴이라
하자. $C_i \subset X$, $i = 1, \ldots, t$를 차원 $1$인 기약
성분들이라 하자. $C_i$에는 유도된 축약 스킴 구조를 주어 스킴으로
본다. $\xi_i \in C_i$를 일반점이라 하자. {\it $C_i$의 $X$ 안에서의
중복도}를 길이
$$
m_i = \text{length}_{\mathcal{O}_{X, \xi_i}} \mathcal{O}_{X, \xi_i}
$$
로 정의한다. $\mathcal{O}_{X, \xi_i}$가 영차원 뇌터 국소환이고 따라서
자기 자신 위에서 유한 길이를 가지므로 이는 잘 정의된다(Algebra,
Proposition \ref{algebra-proposition-dimension-zero-ring}).
추가 정보는 Chow Homology, Section
\ref{chow-section-cycle-of-closed-subscheme}을 보라. 국소 자유층의
차수는 기약 성분들로의 제한에만 의존함이 밝혀진다.

\begin{lemma}
\label{varieties-lemma-degree-in-terms-of-components}
$k$를 체라 하자. $X$를 차원 $\leq 1$인 $k$ 위의 고유 스킴이라 하자.
$\mathcal{E}$를 국소 자유 $\mathcal{O}_X$-가군 중 계수 $n$인 것으로 하자.
그러면
$$
\deg(\mathcal{E}) = \sum m_i \deg(\mathcal{E}|_{C_i})
$$
이다. 여기서 $C_i \subset X$, $i = 1, \ldots, t$는 유도된 축약 스킴
구조를 가지는 차원 $1$의 기약 성분들이고, $m_i$는 $C_i$의 $X$
안에서의 중복도이다.
\end{lemma}

\begin{proof}
$C_i \to \Spec(k)$가 차원 $1$인 고유 사상이므로 명제가 의미가 있음에
유의하자. 아래에서는 $X$의 모든 닫힌 부분스킴이 $k$ 위에서 고유라는
사실을 사용한다. Morphisms, Lemmas
\ref{morphisms-lemma-closed-immersion-proper}와
\ref{morphisms-lemma-composition-proper}를 보라. 열린 부분스킴
$U_i = X \setminus (\bigcup_{j \not = i} C_j)$를 생각하고
$X_i \subset X$를 $U_i$의 스킴론적 폐포라 하자.
$X_i \cap U_i = U_i$가 스킴론적으로 성립하고
$X_i \cap U_j = \emptyset$가 $i \not = j$에 대해 집합론적으로
성립함에 유의하자. 이는 Morphisms, Lemma
\ref{morphisms-lemma-quasi-compact-immersion}의 스킴론적 폐포 설명에서
따른다. 따라서 사상 $X' = \coprod X_i \to X$에 Lemma
\ref{varieties-lemma-degree-birational-pullback}을 적용할 수 있다.
$C_i \subset X_i$가 스킴론적으로 분명하고 $C_i$의 $X_i$ 안에서의
중복도가 $C_i$의 $X$ 안에서의 중복도와 같으므로, 다음 문단에서 다루는
경우로 환원됨을 안다.

\medskip\noindent
$X$가 일반점 $\xi$를 가지는 기약 스킴이라고 가정하자.
$C = X_{red}$의 중복도가 $m$이라 하자. 우리는
$\deg(\mathcal{E}) = m \deg(\mathcal{E}|_C)$임을 보여야 한다.
$\mathcal{I} \subset \mathcal{O}_X$를 닫힌 부분스킴 $C$를 정의하는
아이디얼이라 하자. 최소의 $e \geq 0$을 택하여
$\mathcal{I}^{e + 1} = 0$이 되게 하자(Cohomology of Schemes, Lemma
\ref{coherent-lemma-power-ideal-kills-sheaf}). $e$에 관해 귀납한다.
$e = 0$이면 $X = C$이고 결과는 즉시 성립한다. 그렇지 않으면
$\mathcal{F} = \mathcal{I}^e$로 놓고 이를 연접
$\mathcal{O}_C$-가군으로 본다(Cohomology of Schemes, Lemma
\ref{coherent-lemma-i-star-equivalence}).
$X' \subset X$를 연접 아이디얼 $\mathcal{I}^e$로 잘라낸 닫힌
부분스킴이라 하고 $m'$를 $C$의 $X'$ 안에서의 중복도라 하자. 짧은 완전열
$$
0 \to \mathcal{F} \to \mathcal{O}_X \to \mathcal{O}_{X'} \to 0
$$
의 $\xi$에서의 줄기를 취하면(Algebra, Lemmas
\ref{algebra-lemma-length-additive},
\ref{algebra-lemma-dimension-is-length}, 및
\ref{algebra-lemma-length-independent}을 사용한다) 다음을 얻는다.
$$
m = \text{length}_{\mathcal{O}_{X, \xi}} \mathcal{O}_{X, \xi}
= \dim_{\kappa(\xi)} \mathcal{F}_\xi +
\text{length}_{\mathcal{O}_{X', \xi}} \mathcal{O}_{X', \xi}
= r + m'
$$
여기서 $r$은 $\mathcal{F}$의 $C$ 위의 연접층으로서의 계수이다.
$\mathcal{E}$와 텐서하여 짧은 완전열
$$
0 \to \mathcal{E}|_C \otimes \mathcal{F} \to \mathcal{E} \to
\mathcal{E} \otimes \mathcal{O}_{X'} \to 0
$$
을 얻는다. 귀납가정으로
$\deg(\mathcal{E}|_{X'}) = m' \deg(\mathcal{E}|_C)$이다.
Lemma \ref{varieties-lemma-degree-on-proper-curve}에 의해
$\chi(\mathcal{E}|_C \otimes \mathcal{F}) =
r \deg(\mathcal{E}|_C) + n \chi(\mathcal{F})$이다.
모두 합치면 결과를 얻는다.
\end{proof}

\begin{lemma}
\label{varieties-lemma-degree-tensor-product}
$k$를 체라 하고, $X$를 차원 $\leq 1$인 $k$ 위의 고유 스킴이라 하자.
$\mathcal{E}$와 $\mathcal{V}$를 유한 상수 계수인 국소 자유
$\mathcal{O}_X$-가군들이라 하자. 그러면
$$
\deg(\mathcal{E} \otimes \mathcal{V}) =
\text{rank}(\mathcal{E}) \deg(\mathcal{V}) +
\text{rank}(\mathcal{V}) \deg(\mathcal{E})
$$
\end{lemma}

\begin{proof}
Lemma \ref{varieties-lemma-degree-in-terms-of-components}와 초등 산술에 의해 고유
곡선의 경우로 환원된다. 이 경우는 Lemma
\ref{varieties-lemma-degree-on-proper-curve}에서 따른다.
\end{proof}

\begin{lemma}
\label{varieties-lemma-degree-and-det}
$k$를 체라 하고, $X$를 차원 $\leq 1$인 $k$ 위의 고유 스킴이라 하자.
$\mathcal{E}$를 국소 자유 $\mathcal{O}_X$-가군 중 계수 $n$인 것으로 하자.
그러면
$$
\deg(\mathcal{E}) = \deg(\wedge^n(\mathcal{E})) = \deg(\det(\mathcal{E}))
$$
\end{lemma}

\begin{proof}
Lemma \ref{varieties-lemma-degree-in-terms-of-components}와 초등 산술에 의해 고유
곡선의 경우로 환원된다. 그러면 변형 $f : X' \to X$가 존재하여
$f^*\mathcal{E}$가 가역 가군들을 연속 몫으로 가지는 여과를 갖는다.
Divisors, Lemma \ref{divisors-lemma-filter-after-modification}를 보라.
Lemma \ref{varieties-lemma-degree-birational-pullback}에 의해 $X'$ 위에서 작업해도
된다. 따라서 여과
$$
0 = \mathcal{E}_0 \subset \mathcal{E}_1 \subset \mathcal{E}_2 \subset
\ldots \subset \mathcal{E}_n = \mathcal{E}
$$
가 있고 이것이 국소 자유 $\mathcal{O}_X$-가군들로 이루어지며
$\mathcal{L}_i = \mathcal{E}_i/\mathcal{E}_{i - 1}$가 가역이라고
가정할 수 있다. Modules, Lemma \ref{modules-lemma-det-ses}와 귀납법으로
$\det(\mathcal{E}) = \mathcal{L}_1 \otimes \ldots \otimes \mathcal{L}_n$을
얻는다. 따라서 등식은 Lemma \ref{varieties-lemma-degree-tensor-product}와
가법성(Lemma \ref{varieties-lemma-degree-additive})에서 따른다.
\end{proof}

\begin{lemma}
\label{varieties-lemma-degree-effective-Cartier-divisor}
$k$를 체라 하고, $X$를 차원 $\leq 1$인 $k$ 위의 고유 스킴이라 하자.
$D$를 $X$ 위의 유효 카르티에 인자라 하자. 그러면 $D$는
$\Spec(k)$ 위 차수 $\deg(D) = \dim_k \Gamma(D, \mathcal{O}_D)$인 유한
스킴이다. 국소 자유층 $\mathcal{E}$ 가운데 계수 $n$인 것에 대해
$$
\deg(\mathcal{E}(D)) = n\deg(D) + \deg(\mathcal{E})
$$
이고, 여기서
$\mathcal{E}(D) = \mathcal{E} \otimes_{\mathcal{O}_X} \mathcal{O}_X(D)$이다.
\end{lemma}

\begin{proof}
$D$는 $X$에서 조밀한 곳이 전혀 없으므로(Divisors, Lemma
\ref{divisors-lemma-complement-effective-Cartier-divisor})
$\dim(D) \leq 0$이다. 따라서 Lemma
\ref{varieties-lemma-algebraic-scheme-dim-0}에 의해 $D$는 $k$ 위에서 유한이다.
$k$가 체이므로 사상 $D \to \Spec(k)$는 유한 국소 자유이고 따라서
차수를 가진다(Morphisms, Definition
\ref{morphisms-definition-finite-locally-free}). 이 차수는 보조정리에
진술한 대로 분명히 $\dim_k \Gamma(D, \mathcal{O}_D)$와 같다. Divisors,
Definition
\ref{divisors-definition-invertible-sheaf-effective-Cartier-divisor}에
의해 짧은 완전열
$$
0 \to \mathcal{O}_X \to \mathcal{O}_X(D) \to i_*i^*\mathcal{O}_X(D) \to 0
$$
이 존재한다. 여기서 $i : D \to X$는 닫힌 몰입이다. $\mathcal{E}$와
텐서하면 짧은 완전열
$$
0 \to \mathcal{E} \to \mathcal{E}(D) \to i_*i^*\mathcal{E}(D) \to 0
$$
을 얻는다. 보조정리의 식은 오일러 표수의 가법성(Lemma
\ref{varieties-lemma-euler-characteristic-additive})과 Lemma
\ref{varieties-lemma-chi-tensor-finite}에서 따른다.
\end{proof}

\begin{lemma}
\label{varieties-lemma-divisible}
$k$를 체라 하자. $X$를 $k$ 위의 고유 스킴 가운데
축약이고 연결인 것이라 하자.
$\kappa = H^0(X, \mathcal{O}_X)$라 두자. 그러면 $\kappa/k$는 체의
유한 확대이고, $w = [\kappa : k]$는 다음을 나눈다.
\begin{enumerate}
\item $\deg(\mathcal{E})$. 여기서 임의의 국소 자유
$\mathcal{O}_X$-가군 $\mathcal{E}$를 취한다,
\item $[\kappa(x) : k]$. 여기서 $x \in X$는 임의의 닫힌 점이다, 그리고
\item $\deg(D)$. 여기서 $D \subset X$는 차원이 영인 임의의 닫힌
부분스킴이다.
\end{enumerate}
\end{lemma}

\begin{proof}
$\kappa$에 관한 주장들은 Lemma
\ref{varieties-lemma-proper-geometrically-reduced-global-sections}를 보라.
모든 준연접 $\mathcal{O}_X$-가군에 대해 $k$-벡터 공간
$H^i(X, \mathcal{F})$들은 $\kappa$-벡터 공간이다.
나눗셈 성질들은 이 명제와 정의들에서 쉽게 따른다.
\end{proof}

\begin{lemma}
\label{varieties-lemma-degree-pullback-map-proper-curves}
$k$를 체라 하자. $f : X \to Y$를 $k$ 위의 고유 곡선들 사이의
상수가 아닌 사상이라 하자. $\mathcal{E}$를 국소 자유
$\mathcal{O}_Y$-가군이라 하자. 그러면
$$
\deg(f^*\mathcal{E}) = \deg(X/Y) \deg(\mathcal{E})
$$
\end{lemma}

\begin{proof}
$X$의 $Y$ 위에서의 차수는 Morphisms, Definition
\ref{morphisms-definition-degree}에서 정의된다.
따라서 $f_*\mathcal{O}_X$는 연접 $\mathcal{O}_Y$-가군이며 그 계수는
$\deg(X/Y)$이다. 즉,
$\deg(X/Y) = \dim_{\kappa(\xi)} (f_*\mathcal{O}_X)_\xi$이고, 여기서 $\xi$는
$Y$의 일반점이다.
따라서 다음을 얻는다.
\begin{align*}
\chi(X, f^*\mathcal{E})
& =
\chi(Y, f_*f^*\mathcal{E}) \\
& =
\chi(Y, \mathcal{E} \otimes f_*\mathcal{O}_X) \\
& =
\deg(X/Y) \deg(\mathcal{E}) + n \chi(Y, f_*\mathcal{O}_X) \\
& =
\deg(X/Y) \deg(\mathcal{E}) + n \chi(X, \mathcal{O}_X)
\end{align*}
이는 원하는 바이다. 첫 번째 등식은 $f$가 유한이기 때문이며,
Cohomology of Schemes, Lemma
\ref{coherent-lemma-relative-affine-cohomology}를 보라.
두 번째 등식은 사영 공식에서 따르며, Cohomology, Lemma
\ref{cohomology-lemma-projection-formula}를 보라.
세 번째 등식은 Lemma \ref{varieties-lemma-degree-on-proper-curve}에서 따른다.
\end{proof}

\noindent
다음은 위의 차수에 관한 결과들의 자명하지만 중요한 귀결이다.

\begin{lemma}
\label{varieties-lemma-check-invertible-sheaf-trivial}
$k$를 체라 하자. $X$를 $k$ 위의 고유 곡선이라 하자.
$\mathcal{L}$을 가역 $\mathcal{O}_X$-가군이라 하자.
\begin{enumerate}
\item $\mathcal{L}$이 영이 아닌 단면을 가지면
$\deg(\mathcal{L}) \geq 0$이다.
\item $\mathcal{L}$이 어떤 점에서 소멸하는 영이 아닌 단면 $s$를 가지면
$\deg(\mathcal{L}) > 0$이다.
\item $\mathcal{L}$과 $\mathcal{L}^{-1}$이 영이 아닌 단면들을 가지면
$\mathcal{L} \cong \mathcal{O}_X$이다.
\item $\deg(\mathcal{L}) \leq 0$이고 $\mathcal{L}$이 영이 아닌
단면을 가지면 $\mathcal{L} \cong \mathcal{O}_X$이다.
\item $\mathcal{N} \to \mathcal{L}$이 가역
$\mathcal{O}_X$-가군들의 영이 아닌 사상이면
$\deg(\mathcal{L}) \geq \deg(\mathcal{N})$이고, 등식이 성립하면
이는 동형사상이다.
\end{enumerate}
\end{lemma}

\begin{proof}
$s$를 $\mathcal{L}$의 영이 아닌 단면이라 하자. $X$가 곡선이므로
$s$는 정칙 단면이다. 따라서 유효 카르티에 인자 $D \subset X$와 동형사상
$\mathcal{L} \to \mathcal{O}_X(D)$가 존재하며, 이 동형사상은 $s$를 표준
단면 $1$, 곧 $\mathcal{O}_X(D)$의 표준 단면으로 보낸다. Divisors, Lemma
\ref{divisors-lemma-characterize-OD}를 보라.
그러면 Lemma \ref{varieties-lemma-degree-effective-Cartier-divisor}에 의해
$\deg(\mathcal{L}) = \deg(D)$이다. $\deg(D) \geq 0$이고
$= 0$인 것은 $D = \emptyset$인 것과 동치이므로, 이는 (1)과 (2)를
증명한다. 경우 (3)에서는 $\deg(\mathcal{L}) = 0$이고
$D = \emptyset$임을 알 수 있다. (4)도 마찬가지이다. (5)를 보이려면
가역층
$$
\mathcal{L} \otimes_{\mathcal{O}_X} \mathcal{N}^{\otimes -1} =
\SheafHom_{\mathcal{O}_X}(\mathcal{N}, \mathcal{L})
$$
에 (1)과 (4)를 적용한다. 이 층의 차수는 Lemma
\ref{varieties-lemma-degree-tensor-product}에 의해
$\deg(\mathcal{L}) - \deg(\mathcal{N})$이다.
\end{proof}

\begin{lemma}
\label{varieties-lemma-no-sections-dual-nef}
$k$를 체라 하자. $X$를 $k$ 위의 고유 스킴 가운데 축약이고 연결이며
동차원이고 차원이 $1$인 것이라 하자. $\mathcal{L}$을 가역
$\mathcal{O}_X$-가군이라 하자. $\deg(\mathcal{L}|_C) \leq 0$이라고 하자.
여기서 $C$는 $X$의 임의의 기약 성분이다. 그러면 $H^0(X, \mathcal{L}) = 0$이거나
$\mathcal{L} \cong \mathcal{O}_X$이다.
\end{lemma}

\begin{proof}
$s \in H^0(X, \mathcal{L})$를 영이 아닌 원소라 하자. $X$가 축약이므로
기약 성분 $C$가 $X$ 안에 존재하여 $s|_C \not = 0$이다.
그런데 $s|_C$가 영이 아니면 Lemma
\ref{varieties-lemma-check-invertible-sheaf-trivial}에 의해 $s$는 $C$ 위의 어느
곳에서도 소멸하지 않는다. 그러면 다시 $s$는 $X$의 모든 기약 성분
가운데 $C$와 만나는 것에서 어느 곳에서도 소멸하지 않는다.
$X$가 연결이므로 $s$는 어느 곳에서도 소멸하지 않으며, 보조정리가 따른다.
\end{proof}

\begin{lemma}
\label{varieties-lemma-ample-curve}
$k$를 체라 하자. $X$를 $k$ 위의 고유 곡선이라 하자.
$\mathcal{L}$을 가역 $\mathcal{O}_X$-가군이라 하자.
그러면 $\mathcal{L}$이 풍부한 것은 $\deg(\mathcal{L}) > 0$인 것과 동치이다.
\end{lemma}

\begin{proof}
$\mathcal{L}$이 풍부하면, 어떤 $n > 0$과 단면
$s \in H^0(X, \mathcal{L}^{\otimes n})$이 존재하여 $X_s$는 아핀이다.
$X$는 아핀이 아니므로(그렇지 않으면 Morphisms, Lemma
\ref{morphisms-lemma-finite-proper}에 의해 $X$는 유한일 것이다),
$s$는 어떤 점에서 소멸한다. 따라서 Lemma
\ref{varieties-lemma-check-invertible-sheaf-trivial}에 의해
$\deg(\mathcal{L}^{\otimes n}) > 0$이다.
Lemma \ref{varieties-lemma-degree-tensor-product}에 의해
$\deg(\mathcal{L}) = 1/n\deg(\mathcal{L}^{\otimes n}) > 0$이라고 결론짓는다.

\medskip\noindent
$\deg(\mathcal{L}) > 0$이라고 가정하자. 그러면
$$
\dim_k H^0(X, \mathcal{L}^{\otimes n}) \geq \chi(X, \mathcal{L}^n)
= n\deg(\mathcal{L}) + \chi(X, \mathcal{O}_X)
$$
는 $n$에 대해 선형으로 증가한다. 따라서 닫힌 점들의 임의의 유한 모음
$x_1, \ldots, x_t$을 $X$에서 취하면, 어떤 $n$을 찾아
$\dim_k H^0(X, \mathcal{L}^{\otimes n}) > \sum \dim_k \kappa(x_i)$가 되게
할 수 있다. (힐베르트 영점 정리에 의해 확대체
$\kappa(x_i)/k$가 유한임을 상기하라. 예를 들어 Morphisms, Lemma
\ref{morphisms-lemma-closed-point-fibre-locally-finite-type}를 보라.)
따라서 영이 아닌 $s \in H^0(X, \mathcal{L}^{\otimes n})$을 찾아
$x_1, \ldots, x_t$에서 소멸하게 할 수 있다. 특히
$x_1, \ldots, x_t$를 택하여 $X \setminus \{x_1, \ldots, x_t\}$가
아핀이 되게 하면 $X_s$도 아핀이다
(예를 들어 Properties, Lemma \ref{properties-lemma-affine-cap-s-open}에
의한다. 다만 유한 집합을
$\mathcal{L}|_{X \setminus \{x_1, \ldots, x_t\}}$가 자명해지도록 택하면
이는 즉시 알 수 있다). 결론적으로 어떤 $n > 0$과 영이 아닌 단면
$s \in H^0(X, \mathcal{L}^{\otimes n})$을 찾아 $X_s$가 아핀이 되게 할
수 있다.

\medskip\noindent
아이디얼의 모든 준연접층 $\mathcal{I}$에 대해 어떤 $m > 0$이 존재하여
$H^1(X, \mathcal{I} \otimes \mathcal{L}^{\otimes m})$이 영임을 보이겠다.
그러면 Cohomology of Schemes, Lemma
\ref{coherent-lemma-vanshing-gives-ample}에 의해 증명이 끝난다.
이를 보이기 위해 사상들
$$
\mathcal{I} \xrightarrow{s}
\mathcal{I} \otimes \mathcal{L}^{\otimes n} \xrightarrow{s}
\mathcal{I} \otimes \mathcal{L}^{\otimes 2n} \xrightarrow{s} \ldots
$$
을 생각한다. $\mathcal{I}$는 비틀림이 없으므로 이 사상들은 단사이고
$X_s$ 위에서 동형사상이다. 따라서 여핵들의 $H^1$은 영이다(예를 들어
Cohomology of Schemes, Lemma
\ref{coherent-lemma-coherent-support-dimension-0}에 의한다).
따라서 벡터 공간들의 사상
$$
H^1(X, \mathcal{I}) \to
H^1(X, \mathcal{I} \otimes \mathcal{L}^{\otimes n}) \to
H^1(X, \mathcal{I} \otimes \mathcal{L}^{\otimes 2n}) \to \ldots
$$
들은 전사이다. 한편 $H^1(X, \mathcal{I})$의 차원은 유한하고, 모든 원소는
Cohomology of Schemes, Lemma
\ref{coherent-lemma-section-affine-open-kills-classes}에 의해 궁극적으로
영으로 간다. 따라서 어떤 $e > 0$에 대해
$H^1(X, \mathcal{I} \otimes \mathcal{L}^{\otimes en})$은 영이다.
이로써 증명이 끝난다.
\end{proof}

\begin{lemma}
\label{varieties-lemma-ampleness-in-terms-of-degrees-components}
$k$를 체라 하자. $X$를 차원 $\leq 1$인 $k$ 위의 고유 스킴이라 하자.
$\mathcal{L}$을 가역 $\mathcal{O}_X$-가군이라 하자.
$C_i \subset X$, $i = 1, \ldots, t$를 차원 $1$인 기약 성분들이라 하자.
다음은 동치이다.
\begin{enumerate}
\item $\mathcal{L}$은 풍부하다, 그리고
\item $\deg(\mathcal{L}|_{C_i}) > 0$이다. 여기서 $i = 1, \ldots, t$이다.
\end{enumerate}
\end{lemma}

\begin{proof}
$x_1, \ldots, x_r \in X$를 고립된 닫힌 점들이라 하자.
$x_i = \Spec(\kappa(x_i))$를 스킴으로 생각하자. 스킴의 사상
$$
f : C_1 \amalg \ldots \amalg C_t \amalg x_1 \amalg \ldots \amalg x_r
\longrightarrow X
$$
을 생각하자. 이는 $k$ 위에서 고유인 스킴들의 유한 전사 사상이다
(세부사항은 생략한다). 따라서 $\mathcal{L}$이 풍부한 것은
$f^*\mathcal{L}$이 풍부한 것과 동치이다(Cohomology of Schemes, Lemma
\ref{coherent-lemma-surjective-finite-morphism-ample}).
그러므로 Lemma \ref{varieties-lemma-ample-curve}에 의해 결론을 얻는다.
\end{proof}

\begin{lemma}
\label{varieties-lemma-general-degree-g-line-bundle}
$k$를 대수적으로 닫힌 체라 하자. $X$를 $k$ 위의 고유 곡선이라 하자.
그러면 다음이 존재한다.
\begin{enumerate}
\item 가역 $\mathcal{O}_X$-가군 $\mathcal{L}$로서
$\dim_k H^0(X, \mathcal{L}) = 1$이고 $H^1(X, \mathcal{L}) = 0$인 것, 그리고
\item 가역 $\mathcal{O}_X$-가군 $\mathcal{N}$으로서
$\dim_k H^0(X, \mathcal{N}) = 0$이고 $H^1(X, \mathcal{N}) = 0$인 것.
\end{enumerate}
\end{lemma}

\begin{proof}
닫힌 몰입 $i : X \to \mathbf{P}^n_k$를 택한다
(Lemma \ref{varieties-lemma-dim-1-proper-projective}).
$\mathcal{L} = i^*\mathcal{O}_{\mathbf{P}^n}(d)$로 놓고
$d \gg 0$으로 택하면, 가역층 $\mathcal{L}$이 존재하여
$H^0(X, \mathcal{L}) \not = 0$이고 $H^1(X, \mathcal{L}) = 0$임을 알 수 있다
(소멸에 관해서는 Cohomology of Schemes, Lemma
\ref{coherent-lemma-vanshing-gives-ample}를, 비소멸에 관해서는 그곳의
참고문헌들을 보라). $t = \dim_k H^0(X, \mathcal{L})$에 관한 하강
귀납법으로 (1)의 증명을 끝내겠다. 기초 경우 $t = 1$은 자명하다.
$t > 1$이라고 가정하자.

\medskip\noindent
Lemma \ref{varieties-lemma-dense-smooth-open-variety-over-perfect-field}에서 연구한
비특이점들의 비어 있지 않은 열린 부분집합 $U \subset X$를 취하자.
$s \in H^0(X, \mathcal{L})$를 영이 아닌 원소라 하자. 닫힌 점
$x \in U$가 존재하여 $s$가 $x$에서 소멸하지 않는다.
$\mathcal{I}$를 Lemma \ref{varieties-lemma-regular-point-on-curve}에서와 같은
$i : x \to X$의 아이디얼층이라 하자. 짧은 완전열
$$
0 \to \mathcal{I} \otimes_{\mathcal{O}_X} \mathcal{L} \to
\mathcal{L} \to i_*i^*\mathcal{L} \to 0
$$
을 보자. $H^0(X, i_*i^*\mathcal{L}) = H^0(x, i^*\mathcal{L})$은
차원 $1$이다. 실제로 $x$는 $k$-유리점이다($k$가 대수적으로 닫혀 있다).
$s$가 $x$에서 소멸하지 않으므로
$$
H^0(X, \mathcal{L}) \longrightarrow H^0(X, i_*i^*\mathcal{L})
$$
은 전사이다. 따라서
$\dim_k H^0(X, \mathcal{I} \otimes_{\mathcal{O}_X} \mathcal{L}) = t - 1$이다.
마지막으로 코호몰로지의 긴 완전열은 또한
$H^1(X, \mathcal{I} \otimes_{\mathcal{O}_X} \mathcal{L}) = 0$임을
보여 주며, 이로써 귀납 단계의 증명이 끝난다.

\medskip\noindent
(2)와 같은 가역층을 얻으려면 (1)과 같은 가역층 $\mathcal{L}$을 취하고
앞 문단의 논증을 한 번 더 수행한다.
\end{proof}

\begin{lemma}
\label{varieties-lemma-vanishing-degree-2g-and-1-line-bundle}
$k$를 대수적으로 닫힌 체라 하자. $X$를 $k$ 위의 고유 곡선이라 하자.
$g = \dim_k H^1(X, \mathcal{O}_X)$라 두자. 모든 가역
$\mathcal{O}_X$-가군 $\mathcal{L}$ 가운데 $\deg(\mathcal{L}) \geq 2g - 1$인 것에 대해
$H^1(X, \mathcal{L}) = 0$이다.
\end{lemma}

\begin{proof}
Lemma \ref{varieties-lemma-general-degree-g-line-bundle}의 부분 (2)에서 찾은 가역
가군을 $\mathcal{N}$이라 하자. $\mathcal{N}$의 차수는
$\chi(X, \mathcal{N}) - \chi(X, \mathcal{O}_X) = 0 - (1 - g) = g - 1$이다.
따라서 $\mathcal{L} \otimes \mathcal{N}^{\otimes - 1}$의 차수는
$\deg(\mathcal{L}) - (g - 1) \geq g$이다. 따라서
$\chi(X, \mathcal{L} \otimes \mathcal{N}^{\otimes -1}) \geq g + 1 - g = 1$이다.
그러므로 영이 아닌 전역 단면 $s$가 존재하고, 그 영점 스킴은 유효
카르티에 인자 $D$이며 그 차수는 $\deg(\mathcal{L}) - (g - 1)$이다.
이는 짧은 완전열
$$
0 \to \mathcal{N} \xrightarrow{s} \mathcal{L} \to i_*(\mathcal{L}|_D) \to 0
$$
을 준다. 여기서 $i : D \to X$는 포함 사상이다. 따라서
$H^0(X, \mathcal{L})$은 $H^0(D, \mathcal{L}|_D)$에 동형으로 사상되며,
후자의 차원은 $\deg(\mathcal{L}) - (g - 1)$이다. 결과는 차수의 정의에서
따른다.
\end{proof}

\section{수치적 교차}
\label{varieties-section-num}

\noindent
이 절에서는 고유 스킴 위의 연접층들의 오일러 표수를 다루어 가역
가군들의 수치적 교차수를 얻는다. 주된 도구는 다음 보조정리이다.

\begin{lemma}
\label{varieties-lemma-numerical-polynomial-from-euler}
$k$를 체라 하자. $X$를 $k$ 위의 고유 스킴이라 하자. $\mathcal{F}$를
연접 $\mathcal{O}_X$-가군이라 하자.
$\mathcal{L}_1, \ldots, \mathcal{L}_r$을 가역 $\mathcal{O}_X$-가군들이라 하자.
사상
$$
(n_1, \ldots, n_r) \longmapsto
\chi(X, \mathcal{F} \otimes
\mathcal{L}_1^{\otimes n_1} \otimes \ldots \otimes
\mathcal{L}_r^{\otimes n_r})
$$
은 $n_1, \ldots, n_r$에 관한 수치 다항식이고, 그 전차수는
$\mathcal{F}$의 지지집합의 차원 이하이다.
\end{lemma}

\begin{proof}
$\dim(\text{Supp}(\mathcal{F}))$에 관한 귀납법으로 증명한다.
이 수가 영이면 Lemma \ref{varieties-lemma-chi-tensor-finite}에 의해 함수는 상수이고
그 값은 $\dim_k \Gamma(X, \mathcal{F})$이다.
$\dim(\text{Supp}(\mathcal{F})) > 0$이라고 가정하자.

\medskip\noindent
$\mathcal{F}$가 매장 연관점들을 가지면 Divisors, Lemma
\ref{divisors-lemma-remove-embedded-points}에서 구성한 짧은 완전열
$0 \to \mathcal{K} \to \mathcal{F} \to \mathcal{F}' \to 0$을 생각할 수 있다.
$\mathcal{K}$의 지지집합 차원은 엄밀히 더 작으므로, 귀납 가정에 의해
$\mathcal{K}$에 대해 결과가 성립하며 전차수도 엄밀히 더 작다.
오일러 표수의 가법성(Lemma \ref{varieties-lemma-euler-characteristic-additive})에
의해 $\mathcal{F}'$에 대해 결과를 증명하면 충분하다. 따라서
$\mathcal{F}$가 매장 연관점들을 가지지 않는다고 가정해도 된다.

\medskip\noindent
$i : Z \to X$가 닫힌 몰입이고 $\mathcal{F} = i_*\mathcal{G}$이면,
$X$, $\mathcal{F}$, $\mathcal{L}_1, \ldots, \mathcal{L}_r$에 대한 결과는
$Z$, $\mathcal{G}$, $i^*\mathcal{L}_1, \ldots, i^*\mathcal{L}_r$에 대한
결과와 동치이다(코호몰로지가 일치하기 때문이다. Cohomology of Schemes,
Lemma \ref{coherent-lemma-relative-affine-cohomology}를 보라).
Divisors, Lemma \ref{divisors-lemma-no-embedded-points-endos}를 적용하여
$X$가 매장 성분들을 가지지 않고 $X = \text{Supp}(\mathcal{F})$라고
가정해도 된다.

\medskip\noindent
정칙 유리형 단면 $s$를 $\mathcal{L}_1$에서 택하자. Divisors, Lemma
\ref{divisors-lemma-regular-meromorphic-section-exists}를 보라.
$\mathcal{I} \subset \mathcal{O}_X$를 $s$의 분모 아이디얼이라 하고,
Divisors, Lemma \ref{divisors-lemma-make-maps-regular-section}의 사상들
$$
\mathcal{I}\mathcal{F} \to \mathcal{F},\quad
\mathcal{I}\mathcal{F} \to \mathcal{F} \otimes \mathcal{L}_1
$$
을 생각하자. 이들은 단사이고, 여핵 $\mathcal{Q}$, $\mathcal{Q}'$은
$X = \text{Supp}(\mathcal{F})$의 어디에서도 조밀하지 않은 닫힌
부분스킴들 위에 지지된다. 가역 가군
$\mathcal{L}_1^{\otimes n_1} \otimes \ldots \otimes \mathcal{L}_r^{\otimes n_r}$과
텐서하는 것은 완전하므로, 가법성을 다시 이용하면 다음을 얻는다.
\begin{align*}
&\chi(X, \mathcal{F} \otimes \mathcal{L}_1^{\otimes n_1} \otimes \ldots \otimes
\mathcal{L}_r^{\otimes n_r}) -
\chi(X, \mathcal{F} \otimes \mathcal{L}_1^{\otimes n_1 + 1}
\otimes \ldots \otimes \mathcal{L}_r^{\otimes n_r}) \\
& =
\chi(\mathcal{Q} \otimes \mathcal{L}_1^{\otimes n_1} \otimes \ldots \otimes
\mathcal{L}_r^{\otimes n_r}) -
\chi(\mathcal{Q}' \otimes \mathcal{L}_1^{\otimes n_1} \otimes \ldots \otimes
\mathcal{L}_r^{\otimes n_r})
\end{align*}
따라서 보조정리의 함수 $P(n_1, \ldots, n_r)$는
$$
P(n_1 + 1, n_2, \ldots, n_r) - P(n_1, \ldots, n_r)
$$
가 전차수 $<$ $\mathcal{F}$의 지지집합 차원인 수치 다항식이라는
성질을 갖는다. 물론 대칭성에 의해, 임의의 $i \in \{1, \ldots, r\}$에
대해
$$
P(n_1, \ldots, n_{i - 1}, n_i + 1, n_{i + 1}, \ldots, n_r)
- P(n_1, \ldots, n_r)
$$
에도 같은 것이 성립한다. 간단한 산술 논증으로 $P$는 전차수가
$\dim(\text{Supp}(\mathcal{F}))$ 이하인 수치 다항식임을 알 수 있다.
\end{proof}

\noindent
다음 보조정리는 최고차항의 계수가 대략적으로 지지집합의 일반점들에서
연접 가군이 갖는 길이에만 의존함을 보인다.

\begin{lemma}
\label{varieties-lemma-numerical-polynomial-leading-term}
$k$를 체라 하자. $X$를 $k$ 위의 고유 스킴이라 하자.
$\mathcal{F}$를 연접 $\mathcal{O}_X$-가군이라 하자.
$\mathcal{L}_1, \ldots, \mathcal{L}_r$을 가역 $\mathcal{O}_X$-가군들이라 하자.
$d = \dim(\text{Supp}(\mathcal{F}))$라 두자.
$Z_i \subset X$를 $\text{Supp}(\mathcal{F})$의 차원 $d$인 기약 성분들이라
하자. $\xi_i \in Z_i$를 일반점이라 하고
$m_i = \text{length}_{\mathcal{O}_{X, \xi_i}}(\mathcal{F}_{\xi_i})$로 놓자.
그러면
$$
\chi(X, \mathcal{F} \otimes \mathcal{L}_1^{\otimes n_1} \otimes \ldots \otimes
\mathcal{L}_r^{\otimes n_r}) -
\sum\nolimits_i
m_i\ \chi(Z_i, \mathcal{L}_1^{\otimes n_1} \otimes \ldots \otimes
\mathcal{L}_r^{\otimes n_r}|_{Z_i})
$$
은 $n_1, \ldots, n_r$에 관한 전차수 $< d$인 수치 다항식이다.
\end{lemma}

\begin{proof}
쌍 $(\xi , Z)$를 생각하자. 여기서 $Z \subset X$는 차원 $d$인 정수적
닫힌 부분스킴이고 $\xi$는 그 일반점이다. 그러면 유한 $\mathcal{O}_{X, \xi}$-가군
$\mathcal{F}_\xi$의 지지집합은 $\{\xi\}$에 포함되므로 길이
$m_Z = \text{length}_{\mathcal{O}_{X, \xi}}(\mathcal{F}_\xi)$는 유한이고
(Algebra, Lemma \ref{algebra-lemma-support-point}), $Z = Z_i$가 되는 어떤
$i$가 존재하지 않으면 영이다. 따라서 보조정리의 식은 다음과 같이 쓸 수 있다.
$$
E(\mathcal{F}) =
\chi(X, \mathcal{F} \otimes \mathcal{L}_1^{\otimes n_1} \otimes \ldots \otimes
\mathcal{L}_r^{\otimes n_r}) -
\sum\nolimits
m_Z\ \chi(Z, \mathcal{L}_1^{\otimes n_1} \otimes \ldots \otimes
\mathcal{L}_r^{\otimes n_r}|_Z)
$$
여기서 합은 정수적 닫힌 부분스킴 $Z \subset X$ 가운데 차원 $d$인 것
전체에 걸친다. 대응 $\mathcal{F} \mapsto E(\mathcal{F})$는 짧은 완전열
$0 \to \mathcal{F} \to \mathcal{F}' \to \mathcal{F}'' \to 0$에 대해
가법적이다. 여기서 이 가군들은 연접 $\mathcal{O}_X$-가군들이고 그
지지집합들의 차원은 $\leq d$이다. 이는 오일러 표수의 가법성(Lemma
\ref{varieties-lemma-euler-characteristic-additive})과 길이의 가법성(Algebra, Lemma
\ref{algebra-lemma-length-additive})에서 따른다.
Cohomology of Schemes, Lemma \ref{coherent-lemma-coherent-filter}를 적용하여
연접 부분층들에 의한 여과
$$
0 = \mathcal{F}_0 \subset \mathcal{F}_1 \subset
\ldots \subset \mathcal{F}_m = \mathcal{F}
$$
를 찾자. 여기서 각 $j = 1, \ldots, m$에 대해 정수적 닫힌 부분스킴
$V_j \subset X$와 영이 아닌 아이디얼층
$\mathcal{I}_j \subset \mathcal{O}_{V_j}$가 존재하여
$$
\mathcal{F}_j/\mathcal{F}_{j - 1}
\cong (V_j \to X)_* \mathcal{I}_j
$$
이다. 이로부터 $V_j \subset \text{Supp}(\mathcal{F})$이고 따라서
$\dim(V_j) \leq d$임이 따른다. 위에서 말한 가법성에 의해 각 부분몫
$\mathcal{F}_j/\mathcal{F}_{j - 1}$에 대해 결과를 증명하면 충분하다.
따라서 $\mathcal{F} = (V \to X)_*\mathcal{I}$이고
$V \subset X$가 차원 $\leq d$인 정수적 닫힌 부분스킴이며
$\mathcal{I} \subset \mathcal{O}_V$가 영이 아닌 연접 아이디얼층인 경우만
증명하면 충분하다. $\dim(V) < d$인 경우, 더 일반적으로 $\mathcal{F}$의
지지집합 차원이 $< d$인 경우에는 $E(\mathcal{F})$의 첫 번째 항이 Lemma
\ref{varieties-lemma-numerical-polynomial-from-euler}에 의해 전차수 $< d$이고 두 번째
항은 영이다. $\dim(V) = d$이면 짧은 완전열
$$
0 \to (V \to X)_*\mathcal{I} \to (V \to X)_*\mathcal{O}_V
\to (V \to X)_*(\mathcal{O}_V/\mathcal{I}) \to 0
$$
을 사용할 수 있다. 가운데 층에 대해서는 결과가 성립한다. 실제로 합에
나타나는 유일한 $Z$는 $Z = V$이고 $m_Z = 1$이며,
$$
H^i(X, ((V \to X)_*\mathcal{O}_V) \otimes 
 \mathcal{L}_1^{\otimes n_1} \otimes \ldots \otimes
\mathcal{L}_r^{\otimes n_r}) =
H^i(V,  \mathcal{L}_1^{\otimes n_1} \otimes \ldots \otimes
\mathcal{L}_r^{\otimes n_r}|_V)
$$
이기 때문이다. 이는 사영 공식(Cohomology, Section
\ref{cohomology-section-projection-formula})과 Cohomology of Schemes, Lemma
\ref{coherent-lemma-relative-affine-cohomology}에서 따른다. 따라서 이 경우에는
실제로 $E(\mathcal{F}) = 0$이다. 오른쪽 층에 대해서도 결과가 성립하는데,
그 지지집합의 차원이 $< d$이기 때문이다. 따라서 왼쪽 층에 대해 결과가
성립하고 보조정리가 증명된다.
\end{proof}

\begin{definition}
\label{varieties-definition-intersection-number}
$k$를 체라 하자. $X$를 $k$ 위의 고유 스킴이라 하자.
$i : Z \to X$를 차원 $d$인 닫힌 부분스킴이라 하자.
$\mathcal{L}_1, \ldots, \mathcal{L}_d$를 가역
$\mathcal{O}_X$-가군들이라 하자. {\it 교차수}
$(\mathcal{L}_1 \cdots \mathcal{L}_d \cdot Z)$를 수치 다항식
$$
\chi(X, i_*\mathcal{O}_Z \otimes \mathcal{L}_1^{\otimes n_1} \otimes
\ldots \otimes \mathcal{L}_d^{\otimes n_d}) =
\chi(Z, \mathcal{L}_1^{\otimes n_1} \otimes
\ldots \otimes \mathcal{L}_d^{\otimes n_d}|_Z)
$$
에서 $n_1 \ldots n_d$의 계수로 정의한다. 특별한 경우
$\mathcal{L}_1 = \ldots = \mathcal{L}_d = \mathcal{L}$에는
$(\mathcal{L}^d \cdot Z)$라고 쓴다.
\end{definition}

\noindent
정의에 표시된 등식은 사영 공식(Cohomology, Section
\ref{cohomology-section-projection-formula})과 Cohomology of Schemes, Lemma
\ref{coherent-lemma-relative-affine-cohomology}에서 따른다.
이 교차수들에 관한 몇 가지 보조정리를 증명한다.

\begin{lemma}
\label{varieties-lemma-intersection-number-integer}
Definition \ref{varieties-definition-intersection-number}의 상황에서 교차수
$(\mathcal{L}_1 \cdots \mathcal{L}_d \cdot Z)$는 정수이다.
\end{lemma}

\begin{proof}
차수 $e$인 $n_1, \ldots, n_d$에 관한 임의의 수치 다항식은
$\mathbf{Z}$-선형 결합으로 유일하게 쓸 수 있다. 그 함수들은
${n_1 \choose k_1}{n_2 \choose k_2} \ldots {n_d \choose k_d}$이며
$k_1 + \ldots + k_d \leq e$이다. 이를 $e = d$에 적용한다.
연습문제로 남긴다.
\end{proof}

\begin{lemma}
\label{varieties-lemma-intersection-number-additive}
Definition \ref{varieties-definition-intersection-number}의 상황에서 교차수
$(\mathcal{L}_1 \cdots \mathcal{L}_d \cdot Z)$는 가법적이다. 즉,
$\mathcal{L}_i = \mathcal{L}_i' \otimes \mathcal{L}_i''$이면
$$
(\mathcal{L}_1 \cdots \mathcal{L}_i \cdots \mathcal{L}_d \cdot Z) =
(\mathcal{L}_1 \cdots \mathcal{L}_i' \cdots \mathcal{L}_d \cdot Z) +
(\mathcal{L}_1 \cdots \mathcal{L}_i'' \cdots \mathcal{L}_d \cdot Z)
$$
이다.
\end{lemma}

\begin{proof}
Lemma \ref{varieties-lemma-numerical-polynomial-from-euler}에 의해 함수
$$
(n_1, \ldots, n_{i - 1}, n_i', n_i'', n_{i + 1}, \ldots, n_d)
\mapsto
\chi(Z, \mathcal{L}_1^{\otimes n_1} \otimes
\ldots \otimes (\mathcal{L}_i')^{\otimes n_i'} \otimes
(\mathcal{L}_i'')^{\otimes n_i''} \otimes \ldots \otimes
\mathcal{L}_d^{\otimes n_d}|_Z)
$$
가 전차수 $d$ 이하이고 $d + 1$개 변수에 관한 수치 다항식이기 때문이다.
\end{proof}

\begin{lemma}
\label{varieties-lemma-intersection-number-in-terms-of-components}
Definition \ref{varieties-definition-intersection-number}의 상황에서
$Z_i \subset Z$를 차원 $d$인 기약 성분들이라 하자.
$m_i = \text{length}_{\mathcal{O}_{X, \xi_i}}(\mathcal{O}_{Z, \xi_i})$라 하자.
여기서 $\xi_i \in Z_i$는 일반점이다.
그러면
$$
(\mathcal{L}_1 \cdots \mathcal{L}_d \cdot Z) =
\sum m_i(\mathcal{L}_1 \cdots \mathcal{L}_d \cdot Z_i)
$$
\end{lemma}

\begin{proof}
Lemma \ref{varieties-lemma-numerical-polynomial-leading-term}와 정의들에서 즉시 따른다.
\end{proof}

\begin{lemma}
\label{varieties-lemma-intersection-number-and-pullback}
$k$를 체라 하자. $f : Y \to X$를 $k$ 위의 고유 스킴들 사이의 사상이라
하자. $Z \subset Y$를 차원 $d$인 정수적 닫힌 부분스킴이라 하고,
$\mathcal{L}_1, \ldots, \mathcal{L}_d$를 가역 $\mathcal{O}_X$-가군들이라 하자.
그러면
$$
(f^*\mathcal{L}_1 \cdots f^*\mathcal{L}_d \cdot Z) =
\deg(f|_Z : Z \to f(Z)) (\mathcal{L}_1 \cdots \mathcal{L}_d \cdot f(Z))
$$
이다. 여기서 $\deg(Z \to f(Z))$는 Morphisms, Definition
\ref{morphisms-definition-degree}에서와 같고, $0$이다. 후자는
$\dim(f(Z)) < d$인 경우이다.
\end{lemma}

\begin{proof}
왼쪽 변은 함수에서 $n_1 \ldots n_d$의 계수를 이용하여 계산된다.
$$
\chi(Y, \mathcal{O}_Z \otimes f^*\mathcal{L}_1^{\otimes n_1} \otimes
\ldots \otimes f^*\mathcal{L}_d^{\otimes n_d}) =
\sum (-1)^i
\chi(X, R^if_*\mathcal{O}_Z \otimes
\mathcal{L}_1^{\otimes n_1} \otimes \ldots \otimes
\mathcal{L}_d^{\otimes n_d})
$$
등식은 Lemma \ref{varieties-lemma-euler-characteristic-morphism}와 사영 공식
(Cohomology, Lemma \ref{cohomology-lemma-projection-formula})에서 따른다.
$f(Z)$의 차원이 $< d$이면 오른쪽 변은 Lemma
\ref{varieties-lemma-numerical-polynomial-from-euler}에 의해 전차수 $< d$인
다항식이므로 결과가 성립한다. $\dim(f(Z)) = d$라고 가정하자.
$\xi \in f(Z)$를 일반점이라 하자. 차원론(Lemmas
\ref{varieties-lemma-dimension-locally-algebraic}와
\ref{varieties-lemma-dimension-fibres-locally-algebraic}를 보라)에 의해 $Z$의
일반점은 $Z$의 점 가운데 $\xi$로 사상되는 유일한 점이다.
그러면 $f : Z \to f(Z)$는 $f(Z)$의 비어 있지 않은 열린집합 위에서
유한이다. Morphisms, Lemma \ref{morphisms-lemma-generically-finite}를 보라.
따라서 $\deg(f : Z \to f(Z))$는 정의되고, 실제로 이는
$f_*\mathcal{O}_Z$의 줄기가 $\xi$에서 $\mathcal{O}_{X, \xi}$ 위에 갖는 길이와
같다. 더욱이 $R^if_*\mathcal{O}_X$의 $\xi$에서의 줄기는 $i > 0$일 때
영이다. 방금 $f|_Z$가 $\xi$의 근방에서 유한임을 보았기 때문이다
(따라서 Cohomology of Schemes, Lemma
\ref{coherent-lemma-finite-pushforward-coherent}가 이 소멸을 준다).
따라서 항 $\chi(X, R^if_*\mathcal{O}_Z \otimes
\mathcal{L}_1^{\otimes n_1} \otimes \ldots \otimes
\mathcal{L}_d^{\otimes n_d})$ 가운데 $i > 0$인 것들은 전차수 $< d$이고,
$$
\chi(X, f_*\mathcal{O}_Z \otimes
\mathcal{L}_1^{\otimes n_1} \otimes \ldots \otimes
\mathcal{L}_d^{\otimes n_d})
=
\deg(f : Z \to f(Z)) \chi(f(Z),
\mathcal{L}_1^{\otimes n_1} \otimes \ldots \otimes
\mathcal{L}_d^{\otimes n_d}|_{f(Z)})
$$
이다. 이는 Lemma \ref{varieties-lemma-numerical-polynomial-leading-term}에 의해
전차수 $< d$인 다항식을 법으로 한 등식이다. 원하는 결과가 따른다.
\end{proof}

\begin{lemma}
\label{varieties-lemma-numerical-intersection-effective-Cartier-divisor}
$k$를 체라 하자. $X$를 $k$ 위에서 고유라 하자. $Z \subset X$를
차원 $d$인 닫힌 부분스킴이라 하자. $\mathcal{L}_1, \ldots, \mathcal{L}_d$를
가역 $\mathcal{O}_X$-가군들이라 하자. 유효 카르티에 인자 $D \subset Z$가
존재하여 $\mathcal{L}_1|_Z \cong \mathcal{O}_Z(D)$라고 가정하자. 그러면
$$
(\mathcal{L}_1 \cdots \mathcal{L}_d \cdot Z) =
(\mathcal{L}_2 \cdots \mathcal{L}_d \cdot D)
$$
\end{lemma}

\begin{proof}
$X$를 $Z$로, $\mathcal{L}_i$를 $\mathcal{L}_i|_Z$로 바꾸어도 된다.
따라서 $X = Z$이고 $\mathcal{L}_1 = \mathcal{O}_X(D)$라고 가정해도 된다.
그러면 $\mathcal{L}_1^{-1}$은 $D$의 아이디얼층이고, 짧은 완전열
$$
0 \to \mathcal{L}_1^{\otimes -1} \to \mathcal{O}_X \to \mathcal{O}_D \to 0
$$
을 생각할 수 있다. 다음과 같이 놓자.
$P(n_1, \ldots, n_d) =
\chi(X, \mathcal{L}_1^{\otimes n_1} \otimes \ldots \otimes
\mathcal{L}_d^{\otimes n_d})$
그리고
$Q(n_1, \ldots, n_d) =
\chi(D, \mathcal{L}_1^{\otimes n_1} \otimes \ldots \otimes
\mathcal{L}_d^{\otimes n_d}|_D)$.
가법성에서 다음을 얻는다.
$$
P(n_1, \ldots, n_d) - P(n_1 - 1, n_2, \ldots, n_d) =
Q(n_1, \ldots, n_d)
$$
$P$의 전차수가 $d$ 이하이므로, $n_1 \ldots n_d$의 계수는 $P$에서
취했을 때 $n_2 \ldots n_d$의 계수를 $Q$에서 취한 것과 같다.
\end{proof}

\begin{lemma}
\label{varieties-lemma-ample-positive}
$k$를 체라 하자. $X$를 $k$ 위에서 고유라 하자. $Z \subset X$를 차원
$d$인 닫힌 부분스킴이라 하자. $\mathcal{L}_1, \ldots, \mathcal{L}_d$가
풍부하면 $(\mathcal{L}_1 \cdots \mathcal{L}_d \cdot Z)$는 양수이다.
\end{lemma}

\begin{proof}
$d$에 관한 귀납법으로 증명한다. 경우 $d = 0$은 Lemma
\ref{varieties-lemma-chi-tensor-finite}에서 따른다. $d > 0$이라고 가정하자.
Lemma \ref{varieties-lemma-intersection-number-in-terms-of-components}에 의해
$Z$가 정수적 닫힌 부분스킴이라고 가정해도 된다. 실제로 $X$를 $Z$로,
$\mathcal{L}_i$를 $\mathcal{L}_i|_Z$로 바꾸어 $Z = X$가 차원 $d$인
고유 다양체인 경우로 환원할 수 있다.
Lemma \ref{varieties-lemma-intersection-number-additive}에 의해 $\mathcal{L}_1$을
양의 텐서 거듭제곱으로 바꾸어도 된다. 따라서 영이 아닌 단면
$s \in \Gamma(X, \mathcal{L}_1)$이 존재하여 $X_s$가 아핀이라고 가정해도
된다(여기서는 풍부한 가역층의 정의를 이용한다. Properties, Definition
\ref{properties-definition-ample}을 보라).
$X$는 아핀이 아님을 관찰하자. 고유이면서 아핀이면 유한이기 때문이다
(Morphisms, Lemma \ref{morphisms-lemma-finite-proper}). 이는 $d > 0$과
모순이다. 따라서 $s$의 영점 스킴 $Z(s) \subset X$는 비어 있지 않다.
$X$가 다양체이므로 $s$는 $\mathcal{L}_1$의 정칙 단면이고, 따라서
$Z(s)$는 유효 카르티에 인자이다. 그러므로 $Z(s)$의 여차원은
$1$이고 이는 $X$ 안에서의 여차원이다. 따라서 $Z(s)$의 차원은
$d - 1$이다(여기서는 Divisors, Sections
\ref{divisors-section-effective-Cartier-divisors},
\ref{divisors-section-effective-Cartier-invertible},
\ref{divisors-section-Noetherian-effective-Cartier}의 내용과 Lemma
\ref{varieties-lemma-dimension-locally-algebraic}에서와 같은 차원론을 이용한다).
Lemma \ref{varieties-lemma-numerical-intersection-effective-Cartier-divisor}에 의해
$$
(\mathcal{L}_1 \cdots \mathcal{L}_d \cdot X) =
(\mathcal{L}_2 \cdots \mathcal{L}_d \cdot Z(s))
$$
이다. 귀납법에 의해 오른쪽 변은 양수이며 증명이 끝난다.
\end{proof}

\begin{definition}
\label{varieties-definition-degree}
$k$를 체라 하자. $X$를 $k$ 위의 고유 스킴이라 하자.
$\mathcal{L}$을 풍부한 가역 $\mathcal{O}_X$-가군이라 하자.
임의의 닫힌 부분스킴에 대해 {\it $Z$의 $\mathcal{L}$에 관한 차수}를
$\deg_\mathcal{L}(Z)$로 나타내고, 이를 교차수
$(\mathcal{L}^d \cdot Z)$로 정의한다. 여기서 $d = \dim(Z)$이다.
\end{definition}

\noindent
Lemma \ref{varieties-lemma-ample-positive}에 의해 부분스킴의 차수는 언제나 양의
정수이다. 다음을 이용하면 $\deg_\mathcal{L}(Z) = d$인 것은
$$
\chi(Z, \mathcal{L}^{\otimes n}|_Z) = \frac{d}{\dim(Z)!} n^{\dim(Z)} + l.o.t
$$
인 것과 동치임을 알 수 있다.
$$
(n_1 + \ldots + n_{\dim(Z)})^{\dim(Z)} =
\dim(Z)!\ n_1 \ldots n_{\dim(Z)} + \text{other terms}
$$

\begin{lemma}
\label{varieties-lemma-degree-finite-morphism-in-terms-degrees}
$k$를 체라 하자. $f : Y \to X$를 $k$ 위의 고유 다양체들 사이의 유한
우세 사상이라 하자. $\mathcal{L}$을 풍부한 가역
$\mathcal{O}_X$-가군이라 하자. 그러면
$$
\deg_{f^*\mathcal{L}}(Y) = \deg(f) \deg_\mathcal{L}(X)
$$
이다. 여기서 $\deg(f)$는 Morphisms, Definition
\ref{morphisms-definition-degree}에서와 같다.
\end{lemma}

\begin{proof}
Morphisms, Lemma
\ref{morphisms-lemma-pullback-ample-tensor-relatively-ample}에 의해
$f^*\mathcal{L}$이 풍부하므로 명제는 의미가 있다. 이를 말하고 나면 결과는
Lemma \ref{varieties-lemma-intersection-number-and-pullback}의 특별한 경우이다.
\end{proof}

\noindent
마지막으로 곡선과의 교차수를 Section \ref{varieties-section-divisors-curves}에서
도입한 곡선 위 가역 가군의 차수 개념과 연관시킨다.

\begin{lemma}
\label{varieties-lemma-intersection-numbers-and-degrees-on-curves}
$k$를 체라 하자. $X$를 $k$ 위의 고유 스킴이라 하자.
$Z \subset X$를 차원 $\leq 1$인 닫힌 부분스킴이라 하자.
$\mathcal{L}$을 가역 $\mathcal{O}_X$-가군이라 하자. 그러면
$$
(\mathcal{L} \cdot Z) = \deg(\mathcal{L}|_Z)
$$
이다. 여기서 $\deg(\mathcal{L}|_Z)$는 Definition
\ref{varieties-definition-degree-invertible-sheaf}에서와 같다. $\mathcal{L}$이
풍부하면 $\deg_\mathcal{L}(Z) = \deg(\mathcal{L}|_Z)$이다.
\end{lemma}

\begin{proof}
함수 $n \mapsto \chi(Z, \mathcal{L}|_Z^{\otimes n})$의 차수가 $1$이므로
그 최고차항 계수는 연속한 두 값의 차이라는 사실에서 따른다.
\end{proof}

\begin{proposition}[점근적 리만--로흐]
\label{varieties-proposition-asymptotic-riemann-roch}
$k$를 체라 하자. $X$를 $k$ 위의 고유 스킴 가운데 차원 $d$인 것이라 하자.
$\mathcal{L}$을 풍부한 가역 $\mathcal{O}_X$-가군이라 하자. 그러면
$$
\dim_k \Gamma(X, \mathcal{L}^{\otimes n}) \sim c n^d + l.o.t.
$$
이다. 여기서 $c = \deg_\mathcal{L}(X)/d!$는 양의 상수이다.
\end{proposition}

\begin{proof}
정의들, Lemma \ref{varieties-lemma-ample-positive}, 그리고 Cohomology of Schemes,
Lemma \ref{coherent-lemma-vanshing-gives-ample}의 고차 코호몰로지 소멸에서
따른다.
\end{proof}

\section{매장 차원}
\label{varieties-section-embedding-dimension}

\noindent
매장 차원을 정의하는 방법은 몇 가지 있지만, 대수적으로 닫힌 체 위의
대수적 스킴의 닫힌 점들에 대해서는 모든 정의가 다음 정의와 동치이다.

\begin{definition}
\label{varieties-definition-embed-dim}
$k$를 대수적으로 닫힌 체라 하자. $X$를 국소 대수적 $k$-스킴이라 하고
$x \in X$를 닫힌 점이라 하자. {\it $X$의 $x$에서의 매장 차원}은
$\dim_k \mathfrak m_x/\mathfrak m_x^2$이다.
\end{definition}

\noindent
매장 차원에 관한 사실들을 적는다.
$k, X, x$는 Definition \ref{varieties-definition-embed-dim}에서와 같다고 하자.
\begin{enumerate}
\item $X$의 $x$에서의 매장 차원은 접공간
$T_{X/k, x}$ (Definition \ref{varieties-definition-tangent-space})의
$k$-벡터 공간으로서의 차원이다.
\item $X$의 $x$에서의 매장 차원은 다음 조건을 만족하는 최소 정수
$d \geq 0$이다. 즉, $k$-대수들의 전사
$$
k[[x_1, \ldots, x_d]] \longrightarrow \mathcal{O}_{X, x}^\wedge
$$
가 존재한다.
\item $X$의 $x$에서의 매장 차원은 다음 조건을 만족하는 최소 정수
$d \geq 0$이다. 즉, 열린 근방 $U \subset X$가 $x$를 포함하고 닫힌 몰입
$U \to Y$가 존재하며, 여기서 $Y$는 차원 $d$인 $k$ 위의 매끄러운
다양체이다.
\item $X$의 $x$에서의 매장 차원은 다음 조건을 만족하는 최소 정수
$d \geq 0$이다. 즉, 열린 근방 $U \subset X$가 $x$를 포함하고 비분기 사상
$U \to \mathbf{A}^d_k$가 존재한다.
\item 닫힌 매장 $X \to Y$가 주어지고 $Y$가 $k$ 위에서 매끄럽다면,
$X$의 $x$에서의 매장 차원은 다음 조건을 만족하는 최소 정수 $d \geq 0$이다.
즉, 닫힌 부분스킴 $Z \subset Y$가 존재하여 $X \subset Z$이고,
$Z \to \Spec(k)$는 $x$에서 매끄러우며 $\dim_x(Z) = d$이다.
\end{enumerate}
이 사실들이 필요해지면 정확한 결과를 정식화하고 증명을 제시할 것이다.

\medskip\noindent
대수적으로 닫혀 있지 않은 바탕체 또는 닫혀 있지 않은 점들을 생각하자.
$k$를 체라 하고 $X$를 국소 대수적 $k$-스킴이라 하자.
$x \in X$가 점이면 $X$의 $x$에서의 매장 차원에 관해 몇 가지 선택지가
있다. 즉, 다음 중 하나를 사용할 수 있다.
\begin{enumerate}
\item $\dim_{\kappa(x)}(\mathfrak m_x/\mathfrak m_x^2)$,
\item $\dim_{\kappa(x)}(T_{X/k, x}) =
\dim_{\kappa(x)}(\Omega_{X/k, x} \otimes_{\mathcal{O}_{X, x}} \kappa(x))$
(Lemma \ref{varieties-lemma-tangent-space-cotangent-space}),
\item 다음 조건을 만족하는 최소 정수 $d \geq 0$. 즉, 열린 근방
$U \subset X$가 $x$를 포함하고 닫힌 몰입 $U \to Y$가 존재하며, 여기서
$Y$는 차원 $d$인 $k$ 위의 매끄러운 다양체이다.
\end{enumerate}
표수가 영이면 (1) $=$ (2)이고, 이는 $x$가 닫힌 점일 때 성립한다. 더 일반적으로
$\kappa(x)$가 $k$ 위에서 분리 대수적이면 이것이 성립한다. Lemma
\ref{varieties-lemma-tangent-space-rational-point}를 보라. 기하학적 정의 (3)은
``매장 차원''이라는 말이 불러일으키는 기하학적 직관에 가장 가까워 보인다.
(3)과 (2)가 같은 수를 정의함을 보일 수 있으므로(이는 Lemma
\ref{varieties-lemma-immersion-into-affine}에서 따른다), 이것을 사용하겠다.
우리의 용어에서는 바탕체에 상대적인 매장 차원을 취한다는 점을 명시하겠다.

\begin{definition}
\label{varieties-definition-embedding-dimension}
$k$를 체라 하자. $X$를 국소 대수적 $k$-스킴이라 하자.
$x \in X$를 점이라 하자. {\it $X/k$의 $x$에서의 매장 차원}은
$\dim_{\kappa(x)}(T_{X/k, x})$이다.
\end{definition}

\noindent
$(A, \mathfrak m, \kappa)$가 뇌터 국소환이면 {\it $A$의 매장 차원}을
$\mathfrak m/\mathfrak m^2$의 $\kappa$ 위 차원으로 정의하기도 한다.
위에서 $A$가 체 $k$ 위의 대수로 주어지면
$\dim_\kappa(\Omega_{A/k} \otimes_A \kappa)$를 사용하는 편이 나을 수
있음을 보았다. 이 양을 {\it $A/k$의 매장 차원}이라 부르자.
이 용어를 쓰면 다음을 얻는다.
$$
\text{embed dim of }X/k\text{ at }x =
\text{embed dim of }\mathcal{O}_{X, x}/k =
\text{embed dim of }\mathcal{O}_{X, x}^\wedge/k
$$
여기서 $k, X, x$는 Definition \ref{varieties-definition-embedding-dimension}에서와
같다.

\section{베르티니 정리}
\label{varieties-section-bertini}

\noindent
이 절에서는 다음 형태의 결과들을 증명한다. 매끄러운 사영 다양체 $X$가
체 $k$ 위에 주어지면, 매끄러운 풍부한 인자 $H \subset X$가 존재한다.

\begin{lemma}
\label{varieties-lemma-generate-over-complement}
$k$를 체라 하자. $X$를 $k$ 위의 고유 스킴이라 하자. $\mathcal{L}$을
풍부한 가역 $\mathcal{O}_X$-가군이라 하자. $Z \subset X$를 닫힌
부분스킴이라 하자. 그러면 정수 $n_0$가 존재하여 모든 $n \geq n_0$에
대해 핵 $V_n$을 생각하자. 이는 사상
$\Gamma(X, \mathcal{L}^{\otimes n}) \to \Gamma(Z, \mathcal{L}^{\otimes n}|_Z)$의
핵이며 $\mathcal{L}^{\otimes n}|_{X \setminus Z}$를 생성하고,
표준 사상
$$
X \setminus Z \longrightarrow \mathbf{P}(V_n)
$$
은 $k$ 위의 스킴들의 몰입이다.
\end{lemma}

\begin{proof}
$\mathcal{I} \subset \mathcal{O}_X$를 $Z$의 준연접 아이디얼층이라 하자.
포함
$\mathcal{I} \otimes_{\mathcal{O}_X} \mathcal{L}^{\otimes n} \subset
\mathcal{L}^{\otimes n}$을 통해
$V_n = \Gamma(X, \mathcal{I} \otimes_{\mathcal{O}_X} \mathcal{L}^{\otimes n})$이다.
$n_1$을 택하여 $n \geq n_1$일 때 층
$\mathcal{I} \otimes \mathcal{L}^{\otimes n}$이 전역 생성되게 하자.
Properties, Proposition
\ref{properties-proposition-characterize-ample}을 보라. 그러면
$V_n$은 $\mathcal{L}^{\otimes n}|_{X \setminus Z}$를
$n \geq n_1$에 대해 생성한다.

\medskip\noindent
$n \geq n_1$에 대해 제한 사상을
$\psi_n : V_n \to
\Gamma(X \setminus Z, \mathcal{L}^{\otimes n}|_{X \setminus Z})$로 나타내자.
Constructions, Example \ref{constructions-example-projective-space}에 의해
표준 사상
$$
\varphi = \varphi_{\mathcal{L}^{\otimes n}|_{X \setminus Z}, \psi_n} :
X \setminus Z
\longrightarrow
\mathbf{P}(V_n)
$$
을 얻는다. $n_2$를 택하여 모든 $n \geq n_2$에 대해 가역층
$\mathcal{L}^{\otimes n}$이 $X$ 위에서 매우 풍부하게 하자.
$n_0 = n_1 + n_2$가 조건을 만족한다고 주장한다.

\medskip\noindent
주장을 증명하자. $n \geq n_0$라 하고 $n = n_1 + n'$으로 쓰자.
$x \in X \setminus Z$에 대해 $s_1 \in V_1$을 택하여 $x$에서 소멸하지
않게 할 수 있다. $V' = \Gamma(X, \mathcal{L}^{\otimes n'})$로 놓자.
$n$과 $n'$의 선택에 의해 대응하는 사상
$\varphi' : X \to \mathbf{P}(V')$은 닫힌 몰입이다. 따라서
$s' \in \Gamma(X, \mathcal{L}^{\otimes n'})$을 택하여 $x$에서 소멸하지 않게 하면
$X_{s'} = (\varphi')^{-1}(D_+(s'))$는(Constructions, Lemma
\ref{constructions-lemma-invertible-map-into-proj}를 보라) 아핀이고
$X_{s'} \to D_+(s')$는 닫힌 몰입이다.
그러면 $s = s_1 \otimes s' \in V_n$은 $x$에서 소멸하지 않는다.
$D_+(s) \subset \mathbf{P}(V_n)$가 우리 사영공간에서 이에 대응하는 열린
아핀공간을 나타내면, $\varphi^{-1}(D_+(s)) = X_s \subset X \setminus Z$이다
(위의 참고문헌을 보라). 열린집합 $X_s = X_{s'} \cap X_{s_1}$은 아핀이다.
Properties, Lemma \ref{properties-lemma-affine-cap-s-open}을 보라.
사상 $X_s \to D_+(s)$를 정의하는 환 준동형
$$
\text{Sym}(V)_{(s)} \longrightarrow \mathcal{O}_X(X_s)
$$
을 생각하자. $X_{s'} \to D_+(s')$가 닫힌 몰입이므로 원소들
$$
\frac{s_1 \otimes t'}{s_1 \otimes s'}
$$
의 상들을 생각하자. 여기서 $t' \in V'$이다. 이 상들은
$\mathcal{O}_X(X_{s'}) \to \mathcal{O}_X(X_s)$의 상을 생성한다.
$X_s \to X_{s'}$가 열린 몰입이므로, 이는
$X_s \to D_+(s)$가 아핀 스킴들의 몰입임을 뜻한다(아래를 보라).
따라서 Morphisms, Lemma \ref{morphisms-lemma-check-immersion}에 의해
$\varphi_n$은 몰입이다.

\medskip\noindent
$a : A' \to A$와 $c : B \to A$를 환 준동형들이라 하자.
$\Spec(a)$가 몰입이고 $\Im(a) \subset \Im(c)$라고 하자.
$B' = A' \times_A B$로 놓고 사영들을 $b : B' \to B$와
$c' : B' \to A'$라 하자. 가정에 의해 $c'$은 전사이고 따라서
$\Spec(c')$은 닫힌 몰입이다. 그러므로
$\Spec(c') \circ \Spec(a)$는 몰입이다(Schemes, Lemma
\ref{schemes-lemma-composition-immersion}). 이제 $\Spec(c)$는 몰입
$\Spec(c') \circ \Spec(a) = \Spec(b) \circ \Spec(c)$을 분해하므로
몰입이어야 한다. Morphisms, Lemma
\ref{morphisms-lemma-immersion-permanence}를 보라.
\end{proof}

\begin{situation}
\label{varieties-situation-family-divisors}
$k$를 체라 하고,
$X$를 $k$ 위의 스킴이라 하고,
$\mathcal{L}$을 가역 $\mathcal{O}_X$-가군이라 하고,
$V$를 유한 차원 $k$-벡터 공간이라 하고,
$\psi : V \to \Gamma(X, \mathcal{L})$를 $k$-선형 사상이라 하자.
$\dim(V) = r$라 하고 기저 $v_1, \ldots, v_r$가 $V$에 주어졌다고 하자.
그러면 ``보편 인자''
$$
H_{univ} = Z(s_{univ}) \subset  \mathbf{A}^r \times_k X
$$
를 얻는다. 이는 단면
$$
s_{univ} = \sum\nolimits_{i = 1, \ldots, r} x_i \psi(v_i) \in
\Gamma(\mathbf{A}^r \times_k X, \text{pr}_2^*\mathcal{L})
$$
의 영점 스킴이다(Divisors, Definition
\ref{divisors-definition-zero-scheme-s}). 체 확대 $k'/k$에 대해
$k'$-점 $v \in \mathbf{A}^r_k(k')$들은 벡터
$(a_1, \ldots, a_r)$에 대응하며 그 성분들은 $k'$의 원소이다. 따라서 한편으로 $v$를 원소
$v = \sum_{i = 1, \ldots, r} a_i v_i \in V \otimes_k k'$로 생각할 수
있고, 다른 한편으로 $v$에 단면
$$
\psi(v) = \sum\nolimits_{i = 1, \ldots, r} a_i \psi(v_i) \in
\Gamma(X_{k'}, \mathcal{L}|_{X_{k'}})
$$
을 대응시킬 수 있다. 이 표기에서 $H_{univ}$의
$v \in V \otimes k'$ 위 올은 $\psi(v)$의 영점 스킴임이 분명하다. 식으로 쓰면
$$
H_v = H_{univ, v} = Z(\psi(v))
$$
이다. 표시한 대로 이 공통값을 $H_v$로 나타내겠다. 마지막으로 이 상황에서
$P$를 벡터 $v \in V \otimes_k k'$의 성질이라 하자. 여기서 $k'/k$는
임의의 체 확대이다\footnote{예를 들어 $H_v$가 $k'$ 위에서 매끄럽다는 조건이나
$k'$ 위에서 기하적으로 기약이라는 조건을 생각할 수 있다.}. $P$가
{\it 일반적으로 성립한다}고 말하자. 여기서 대상은
$v \in V \otimes_k k'$이며, 그 뜻은 비어 있지 않은 자리스키 열린집합
$U \subset \mathbf{A}^r_k$가 존재하여, $v$가 $k'$-점으로서 $U$에 대응하고
임의의 $k'/k$에 대해 이 조건이 성립할 때 $P(v)$가 성립한다는 것이다.
\end{situation}

\begin{lemma}
\label{varieties-lemma-bertini}
Situation \ref{varieties-situation-family-divisors}에서 다음을 가정하자.
\begin{enumerate}
\item $X$는 $k$ 위에서 매끄럽다,
\item $\psi : V \to \Gamma(X, \mathcal{L})$의 상은 $\mathcal{L}$을 생성한다,
\item 대응하는 사상
$\varphi_{\mathcal{L}, \psi} : X \to \mathbf{P}(V)$는 몰입이다.
\end{enumerate}
그러면 일반적인 $v \in V \otimes_k k'$에 대해 스킴 $H_v$는
$k'$ 위에서 매끄럽다.
\end{lemma}

\begin{proof}
($X$는 사영공간의 국소 닫힌 부분스킴이므로 분리이고 유한형임을
관찰한다.) 보조정리의 진술 위에서 도입한 표기를 사용하자.
사영들을 생각한다.
$$
\xymatrix{
\mathbf{A}^r_k \times_k X \ar[d] &
H_{univ} \ar[l] \ar[ld]^p \ar[r] \ar[rd]_q &
\mathbf{A}^r_k \times_k X \ar[d] \\
X & &
\mathbf{A}^r_k
}
$$
$\Sigma \subset H_{univ}$를 사상 $q : H_{univ} \to \mathbf{A}^r_k$의
특이 궤적, 즉 $q$가 매끄럽지 않은 점들의 집합이라 하자. 사상의 매끄러운
궤적은 정의상 열려 있으므로 $\Sigma$는 닫혀 있다. 매끄러운 사상의 올은
매끄러우므로 $q(\Sigma)$가 $\mathbf{A}^r_k$의 진부분 닫힌 부분집합에
포함됨을 증명하면 충분하다. $\Sigma$는 (유도된 축약 스킴 구조를 주면)
$k$ 위의 유한형 스킴이므로 $\dim(\Sigma) < r$임을 증명하면 충분하다.
이는 Lemma \ref{varieties-lemma-dimension-fibres-locally-algebraic}에서 따른다.
$k$를 더 큰 체로 바꾸어도 차원은 변하지 않으므로(Morphisms, Lemma
\ref{morphisms-lemma-dimension-fibre-after-base-change}), $k$가 대수적으로
닫혀 있다고 가정하며 그렇게 하겠다. 차원론(Lemma
\ref{varieties-lemma-dimension-fibres-locally-algebraic})에 의해, 닫힌
$x \in X \setminus Z$에 대해
$p^{-1}(\{x\}) \cap \Sigma$의 차원이 $< r - \dim(X')$임을 증명하면
충분하다. 여기서 $X'$은 $X$의 유일한 기약 성분 가운데 $x$를 포함하는 것이다.
$X$가 $k$ 위에서 매끄럽고 $x$가 닫힌 점이므로
$\dim(X') = \dim \mathfrak m_x/\mathfrak m_x^2$이다
(Morphisms, Lemma \ref{morphisms-lemma-smooth-omega-finite-locally-free}와
Algebra, Lemma \ref{algebra-lemma-rank-omega}). 따라서 모든 닫힌
$x \in X$에 대해
$$
\dim p^{-1}(x) \cap \Sigma < r - \dim \mathfrak m_x/\mathfrak m_x^2
$$
이면 원하는 결과를 얻는다.

\medskip\noindent
$V$가 $\mathcal{L}$을 전역 생성하므로, 기약 성분 $X'$을 $X$에서 택할 때마다
$\mathbf{A}^r$의 비어 있지 않은 자리스키 열린집합이 존재하여 그 위의
$q$의 올들은 $X'$을 포함하지 않는다. (예를 들어 $x' \in X'$이 닫힌
점이면, 벡터 $v \in V$들 가운데 $\psi(v)$가 $x'$에서 소멸하지 않는 것에 대응하는
열린집합을 취할 수 있다. 이 열린집합은 $\mathbf{A}^r_k$ 안의 초평면의
여집합일 것이다.) $U \subset \mathbf{A}^r$을 이 열린집합들의 (비어 있지
않은) 교집합이라 하자. 그러면 $q^{-1}(U) \to U$의 올들은
$U \times_k X \to U$의 올들 위의 유효 카르티에 인자들이다(정수적 스킴
위의 가역 가군의 영이 아닌 단면은 정칙 단면이기 때문이다). 따라서 사상
$q^{-1}(U) \to U$는 Divisors, Lemma
\ref{divisors-lemma-fibre-Cartier}에 의해 평탄하다. 그러므로 닫힌
$x \in X$와 $v \in V = \mathbf{A}^r_k(k)$에 대해,
$(x, v) \in H_{univ}$, 즉 $x \in H_v$이면, $q$가 $(x, v)$에서
매끄러운 것은 올 $H_v$가 $x$에서 매끄러운 것과 동치이다. Morphisms,
Lemma \ref{morphisms-lemma-smooth-at-point}를 보라.

\medskip\noindent
$\psi(v)_x$라는 상을 생각하자. 이는 줄기 $\mathcal{L}_x$ 안의 상이며,
해당 단면은 $v \in V$에 대응한다. 다음이 성립한다.
$$
x \in H_v \Leftrightarrow \psi(v)_x \in \mathfrak m_x\mathcal{L}_x
$$
이것이 참이면 다음이 성립한다.
$$
H_v\text{ singular at }x \Leftrightarrow
\psi(v)_x \in \mathfrak m_x^2\mathcal{L}_x
$$
실제로 $\psi(v)_x$가 $\mathfrak m_x^2\mathcal{L}_x$에 포함되지 않는다
$\Leftrightarrow$
$H_v \subset X$의 $x$에서의 국소 방정식이 $\mathfrak m_x^2$에 포함되지 않는다
$\Leftrightarrow$
$\mathcal{O}_{H_v, x}$가 정칙이다
(Algebra, Lemma \ref{algebra-lemma-regular-ring-CM})
$\Leftrightarrow$
$H_v$가 $x$에서 $k$ 위로 매끄럽다
(Algebra, Lemma \ref{algebra-lemma-separable-smooth}).
따라서 $p^{-1}(x) \cap \Sigma$의 닫힌 점들은 $v \in V$ 가운데
$\psi(v)_x \in \mathfrak m_x^2\mathcal{L}_x$인 것들에 대응한다.
그러나 $\varphi_{\mathcal{L}, \psi}$가 몰입이므로 사상
$$
V \longrightarrow \mathcal{L}_x/\mathfrak m_x^2\mathcal{L}_x
$$
은 전사이다(작은 세부사항은 생략한다). 위의 논의에 의해 궤적
$p^{-1}(x) \cap \Sigma$를 $V$의 부분공간으로 보았을 때 그 닫힌 점들은 이 사상의
핵이고, 따라서 그 차원은 원하는 대로
$r - \dim \mathfrak m_x/\mathfrak m_x^2 - 1$이다.
\end{proof}

\section{엔리케스--세베리--자리츠키}
\label{varieties-section-vanishing-negative}

\noindent
이 절에서는 ``매우 음인'' 가역 가군과 텐서하면 낮은 차수의 코호몰로지가
소멸한다는 형태의 몇 가지 결과를 증명한다. 또한 차원 $\geq 2$인 정규
고유 스킴의 초곡면 단면이 연결임을 도출한다.

\begin{lemma}
\label{varieties-lemma-vanishin-h0-negative}
$k$를 체라 하자. $X$를 $k$ 위의 고유 스킴이라 하자. $\mathcal{L}$을
풍부한 가역 $\mathcal{O}_X$-가군이라 하자. $\mathcal{F}$를 연접
$\mathcal{O}_X$-가군이라 하자. $\text{Ass}(\mathcal{F})$가 닫힌 점을
하나도 포함하지 않으면
$\Gamma(X, \mathcal{F} \otimes_{\mathcal{O}_X} \mathcal{L}^{\otimes n}) = 0$이다.
이는 $n \ll 0$에 대해 성립한다.
\end{lemma}

\begin{proof}
연접 $\mathcal{O}_X$-가군 $\mathcal{F}$에 대해 $\mathcal{P}(\mathcal{F})$를
다음 성질이라 하자. 어떤 $n_0 \in \mathbf{Z}$가 존재하여 $n \leq n_0$일
때마다 모든 단면 $s$를 생각하면, 이는
$\mathcal{F} \otimes_{\mathcal{O}_X} \mathcal{L}^{\otimes n}$의 단면이고 그
지지집합은 닫힌 점들로만 이루어진다.
$\text{Ass}(\mathcal{F}) =
\text{Ass}(\mathcal{F} \otimes_{\mathcal{O}_X} \mathcal{L}^{\otimes n})$이므로
$\mathcal{P}$가 모든 연접 가군에 대해 성립하며 이 가군들은 $X$ 위에
있음을 증명하면 충분하다.
이를 위해 Cohomology of Schemes, Lemma
\ref{coherent-lemma-property-higher-rank-cohomological}의 조건 (1), (2), (3)이
성립함을 증명하겠다.

\medskip\noindent
조건 (1)을 보이기 위해
$$
0 \to \mathcal{F}_1 \to \mathcal{F} \to \mathcal{F}_2 \to 0
$$
이 연접 $\mathcal{O}_X$-가군들의 짧은 완전열이고, $\mathcal{P}$가
$\mathcal{F}_i$에 대해 성립하며 $i = 1, 2$라고 하자. 찾은 경계값들을
$n_1, n_2$라 하자. $\mathcal{F}'_2 \subset \mathcal{F}_2$를 지지집합이
닫힌 점들의 유한 집합인 최대 연접 부분가군이라 하자.
$\mathcal{I} \subset \mathcal{O}_X$를 $\mathcal{F}'_2$의 소멸자라 하자.
$\mathcal{L}$이 풍부하므로 어떤 $e > 0$을 찾아
$\mathcal{I} \otimes_{\mathcal{O}_X} \mathcal{L}^{\otimes e}$가 전역 생성되게
할 수 있다. $n_0 = \min(n_2, n_1 - e)$로 놓자. $n \leq n_0$라 하고,
$t$를 $\mathcal{F} \otimes \mathcal{L}^{\otimes n}$의 전역 단면이라 하자.
$t$의 상을 $\mathcal{F}_2 \otimes \mathcal{L}^{\otimes n}$에서 취하면
$\mathcal{F}'_2 \otimes \mathcal{L}^{\otimes n}$에 들어간다. 실제로
$n \leq n_2$이기 때문이다.
따라서 임의의
$s \in \Gamma(X, \mathcal{I} \otimes_{\mathcal{O}_X} \mathcal{L}^{\otimes e})$에
대해 곱 $t \otimes s$는
$\mathcal{F}_1 \otimes \mathcal{L}^{\otimes n + e}$에 속한다.
따라서 $t \otimes s$의 지지집합은 $\text{Ass}(\mathcal{F}_1)$에 있는
닫힌 점들의 유한 집합에 포함된다. 실제로 $n + e \leq n_1$이다.
$e$의 선택에 의해 $s$를 택하여 $\mathcal{F}'_2$의 지지집합에 속하지
않는 임의의 점에서 가역이 되게 할 수 있다. 따라서 $t$의 지지집합은
$\text{Ass}(\mathcal{F}_1)$에 있는 닫힌 점들의 유한 집합과
$\text{Ass}(\mathcal{F}_2)$에 있는 닫힌 점들의 유한 집합의 합집합에
포함된다. 이로써 조건 (1)의 증명이 끝난다.

\medskip\noindent
조건 (2)는 즉시 성립한다.

\medskip\noindent
조건 (3)을 위해 $\mathcal{G} = \mathcal{O}_Z$로 택한다. 이 경우 $Z$가
$X$의 닫힌 점이면 보일 것이 없다. $\dim(Z) > 0$이면
$\Gamma(Z, \mathcal{L}^{\otimes n}|_Z) = 0$임을 $n < 0$에 대해 보이겠다.
실제로 $s$를 $\mathcal{L}|_Z$의 음의 거듭제곱의 영이 아닌 단면이라 하자.
$t$를 $\mathcal{L}|_Z$의 양의 거듭제곱의 영이 아닌 단면으로 택하자
($\mathcal{L}$이 풍부하므로 가능하다. Properties, Proposition
\ref{properties-proposition-characterize-ample}을 보라). 그러면
$s^{\deg(t)} \otimes t^{\deg(s)}$는 $\mathcal{O}_Z$의 영이 아닌 전역
단면이고($Z$가 정수적이기 때문이다), 따라서 단원이다(Lemma
\ref{varieties-lemma-proper-geometrically-reduced-global-sections}). 이는 $t$가
$\mathcal{L}$의 양의 거듭제곱을 자명화하는 단면임을 뜻한다. 따라서 함수
$n \mapsto \dim_k \Gamma(X, \mathcal{L}^{\otimes n})$는 양의 정수들의
무한 집합에서 유계이다. 그러나 $\dim(Z) > 0$이므로 이는 점근적
리만--로흐(Proposition \ref{varieties-proposition-asymptotic-riemann-roch})와 모순이다.
\end{proof}

\begin{lemma}[엔리케스--세베리--자리츠키]
\label{varieties-lemma-vanishin-h1-negative}
$k$를 체라 하자. $X$를 $k$ 위의 고유 스킴이라 하자. $\mathcal{L}$을
풍부한 가역 $\mathcal{O}_X$-가군이라 하자. $\mathcal{F}$를 연접
$\mathcal{O}_X$-가군이라 하자. 닫힌 $x \in X$에 대해
$\text{depth}(\mathcal{F}_x) \geq 2$라고 가정하자. 그러면
$H^1(X, \mathcal{F} \otimes_{\mathcal{O}_X} \mathcal{L}^{\otimes m}) = 0$이다.
이는 $m \ll 0$에 대해 성립한다.
\end{lemma}

\begin{proof}
닫힌 몰입 $i : X \to \mathbf{P}^n_k$를 택하여
$i^*\mathcal{O}(1) \cong \mathcal{L}^{\otimes e}$가 어떤 $e > 0$에 대해
성립하게 하자(Morphisms, Lemma
\ref{morphisms-lemma-finite-type-over-affine-ample-very-ample}를 보라).
그러면 $\mathbf{P}^n_k$ 위의
$$
\mathcal{G} =
i_*(\mathcal{F} \oplus \mathcal{F} \otimes \mathcal{L} \oplus \ldots
\oplus \mathcal{F} \otimes \mathcal{L}^{\otimes e - 1})
\quad\text{and}\quad \mathcal{O}(1)
$$
에 대해 보조정리를 증명하면 충분하다. 실제로
$$
H^1(\mathbf{P}^n_k, \mathcal{G}(m)) =
\bigoplus\nolimits_{j = 0, \ldots, e - 1}
H^1(X, \mathcal{F} \otimes \mathcal{L}^{\otimes j + me})
$$
이다. 이는 Cohomology of Schemes, Lemma
\ref{coherent-lemma-relative-affine-cohomology}에서 따른다. 또한
$y \in \mathbf{P}^n_k$가 닫힌 점이면
$\text{depth}(\mathcal{G}_y) = \infty$이고 이는 $y \not \in i(X)$일 때이며,
$\text{depth}(\mathcal{G}_y) = \text{depth}(\mathcal{F}_x)$이고 이는
$y = i(x)$일 때이다. 실제로 이
경우 $\mathcal{G}_y \cong \mathcal{F}_x^{\oplus e}$이고 이는
$\mathcal{O}_{\mathbf{P}^n_k, x}$ 위의 가군이므로, 예를 들어 Algebra,
Lemma \ref{algebra-lemma-depth-goes-down-finite}를 사용하여 등식을 얻을 수
있다.

\medskip\noindent
$X = \mathbf{P}^n_k$이고 $\mathcal{L} = \mathcal{O}(1)$이며 $k$가
무한이라고 가정하자. 완전열
$$
0 \to \mathcal{F}(-1) \xrightarrow{s} \mathcal{F} \to \mathcal{G} \to 0
$$
을 결정하는 $s \in H^0(\mathbf{P}^1_k, \mathcal{O}(1))$을 택하자.
Lemma \ref{varieties-lemma-exact-sequence-induction}에서와 같다. 사상
$\mathcal{F}(-1) \to \mathcal{F}$는 아핀 국소적으로 $\mathcal{F}$ 위의
영인자가 아닌 원소로 곱하여 주어지므로, 닫힌
$x \in \mathbf{P}^n_k$에 대해 $\text{depth}(\mathcal{G}_x) \geq 1$이다.
Algebra, Lemma \ref{algebra-lemma-depth-drops-by-one}을 보라. 따라서 Lemma
\ref{varieties-lemma-vanishin-h0-negative}에 의해 $H^0(\mathcal{G}(m)) = 0$이고,
이는 $m \ll 0$일 때 성립한다. 텐서한 뒤 코호몰로지의 긴 완전열을 보면
(Remark \ref{varieties-remark-exact-sequence-induction-cohomology}를 보라), 수열
$$
\dim H^1(\mathbf{P}^n_k, \mathcal{F}(m))
$$
은 $m \leq m_0$일 때 안정되며, 여기서 $m_0$는 어떤 정수이다. 이 공간들의
공통 차원을 $N$이라 하자. 이는 $m \leq m_0$일 때의 공통 차원이다.
$N = 0$임을 보여야 한다.

\medskip\noindent
$d > 0$이고 $m \leq m_0$일 때 쌍선형 사상
$$
H^0(\mathbf{P}^n_k, \mathcal{O}(d)) \times
H^1(\mathbf{P}^n_k, \mathcal{F}(m - d))
\longrightarrow
H^1(\mathbf{P}^n_k, \mathcal{F}(m))
$$
을 생각하자. 선형대수에 의해 여차원 $\leq N^2$인 부분공간
$V_m \subset H^0(\mathbf{P}^n_k, \mathcal{O}(d))$가 존재하여
$s' \in V_m$로 곱하면
$H^1(\mathbf{P}^n_k, \mathcal{F}(m - d))$를 소멸시킨다.
$m' < m \leq m_0$일 때 도표
$$
\xymatrix{
H^0(\mathbf{P}^n_k, \mathcal{O}(d)) \times
H^1(\mathbf{P}^n_k, \mathcal{F}(m' - d)) \ar[r] \ar[d]^{1 \times s^{m' - m}} &
H^1(\mathbf{P}^n_k, \mathcal{F}(m')) \ar[d]^{s^{m' - m}}\\
H^0(\mathbf{P}^n_k, \mathcal{O}(d)) \times
H^1(\mathbf{P}^n_k, \mathcal{F}(m - d)) \ar[r] &
H^1(\mathbf{P}^n_k, \mathcal{F}(m))
}
$$
는 가환하며 수직 사상들은 동형사상이다. 따라서 $V_m = V$는
$m \leq m_0$에 무관하다. $x \in \text{Ass}(\mathcal{F})$에 대해
$Z = \overline{\{x\}}$로 놓자. $d$가 충분히 크면 선형 사상
$$
H^0(\mathbf{P}^n_k, \mathcal{O}(d)) \to H^0(Z, \mathcal{O}(d)|_Z)
$$
의 계수는 $> N^2$이다. 이는 $\dim(Z) \geq 1$이기 때문이다(예를 들어
점근적 리만--로흐와 $\mathcal{O}(1)$의 풍부함에서 따른다. 세부사항은
생략한다). 따라서 $s' \in V$를 찾아 $s'$가 $\mathcal{F}$의 어느
연관점에서도 소멸하지 않게 할 수 있다(연관점들의 집합이 유한임을
사용한다). 그러면
$$
0 \to \mathcal{F}(-d) \xrightarrow{s'} \mathcal{F} \to \mathcal{G}' \to 0
$$
을 얻고, 앞에서와 같이 $s'$로 곱하는 사상이
$H^1(\mathbf{P}^n_k, \mathcal{F}(m - d))$ 위에서 단사라고 다시 결론짓는다.
이는 $m \ll 0$에 대해 성립한다. 그런데 이는 $s'$의 선택과 모순이며,
$N = 0$인 경우만 예외이다. 원하는 바를 얻었다.

\medskip\noindent
$k$가 유한인 경우를 아직 다루어야 한다. 이 경우 $K/k$를 임의의 무한
대수적 체 확대라 하자. $\mathcal{F}_K$와 $\mathcal{L}_K$로
$\mathcal{F}$와 $\mathcal{L}$을
$X_K = \Spec(K) \times_{\Spec(k)} X$에 당겨 올린 것들을 나타내자. 그러면
$$
H^1(X_K, \mathcal{F}_K \otimes \mathcal{L}_K^{\otimes m}) =
H^1(X, \mathcal{F} \otimes \mathcal{L}^{\otimes m}) \otimes_k K
$$
이다. 이는 Cohomology of Schemes, Lemma
\ref{coherent-lemma-flat-base-change-cohomology}에서 따른다. 한편 닫힌 점
$x_K$가 $X_K$에 있으면 닫힌 점 $x$로 사상되고, 이 점은 $X$에 있다.
$K/k$가 대수적 확대이기
때문이다. 환 준동형 $\mathcal{O}_{X, x} \to \mathcal{O}_{X_K, x_K}$는
평탄하다(Lemma \ref{varieties-lemma-change-fields-flat}). 따라서
$$
\text{depth}(\mathcal{F}_{x_K}) =
\text{depth}(\mathcal{F}_x \otimes_{\mathcal{O}_{X, x}}
\mathcal{O}_{X_K, x_K}) \geq
\text{depth}(\mathcal{F}_x)
$$
이다. Algebra, Lemma \ref{algebra-lemma-apply-grothendieck-module}에
의한다(사실 여기서는 등식이 성립하지만 필요하지 않다). 따라서 $k$ 위의
결과는 무한체 $K$ 위의 결과에서 따르며 증명이 끝난다.
\end{proof}

\begin{lemma}
\label{varieties-lemma-connectedness-ample-divisor}
$k$를 체라 하자. $X$를 $k$ 위의 고유 스킴이라 하자. $\mathcal{L}$을
풍부한 가역 $\mathcal{O}_X$-가군이라 하자.
$s \in \Gamma(X, \mathcal{L})$라 하자. 다음을 가정하자.
\begin{enumerate}
\item $s$는 정칙 단면이다
(Divisors, Definition \ref{divisors-definition-regular-section}),
\item 모든 닫힌 점 $x \in X$에 대해
$\text{depth}(\mathcal{O}_{X, x}) \geq 2$이다, 그리고
\item $X$는 연결이다.
\end{enumerate}
그러면 영점 스킴 $Z(s)$, 즉 $s$의 영점 스킴은 연결이다.
\end{lemma}

\begin{proof}
$s$가 정칙 단면이므로 $s^n \in \Gamma(X, \mathcal{L}^{\otimes n})$도
모든 $n > 1$에 대해 정칙 단면이다. 더욱이 포함
사상 $Z(s) \to Z(s^n)$은 바탕 위상공간들 사이의 전단사이다. 따라서
$Z(s)$가 비연결이면 $Z(s^n)$도 비연결이다. 이제 표준 짧은 완전열
$$
0 \to \mathcal{L}^{\otimes -n} \xrightarrow{s^n}
\mathcal{O}_X \to \mathcal{O}_{Z(s^n)} \to 0
$$
을 생각하자. $k$-대수 $R_n = \Gamma(X, \mathcal{O}_{Z(s^n)})$를 생각하자.
$Z(s)$가 비연결, 즉 $Z(s^n)$이 비연결이면 다음 두 경우가 있다.
$R_n$이 영인 경우는 $Z(s^n) = \emptyset$인 경우이고,
$R_n$이 자명하지 않은 멱등원을 포함하는 경우는
$Z(s^n) = U \amalg V$이고 $U, V \subset Z(s^n)$가 열리고 비어 있지 않은
경우이다(독자는 Lemma
\ref{varieties-lemma-proper-geometrically-reduced-global-sections}를 참고할 수 있다).
따라서 사상 $\Gamma(X, \mathcal{O}_X) \to R_n$은 동형사상일 수 없다.
따라서 $H^0(X, \mathcal{L}^{\otimes -n})$ 또는
$H^1(X, \mathcal{L}^{\otimes -n})$이 무한히 많은 양의 $n$에 대해 영이
아니다. 이는 Lemma
\ref{varieties-lemma-vanishin-h0-negative} 또는
\ref{varieties-lemma-vanishin-h1-negative}와 모순이며 증명이 끝난다.
\end{proof}
